% Generated mechanically from frozen input inputs/p01/ja/categories.tex.
% Do not edit this generated file; edit the bound locale source instead.


% OK, start here.
%

\setcounter{chapter}{3}
\chapter{圏論}



\phantomsection
\label{categories-section-phantom}


\section{序論}
\label{categories-section-introduction}

\noindent
圏は \cite{GenEqui} において初めて導入された。
圏からなる圏（これは真の類である）は $2$-圏である。同様に、スタックからなる圏も
$2$-圏をなす。圏については既に知っているが $2$-圏については知らない読者は、
後に現れる形式的定義への導入として
節 \ref{categories-section-formal-cat-cat}
を読まれたい。

\section{定義}
\label{categories-section-definition-categories}

\noindent
記法を定める意味も兼ねて、定義を振り返る。

\begin{definition}
\label{categories-definition-category}
{\it 圏} $\mathcal{C}$ とは、次のデータからなるものである：
\begin{enumerate}
\item 対象の集合 $\Ob(\mathcal{C})$。
\item 各対 $x, y \in \Ob(\mathcal{C})$ に対する射の集合
$\Mor_\mathcal{C}(x, y)$。
\item 各三つ組 $x, y, z\in \Ob(\mathcal{C})$ に対する合成写像
$ \Mor_\mathcal{C}(y, z) \times \Mor_\mathcal{C}(x, y)
\to \Mor_\mathcal{C}(x, z) $。これは $(\phi, \psi) \mapsto
\phi \circ \psi$ と表す。
\end{enumerate}
これらのデータは次の規則を満たすものとする：
\begin{enumerate}
\item 各元 $x\in \Ob(\mathcal{C})$ に対して射
$\text{id}_x\in \Mor_\mathcal{C}(x, x)$ が存在し、合成が意味をもつ限り
$\text{id}_x \circ \phi = \phi$ および $\psi \circ \text{id}_x = \psi $
が成り立つ。
\item 合成は結合的である。すなわち、合成が意味をもつ限り
$(\phi \circ \psi) \circ \chi =
\phi \circ ( \psi \circ \chi)$ が成り立つ。
\end{enumerate}
\end{definition}


\noindent
すべての射集合 $\Mor_\mathcal{C}(x, y)$ は互いに素であると仮定するのが
通例である。そうすれば、射 $\phi : x \to y$ は一意な{\it 始域} $x$ と
一意な{\it 終域} $y$ をもつ。この仮定は厳密には必要でないが、仮定しない
場合には上の条件 (2) の定式化に注意を要する。この仮定は便利なので、以下でも
しばしば採用する。このとき、$\phi$ の始域が $\psi$ の終域に等しければ
$\phi$ と $\psi$ は{\it 合成可能}であるといい、その場合に
$\phi \circ \psi$ が定義される。同値な定義として、圏を
$(\text{Ob}, \text{Arrows}, s, t, \circ)$ という五つ組、すなわち対象の集合、
射（矢印）の集合、始域、終域および合成であって、長い公理の列を満たすものとして
定義することもできる。以下では折に触れてこの見方も用いる。

\begin{remark}
\label{categories-remark-big-categories}
大きな圏。文献によっては、対象が真の類をなす圏も許される。本書でもこれを
許すが、以下に列挙する場合に限る（この一覧は必要に応じて更新する）。特に
「$\mathcal{C}$ を圏とする」と述べたときは、$\Ob(\mathcal{C})$ が集合である
ことを含意する。
\begin{enumerate}
\item 集合の圏 $\textit{Sets}$。
\item アーベル群の圏 $\textit{Ab}$。
\item 群の圏 $\textit{Groups}$。
\item 群 $G$ が与えられたとき、左 $G$-作用をもつ集合の圏
$G\textit{-Sets}$。
\item 環 $R$ が与えられたとき、$R$-加群の圏 $\text{Mod}_R$。
\item 体 $k$ が与えられたとき、$k$ 上のベクトル空間の圏。
\item 環の圏。
\item 分割冪環の圏。分割冪代数、節
\ref{dpa-section-divided-power-rings} を参照されたい。
\item スキームの圏。
\item 位相空間の圏 $\textit{Top}$。
\item 位相空間 $X$ が与えられたとき、$X$ 上の集合の前層の圏
$\textit{PSh}(X)$。
\item 位相空間 $X$ が与えられたとき、$X$ 上の集合の層の圏
$\Sh(X)$。
\item 位相空間 $X$ が与えられたとき、$X$ 上のアーベル群の前層の圏
$\textit{PAb}(X)$。
\item 位相空間 $X$ が与えられたとき、$X$ 上のアーベル群の層の圏
$\textit{Ab}(X)$。
\item 小さい圏 $\mathcal{C}$ が与えられたとき、$\mathcal{C}$ から
$\textit{Sets}$ への関手の圏。
\item 圏 $\mathcal{C}$ が与えられたとき、$\mathcal{C}$ 上の集合の前層の圏。
\item サイト $\mathcal{C}$ が与えられたとき、$\mathcal{C}$ 上の集合の層の圏。
\end{enumerate}
これらをここで列挙する理由の一つは、「大きな」圏の「集まり」のようなものを
扱うことを避けるためである。それはすべての類の集まりを扱うようなものであり、
著者の考えでは明らかにメタ数学的対象である。
\end{remark}

\begin{remark}
\label{categories-remark-unique-identity}
定義から直ちに、$\mathcal{A}$ の対象 $x$ の任意の二つの恒等射は一致する。
したがって、$x$ の{\it その}恒等射 $\text{id}_x$ と呼ぶことができ、以下でも
そう呼ぶ。
\end{remark}

\begin{definition}
\label{categories-definition-isomorphism}
射 $\phi : x \to y$ が圏 $\mathcal{C}$ の{\it 同型射}であるとは、射
$\psi : y \to x$ で
$\phi \circ \psi = \text{id}_y$ および
$\psi \circ \phi = \text{id}_x$ を満たすものが存在することをいう。
\end{definition}

\noindent
同型射 $\phi$ は{\it 可逆}射とも呼ばれる。定義における射 $\psi$ を
{\it 逆射}といい、$\phi^{-1}$ と表す。逆射は、存在すれば一意である。また、
圏 $\mathcal{A}$ の対象 $x$ が与えられたとき、
$\Mor_\mathcal{A}(x, x)$ の可逆元全体 $\text{Aut}_\mathcal{A}(x)$ は
合成に関して群をなす。この群を $\mathcal{A}$ における $x$ の
{\it 自己同型群}という。

\begin{definition}
\label{categories-definition-groupoid}
すべての射が同型射である圏を{\it グルーポイド}という。
\end{definition}

\begin{example}
\label{categories-example-group-groupoid}
群 $G$ から、ただ一つの対象 $x$ と射 $\Mor(x, x) = G$ をもち、合成則が
$G$ の群演算で与えられるグルーポイドが得られる。ただ一つの対象をもつ
グルーポイドはすべてこの形である。
\end{example}

\begin{example}
\label{categories-example-set-groupoid}
集合 $C$ から次のようにグルーポイド $\mathcal{C}$ が得られる。対象は
$\Ob(\mathcal{C}) := C$ とし、射については、$x\neq y$ なら
$\Mor(x, y)$ を空集合、$x = y$ なら $\{\text{id}_x\}$ とする。
\end{example}

\begin{definition}
\label{categories-definition-functor}
二つの圏 $\mathcal{A}, \mathcal{B}$ の間の{\it 関手}
$F : \mathcal{A} \to \mathcal{B}$ とは、次のデータからなるものである：
\begin{enumerate}
\item 写像 $F : \Ob(\mathcal{A}) \to \Ob(\mathcal{B})$。
\item 各 $x, y \in \Ob(\mathcal{A})$ に対する写像
$F : \Mor_\mathcal{A}(x, y) \to \Mor_\mathcal{B}(F(x), F(y))$。
これは $\phi \mapsto F(\phi)$ と表す。
\end{enumerate}
これらのデータは合成および恒等射と次のように両立しなければならない：
$\mathcal{A}$ の射の合成可能な対 $(\phi, \psi)$ に対して
$F(\phi \circ \psi) = F(\phi) \circ F(\psi)$ が成り立ち、また
$F(\text{id}_x) = \text{id}_{F(x)}$ が成り立つ。
\end{definition}

\noindent
各圏 $\mathcal{A}$ は{\it 恒等}関手 $\text{id}_\mathcal{A}$ をもつ。
さらに、関手 $G : \mathcal{B} \to \mathcal{C}$ と関手
$F : \mathcal{A} \to \mathcal{B}$ が与えられたとき、明らかな仕方で
定義される{\it 合成}関手 $G \circ F : \mathcal{A} \to \mathcal{C}$ がある。

\begin{definition}
\label{categories-definition-faithful}
$F : \mathcal{A} \to \mathcal{B}$ を関手とする。
\begin{enumerate}
\item 任意の対象 $x, y \in \Ob(\mathcal{A})$ に対して写像
$$
F : \Mor_\mathcal{A}(x, y) \to \Mor_\mathcal{B}(F(x), F(y))
$$
が単射であるとき、$F$ は{\it 忠実}であるという。
\item これらの写像がすべて全単射であるとき、$F$ は
{\it 充満忠実}であるという。
\item
任意の対象 $y \in \Ob(\mathcal{B})$ に対して、$F(x)$ が
$\mathcal{B}$ において $y$ と同型になるような対象
$x \in \Ob(\mathcal{A})$ が存在するとき、関手 $F$ は
{\it 本質的全射}であるという。
\end{enumerate}
\end{definition}

\begin{definition}
\label{categories-definition-subcategory}
圏 $\mathcal{B}$ の{\it 部分圏}とは、対象および射がそれぞれ
$\mathcal{B}$ の対象および射の部分集合をなし、$\mathcal{A}$ における
始域、終域および合成が $\mathcal{B}$ のものと一致し、さらに各対象の恒等射が $\mathcal{B}$ に
おける恒等射と一致するような圏 $\mathcal{A}$ のことである。すべての
$x, y \in \Ob(\mathcal{A})$ に対して $\Mor_\mathcal{A}(x, y)
= \Mor_\mathcal{B}(x, y)$ が成り立つとき、$\mathcal{A}$ は
$\mathcal{B}$ の{\it 充満部分圏}であるという。$\mathcal{A}$ が充満部分圏で
あり、かつ $x \in \Ob(\mathcal{A})$ が与えられたとき、$x$ と同型な $\mathcal{B}$ の対象が
すべて $\mathcal{A}$ にも属するとき、$\mathcal{A}$ は $\mathcal{B}$ の
{\it 厳密充満}部分圏であるという。
\end{definition}

\noindent
$\mathcal{A} \subset \mathcal{B}$ が部分圏ならば、恒等写像は
$\mathcal{A}$ から $\mathcal{B}$ への関手である。さらに、部分圏
$\mathcal{A} \subset \mathcal{B}$ が充満であることと、包含関手が
充満忠実であることは同値である。また、圏 $\mathcal{B}$ が与えられたとき、
$\mathcal{B}$ の充満部分圏全体の
集合は、$\Ob(\mathcal{B})$ の部分集合全体の集合と同じである。

\begin{remark}
\label{categories-remark-functor-into-sets}
$\mathcal{A}$ を圏とする。集合の圏が「大きい」にもかかわらず、
$\mathcal{A}$ から $\textit{Sets}$ への関手 $F$ は数学的対象である
（すなわち、それは集合であって、類でも集合論の論理式でもない。集合論、
節 \ref{sets-section-sets-everything} を参照されたい）。実際、対象における
$F$ の値域 $F(\Ob(\mathcal{A}))$ は集合である。そこで、$F$ を
$\mathcal{A}$ と、対象が $F(\Ob(\mathcal{A}))$ の元であるような
集合の圏の充満部分圏との間の関手とみなすことができる。
\end{remark}

\begin{example}
\label{categories-example-group-homomorphism-functor}
群の準同型 $p : G\to H$ は、例 \ref{categories-example-group-groupoid} の対応する
グルーポイドの間の関手を与える。この関手が忠実（それぞれ充満忠実）である
ことと、$p$ が単射（それぞれ同型）であることは同値である。
\end{example}

\begin{example}
\label{categories-example-category-over-X}
圏 $\mathcal{C}$ と対象 $X\in \Ob(\mathcal{C})$ が与えられたとき、
{\it $X$ 上の対象の圏} $\mathcal{C}/X$ を次のように定義する。
$\mathcal{C}/X$ の対象は、ある $Y\in \Ob(\mathcal{C})$ に対する射
$Y\to X$ である。対象 $Y\to X$ と $Y'\to X$ の間の射は、明らかな図式を
可換にする $\mathcal{C}$ の射 $Y\to Y'$ である。射を単に忘れる関手
$p_X : \mathcal{C}/X\to \mathcal{C}$ がある。また、$\mathcal{C}$ の射
$f : X'\to X$ が与えられれば、$f$ との合成によって誘導される関手
$F : \mathcal{C}/X' \to \mathcal{C}/X$ があり、
$p_X\circ F = p_{X'}$ が成り立つ。
\end{example}

\begin{example}
\label{categories-example-category-under-X}
圏 $\mathcal{C}$ と対象 $X\in \Ob(\mathcal{C})$ が与えられたとき、
{\it $X$ の下の対象の圏} $X/\mathcal{C}$ を次のように定義する。
$X/\mathcal{C}$ の対象は、ある $Y\in \Ob(\mathcal{C})$ に対する射
$X\to Y$ である。対象 $X\to Y$ と $X\to Y'$ の間の射は、明らかな図式を
可換にする $\mathcal{C}$ の射 $Y\to Y'$ である。射を単に忘れる関手
$p_X : X/\mathcal{C}\to \mathcal{C}$ がある。また、$\mathcal{C}$ の射
$f : X'\to X$ が与えられれば、$f$ との合成によって誘導される関手
$F : X/\mathcal{C} \to X'/\mathcal{C}$ があり、
$p_{X'}\circ F = p_X$ が成り立つ。
\end{example}

\begin{definition}
\label{categories-definition-transformation-functors}
$F, G : \mathcal{A} \to \mathcal{B}$ を関手とする。
{\it 自然変換}、または{\it 関手の射} $t : F \to G$ とは、
次を満たす族 $\{t_x\}_{x\in \Ob(\mathcal{A})}$ のことである：
\begin{enumerate}
\item $t_x : F(x) \to G(x)$ は圏 $\mathcal{B}$ の射であり、
\item $\mathcal{A}$ の各射 $\phi : x \to y$ に対して、次の図式が可換である：
$$
\xymatrix{
F(x) \ar[r]^{t_x} \ar[d]_{F(\phi)} & G(x) \ar[d]^{G(\phi)} \\
F(y) \ar[r]^{t_y} & G(y) }
$$
\end{enumerate}
\end{definition}

\noindent
$t$ が $F$ から $G$ への射であることを示すために、図式
$$
\xymatrix{
\mathcal{A}
\rtwocell^F_G{t}
&
\mathcal{B}
}
$$
を用いることもある。

\medskip\noindent
各関手 $F$ には{\it 恒等}変換 $\text{id}_F : F \to F$ が付随する。
さらに、関手の射 $t : F \to G$ と関手の射 $s : E \to F$ が
与えられたとき、その{\it 合成} $t \circ s$ を
$$
(t \circ s)_x = t_x \circ s_x : E(x) \to G(x)
$$
（$x \in \Ob(\mathcal{A})$）により定義する。これが実際に $E$ から $G$ への
関手の射であることは容易に確かめられる。このようにして、圏
$\mathcal{A}$ と $\mathcal{B}$ が与えられると新たな圏、すなわち
$\mathcal{A}$ と $\mathcal{B}$ の間の関手からなる圏が得られる。

\begin{remark}
\label{categories-remark-functors-sets-sets}
$\mathcal{A}$ が「大きな」圏の場合には、同じことが成り立たない例がここに
現れる。たとえば関手 $\textit{Sets} \to \textit{Sets}$ を考えよう。
現在の定義では、このような関手は集合ではなく類である。言い換えれば、これは
集合論の論理式（いくつかの変数を指定した集合に等しいものとする）によって
与えられる。数学的対象の定義の一部として集合論の可能な論理式すべてを考えるのは
得策でない。同じ問題は、たとえばスキームの圏上の層を考えるときにも生じる。
この点には後で立ち返る。
\end{remark}

\begin{definition}
\label{categories-definition-equivalence-categories}
{\it 圏同値} $F : \mathcal{A} \to \mathcal{B}$ とは、関手
$G : \mathcal{B} \to \mathcal{A}$ であって、合成 $F \circ G$ および
$G \circ F$ がそれぞれ恒等関手 $\text{id}_\mathcal{B}$ および
$\text{id}_\mathcal{A}$ と同型になるものが存在するような関手のことである。
この場合、$G$ を $F$ の{\it 擬逆関手}という。
\end{definition}

\begin{lemma}
\label{categories-lemma-construct-quasi-inverse}
$F : \mathcal{A} \to \mathcal{B}$ を充満忠実関手とする。各
$X \in \Ob(\mathcal{B})$ に対して、$\mathcal{A}$ の対象 $j(X)$ と同型射
$i_X : X \to F(j(X))$ が与えられていると仮定する。このとき、一意な関手
$j : \mathcal{B} \to \mathcal{A}$ が存在し、$j$ は対象に対するこの対応を
拡張する。また、同型射 $i_X$ は関手の同型
$\text{id}_\mathcal{B} \to F \circ j$ を定める。
さらに、$j$ と $F$ は互いに擬逆な圏同値である。
\end{lemma}

\begin{proof}
$j : \mathcal{B} \to \mathcal{A}$ の構成は二段階からなる。まず、
$X \in \mathcal{B}$ に $j(X)$ を対応させる写像
$j : \Ob(\mathcal{B}) \to \Ob(\mathcal{A})$ を定義する。次に、
$X,Y \in \Ob(\mathcal{B})$ かつ $\phi : X \to Y$ のとき、
$\phi' := i_Y \circ\phi\circ i_X^{-1}$ を考える。$F$ は充満忠実なので、
$F(\varphi) = \phi'$ を満たす一意な $\varphi$ が存在する。
$j(\phi) = \varphi$ と定める。$j$ が関手であることの確認は省略する。
構成により図式
$$
\xymatrix{
X \ar[r]_-{i_X} \ar[d]_{\phi}
&
F(j(X)) \ar[d]^{F\circ j(\phi)}
\\
Y \ar[r]^-{i_Y}
&
F(j(Y))
}
$$
は可換である。したがって、各 $i_X$ が同型射であることから、
$\{i_X\}_X$ は関手の同型 $\text{id}_\mathcal{B} \to F\circ j$ である。
結論を得るには、さらに $j \circ F$ が $\text{id}_\mathcal{A}$ と同型である
ことを示さなければならない。しかし $F$ は充満忠実なので、これを示すには
$F$ を後合成した後で示せば十分である。すなわち、
$F \circ j \circ F$ が $F \circ \text{id}_\mathcal{A}$ と同型であることを
示せば十分である（細部は省略する）。これは
$F \circ j \cong \text{id}_\mathcal{B}$ から明らかである。
\end{proof}




\begin{lemma}
\label{categories-lemma-equivalence-categories}
関手が圏同値であるための必要十分条件は、その関手が充満忠実かつ本質的全射で
あることである。
\end{lemma}

\begin{proof}
$F : \mathcal{A} \to \mathcal{B}$ を本質的全射かつ充満忠実とする。
約束によりすべての圏は小さく、また $F$ は本質的全射なので、選択公理を用いて
各 $X \in \Ob(\mathcal{B})$ に対し、$\mathcal{A}$ の対象 $j(X)$ と同型射
$i_X : X \to F(j(X))$ を選ぶことができる。そこで、$F$ が充満忠実である
ことを用いて補題 \ref{categories-lemma-construct-quasi-inverse} を適用する。
\end{proof}

\begin{definition}
\label{categories-definition-product-category}
$\mathcal{A}$、$\mathcal{B}$ を圏とする。{\it 直積圏}
$\mathcal{A} \times \mathcal{B}$ を、対象が
$\Ob(\mathcal{A} \times \mathcal{B}) =
\Ob(\mathcal{A}) \times \Ob(\mathcal{B})$
であり、射が
$$
\Mor_{\mathcal{A} \times \mathcal{B}}((x, y), (x', y'))
:=
\Mor_\mathcal{A}(x, x')\times
\Mor_\mathcal{B}(y, y').
$$
である圏として定義する。合成は成分ごとに定義する。
\end{definition}

\section{反対圏と米田の補題}
\label{categories-section-opposite}

\begin{definition}
\label{categories-definition-opposite}
圏 $\mathcal{C}$ が与えられたとき、その{\it 反対圏}
$\mathcal{C}^{opp}$ とは、$\mathcal{C}$ と同じ対象をもち、すべての射を
逆向きにした圏のことである。
\end{definition}

\noindent
言い換えれば、
$\Mor_{\mathcal{C}^{opp}}(x, y) = \Mor_\mathcal{C}(y, x)$
である。$\mathcal{C}^{opp}$ における合成は、向きが逆であることを除けば
$\mathcal{C}$ における合成と同じである。すなわち、
$\phi : y \to z$ と $\psi : x \to y$ が $\mathcal{C}^{opp}$ の射、
言い換えれば $\mathcal{C}$ における矢印 $z \to y$ と $y \to x$ ならば、
$\phi \circ^{opp} \psi$ は $\mathcal{C}$ における合成
$z \to y \to x$ に対応する $\mathcal{C}^{opp}$ の射 $x \to z$ である。

\begin{definition}
\label{categories-definition-contravariant}
$\mathcal{C}$、$\mathcal{S}$ を圏とする。$\mathcal{C}$ から
$\mathcal{S}$ への{\it 反変}関手 $F$ とは、関手
$\mathcal{C}^{opp}\to \mathcal{S}$ のことである。
\end{definition}

\noindent
具体的には、反変関手 $F$ は写像 $F : \Ob(\mathcal{C}) \to
\Ob(\mathcal{S})$ と、$\mathcal{C}$ の各射 $\psi : x \to y$ に対する射
$F(\psi) : F(y) \to F(x)$ によって与えられる。これらは、別の射
$\phi : y \to z$ が与えられたとき、$F(z) \to F(x)$ の射として
$F(\phi \circ \psi) = F(\psi) \circ F(\phi)$ を満たさなければならない
（順序が逆になることに注意されたい）。

\begin{definition}
\label{categories-definition-presheaf}
$\mathcal{C}$ を圏とする。
\begin{enumerate}
\item {\it $\mathcal{C}$ 上の集合の前層}、または単に{\it 前層}とは、
$\mathcal{C}$ から $\textit{Sets}$ への反変関手 $F$ のことである。
\item 前層の圏を $\textit{PSh}(\mathcal{C})$ と表す。
\end{enumerate}
\end{definition}

\noindent
もちろん、前層の圏は真の類である。

\begin{example}
\label{categories-example-hom-functor}
点関手。任意の $U\in \Ob(\mathcal{C})$ に対して、反変関手
$$
\begin{matrix}
h_U & : & \mathcal{C}
&
\longrightarrow
&
\textit{Sets} \\
& &
X
&
\longmapsto
&
\Mor_\mathcal{C}(X, U)
\end{matrix}
$$
があり、対象 $X$ を集合 $\Mor_\mathcal{C}(X, U)$ に送る。言い換えれば、
$h_U$ は前層である。射 $f : X\to Y$ が与えられたとき、対応する写像
$h_U(f) : \Mor_\mathcal{C}(Y, U)\to \Mor_\mathcal{C}(X, U)$ は
$\phi$ を $\phi\circ f$ に送る。この前層は常に
$h_U : \mathcal{C}^{opp} \to \textit{Sets}$ と表す。これは $U$ に付随する
{\it 表現可能前層}と呼ばれる。$\mathcal{C}$ がスキームの圏である場合、
この関手を $U$ の\emph{点関手}と呼ぶこともある。
\end{example}

\noindent
$\mathcal{C}$ の射 $\phi : U \to V$ が与えられると、射 $U \to V$ との
合成によって定義される関手の自然変換 $h(\phi) : h_U \to h_V$ が得られる。
これは $\mathcal{C}$ における射の合成を関手の変換の合成に移す。
言い換えれば、関手
$$
h :
\mathcal{C}
\longrightarrow
\textit{PSh}(\mathcal{C})
$$
が得られる。終域は「大きな」圏であることに注意されたい。注意
\ref{categories-remark-big-categories} を参照されたい。一方、$h$ は実際の数学的対象
（すなわち集合）である。注意 \ref{categories-remark-functor-into-sets} と比較されたい。

\begin{lemma}[米田の補題]
\label{categories-lemma-yoneda}
\begin{reference}
ある形では \cite{Yoneda-homology} に現れ、Grothendieck は
\cite{Gr-II} において一般化された形で用いた。
\end{reference}
$U, V \in \Ob(\mathcal{C})$ とする。任意の関手の射
$s : h_U \to h_V$ に対して、一意な射 $\phi : U \to V$ で
$h(\phi) = s$ を満たすものが存在する。言い換えれば、関手 $h$ は充満忠実で
ある。より一般に、任意の反変関手 $F$ と $\mathcal{C}$ の任意の対象 $U$ に
対して、自然な全単射
$$
\Mor_{\textit{PSh}(\mathcal{C})}(h_U, F) \longrightarrow F(U),
\quad
s \longmapsto s_U(\text{id}_U).
$$
がある。
\end{lemma}

\begin{proof}
第一の主張については、
$\phi = s_U(\text{id}_U) \in \Mor_\mathcal{C}(U, V)$
とおけばよい。第二の主張については、$\xi \in F(U)$ が与えられたとき、
$h_U(V) = \Mor_\mathcal{C}(V, U)$ の元 $f : V \to U$ を
$F(f)(\xi)$ に送る写像 $s_V : h_U(V) \to F(V)$ によって $s$ を定義する。
\end{proof}

\begin{definition}
\label{categories-definition-representable-functor}
反変関手 $F : \mathcal{C}\to \textit{Sets}$ が{\it 表現可能}であるとは、
それが $\mathcal{C}$ のある対象 $U$ の点関手 $h_U$ と同型であることをいう。
\end{definition}

\noindent
$\mathcal{C}$ を圏とし、
$F : \mathcal{C}^{opp} \to \textit{Sets}$ を表現可能関手とする。
$\mathcal{C}$ の対象 $U$ と同型 $s : h_U \to F$ を選ぶ。米田の補題により、
対 $(U, s)$ は一意な同型を除いて一意である。対象 $U$ を $F$ を
{\it 表現する}対象という。米田の補題により、変換 $s$ は一意な元
$\xi \in F(U)$ に対応する。この元を{\it 普遍対象}という。これは、各
$V \in \Ob(\mathcal{C})$ に対して写像
$$
\Mor_\mathcal{C}(V, U) \longrightarrow F(V),\quad
(f : V \to U) \longmapsto F(f)(\xi)
$$
が全単射になるという性質をもつ。したがって $\xi$ は、$F(V)$ の各元が、
一意な $\mathcal{C}$ の射 $f : V \to U$ に対する $F(f)$ による $\xi$ の
像に等しいという意味で普遍的である。







\section{二対象の積}
\label{categories-section-products-pairs}

\begin{definition}
\label{categories-definition-products}
$x, y\in \Ob(\mathcal{C})$ とする。$x$ と $y$ の{\it 積}とは、対象
$x \times y \in \Ob(\mathcal{C})$ と射
$p\in \Mor_{\mathcal C}(x \times y, x)$ および
$q\in\Mor_{\mathcal C}(x \times y, y)$ であって、次の普遍性を満たすものの
ことである。任意の $w\in \Ob(\mathcal{C})$ と射
$\alpha \in \Mor_{\mathcal C}(w, x)$ および
$\beta \in \Mor_\mathcal{C}(w, y)$ に対して、図式
$$
\xymatrix{
w \ar[rrrd]^\beta \ar@{-->}[rrd]_\gamma \ar[rrdd]_\alpha & & \\
& & x \times y \ar[d]_p \ar[r]_q & y \\
& & x &
}
$$
を可換にする一意な $\gamma\in \Mor_{\mathcal C}(w, x \times y)$ が存在する。
\end{definition}

\noindent
積が存在すれば、それは一意な同型を除いて一意である。実際、定義は
$x \times y$ が $\mathcal{C}$ の対象であって、$w$ に関して関手的に
$$
h_{x \times y}(w) = h_x(w) \times h_y(w)
$$
を満たすことを要求しているので、これは米田の補題から従う。言い換えれば、
積 $x \times y$ は関手 $w \mapsto h_x(w) \times h_y(w)$ を表現する対象である。

\begin{definition}
\label{categories-definition-has-products-of-pairs}
任意の $x, y \in \Ob(\mathcal{C})$ に対して積 $x \times y$ が存在するとき、
圏 $\mathcal{C}$ は{\it 二対象の積をもつ}という。
\end{definition}

\noindent
この用語は、通常は別の意味をもつ「積をもつ」または「有限積をもつ」という
概念と区別するために用いる（特に後二者は常に終対象の存在を含意する）。





\section{二対象の余積}
\label{categories-section-coproducts-pairs}

\begin{definition}
\label{categories-definition-coproducts}
$x, y \in \Ob(\mathcal{C})$ とする。$x$ と $y$ の{\it 余積}、または
{\it 融合和}とは、対象 $x \amalg y \in \Ob(\mathcal{C})$ と射
$i \in \Mor_{\mathcal C}(x, x \amalg y)$ および
$j \in \Mor_{\mathcal C}(y, x \amalg y)$ であって、次の普遍性を満たすものの
ことである。任意の $w \in \Ob(\mathcal{C})$ と射
$\alpha \in \Mor_{\mathcal C}(x, w)$ および
$\beta \in \Mor_\mathcal{C}(y, w)$ に対して、図式
$$
\xymatrix{
& y \ar[d]^j \ar[rrdd]^\beta \\
x \ar[r]^i \ar[rrrd]_\alpha & x \amalg y \ar@{-->}[rrd]^\gamma \\
& & & w
}
$$
を可換にする一意な $\gamma \in \Mor_{\mathcal C}(x \amalg y, w)$ が存在する。
\end{definition}

\noindent
余積が存在すれば、それは一意な同型を除いて一意である。実際、定義は
$x \amalg y$ が $\mathcal{C}$ の対象であって、$w$ に関して関手的に
$$
\Mor_\mathcal{C}(x \amalg y, w) =
\Mor_\mathcal{C}(x, w) \times \Mor_\mathcal{C}(y, w)
$$
を満たすことを要求しているので、これは反対圏に適用した米田の補題から従う。

\begin{definition}
\label{categories-definition-has-coproducts-of-pairs}
任意の $x, y \in \Ob(\mathcal{C})$ に対して余積 $x \amalg y$ が存在するとき、
圏 $\mathcal{C}$ は{\it 二対象の余積をもつ}という。
\end{definition}

\noindent
この用語は、通常は別の意味をもつ「余積をもつ」または「有限余積をもつ」という
概念と区別するために用いる（特に後二者は常に $\mathcal{C}$ における始対象の
存在を含意する）。

\section{ファイバー積}
\label{categories-section-fibre-products}

\begin{definition}
\label{categories-definition-fibre-products}
$x, y, z\in \Ob(\mathcal{C})$、
$f\in \Mor_\mathcal{C}(x, y)$ および
$g\in \Mor_{\mathcal C}(z, y)$ とする。$f$ と $g$ の
{\it ファイバー積}とは、対象 $x \times_y z\in \Ob(\mathcal{C})$ と射
$p \in \Mor_{\mathcal C}(x \times_y z, x)$ および
$q \in \Mor_{\mathcal C}(x \times_y z, z)$ であって、図式
$$
\xymatrix{
x \times_y z \ar[r]_q \ar[d]_p & z \ar[d]^g \\
x \ar[r]^f & y
}
$$
を可換にし、さらに次の普遍性を満たすもののことである。任意の
$w\in \Ob(\mathcal{C})$ と射 $\alpha \in \Mor_{\mathcal C}(w, x)$ および
$\beta \in \Mor_\mathcal{C}(w, z)$ で $f \circ \alpha = g \circ \beta$ を
満たすものに対して、図式
$$
\xymatrix{
w \ar[rrrd]^\beta \ar@{-->}[rrd]_\gamma \ar[rrdd]_\alpha & & \\
& & x \times_y z \ar[d]^p \ar[r]_q & z \ar[d]^g \\
& & x \ar[r]^f & y
}
$$
を可換にする一意な $\gamma \in \Mor_{\mathcal C}(w, x \times_y z)$ が存在する。
\end{definition}

\noindent
ファイバー積が存在すれば、それは一意な同型を除いて一意である。実際、定義は
$x \times_y z$ が $\mathcal{C}$ の対象であって、$w$ に関して関手的に
$$
h_{x \times_y z}(w) = h_x(w) \times_{h_y(w)} h_z(w)
$$
を満たすことを要求しているので、これは米田の補題から従う。言い換えれば、
ファイバー積 $x \times_y z$ は関手
$w \mapsto h_x(w) \times_{h_y(w)} h_z(w)$ を表現する対象である。

\begin{definition}
\label{categories-definition-cartesian}
圏における可換図式
$$
\xymatrix{
w \ar[r] \ar[d] &
z \ar[d] \\
x \ar[r] &
y
}
$$
が{\it カルテジアン}であるとは、$w$ と射 $w \to x$ および $w \to z$ が、
射 $x \to y$ と $z \to y$ のファイバー積をなすことをいう。
\end{definition}

\begin{definition}
\label{categories-definition-has-fibre-products}
任意の $f\in \Mor_{\mathcal C}(x, y)$ と
$g\in \Mor_{\mathcal C}(z, y)$ に対してファイバー積が存在するとき、圏
$\mathcal{C}$ は{\it ファイバー積をもつ}という。
\end{definition}

\begin{definition}
\label{categories-definition-representable-morphism}
圏 $\mathcal{C}$ の射 $f : x \to y$ が{\it 表現可能}であるとは、
$\mathcal{C}$ の各射 $z \to y$ に対してファイバー積
$x \times_y z$ が存在することをいう。
\end{definition}

\begin{lemma}
\label{categories-lemma-composition-representable}
$\mathcal{C}$ を圏とする。$f : x \to y$ と $g : y \to z$ が
表現可能であるとする。このとき $g \circ f : x \to z$ は表現可能である。
\end{lemma}

\begin{proof}
$t \in \Ob(\mathcal C)$ および
$ \varphi \in \Mor_{\mathcal C}(t,z) $ とする。$g$ と $f$ は表現可能なので、
定義 \ref{categories-definition-fibre-products} の普遍性をもつ可換図式
$$
\xymatrix{
y \times_z t \ar[r]_q \ar[d]_p &
t \ar[d]^{\varphi} \\
y \ar[r]^{g} & z
}
\quad\quad
\xymatrix{
x \times_y (y\! \times_z\! t) \ar[r]_{q'} \ar[d]_{p'} &
y \times_z t \ar[d]^p \\
x \ar[r]^f & y
}
$$
が得られる。射 $q \circ q' : x \times_z t \to t$ および
$p' : x \times_z t \to x$ とともに
$x \times_z t = x \times_y (y \times_z t)$ がファイバー積であると主張する。
まず、上の図式の可換性から
$\varphi \circ q \circ q' = g \circ f \circ p'$ が従う。
普遍性を確かめるため、$w \in \Ob(\mathcal C)$ とし、
$\alpha : w \to x$ および $\beta : w \to t$ が
$\varphi \circ \beta = g \circ f \circ \alpha$ を満たす射であるとする。
ファイバー積の定義により、図式
$$
\xymatrix{
w \ar[rrrd]^\beta \ar@{-->}[rrd]_\delta \ar[rrdd]_{f\circ\alpha} & & \\
& & y \times_z t \ar[d]_p \ar[r]_q & t \ar[d]^{\varphi} \\
& & y \ar[r]^{g} & z
}
$$
および
$$
\xymatrix{
w \ar[rrrd]^\delta \ar@{-->}[rrd]_\gamma \ar[rrdd]_{\alpha} & & \\
& & x \times_y (y\!\times_z\! t) \ar[d]_{p'} \ar[r]_{q'} &
y \times_z t \ar[d]^{p} \\
& & x \ar[r]^{f} & y
}
$$
を可換にする一意な射 $\delta$ および $\gamma$ が存在する。このとき
$\gamma$ は図式
$$
\xymatrix{
w \ar[rrrd]^\beta \ar@{-->}[rrd]_\gamma \ar[rrdd]_{\alpha} & & \\
& & x \times_z t \ar[d]_{p'} \ar[r]_{q\circ q'} & t \ar[d]^{\varphi} \\
& & x \ar[r]^{g\circ f} & z
}
$$
を可換にする。その一意性を示すため、
$ q\circ q'\circ\gamma' = \beta $ および $ p'\circ \gamma' = \alpha $ を
満たす $\gamma'$ を考える。$\gamma$ の一意性により、結論を得るには
$ q'\circ\gamma' = \delta $ および $ p'\circ\gamma' = \alpha $ を示せば
よい。第二の等式は仮定した。第一の等式には、先に構成した射の一意性も用いる。
$\delta$ は $ q\circ\delta = \beta $ および
$ p\circ\delta = f\circ\alpha $ を満たす唯一の射である。既に
$ q\circ (q'\circ\gamma') = \beta $ を仮定している。さらに、
ファイバー積の定義から $ f\circ p' = p\circ q' $ である。したがって
$$
p\circ (q'\circ\gamma') =
(p\circ q')\circ\gamma' =
(f\circ p')\circ\gamma' =
f\circ (p'\circ\gamma') = f\circ\alpha.
$$
よって $ q'\circ\gamma' = \delta $ であり、証明が完了する。
\end{proof}

\begin{lemma}
\label{categories-lemma-base-change-representable}
$\mathcal{C}$ を圏とする。$f : x \to y$ が表現可能であるとし、
$y' \to y$ を $\mathcal{C}$ の射とする。このとき射
$x' := x \times_y y' \to y'$ も表現可能である。
\end{lemma}

\begin{proof}
$z \to y'$ を射とする。ファイバー積 $x' \times_{y'} z$ は関手
\begin{eqnarray*}
w & \mapsto & h_{x'}(w)\times_{h_{y'}(w)} h_z(w) \\
& = & (h_x(w) \times_{h_y(w)} h_{y'}(w)) \times_{h_{y'}(w)} h_z(w) \\
& = & h_x(w) \times_{h_y(w)} h_z(w)
\end{eqnarray*}
を表現するはずであり、この関手は仮定により表現可能である。
\end{proof}

\section{ファイバー積の例}
\label{categories-section-example-fibre-products}

\noindent
この節ではファイバー積の例を列挙し、それらを記述する。

\medskip\noindent
最初のまったく自明な例として、集合の圏はファイバー積をもち、したがって
すべての射が表現可能であることに注意する。実際、$f : X \to Y$ と
$g : Z \to Y$ が集合の写像ならば、$X \times_Y Z$ を、$f(x) = g(z)$ を
満たす対 $(x, z)$ からなる $X \times Z$ の部分集合として定義する。
射 $p : X \times_Y Z \to X$ および $q : X \times_Y Z \to Z$ は、それぞれ
射影 $(x, z) \mapsto x$ および $(x, z) \mapsto z$ である。最後に、
$\alpha : W \to X$ と $\beta : W \to Z$ が
$f \circ \alpha = g \circ \beta$ を満たす射ならば、写像
$W \to X \times Z$、$w\mapsto (\alpha(w), \beta(w))$ の像は明らかに
望みどおり $X \times_Y Z$ に含まれる。

\medskip\noindent
対象がある種の代数構造を備えた集合である多くの圏では、台集合のファイバー積が
その圏におけるファイバー積にもなる。たとえば、上の $X$、$Y$、$Z$ が群で、
$f$、$g$ が群準同型であると仮定する。このとき集合論的ファイバー積
$X \times_Y Z$ は、二つの対の積を
$(x, z) \cdot (x', z') = (xx', zz')$ と定めるだけで群の構造を受け継ぐ。
同じ議論が成り立つ圏を列挙する。
\begin{enumerate}
\item 群の圏 $\textit{Groups}$。
\item 固定した群 $G$ に対し、左 $G$-作用を備えた集合の圏
$G\textit{-Sets}$。
\item 環の圏。
\item 環 $R$ が与えられたときの $R$-加群の圏。
\end{enumerate}




\section{ファイバー積と表現可能性}
\label{categories-section-representable-map-presheaves}

\noindent
この節では、ある圏から集合の圏への反変関手の圏におけるファイバー積を
具体的に調べる。この議論は、後にサイト、前層および層を論じる際に置き換えられる。
このような関手の間の「表現可能射」という概念は、それ自体にも関心がある。

\begin{lemma}
\label{categories-lemma-fibre-product-presheaves}
$\mathcal{C}$ を圏とする。$F, G, H : \mathcal{C}^{opp} \to \textit{Sets}$ を
関手とし、$a : F \to G$ および $b : H \to G$ を関手の変換とする。
このとき、圏 $\textit{PSh}(\mathcal{C})$ におけるファイバー積
$F \times_{a, G, b} H$ は存在し、$\mathcal{C}$ の任意の対象 $X$ に対して
$$
(F \times_{a, G, b} H)(X) =
F(X) \times_{a_X, G(X), b_X} H(X)
$$
により与えられる。
\end{lemma}

\begin{proof}
省略する。
\end{proof}

\noindent
特別な場合として、射 $a : F \to G$、対象 $U \in \Ob(\mathcal{C})$ および
元 $\xi \in G(U)$ が与えられていると仮定する。米田の補題
\ref{categories-lemma-yoneda} により、これは変換 $\xi : h_U \to G$ を与える。
この場合のファイバー積は
$$
(h_U \times_{\xi, G, a} F)(X) =
\{ (f, \xi') \mid f : X \to U, \ \xi' \in F(X), \ G(f)(\xi) = a_X(\xi')\}
$$
によって記述される。$F$ と $G$ も表現可能ならば、これはファイバー積が
存在するときにそれを表現する関手である。節 \ref{categories-section-fibre-products} を
参照されたい。定義 \ref{categories-definition-representable-morphism} との類似から、
表現可能変換の概念を次のように定義する。

\begin{definition}
\label{categories-definition-representable-map-presheaves}
$\mathcal{C}$ を圏とし、$F, G : \mathcal{C}^{opp} \to \textit{Sets}$ を
関手とする。射 $a : F \to G$ が{\it 表現可能}である、または
{\it $F$ は $G$ 上相対的に表現可能}であるとは、任意の
$U \in \Ob(\mathcal{C})$ と任意の $\xi \in G(U)$ に対して関手
$h_U \times_{\xi, G, a} F$ が表現可能であることをいう。
\end{definition}

\begin{lemma}
\label{categories-lemma-representable-over-representable}
$\mathcal{C}$ を圏とし、$a : F \to G$ を $\mathcal{C}$ から
$\textit{Sets}$ への反変関手の射とする。$a$ が表現可能であり、かつ $G$ が
表現可能関手ならば、$F$ は表現可能である。
\end{lemma}

\begin{proof}
省略する。
\end{proof}

\begin{lemma}
\label{categories-lemma-representable-diagonal}
$\mathcal{C}$ を圏とし、$F : \mathcal{C}^{opp} \to \textit{Sets}$ を関手と
する。$\mathcal{C}$ は二対象の積およびファイバー積をもつと仮定する。
次の条件は同値である：
\begin{enumerate}
\item 対角射 $\Delta : F \to F \times F$ は表現可能である。
\item $\mathcal{C}$ の任意の $U$ と任意の $\xi \in F(U)$ に対して、
写像 $\xi : h_U \to F$ は表現可能である。
\item $\mathcal{C}$ の任意の対 $U, V$ と任意の
$\xi \in F(U)$、$\xi' \in F(V)$ に対して、ファイバー積
$h_U \times_{\xi, F, \xi'} h_V$ は表現可能である。
\end{enumerate}
\end{lemma}

\begin{proof}
引き続き米田の補題を用いて、$F(U)$ を関手の変換 $h_U \to F$ と同一視する。

\medskip\noindent
(2) と (3) の同値性。$U, \xi, V, \xi'$ を (3) のとおりとする。(2) と (3) は
どちらも、ちょうど $h_U \times_{\xi, F, \xi'} h_V$ が表現可能であることを
述べている。相違は、主張 (3) が $U$ と $V$ に関して対称であるのに対し、
(2) はそうでないことだけである。

\medskip\noindent
条件 (1) を仮定する。$U, \xi, V, \xi'$ を (3) のとおりとする。
$h_U \times h_V = h_{U \times V}$ は表現可能であることに注意する。
積 $\xi \times \xi' : h_U \times h_V \to F \times F$ に対応する写像を
$\eta : h_{U \times V} \to F \times F$ と表す。仮定により、ファイバー積
$F \times_{\Delta, F \times F, \eta} h_{U \times V}$ は表現可能である。
これは、$W \in \Ob(\mathcal{C})$ と射
$W \to U$、$W \to V$ および $h_W \to F$ が存在し、図式
$$
\xymatrix{
h_W \ar[d] \ar[r] & h_U \times h_V \ar[d]^{\xi \times \xi'} \\
F \ar[r] & F \times F
}
$$
がカルテジアンになることを意味する。補題
\ref{categories-lemma-fibre-product-presheaves} におけるファイバー積の明示的記述を
用いれば、これが望みどおり
$h_W = h_U \times_{\xi, F, \xi'} h_V$ を含意することが分かる。

\medskip\noindent
同値な条件 (2) および (3) を仮定する。$U$ を $\mathcal{C}$ の対象とし、
$(\xi, \xi') \in (F \times F)(U)$ とする。(3) によりファイバー積
$h_U \times_{\xi, F, \xi'} h_U$ は表現可能である。対象 $W$ と同型
$h_W \to h_U \times_{\xi, F, \xi'} h_U$ を選ぶ。二つの射影
$\text{pr}_i : h_U \times_{\xi, F, \xi'} h_U \to h_U$ は、米田により射
$p_i : W \to U$ に対応する。
$W' = W \times_{(p_1, p_2), U \times U} U$ とおく。
$$
h_{W'} = h_W \times_{h_U \times h_U} h_U
= (h_U \times_{\xi, F, \xi'} h_U) \times_{h_U \times h_U} h_U
= F \times_{F \times F} h_U.
$$
であるから、$W'$ が $F \times_{\Delta, F \times F} h_U$ を表現することは
形式的に示される。したがって $\Delta$ は表現可能であり、証明が完了する。
\end{proof}

\section{押し出し}
\label{categories-section-pushouts}

\noindent
ファイバー積の双対概念は押し出しである。

\begin{definition}
\label{categories-definition-pushouts}
$x, y, z\in \Ob(\mathcal{C})$、
$f\in \Mor_\mathcal{C}(y, x)$ および
$g\in \Mor_{\mathcal C}(y, z)$ とする。$f$ と $g$ の{\it 押し出し}とは、
対象 $x\amalg_y z\in \Ob(\mathcal{C})$ と射
$p\in \Mor_{\mathcal C}(x, x\amalg_y z)$ および
$q\in\Mor_{\mathcal C}(z, x\amalg_y z)$ であって、図式
$$
\xymatrix{
y \ar[r]_g \ar[d]_f & z \ar[d]^q \\
x \ar[r]^p & x\amalg_y z
}
$$
を可換にし、さらに次の普遍性を満たすもののことである。任意の
$w\in \Ob(\mathcal{C})$ と射 $\alpha \in \Mor_{\mathcal C}(x, w)$ および
$\beta \in \Mor_\mathcal{C}(z, w)$ で
$\alpha \circ f = \beta \circ g$ を満たすものに対して、図式
$$
\xymatrix{
y \ar[r]_g \ar[d]_f & z \ar[d]^q \ar[rrdd]^\beta & & \\
x \ar[r]^p \ar[rrrd]^\alpha & x \amalg_y z \ar@{-->}[rrd]^\gamma & & \\
& & & w
}
$$
を可換にする一意な $\gamma\in \Mor_{\mathcal C}(x\amalg_y z, w)$ が存在する。
\end{definition}

\noindent
三つ組 $(x\amalg_y z, p, q)$ が存在するならば、一意な同型を除いて一意である
ことを直接の議論で証明するのは容易である。別の方法として、押し出しを反対圏に
おけるファイバー積とみなせば、節 \ref{categories-section-fibre-products} の議論から
この一意性が直ちに得られる。

\begin{definition}
\label{categories-definition-cocartesian}
圏における可換図式
$$
\xymatrix{
y \ar[r] \ar[d] & z \ar[d] \\
x \ar[r] & w
}
$$
が{\it 余カルテジアン}であるとは、$w$ と射 $x \to w$ および $z \to w$ が、
射 $y \to x$ と $y \to z$ の押し出しをなすことをいう。
\end{definition}


\section{等化子}
\label{categories-section-equalizers}

\begin{definition}
\label{categories-definition-equalizers}
$X$、$Y$ を圏 $\mathcal{C}$ の対象とし、$a, b : X \to Y$ を射とする。
射 $e : Z \to X$ が対 $(a, b)$ の{\it 等化子}であるとは、
$a \circ e = b \circ e$ が成り立ち、かつ $(Z, e)$ が次の普遍性を満たすことを
いう。$\mathcal{C}$ の任意の射 $t : W \to X$ で
$a \circ t = b \circ t$ を満たすものに対して、$t = e \circ s$ を満たす
一意な射 $s : W \to Z$ が存在する。
\end{definition}

\noindent
上のファイバー積の場合と同様に、等化子が存在すれば、それは一意な同型を除いて
一意である。この定義は、射が $2$ 本より多い場合にも直ちに一般化できる。

\section{余等化子}
\label{categories-section-coequalizers}

\begin{definition}
\label{categories-definition-coequalizers}
$X$、$Y$ を圏 $\mathcal{C}$ の対象とし、$a, b : X \to Y$ を射とする。
射 $c : Y \to Z$ が対 $(a, b)$ の{\it 余等化子}であるとは、
$c \circ a = c \circ b$ が成り立ち、かつ $(Z, c)$ が次の普遍性を満たすことを
いう。$\mathcal{C}$ の任意の射 $t : Y \to W$ で
$t \circ a = t \circ b$ を満たすものに対して、$t = s \circ c$ を満たす
一意な射 $s : Z \to W$ が存在する。
\end{definition}

\noindent
上の押し出しの場合と同様に、余等化子が存在すれば、それは一意な同型を除いて
一意である。このことは反対圏を考えれば等化子の一意性から従う。この定義は、
射が $2$ 本より多い場合にも直ちに一般化できる。

\section{始対象と終対象}
\label{categories-section-initial-final}

\begin{definition}
\label{categories-definition-initial-final}
$\mathcal{C}$ を圏とする。
\begin{enumerate}
\item 圏 $\mathcal{C}$ の対象 $x$ が{\it 始対象}であるとは、$\mathcal{C}$ の
各対象 $y$ に対して射 $x \to y$ がちょうど一つ存在することをいう。
\item 圏 $\mathcal{C}$ の対象 $x$ が{\it 終対象}であるとは、$\mathcal{C}$ の
各対象 $y$ に対して射 $y \to x$ がちょうど一つ存在することをいう。
\end{enumerate}
\end{definition}

\noindent
集合の圏では、空集合 $\emptyset$ が始対象であり、実際これが唯一の始対象である。
また、任意の{\it 一元集合}、すなわち元を一つもつ集合は終対象である
（したがって終対象は一意ではない）。




\section{モノ射とエピ射}
\label{categories-section-mono-epi}

\begin{definition}
\label{categories-definition-mono-epi}
$\mathcal{C}$ を圏とし、$f : X \to Y$ を $\mathcal{C}$ の射とする。
\begin{enumerate}
\item 任意の対象 $W$ と射の任意の対 $a, b : W \to X$ に対し、
$f \circ a = f \circ b$ ならば $a = b$ が成り立つとき、$f$ は
{\it モノ射}であるという。
\item 任意の対象 $W$ と射の任意の対 $a, b : Y \to W$ に対し、
$a \circ f = b \circ f$ ならば $a = b$ が成り立つとき、$f$ は
{\it エピ射}であるという。
\end{enumerate}
\end{definition}

\begin{example}
\label{categories-example-mono-epi-sets}
集合の圏では、モノ射は単射に対応し、エピ射は全射に対応する。
\end{example}

\begin{lemma}
\label{categories-lemma-characterize-mono-epi}
$\mathcal{C}$ を圏とし、$f : X \to Y$ を $\mathcal{C}$ の射とする。このとき
\begin{enumerate}
\item $f$ がモノ射であるための必要十分条件は、$X$ がファイバー積
$X \times_Y X$ であることであり、
\item $f$ がエピ射であるための必要十分条件は、$Y$ が押し出し
$Y \amalg_X Y$ であることである。
\end{enumerate}
\end{lemma}

\begin{proof}
$f$ がモノ射であると仮定する。$W$ を $\mathcal C$ の対象とし、
$\alpha, \beta \in \Mor_\mathcal C(W,X)$ が
$f\circ \alpha = f\circ \beta$ を満たすとする。$f$ はモノ射なので
$\alpha = \beta$ である。さらに可換図式
$$
\xymatrix{
X \ar[r]^{\text{id}_X} \ar[d]_{\text{id}_X} & X \ar[d]^{f} \\
X \ar[r]^{f} & Y
}
$$
は $\gamma := \alpha = \beta$ として普遍性を満たす。したがって $X$ は
実際にファイバー積 $X\times_Y X$ である。

\medskip\noindent
$X \times_Y X \cong X $ と仮定する。図式
$$
\xymatrix{
X \ar[r]^{\text{id}_X} \ar[d]_{\text{id}_X} & X \ar[d]^{f} \\
X \ar[r]^{f} & Y }
$$
は可換である。$W \in \Ob(\mathcal C)$ と
$\alpha, \beta : W \to X$ が $f \circ \alpha = f \circ \beta$ を満たせば、
$$
\gamma = \text{id}_X\circ\gamma = \alpha = \beta
$$
を満たす一意な $\gamma$ が存在するので、$\alpha = \beta$ が従う。

\medskip\noindent
第二の主張の証明もまったく同じであり、押し出し
$Y\amalg_X Y = Y$ を用いればよい。
\end{proof}

\section{極限と余極限}
\label{categories-section-limits}

\noindent
$\mathcal{C}$ を圏とする。$\mathcal{C}$ における \textit{図式} とは、単に関手
$M : \mathcal{I} \to \mathcal{C}$ のことである。$\mathcal{I}$ を
\textit{添字圏} といい、また $M$ を $\mathcal{I}$-図式ともいう。
$\mathcal{I}$ の対象 $i$ の像を $M_i$ と書く。したがって、
$\mathcal{I}$ の射 $\phi : i \to i'$ に対して
$M(\phi) : M_i \to M_{i'}$ である。

\begin{definition}
\label{categories-definition-limit}
圏 $\mathcal{C}$ における $\mathcal{I}$-図式 $M$ の \textit{極限} とは、
$\mathcal{C}$ の対象 $\lim_\mathcal{I} M$ と射
$p_i : \lim_\mathcal{I} M \to M_i$ であって、次を満たすものをいう。
\begin{enumerate}
\item $\mathcal{I}$ の射 $\phi : i \to i'$ に対して
$p_{i'} = M(\phi) \circ p_i$ である。
\item $\mathcal{C}$ の任意の対象 $W$ と、$i \in \Ob(\mathcal{I})$ で
添字付けられた射の族 $q_i : W \to M_i$ が与えられ、
$\mathcal{I}$ のすべての $\phi : i \to i'$ に対して
$q_{i'} = M(\phi) \circ q_i$ が成り立つとする。このとき、
各 $\mathcal{I}$ の対象 $i$ について $q_i = p_i \circ q$ を満たす射
$q : W \to \lim_\mathcal{I} M$ が一意に存在する。
\end{enumerate}
\end{definition}

\noindent
極限 $(\lim_\mathcal{I} M, (p_i)_{i\in \Ob(\mathcal{I})})$ は、存在すれば、
定義の一意性の要請により一意な同型を除いて一意である。
二対象の積、ファイバー積、および等化子は極限の例である。
空図式の極限は $\mathcal{C}$ の終対象である。集合の圏ではすべての極限が存在する。
双対概念が余極限である。

\begin{definition}
\label{categories-definition-colimit}
圏 $\mathcal{C}$ における $\mathcal{I}$-図式 $M$ の \textit{余極限} とは、
$\mathcal{C}$ の対象 $\colim_\mathcal{I} M$ と射
$s_i : M_i \to \colim_\mathcal{I} M$ であって、次を満たすものをいう。
\begin{enumerate}
\item $\mathcal{I}$ の射 $\phi : i \to i'$ に対して
$s_i = s_{i'} \circ M(\phi)$ である。
\item $\mathcal{C}$ の任意の対象 $W$ と、$i \in \Ob(\mathcal{I})$ で
添字付けられた射の族 $t_i : M_i \to W$ が与えられ、
$\mathcal{I}$ のすべての $\phi : i \to i'$ に対して
$t_i = t_{i'} \circ M(\phi)$ が成り立つとする。このとき、
各 $\mathcal{I}$ の対象 $i$ について $t_i = t \circ s_i$ を満たす射
$t : \colim_\mathcal{I} M \to W$ が一意に存在する。
\end{enumerate}
\end{definition}

\noindent
余極限 $(\colim_\mathcal{I} M, (s_i)_{i\in \Ob(\mathcal{I})})$ は、存在すれば、
定義の一意性の要請により一意な同型を除いて一意である。
二対象の余積、押し出し、および余等化子は余極限の例である。
空図式の余極限は $\mathcal{C}$ の始対象である。
集合の圏ではすべての余極限が存在する。

\begin{remark}
\label{categories-remark-diagram-small}
(余)極限の添字圏が対象の真の類をもつことは決して許さない。
本プロジェクトでは、これはその添字圏が
Remark \ref{categories-remark-big-categories} に列挙された圏のいずれでもないことを意味する。
\end{remark}

\begin{remark}
\label{categories-remark-limit-colim}
$\mathcal{I}$ を添字とする形の代わりに、しばしば
$\lim_i M_i$, $\colim_i M_i$,
$\lim_{i\in \mathcal{I}} M_i$, または $\colim_{i\in \mathcal{I}} M_i$ と書く。
この記法と、後の Section \ref{categories-section-limit-sets} における
集合の極限および余極限の記述を用いれば、次のように述べられる。
$M : \mathcal{I} \to \mathcal{C}$ を図式とする。
\begin{enumerate}
\item 対象 $\lim_i M_i$ は、存在すれば、次の性質を満たす。
$$
\Mor_\mathcal{C}(W, \lim_i M_i)
=
\lim_i \Mor_\mathcal{C}(W, M_i)
$$
ここで右辺の極限は集合の圏でとる。
\item 対象 $\colim_i M_i$ は、存在すれば、次の性質を満たす。
$$
\Mor_\mathcal{C}(\colim_i M_i, W)
=
\lim_{i \in \mathcal{I}^{opp}} \Mor_\mathcal{C}(M_i, W)
$$
ここで右辺は、集合の圏に値をもつ反対圏上の極限である。
\end{enumerate}
米田の補題（およびその双対）により、この公式は極限と余極限をそれぞれ完全に決定する。
\end{remark}

\begin{remark}
\label{categories-remark-cones-and-cocones}
$M : \mathcal{I} \to \mathcal{C}$ を図式とする。この状況で $M$ の
\textit{錐} とは、対象 $W$ と射の族
$q_i : W \to M_i$, $i \in \Ob(\mathcal{I})$ であって、
$\mathcal{I}$ のすべての射 $\phi : i \to i'$ に対して図式
$$
\xymatrix{
& W \ar[dl]_{q_i} \ar[dr]^{q_{i'}} \\
M_i \ar[rr]^{M(\phi)} & & M_{i'}
}
$$
が可換となるものをいう。錐全体は、射の明らかな概念によって圏をなす。
明らかに、$M$ の極限は、存在すれば錐の圏の終対象である。
双対的に、$M$ の \textit{余錐} とは、対象 $W$ と射の族
$t_i : M_i \to W$ であって、$\mathcal{I}$ のすべての射
$\phi : i \to i'$ に対して図式
$$
\xymatrix{
M_i \ar[rr]^{M(\phi)} \ar[dr]_{t_i} & & M_{i'} \ar[dl]^{t_{i'}} \\
& W
}
$$
が可換となるものをいう。余錐全体も、射の明らかな概念によって圏をなす。
上と同様に、$M$ の余極限が存在するための必要十分条件は、余錐の圏が始対象をもつことである。
\end{remark}

\noindent
極限と余極限の概念の応用として、積と余積を定義する。

\begin{definition}
\label{categories-definition-product}
$I$ を集合とし、各 $i \in I$ に対して圏 $\mathcal{C}$ の対象
$M_i$ が与えられているとする。\textit{積} $\prod_{i\in I} M_i$ とは、
定義により $\lim_\mathcal{I} M$（存在する場合）である。ここで
$\mathcal{I}$ は、$I$ の元を対象とし、射として恒等射のみをもつ圏である。
\end{definition}

\noindent
重要な特別の場合は $I = \emptyset$ であり、このとき積は圏の終対象である。
射 $p_i : \prod M_i \to M_i$ を \textit{射影射} という。

\begin{definition}
\label{categories-definition-coproduct}
$I$ を集合とし、各 $i \in I$ に対して圏 $\mathcal{C}$ の対象
$M_i$ が与えられているとする。\textit{余積} $\coprod_{i\in I} M_i$ とは、
定義により $\colim_\mathcal{I} M$（存在する場合）である。ここで
$\mathcal{I}$ は、$I$ の元を対象とし、射として恒等射のみをもつ圏である。
\end{definition}

\noindent
重要な特別の場合は $I = \emptyset$ であり、このとき余積は圏の始対象である。
余積には射 $M_i \to \coprod M_i$ が備わっていることに注意する。
これらは \textit{余射影} と呼ばれることがある。

\begin{lemma}
\label{categories-lemma-functorial-colimit}
$M : \mathcal{I} \to \mathcal{C}$ と
$N : \mathcal{J} \to \mathcal{C}$ を、余極限が存在する図式とする。
$H : \mathcal{I} \to \mathcal{J}$ を関手とし、
$t : M \to N \circ H$ を関手の変換とする。このとき一意な射
$$
\theta :
\colim_\mathcal{I} M
\longrightarrow
\colim_\mathcal{J} N
$$
が存在し、すべての図式
$$
\xymatrix{
M_i \ar[d]_{t_i} \ar[r]
&
\colim_\mathcal{I} M \ar[d]^{\theta}
\\
N_{H(i)} \ar[r]
&
\colim_\mathcal{J} N
}
$$
は可換である。
\end{lemma}

\begin{proof}
省略する。
\end{proof}

\begin{lemma}
\label{categories-lemma-functorial-limit}
$M : \mathcal{I} \to \mathcal{C}$ と
$N : \mathcal{J} \to \mathcal{C}$ を、極限が存在する図式とする。
$H : \mathcal{I} \to \mathcal{J}$ を関手とし、
$t : N \circ H \to M$ を関手の変換とする。このとき一意な射
$$
\theta :
\lim_\mathcal{J} N
\longrightarrow
\lim_\mathcal{I} M
$$
が存在し、すべての図式
$$
\xymatrix{
\lim_\mathcal{J} N \ar[d]^{\theta} \ar[r]
&
N_{H(i)} \ar[d]_{t_i}
\\
\lim_\mathcal{I} M \ar[r]
&
M_i
}
$$
は可換である。
\end{lemma}

\begin{proof}
省略する。
\end{proof}


\begin{lemma}
\label{categories-lemma-colimits-commute}
$\mathcal{I}$ と $\mathcal{J}$ を添字圏とし、
$M : \mathcal{I} \times \mathcal{J} \to \mathcal{C}$ を関手とする。
すべての $i$ に対して $M_{i, \infty} = \colim_j M_{i,j}$ が存在すると仮定する。
このとき、得られる関手 $M_{-, \infty} : \mathcal{I} \to \mathcal{C}$ が
余極限をもつための必要十分条件は $M$ が余極限をもつことであり、その場合、
両者の余極限は一致する。特に、表示された余極限がすべて存在すれば
$$
\colim_i \colim_j M_{i, j}
=
\colim_{i, j} M_{i, j}
=
\colim_j \colim_i M_{i, j}
$$
である。極限についても同様である。
\end{lemma}

\begin{proof}
省略する。
\end{proof}

\begin{lemma}
\label{categories-lemma-limits-products-equalizers}
$M : \mathcal{I} \to \mathcal{C}$ を図式とする。
$I = \Ob(\mathcal{I})$ および $A = \text{Arrows}(\mathcal{I})$ と書き、
始域写像と終域写像を $s, t : A \to I$ とする。
$\prod_{i \in I} M_i$ と $\prod_{a \in A} M_{t(a)}$ が存在すると仮定する。
さらに、
$$
\xymatrix{
\prod_{i \in I} M_i
\ar@<1ex>[r]^\phi \ar@<-1ex>[r]_\psi
&
\prod_{a \in A} M_{t(a)}
}
$$
の等化子が存在すると仮定する。ここで射は、それぞれの成分により
$p_a \circ \psi = M(a) \circ p_{s(a)}$ および
$p_a \circ \phi = p_{t(a)}$ と定められる。この等化子が図式の極限である。
\end{lemma}

\begin{proof}
省略する。
\end{proof}


\begin{lemma}
\label{categories-lemma-colimits-coproducts-coequalizers}
\begin{slogan}
すべての余積と余等化子が存在すれば、すべての余極限が存在する。
\end{slogan}
$M : \mathcal{I} \to \mathcal{C}$ を図式とする。
$I = \Ob(\mathcal{I})$ および $A = \text{Arrows}(\mathcal{I})$ と書き、
始域写像と終域写像を $s, t : A \to I$ とする。
$\coprod_{i \in I} M_i$ と $\coprod_{a \in A} M_{s(a)}$ が存在すると仮定する。
さらに、
$$
\xymatrix{
\coprod_{a \in A} M_{s(a)}
\ar@<1ex>[r]^\phi \ar@<-1ex>[r]_\psi
&
\coprod_{i \in I} M_i
}
$$
の余等化子が存在すると仮定する。ここで射は、それぞれの成分により
次のように定められる。成分 $M_{s(a)}$ は $\psi$ により、射 $M(a)$ を介して
成分 $M_{t(a)}$ へ写る。成分 $M_{s(a)}$ は $\phi$ により、恒等射を介して
成分 $M_{s(a)}$ へ写る。この余等化子が図式の余極限である。
\end{lemma}

\begin{proof}
省略する。
\end{proof}

\section{集合の圏における極限と余極限}
\label{categories-section-limit-sets}

\noindent
$\textit{Sets}$ では極限と余極限が存在するだけでなく、それらを容易に記述できる。
実際、$M : \mathcal{I} \to \textit{Sets}$, $i \mapsto M_i$ を集合の図式とし、
$I = \Ob(\mathcal{I})$ と書く。極限は
$$
\lim_\mathcal{I} M
=
\{
(m_i)_{i\in I} \in \prod\nolimits_{i\in I} M_i
\mid
\forall \phi : i \to i' \text{ in }\mathcal{I},
M(\phi)(m_i) = m_{i'}
\}.
$$
と記述される。したがって、極限の元を、すべての集合 $M_i$ の元からなる
整合的な系と考える。

\medskip\noindent
一方、余極限は
$$
\colim_\mathcal{I} M
=
(\coprod\nolimits_{i\in I} M_i)/\sim
$$
である。ここで同値関係 $\sim$ は、$m_i \in M_i$,
$m_{i'} \in M_{i'}$ であり、ある $\phi : i \to i'$ に対して
$M(\phi)(m_i) = m_{i'}$ となるとき $m_i \sim m_{i'}$ とおくことで
生成される同値関係である。言い換えると、$m_i \in M_i$ と
$m_{i'} \in M_{i'}$ が同値であるとは、$\mathcal{I}$ に射の鎖
$$
\xymatrix{
&
i_1 \ar[ld] \ar[rd] & &
i_3 \ar[ld] & &
i_{2n-1} \ar[rd] & \\
i = i_0 & &
i_2 & &
\ldots & &
i_{2n} = i'
}
$$
が存在し、元 $m_{i_j} \in M_{i_j}$ が存在して、それらが上の
$\mathcal{I}$ の射から誘導される写像
$M_{i_{2k-1}} \to M_{i_{2k-2}}$ および
$M_{i_{2k-1}} \to M_{i_{2k}}$ のもとで互いに写り合うことをいう。

\medskip\noindent
これは扱いやすい種類の対象ではない。しかし図式がフィルター付きなら、
はるかに容易に記述できる。Section \ref{categories-section-directed-colimits} でこれを説明する。



\section{連結極限}
\label{categories-section-connected-limits}

\noindent
添字圏が連結であるとき、(余)極限を連結であるという。

\begin{definition}
\label{categories-definition-category-connected}
$x \sim y \Leftrightarrow \Mor_\mathcal{I}(x, y) \not = \emptyset$ により
生成される同値関係がちょうど一つの同値類をもつとき、圏 $\mathcal{I}$ は
\textit{連結} であるという。
\end{definition}

\noindent
ここでは、連結空間は空でないとする
Topology, Definition \ref{topology-definition-connected-components} の規約に従う。
次の結果は、ある漠然とした意味で連結極限を特徴づける。

\begin{lemma}
\label{categories-lemma-connected-limit-over-X}
$\mathcal{C}$ を圏とし、$X$ を $\mathcal{C}$ の対象とする。
$M : \mathcal{I} \to \mathcal{C}/X$ を、$X$ 上の対象の圏における図式とする。
添字圏 $\mathcal{I}$ が連結であり、$M$ の極限が $\mathcal{C}/X$ に存在するならば、
合成 $\mathcal{I} \to \mathcal{C}/X \to \mathcal{C}$ の極限も存在し、
同じ対象で与えられる。
\end{lemma}

\begin{proof}
$L \to X$ を $\mathcal{C}/X$ における極限を表現する対象とする。関手
$$
W \longmapsto \lim_i \Mor_\mathcal{C}(W, M_i).
$$
を考える。右辺の集合の元を $(\varphi_i)$ とする。各 $M_i$ には射
$s_i : M_i \to X$ が備わっているから、射
$f_i = s_i \circ \varphi_i : W \to X$ を得る。しかし $\mathcal{I}$ は
連結なので、すべての $f_i$ は等しい。また $\mathcal{I}$ は空でないので、
少なくとも一つの $f_i$ が存在する。したがって、この共通の写像 $W \to X$ は
$W$ に $\mathcal{C}/X$ の対象としての構造を与え、$(\varphi_i)$ は
$\lim_i \Mor_{\mathcal{C}/X}(W, M_i)$ の元を定める。よって、各 $\varphi_i$ が
$\phi$ と $L \to M_i$ の合成となるような一意な射 $\phi : W \to L$ を得る。
\end{proof}

\begin{lemma}
\label{categories-lemma-connected-colimit-under-X}
$\mathcal{C}$ を圏とし、$X$ を $\mathcal{C}$ の対象とする。
$M : \mathcal{I} \to X/\mathcal{C}$ を、$X$ の下の対象の圏における図式とする。
添字圏 $\mathcal{I}$ が連結であり、$M$ の余極限が $X/\mathcal{C}$ に存在するならば、
合成 $\mathcal{I} \to X/\mathcal{C} \to \mathcal{C}$ の余極限も存在し、
同じ対象で与えられる。
\end{lemma}

\begin{proof}
省略する。ヒント：この補題は Lemma \ref{categories-lemma-connected-limit-over-X} の双対である。
\end{proof}

\section{共終圏と共始圏}
\label{categories-section-cofinal}

\noindent
文献では、次の定義の cofinal の代わりに ``final'' という語を用いることがある。

\begin{definition}
\label{categories-definition-cofinal}
$H : \mathcal{I} \to \mathcal{J}$ を圏の間の関手とする。次を満たすとき、
\textit{$\mathcal{I}$ は $\mathcal{J}$ において共終である}、または
$H$ は \textit{共終} であるという。
\begin{enumerate}
\item すべての $y \in \Ob(\mathcal{J})$ に対して、
$x \in \Ob(\mathcal{I})$ と射 $y \to H(x)$ が存在する。
\item $y \in \Ob(\mathcal{J})$, $x, x' \in \Ob(\mathcal{I})$、および射
$y \to H(x)$, $y \to H(x')$ が与えられたとする。このとき、
$\mathcal{I}$ に射の列
$$
x = x_0 \leftarrow x_1 \rightarrow x_2 \leftarrow x_3 \rightarrow \ldots
\rightarrow x_{2n} = x'
$$
と、$\mathcal{J}$ に射 $y \to H(x_i)$ が存在し、
$y \to H(x_0)$ と $y \to H(x_{2n})$ は与えられた射であり、
$k = 0, \ldots, n - 1$ に対して図式
$$
\xymatrix{
& y \ar[ld] \ar[d] \ar[rd] \\
H(x_{2k}) & H(x_{2k + 1}) \ar[l] \ar[r] & H(x_{2k + 2})
}
$$
が可換である。
\end{enumerate}
言い換えると、$\mathcal{J}$ の対象 $y$ を固定し、$x$ が $\mathcal{I}$ の対象、
$b : y \to H(x)$ が射であるような対 $(x, b)$ の集合 $S$ を考える。
射 $a : x \to x'$ が存在して $b' = H(a) \circ b$ となるとき
$(x, b) \sim (x', b')$ とおくことにより $S$ 上に生成される同値関係を考える。
このとき $S$ はちょうど一つの同値類からなることを要請する。
\end{definition}

\begin{lemma}
\label{categories-lemma-cofinal}
$H : \mathcal{I} \to \mathcal{J}$ を圏の関手とし、
$\mathcal{I}$ が $\mathcal{J}$ において共終であると仮定する。このとき、任意の図式
$M : \mathcal{J} \to \mathcal{C}$ に対して、いずれか一方が存在すれば標準同型
$$
\colim_\mathcal{I} M \circ H
=
\colim_\mathcal{J} M
$$
が成り立つ。
\end{lemma}

\begin{proof}
省略する。
\end{proof}

\begin{definition}
\label{categories-definition-initial}
$H : \mathcal{I} \to \mathcal{J}$ を圏の間の関手とする。次を満たすとき、
\textit{$\mathcal{I}$ は $\mathcal{J}$ において共始である}、または
$H$ は \textit{共始} であるという。
\begin{enumerate}
\item すべての $y \in \Ob(\mathcal{J})$ に対して、
$x \in \Ob(\mathcal{I})$ と射 $H(x) \to y$ が存在する。
\item 任意の $y \in \Ob(\mathcal{J})$, $x , x' \in \Ob(\mathcal{I})$ と、
$\mathcal{J}$ の射 $H(x) \to y$, $H(x') \to y$ に対して、
$\mathcal{I}$ に射の列
$$
x = x_0 \leftarrow x_1 \rightarrow x_2 \leftarrow x_3 \rightarrow \ldots
\rightarrow x_{2n} = x'
$$
と、$\mathcal{J}$ に射 $H(x_i) \to y$ が存在し、
$k = 0, \ldots, n - 1$ に対して図式
$$
\xymatrix{
H(x_{2k}) \ar[rd] &
H(x_{2k + 1}) \ar[l] \ar[r] \ar[d] &
H(x_{2k + 2}) \ar[ld] \\
& y
}
$$
が可換である。
\end{enumerate}
\end{definition}

\noindent
これは共終関手の双対概念にほかならない。

\begin{lemma}
\label{categories-lemma-initial}
$H : \mathcal{I} \to \mathcal{J}$ を圏の関手とし、
$\mathcal{I}$ が $\mathcal{J}$ において共始であると仮定する。
このとき、任意の図式 $M : \mathcal{J} \to \mathcal{C}$ に対して、
いずれか一方が存在すれば標準同型
$$
\lim_\mathcal{I} M \circ H = \lim_\mathcal{J} M
$$
が成り立つ。
\end{lemma}

\begin{proof}
省略する。
\end{proof}

\begin{lemma}
\label{categories-lemma-colimit-constant-connected-fibers}
$F : \mathcal{I} \to \mathcal{I}'$ を関手とする。次を仮定する。
\begin{enumerate}
\item $\mathcal{I}$ の $\mathcal{I}'$ 上のファイバー圏
（Definition \ref{categories-definition-fibre-category} を参照）はすべて連結である。
\item $\mathcal{I}'$ のすべての射 $\alpha' : x' \to y'$ に対して、
$F(\alpha) = \alpha'$ を満たす $\mathcal{I}$ の射 $\alpha : x \to y$ が存在する。
\end{enumerate}
このとき、任意の図式 $M : \mathcal{I}' \to \mathcal{C}$ に対して、
余極限 $\colim_\mathcal{I} M \circ F$ が存在するための必要十分条件は
$\colim_{\mathcal{I}'} M$ が存在することであり、その場合これらの余極限は一致する。
\end{lemma}

\begin{proof}
$\mathcal{I}$ が $\mathcal{I}'$ において共終であることを示して
Lemma \ref{categories-lemma-cofinal} を適用すれば証明できる。しかし、次のように直接証明することもできる。
$\mathcal{C}$ の任意の対象 $T$ に対して
$$
\lim_{\mathcal{I}^{opp}} \Mor_\mathcal{C}(M_{F(i)}, T)
=
\lim_{(\mathcal{I}')^{opp}} \Mor_\mathcal{C}(M_{i'}, T)
$$
となることを示せば十分である。右辺の元を
$(g_{i'})_{i' \in \Ob(\mathcal{I}')}$ とすれば、$f_i = g_{F(i)}$ とおくことにより
左辺の元 $(f_i)_{i \in \Ob(\mathcal{I})}$ を得る。逆に、左辺の元を
$(f_i)_{i \in \Ob(\mathcal{I})}$ とする。各（連結な）ファイバー圏
$\mathcal{I}_{i'}$ 上で関手 $M \circ F$ は値 $M_{i'}$ をもつ定値関手である。
したがって、$F(i) = i'$ を満たす $i \in \Ob(\mathcal{I})$ に対する射
$f_i$ はすべて同じであり、良定義な射 $g_{i'} : M_{i'} \to T$ を定める。
仮定 (2) により、族 $(g_{i'})_{i' \in \Ob(\mathcal{I}')}$ は右辺の元を定める。
\end{proof}

\begin{lemma}
\label{categories-lemma-product-with-connected}
$\mathcal{I}$ と $\mathcal{J}$ を圏とし、射影を
$p : \mathcal{I} \times \mathcal{J} \to \mathcal{J}$ と書く。
$\mathcal{I}$ が連結ならば、図式 $M : \mathcal{J} \to \mathcal{C}$ に対して、
余極限 $\colim_\mathcal{J} M$ が存在するための必要十分条件は
$\colim_{\mathcal{I} \times \mathcal{J}} M \circ p$ が存在することであり、
その場合これらの余極限は等しい。
\end{lemma}

\begin{proof}
これは Lemma \ref{categories-lemma-colimit-constant-connected-fibers} の特別な場合である。
\end{proof}

\section{有限極限と有限余極限}
\label{categories-section-finite-limits}

\noindent
添字圏が有限、すなわち対象も射も有限個しかもたない (余)極限を
\textit{有限} (余)極限という。添字圏が空でない (余)極限を
\textit{非空} という。添字圏が連結である (余)極限を
\textit{連結} という。Definition \ref{categories-definition-category-connected} を参照されたい。
実は、有限添字圏だけで「十分」である。

\begin{lemma}
\label{categories-lemma-finite-diagram-category}
$\mathcal{I}$ を次の条件を満たす圏とする。
\begin{enumerate}
\item $\Ob(\mathcal{I})$ は有限である。
\item $\mathcal{I}$ のすべての射が合成
$f_{j_1} \circ f_{j_2} \circ \ldots \circ f_{j_k}$ として書けるような、
有限個の射 $f_1, \ldots, f_m \in \text{Arrows}(\mathcal{I})$ が存在する。
\end{enumerate}
このとき、次を満たす関手 $F : \mathcal{J} \to \mathcal{I}$ が存在する。
\begin{enumerate}
\item[(a)] $\mathcal{J}$ は有限圏である。
\item[(b)] 任意の図式 $M : \mathcal{I} \to \mathcal{C}$ に対して、
$\mathcal{I}$ 上の $M$ の (余)極限が存在するための必要十分条件は、
$\mathcal{J}$ 上の $M \circ F$ の (余)極限が存在することであり、
この場合、両 (余)極限は標準的に同型である。
\end{enumerate}
さらに、$\mathcal{J}$ が連結（それぞれ非空）であるための必要十分条件は、
$\mathcal{I}$ がそうであることである。
\end{lemma}

\begin{proof}
$\Ob(\mathcal{I}) = \{x_1, \ldots, x_n\}$ とする。
$f_j : x_{s(j)} \to x_{t(j)}$ となるような関数
$s, t : \{1, \ldots, m\} \to \{1, \ldots, n\}$ を考える。
$\Ob(\mathcal{J}) = \{y_1, \ldots, y_n, z_1, \ldots, z_n\}$ とおく。
恒等射のほかに、射 $g_j : y_{s(j)} \to z_{t(j)}$,
$j = 1, \ldots, m$ と、射 $h_i : y_i \to z_i$,
$i = 1, \ldots, n$ を導入する。$\mathcal{J}$ のすべての非恒等射は
$y$ から $z$ へ向かうので、定義すべき合成も、確認すべき結合律もない。
$F(y_i) = F(z_i) = x_i$ とおき、$F(g_j) = f_j$ および
$F(h_i) = \text{id}_{x_i}$ とおく。$F$ が関手であることは明らかである。
$\mathcal{J}$ が有限であることも明らかである。
$\mathcal{J}$ が連結、あるいは非空であるための必要十分条件が
$\mathcal{I}$ について同じことが成り立つことであるのも明らかである。

\medskip\noindent
$M : \mathcal{I} \to \mathcal{C}$ を図式とする。
$\mathcal{C}$ の対象 $W$ と、Definition \ref{categories-definition-limit} におけるような射
$q_i : W \to M(x_i)$ を考える。このとき、
$q_i : W \to M(F(y_i)) = M(F(z_i)) = M(x_i)$ とすることにより、
図式 $M \circ F$ に対して Definition \ref{categories-definition-limit} におけるような
写像の族を得る。逆に、図式 $M \circ F$ に対して
Definition \ref{categories-definition-limit} におけるような写像
$qy_i : W \to M(F(y_i))$ および $qz_i : W \to M(F(z_i))$ が
与えられているとする。
$$
M(F(h_i)) = \text{id} : M(F(y_i)) = M(x_i) \longrightarrow M(x_i) = M(F(z_i))
$$
なので、すべての $i$ に対して $qy_i = qz_i$ である。この共通の値を
$q_i$ とおく。$q_{s(j)} = qy_{s(j)}$ と $q_{t(j)} = qz_{t(j)}$ の
射 $M(f_j)$ に関する整合性は、族 $q_i$ が $\mathcal{I}$ のすべての射と
整合することを保証する。実際、仮定によりそのような射はすべて射 $f_j$ の合成である。
したがって標準全単射
$$
\lim_{B \in \Ob(\mathcal{J})} \Mor_\mathcal{C}(W, M(F(B)))
=
\lim_{A \in \Ob(\mathcal{I})} \Mor_\mathcal{C}(W, M(A))
$$
を得る。これは補題の極限に関する主張を含意する。余極限に関する主張も
同じ方法で証明される（証明は省略する）。
\end{proof}

\begin{lemma}
\label{categories-lemma-fibre-products-equalizers-exist}
$\mathcal{C}$ を圏とする。次は同値である。
\begin{enumerate}
\item $\mathcal{C}$ に連結有限極限が存在する。
\item $\mathcal{C}$ に等化子とファイバー積が存在する。
\end{enumerate}
\end{lemma}

\begin{proof}
等化子とファイバー積は連結有限極限なので、(1) から (2) が従う。
逆を示すため、$\mathcal{I}$ を有限連結添字圏とする。
Lemma \ref{categories-lemma-finite-diagram-category} の証明で構成された添字圏の関手を
$F : \mathcal{J} \to \mathcal{I}$ とする。すると $\mathcal{I}$ を
$\mathcal{J}$ で置き換えてよい。したがって、
$\Ob(\mathcal{I}) = \{x_1, \ldots, x_n\} \amalg \{y_1, \ldots, y_m\}$,
$n, m \geq 1$ であり、$\mathcal{I}$ のすべての非恒等射は、ある $i$, $j$ に対する
射 $f : x_i \to y_j$ であると仮定してよい。

\medskip\noindent
$n > 1$ と仮定する。$\mathcal{I}$ は連結なので、添字 $i_1, i_2$, $j_0$ と射
$a : x_{i_1} \to y_{j_0}$, $b : x_{i_2} \to y_{j_0}$ が存在する。圏
$$
\mathcal{I}' =
\{x\} \amalg \{x_1, \ldots, \hat x_{i_1}, \ldots, \hat x_{i_2}, \ldots x_n\}
\amalg \{y_1, \ldots, y_m\}
$$
を、
$$
\Mor_{\mathcal{I}'}(x, y_j) = \Mor_\mathcal{I}(x_{i_1}, y_j)
\amalg \Mor_\mathcal{I}(x_{i_2}, y_j)
$$
とし、それ以外の射集合は $\mathcal{I}$ と同じものとして定める。任意の関手
$M : \mathcal{I} \to \mathcal{C}$ に対して、
$$
M'(x) = M(x_{i_1}) \times_{M(a), M(y_{j_0}), M(b)} M(x_{i_2})
$$
とおき、例えば $f : x_{i_1} \to y_j$ に対応する射
$f' : x \to y_j$ に対して $M'(f) = M(f) \circ \text{pr}_1$ とおくことにより、
関手 $M' : \mathcal{I}' \to \mathcal{C}$ を構成できる。このとき、関手 $M$ が
極限をもつための必要十分条件は、関手 $M'$ が極限をもつことである
（証明は省略する）。したがって帰納法により $n = 1$ の場合に帰着する。

\medskip\noindent
$n = 1$ ならば、任意の $M : \mathcal{I} \to \mathcal{C}$ の極限は、
射 $x_1 \to y_j$ の対の逐次的な等化子であるから、仮定により存在する。
\end{proof}

\begin{lemma}
\label{categories-lemma-almost-finite-limits-exist}
$\mathcal{C}$ を圏とする。次は同値である。
\begin{enumerate}
\item $\mathcal{C}$ に非空有限極限が存在する。
\item $\mathcal{C}$ に二対象の積と等化子が存在する。
\item $\mathcal{C}$ に二対象の積とファイバー積が存在する。
\end{enumerate}
\end{lemma}

\begin{proof}
二対象の積、ファイバー積、および等化子は、添字圏が空でない極限なので、
(1) は (2) と (3) の両方を含意する。(2) を仮定する。このとき有限非空積と
等化子が存在する。よって Lemma \ref{categories-lemma-limits-products-equalizers} により
有限非空極限が存在し、(1) が成り立つ。(3) を仮定する。
$a, b : A \to B$ を $\mathcal{C}$ の射とすれば、$a, b$ の等化子は
$$
(A \times_{a, B, b} A)\times_{(\text{pr}_1, \text{pr}_2), A \times A, \Delta} A.
$$
である。したがって (3) は (2) を含意し、補題が証明された。
\end{proof}

\begin{lemma}
\label{categories-lemma-finite-limits-exist}
$\mathcal{C}$ を圏とする。次は同値である。
\begin{enumerate}
\item $\mathcal{C}$ に有限極限が存在する。
\item 有限積と等化子が存在する。
\item 圏に終対象があり、ファイバー積が存在する。
\end{enumerate}
\end{lemma}

\begin{proof}
有限積、ファイバー積、等化子、および終対象は有限添字圏上の極限なので、
(1) は (2) と (3) の両方を含意する。上の
Lemma \ref{categories-lemma-limits-products-equalizers} により (2) は (1) を含意する。
(3) を仮定する。積 $A \times B$ は終対象上のファイバー積であることに注意する。
$a, b : A \to B$ を $\mathcal{C}$ の射とすれば、$a, b$ の等化子は
$$
(A \times_{a, B, b} A)\times_{(\text{pr}_1, \text{pr}_2), A \times A, \Delta} A.
$$
である。したがって (3) は (2) を含意し、補題が証明された。
\end{proof}

\begin{lemma}
\label{categories-lemma-push-outs-coequalizers-exist}
$\mathcal{C}$ を圏とする。次は同値である。
\begin{enumerate}
\item $\mathcal{C}$ に連結有限余極限が存在する。
\item $\mathcal{C}$ に余等化子と押し出しが存在する。
\end{enumerate}
\end{lemma}

\begin{proof}
省略する。ヒント：これは Lemma \ref{categories-lemma-fibre-products-equalizers-exist} の双対である。
\end{proof}

\begin{lemma}
\label{categories-lemma-almost-finite-colimits-exist}
$\mathcal{C}$ を圏とする。次は同値である。
\begin{enumerate}
\item $\mathcal{C}$ に非空有限余極限が存在する。
\item $\mathcal{C}$ に二対象の余積と余等化子が存在する。
\item $\mathcal{C}$ に二対象の余積と押し出しが存在する。
\end{enumerate}
\end{lemma}

\begin{proof}
省略する。ヒント：これは Lemma \ref{categories-lemma-almost-finite-limits-exist} の双対である。
\end{proof}

\begin{lemma}
\label{categories-lemma-colimits-exist}
$\mathcal{C}$ を圏とする。次は同値である。
\begin{enumerate}
\item $\mathcal{C}$ に有限余極限が存在する。
\item $\mathcal{C}$ に有限余積と余等化子が存在する。
\item 圏に始対象があり、押し出しが存在する。
\end{enumerate}
\end{lemma}

\begin{proof}
省略する。ヒント：これは Lemma \ref{categories-lemma-finite-limits-exist} の双対である。
\end{proof}

\section{フィルター付き余極限}
\label{categories-section-directed-colimits}

\noindent
フィルター付き図式上の余極限は、計算も記述もより容易である。まず定義を与える。

\begin{definition}
\label{categories-definition-directed}
図式 $M : \mathcal{I} \to \mathcal{C}$ が \textit{有向}、または
\textit{フィルター付き} であるとは、次の条件が成り立つことをいう。
\begin{enumerate}
\item 圏 $\mathcal{I}$ は少なくとも一つの対象をもつ。
\item $\mathcal{I}$ の任意の二対象 $x, y$ に対して、対象 $z$ と射
$x \to z$, $y \to z$ が存在する。
\item $\mathcal{I}$ の任意の二対象 $x, y$ と、$\mathcal{I}$ の任意の二射
$a, b : x \to y$ に対して、$\mathcal{C}$ の射として
$M(c \circ a) = M(c \circ b)$ を満たす $\mathcal{I}$ の射
$c : y \to z$ が存在する。
\end{enumerate}
添字圏 $\mathcal{I}$ が \textit{有向}、または \textit{フィルター付き}
であるとは、$\text{id} : \mathcal{I} \to \mathcal{I}$ がフィルター付きであることをいう
（言い換えると、上の (3) から $M$ を消去する）。
\end{definition}

\noindent
フィルター付き添字圏をもつ図式はすべてフィルター付きであり、通常、
フィルター付き余極限はこのように現れる。実際、
$M : \mathcal{I} \to \mathcal{C}$ がフィルター付き図式ならば、
$M$ を $\mathcal{I} \to \mathcal{I}' \to \mathcal{C}$ と分解できる。
ここで $\mathcal{I}'$ はフィルター付き添字圏であり\footnote{具体的には、
$\mathcal{I}'$ の対象を $\mathcal{I}$ と同じものとし、
$\Mor_{\mathcal{I}'}(x, y)$ を $\Mor_\mathcal{I}(x, y)$ の、
$M(a) = M(b)$ となる $a, b : x \to y$ を同一視する同値関係による商とする。}、
$\colim_\mathcal{I} M$ が存在するための必要十分条件は
$\colim_{\mathcal{I}'} M'$ が存在することであり、その場合、両余極限は標準的に同型である。

\medskip\noindent
$M : \mathcal{I} \to \textit{Sets}$ をフィルター付き図式とする。この場合、公式
$$
\colim_\mathcal{I} M
=
(\coprod\nolimits_{i\in I} M_i)/\sim
$$
における同値関係を、単に
$$
m_i \sim m_{i'}
\Leftrightarrow
\exists i'', \phi : i \to i'', \phi': i' \to i'',
M(\phi)(m_i) = M(\phi')(m_{i'}).
$$
と記述できる。言い換えると、二つの元が余極限において等しいための必要十分条件は、
それらが「やがて等しくなる」ことである。

\begin{lemma}
\label{categories-lemma-directed-commutes}
$\mathcal{I}$ と $\mathcal{J}$ を添字圏とし、$\mathcal{I}$ はフィルター付き、
$\mathcal{J}$ は有限であると仮定する。
$M : \mathcal{I} \times \mathcal{J} \to \textit{Sets}$,
$(i, j) \mapsto M_{i, j}$ を集合の図式からなる図式とする。このとき
$$
\colim_i \lim_j M_{i, j}
=
\lim_j \colim_i M_{i, j}.
$$
である。特に、$\mathcal{I}$ 上の余極限は、集合の有限積、ファイバー積、
および等化子と交換する。
\end{lemma}

\begin{proof}
省略する。実際、ある圏がフィルター付きであるための必要十分条件が、
その圏上の余極限が有限極限（集合の圏へのもの）と交換することであると証明するのは、
楽しい演習である。
\end{proof}

\noindent
$\mathcal{J}$ が無限の場合のこの補題に対する反例を与える。
$\mathcal{I}$ を $\mathbf{N} = \{1, 2, 3, \ldots\}$ からなり、
$i \leq i'$ のときに一意な射 $i \to i'$ をもつ圏とする。
$\mathcal{J}$ を離散圏 $\mathbf{N} = \{1, 2, 3, \ldots\}$
（射は恒等射のみ）とする。$M_{i, j} = \{1, 2, \ldots, i\}$ とし、
$i \leq i'$ のときの写像 $M_{i, j} \to M_{i', j}$ は明らかな包含写像とする。
この場合、$\colim_i M_{i, j} = \mathbf{N}$ であり、したがって
$$
\lim_j \colim_i M_{i, j}
=
\prod\nolimits_j \mathbf{N}
=
\mathbf{N}^\mathbf{N}
$$
である。一方、$\lim_j M_{i, j} = \prod\nolimits_j M_{i, j}$ なので
$$
\colim_i \lim_j M_{i, j}
=
\bigcup\nolimits_i \{1, 2, \ldots, i\}^{\mathbf{N}}
$$
となり、これはもう一方の極限より小さい。

\begin{lemma}
\label{categories-lemma-cofinal-in-filtered}
\begin{reference}
\cite[Proposition 3.2.4]{KS}
\end{reference}
$\mathcal{I}$ を圏とし、$\mathcal{J}$ を充満部分圏とする。
$\mathcal{I}$ はフィルター付きであると仮定する。さらに、$\mathcal{I}$ の
任意の対象 $i$ に対して、$\mathcal{J}$ のある対象 $j$ への射
$i \to j$ が存在すると仮定する。このとき $\mathcal{J}$ は
フィルター付きであり、$\mathcal{I}$ において共終である。
\end{lemma}

\begin{proof}
省略する。関係する概念についての楽しい演習である。
\end{proof}

\noindent
添字圏に課しうる条件を、より細かく制御する必要がある場合がある。
そこで、要請しうる事柄についていくつかの補題を加える。

\begin{lemma}
\label{categories-lemma-preserve-products}
$\mathcal{I}$ を添字圏、すなわち圏とする。$\mathcal{I}$ の任意の二対象
$x, y$ に対して、対象 $z$ と射 $x \to z$, $y \to z$ が存在すると仮定する。
このとき次が成り立つ。
\begin{enumerate}
\item $M$ と $N$ が $\mathcal{I}$ 上の集合の図式ならば、
$\colim (M_i \times N_i) \to \colim M_i \times \colim N_i$ は全射である。
\item 一般には、$\mathcal{I}$ 上の集合の図式の余極限は、有限非空積と交換しない。
\end{enumerate}
\end{lemma}

\begin{proof}
(1) の証明。$\colim M_i \times \colim N_i$ の元を
$(\overline{m}, \overline{n})$ とする。ある $x, y \in \Ob(\mathcal{I})$ に対して
$m \in M_x$, $n \in N_y$ を、$m$ が $\overline{m}$ に写り、$n$ が
$\overline{n}$ に写るように選べる。Section \ref{categories-section-limit-sets} を参照されたい。
$\mathcal{I}$ において $a : x \to z$ と $b : y \to z$ を選ぶ。このとき
$(M(a)(m), N(b)(n))$ は $(M \times N)_z$ の元であり、
$\colim (M_i \times N_i)$ におけるその像は、望みどおり
$(\overline{m}, \overline{n})$ に写る。

\medskip\noindent
(2) の証明。$G$ を非自明群とし、$\mathcal{I}$ を自己準同型モノイドが $G$ である
一対象圏とする。このとき $\mathcal{I}$ は補題で述べた条件を自明に満たす。
$G$ を自身に平行移動で作用させ、$G$-集合 $G$ を集合値の
$\mathcal{I}$-図式とみなす。このとき
$$
\colim_\mathcal{I} G \times \colim_\mathcal{I} G \cong G/G \times G/G
$$
は
$$
\colim_\mathcal{I} (G \times G) \cong (G \times G)/G
$$
と同型ではない。この例は、有限積だけに関心がある場合でも、
Lemma \ref{categories-lemma-directed-commutes} の追加条件を単に取り除くことはできないと示している。
\end{proof}

\begin{lemma}
\label{categories-lemma-colimits-abelian-as-sets}
$\mathcal{I}$ を添字圏、すなわち圏とする。$\mathcal{I}$ の任意の二対象
$x, y$ に対して、対象 $z$ と射 $x \to z$, $y \to z$ が存在すると仮定する。
$M : \mathcal{I} \to \textit{Ab}$ を $\mathcal{I}$ 上のアーベル群の図式とする。
このとき、集合の圏における $M$ の余極限から、アーベル群の圏における
$M$ の余極限への写像は全射である。
\end{lemma}

\begin{proof}
集合の圏における余極限は、関係による非交和 $\coprod M_i$ の商であった。
Section \ref{categories-section-limit-sets} を参照されたい。同様に、アーベル群の圏における
余極限は直和 $\bigoplus M_i$ の商である。補題の仮定は、
$i, j \in \Ob(\mathcal{I})$, $m \in M_i$, $n \in M_j$ が与えられたとき、
対象 $k$ と射 $a : i \to k$, $b : j \to k$ を選べることを意味する。
したがって $m + n$ は、$M_k$ の元 $M(a)(m) + M(b)(n)$ によって
余極限で表現される。ゆえに $\coprod M_i$ は余極限へ全射となる。
\end{proof}

\begin{lemma}
\label{categories-lemma-split-into-connected}
$\mathcal{I}$ を添字圏、すなわち圏とする。$\mathcal{I}$ において、
任意の実線部分が与えられた図式
$$
\xymatrix{
x \ar[d] \ar[r] & y \ar@{..>}[d] \\
z \ar@{..>}[r] & w
}
$$
に対し、図式を可換にする対象 $w$ と点線の射が存在すると仮定する。
このとき $\mathcal{I}$ は空であるか、または同じ性質をもつ非空な連結圏の
非空な非交和である。
\end{lemma}

\begin{proof}
$\mathcal{I}$ が空圏ならば補題は成り立つ。そうでなければ、
$\mathcal{I}$ の対象上の関係を、ある $z$ と射 $x \to z$, $y \to z$
が存在するとき $x \sim y$ と定める。補題の仮定により、これは同値関係である。
したがって $\Ob(\mathcal{I})$ は同値類の非交和である。
これらの同値類に対応する充満部分圏を $\mathcal{I}_j$ とする。
すると $\mathcal{I} = \coprod \mathcal{I}_j$ であり、望みどおり各
$\mathcal{I}_j$ は非空である。
\end{proof}

\begin{lemma}
\label{categories-lemma-preserve-injective-maps}
$\mathcal{I}$ を添字圏、すなわち圏とする。$\mathcal{I}$ において、
任意の実線部分が与えられた図式
$$
\xymatrix{
x \ar[d] \ar[r] & y \ar@{..>}[d] \\
z \ar@{..>}[r] & w
}
$$
に対し、図式を可換にする対象 $w$ と点線の射が存在すると仮定する。
このとき次が成り立つ。
\begin{enumerate}
\item $\mathcal{I}$ 上の集合の図式の単射な射 $M \to N$ は、集合の単射
$\colim M_i \to \colim N_i$ を誘導する。
\item 一般には、アーベル群の図式とその余極限について同じことは成り立たない。
\end{enumerate}
\end{lemma}

\begin{proof}
$\mathcal{I}$ が空圏ならば補題は成り立つ。そこで $\mathcal{I}$ は非空と仮定してよい。
この場合、各 $\mathcal{I}_j$ が非空で同じ性質を満たすように
$\mathcal{I} = \coprod \mathcal{I}_j$ と書ける。Lemma
\ref{categories-lemma-split-into-connected} を参照されたい。
$\colim_\mathcal{I} M = \coprod_j \colim_{\mathcal{I}_j} M|_{\mathcal{I}_j}$
なので、(1) の証明は連結な場合に帰着する。

\medskip\noindent
$\mathcal{I}$ は連結で、$M \to N$ は単射、すなわちすべての写像
$M_i \to N_i$ が単射であると仮定する。$M_i$ を $M_i \to N_i$ の像と同一視し、
すなわち $M_i$ を $N_i$ の部分集合とみなす。
以下では Section \ref{categories-section-limit-sets} で与えた余極限の記述を、
断りなく用いる。$s, s' \in \colim M_i$ が $\colim N_i$ の同じ元へ写るとする。
$s$ は $M_i$ の元 $m$ から生じ、$s'$ は $M_{i'}$ の元 $m'$ から生じるとする。
このとき、$\mathcal{I}$ の対象の列
$i = i_0, i_1, \ldots, i_n = i'$ と射
$$
\xymatrix{
&
i_1 \ar[ld] \ar[rd] & &
i_3 \ar[ld] & &
i_{2n-1} \ar[rd] & \\
i = i_0 & &
i_2 & &
\ldots & &
i_{2n} = i'
}
$$
および元 $n_{i_j} \in N_{i_j}$ を見いだせる。これらの元は、上の
$\mathcal{I}$ の射から誘導される写像 $N_{i_{2k-1}} \to N_{i_{2k-2}}$ と
$N_{i_{2k-1}} \to N_{i_{2k}}$ のもとで互いに写り合い、
$n_{i_0} = m$ かつ $n_{i_{2n}} = m'$ である。これが $s = s'$ を含意することを
$n$ に関する帰納法で示す。$n = 0$ の場合は自明である。
$n \geq 1$ と仮定する。$\mathcal{I}$ に関する仮定を用いると、可換図式
$$
\xymatrix{
& i_1 \ar[ld] \ar[rd] \\
i_0 \ar[rd] & & i_2 \ar[ld] \\
& w
}
$$
を得る。$m$ と $n_{i_2}$ はともに元 $n_{i_1}$ の像なので、$N_w$ の同じ元へ写る。
特に、この元は $m'' \in M_w$ であり、$\colim M_i$ において $s$ と同じ元を与える。
そこで鎖
$$
\xymatrix{
&
i_3 \ar[ld] \ar[rd] & &
i_5 \ar[ld] & &
i_{2n-1} \ar[rd] & \\
w & &
i_4 & &
\ldots & &
i_{2n} = i'
}
$$
と、$j \geq 3$ に対する元 $n_{i_j}$ が得られ、これは最初の鎖より短い。
これで帰納段階が示され、(1) の証明が完了する。

\medskip\noindent
$G$ を群とし、$\mathcal{I}$ を自己準同型モノイドが $G$ である一対象圏とする。
このとき、$g_1, g_2 \in G$ が与えられれば $h_1 g_1 = h_2 g_2$ を満たす
$h_1, h_2 \in G$ を見いだせるので、$\mathcal{I}$ は補題の条件を満たす。
$\textit{Ab}$ における $\mathcal{I}$ 上の図式 $M$ は、$G$-作用をもつアーベル群
$M$ と同じものであり、$\colim_\mathcal{I} M$ は $M$ の余不変量 $M_G$ である。
$G$ を位数 $2$ の群とし、$M = \mathbf{Z}/2\mathbf{Z}$ には自明に作用させ、これを
$N = \mathbf{Z}/2\mathbf{Z} \times \mathbf{Z}/2\mathbf{Z}$ の第一直和因子へ写す。
ただし $G$ の非自明元は後者に
$(x, y) \mapsto (x + y, y)$ により作用する。このとき $M_G \to N_G$ は零写像である。
\end{proof}

\begin{lemma}
\label{categories-lemma-split-into-directed}
$\mathcal{I}$ を添字圏、すなわち圏とする。次を仮定する。
\begin{enumerate}
\item $\mathcal{I}$ の任意の二射 $a : w \to x$ と $b : w \to y$ に対し、
対象 $z$ と射 $c : x \to z$, $d : y \to z$ が存在して
$c \circ a = d \circ b$ を満たす。
\item 任意の二射 $a, b : x \to y$ に対し、$c \circ a = c \circ b$ を満たす射
$c : y \to z$ が存在する。
\end{enumerate}
このとき $\mathcal{I}$ は、（空でもよい）互いに素なフィルター付き添字圏
$\mathcal{I}_j$ の和である。
\end{lemma}

\begin{proof}
$\mathcal{I}$ が空圏ならば補題は成り立つ。そうでなければ、
$\mathcal{I}$ の対象上の関係を、ある $z$ と射 $x \to z$, $y \to z$
が存在するとき $x \sim y$ と定める。補題の第一の仮定により、
これは同値関係である。したがって $\Ob(\mathcal{I})$ は同値類の非交和である。
これらの同値類に対応する充満部分圏を $\mathcal{I}_j$ とする。
残りは定義から明らかである。
\end{proof}

\begin{lemma}
\label{categories-lemma-almost-directed-commutes-equalizers}
$\mathcal{I}$ を、上の Lemma \ref{categories-lemma-split-into-directed} の仮定を満たす添字圏とする。
このとき $\mathcal{I}$ 上の余極限は、集合のファイバー積および等化子
（より一般には有限連結極限）と交換する。
\end{lemma}

\begin{proof}
Lemma \ref{categories-lemma-split-into-directed} により、各 $\mathcal{I}_j$ が
フィルター付きとなるように $\mathcal{I} = \coprod \mathcal{I}_j$ と書ける。
Lemma \ref{categories-lemma-directed-commutes} により、$\mathcal{I}_j$ の余極限は
等化子およびファイバー積と交換する。したがって、集合の圏において等化子と
ファイバー積が余積（空な余積を含む）と交換することを示せば十分である。
言い換えると、集合 $J$、集合 $A_j, B_j, C_j$、および $j \in J$ に対する集合写像
$A_j \to B_j$, $C_j \to B_j$ が与えられたとき、
$$
(\coprod\nolimits_{j \in J} A_j)
\times_{(\coprod\nolimits_{j \in J} B_j)}
(\coprod\nolimits_{j \in J} C_j)
=
\coprod\nolimits_{j \in J} A_j \times_{B_j} C_j
$$
を、また $a_j, a'_j : A_j \to B_j$ が与えられたとき、
$$
\text{Equalizer}(
\coprod\nolimits_{j \in J} a_j,
\coprod\nolimits_{j \in J} a'_j)
=
\coprod\nolimits_{j \in J}
\text{Equalizer}(a_j, a'_j)
$$
を示せばよい。これは $J = \emptyset$ の場合にも成り立つ。詳細は省略する。
\end{proof}


\section{余フィルター付き極限}
\label{categories-section-codirected-limits}

\noindent
余フィルター付き図式上の極限は、より容易に計算または記述できる。
まず定義を与える。

\begin{definition}
\label{categories-definition-codirected}
図式 $M : \mathcal{I} \to \mathcal{C}$ が次の条件を満たすとき、
これを {\it 余有向} または {\it 余フィルター付き} という。
\begin{enumerate}
\item 圏 $\mathcal{I}$ は少なくとも一つの対象をもつ。
\item $\mathcal{I}$ の任意の二対象 $x, y$ に対して、対象 $z$ と射
$z \to x$, $z \to y$ が存在する。
\item $\mathcal{I}$ の任意の二対象 $x, y$ および $\mathcal{I}$ の任意の二射
$a, b : x \to y$ に対して、$\mathcal{C}$ の射として
$M(a \circ c) = M(b \circ c)$ を満たす $\mathcal{I}$ の射
$c : w \to x$ が存在する。
\end{enumerate}
添字圏 $\mathcal{I}$ について、$\text{id} : \mathcal{I} \to \mathcal{I}$ が
余フィルター付きであるとき、この添字圏は {\it 余有向} または
{\it 余フィルター付き} であるという（言い換えれば、上の (3) から
$M$ を取り除く）。
\end{definition}

\noindent
余フィルター付き添字圏をもつ任意の図式は余フィルター付きである。
通常この状況はこのようにして現れる。

\medskip\noindent
集合の余フィルター付き極限が一般の極限より「容易」である理由の例として、
有限非空集合からなる余フィルター付き図式の極限は非空であるという事実を挙げる
（Lemma \ref{categories-lemma-nonempty-limit}）。この結果は、有限非空集合の一般の極限には
成り立たない。


\section{前順序集合上の極限と余極限}
\label{categories-section-posets-limits}

\noindent
図式の特別な場合として、前順序集合上の系がある。

\begin{definition}
\label{categories-definition-directed-set}
$I$ を集合とし、$\leq$ を $I$ 上の二項関係とする。
\begin{enumerate}
\item $\leq$ が推移的（$i \leq j$ かつ $j \leq k$ ならば
$i \leq k$）かつ反射的（すべての $i \in I$ に対して $i \leq i$）
であるとき、これを {\it 前順序} という。
\item 前順序を備えた集合を {\it 前順序集合} という。
\item 前順序集合 $(I, \leq)$ が {\it 有向集合} であるとは、
$I$ が空でなく、$\forall i, j \in I$ に対して、
$i \leq k, j \leq k$ を満たす $k \in I$ が存在することをいう。
\item $\leq$ が反対称的（$i \leq j$ かつ $j \leq i$ ならば
$i = j$）な前順序であるとき、これを {\it 半順序} という。
\item 半順序を備えた集合を {\it 半順序集合} という。
\item 順序が半順序である有向集合を {\it 有向半順序集合} という。
\end{enumerate}
\end{definition}

\noindent
前順序集合について述べる際には、記号から $\leq$ を省くのが慣例である。
すなわち、前順序集合 $(I, \leq)$ ではなく前順序集合 $I$ と呼ぶ。
前順序集合 $I$ が与えられたとき、記号 $\geq$ を、すべての
$i, j \in I$ に対する規則 $i \geq j \Leftrightarrow j \leq i$ によって定める。
英語の「partially ordered set」は「poset」と略されることがある。

\medskip\noindent
前順序集合 $I$ から次のように圏を構成できる。対象は $I$ の元であり、
$i \leq i'$ ならば射 $i \to i'$ がただ一つ存在し、そうでなければ存在しない。
逆に、任意の二対象間に高々一つの射しかない圏 $\mathcal{C}$ が与えられると、
集合 $\Ob(\mathcal{C})$ には規則
$x \leq y \Leftrightarrow \Mor_\mathcal{C}(x, y) \not = \emptyset$
で定まる前順序が入る。

\begin{definition}
\label{categories-definition-system-over-poset}
$(I, \leq)$ を前順序集合、$\mathcal{C}$ を圏とする。
\begin{enumerate}
\item $\mathcal{C}$ における $I$ 上の {\it 系}、または
$\mathcal{C}$ における $I$ 上の {\it 帰納系} とは、$\mathcal{C}$ の対象
$M_i$ と、各 $i \leq i'$ に対する射 $f_{ii'} : M_i \to M_{i'}$ であって、
$f_{ii} = \text{id}$ を満たし、$i \leq i' \leq i''$ のとき
$f_{ii''} = f_{i'i''} \circ f_{i i'}$ を満たすものをいう。
\item $\mathcal{C}$ における $I$ 上の {\it 逆系}、または
$\mathcal{C}$ における $I$ 上の {\it 射影系} とは、$\mathcal{C}$ の対象
$M_i$ と、各 $i' \leq i$ に対する射 $f_{ii'} : M_i \to M_{i'}$ であって、
$f_{ii} = \text{id}$ を満たし、$i'' \leq i' \leq i$ のとき
$f_{ii''} = f_{i'i''} \circ f_{i i'}$ を満たすものをいう
（不等号の向きが逆であることに注意されたい）。
\end{enumerate}
これを表すため、$(M_i, f_{ii'})$ は $I$ 上の（逆）系であるという。
写像 $f_{ii'}$ は {\it 遷移写像} と呼ばれることがある。
\end{definition}

\noindent
言い換えると、$I$ 上の系とは図式 $M : \mathcal{I} \to \mathcal{C}$ にほかならない。
ここで $\mathcal{I}$ は上で $I$ に対応させた圏、すなわち対象が $I$ の元で、
$i \leq i'$ であるとき、かつそのときに限って、$\mathcal{I}$ に一意な射
$i \to i'$ が存在する圏である。逆系は図式
$M : \mathcal{I}^{opp} \to \mathcal{C}$ である。
この観点からは、$I$ 上の任意の（逆）系の（余）極限を取ることもできる。
しかし、慣例として {\it $I$ 上の系については余極限のみ} を取り、
{\it $I$ 上の逆系については極限のみ} を取る。
より正確には、$I$ 上の系 $(M_i, f_{ii'})$ が与えられたとき、
系 $(M_i, f_{ii'})$ の余極限を
$$
\colim_{i \in I} M_i = \colim_\mathcal{I} M,
$$
すなわち対応する図式の余極限として定める。
$I$ 上の逆系 $(M_i, f_{ii'})$ が与えられたとき、
逆系 $(M_i, f_{ii'})$ の極限を
$$
\lim_{i \in I} M_i = \lim_{\mathcal{I}^{opp}} M,
$$
すなわち対応する図式の極限として定める。

\begin{remark}
\label{categories-remark-preorder-versus-partial-order}
$I$ を前順序集合とする。$I$ から標準的な半順序集合 $\overline{I}$ と
順序保存写像 $\pi : I \to \overline{I}$ を構成できる。
実際、$I$ 上の同値関係 $\sim$ を規則
$$
i \sim j \Leftrightarrow (i \leq j\text{ and }j \leq i).
$$
によって定める。$\overline{I} = I/\sim$ と置き、
$\pi : I \to \overline{I}$ を商写像とする。最後に $\overline{I}$ には、
$\pi(i) \leq \pi(j) \Leftrightarrow i \leq j$ を満たす一意な半順序が入る。
$I$ が有向集合ならば、$\overline{I}$ は有向半順序集合であることに注意されたい。
$\overline{I}$ 上の（逆）系 $N$ が与えられると、
$M_i = N_{\pi(i)}$ と置くことにより $I$ 上の（逆）系 $M$ を得る。
この構成は、$I$ および $\overline{I}$ 上の逆系の圏の間の関手を定める。
実際、これは同値である。その理由は、$i \sim j$ ならば、$I$ 上の任意の系
$M$ に対して、写像 $M_i \to M_j$ と $M_j \to M_i$ が互いに逆な同型だからである。
より正確には、$\pi$ の切断 $s : \overline{I} \to I$ を選ぶと、
上の関手の準逆は $M$ を、$N_{\overline{i}} = M_{s(\overline{i})}$ を満たす
$N$ へ送る。最後に、この対応は系の余極限と両立する。
$M$ と $N$ が上のように関係し、
$\colim_{\overline{I}} N$ または $\colim_I M$ のいずれかが存在すれば、
他方も存在し、$\colim_{\overline{I}} N = \colim_I M$ である。
逆系および逆系の極限についても同様の結果が成り立つ。
\end{remark}

\noindent
Remark \ref{categories-remark-preorder-versus-partial-order} の要点は、系の余極限または
逆系の極限を計算する際、前順序を常に半順序と仮定してよいということである。

\begin{definition}
\label{categories-definition-directed-system}
$I$ を前順序集合とする。系（それぞれ逆系）$(M_i, f_{ii'})$ が
{\it 有向系}（それぞれ {\it 有向逆系}）であるとは、$I$ が有向集合
（Definition \ref{categories-definition-directed-set}）、すなわち $I$ が非空で、
すべての $i_1, i_2 \in I$ に対して、$i\in I$ で
$i_1 \leq i$ かつ $i_2 \leq i$ を満たすものが存在することをいう。
\end{definition}

\noindent
この場合、余極限は（残念ながら）「直極限」と呼ばれることがある。
ここではこの用語を用いない。次の意味で、フィルター付き圏上の図式は
有向系より一般的なものではない。

\begin{lemma}
\label{categories-lemma-directed-category-system}
$\mathcal{I}$ をフィルター付き添字圏とする。次の性質をもつ有向集合 $I$ と、
$\mathcal{I}$ における $I$ 上の系 $(x_i, \varphi_{ii'})$ が存在する。
\begin{enumerate}
\item 任意の圏 $\mathcal{C}$ と、$\mathcal{C}$ に値をとる任意の図式
$M : \mathcal{I} \to \mathcal{C}$ に対し、対応する $I$ 上の系を
$(M(x_i), M(\varphi_{ii'}))$ と書く。$\colim_{i \in I} M(x_i)$ が存在すれば
$\colim_\mathcal{I} M$ も存在し、Lemma \ref{categories-lemma-functorial-colimit} の変換
$$
\theta :
\colim_{i \in I} M(x_i)
\longrightarrow
\colim_\mathcal{I} M
$$
は同型である。
\item 任意の圏 $\mathcal{C}$ と、$\mathcal{C}$ における任意の図式
$M : \mathcal{I}^{opp} \to \mathcal{C}$ に対し、対応する $I$ 上の逆系を
$(M(x_i), M(\varphi_{ii'}))$ と書く。$\lim_{i \in I} M(x_i)$ が存在すれば
$\lim_{\mathcal{I}^{opp}} M$ も存在し、Lemma \ref{categories-lemma-functorial-limit} の変換
$$
\theta :
\lim_{\mathcal{I}^{opp}} M
\longrightarrow
\lim_{i \in I} M(x_i)
$$
は同型である。
\end{enumerate}
\end{lemma}

\begin{proof}
Definition \ref{categories-definition-system-over-poset} の後で説明したように、
前順序集合を圏、系を関手とみなせる。証明を通して、この二つの観点を自由に行き来する。
第一の主張を示すため、有向集合に対応する圏 $\mathcal{I}_0$\footnote{実際、
この構成は有向半順序集合を与える。} と共終関手
$M_0 : \mathcal{I}_0 \to \mathcal{I}$ を構成する。
すると Lemma \ref{categories-lemma-cofinal} により、図式
$M : \mathcal{I} \to \mathcal{C}$ の余極限は図式
$M \circ M_0 : \mathcal{I}_0 \to \mathcal{C}$ の余極限と一致し、
主張が従う。第二の主張は第一の双対であり、$\mathcal{C}$ における極限を
$\mathcal{C}^{opp}$ における余極限と解釈して証明できる。詳細は省略する。

\medskip\noindent
圏 $\mathcal{F}$ が {\em 有限生成} であるとは、$\mathcal{F}$ の射からなる
有限集合 $F$ が存在して、$\mathcal{F}$ の各射を $F$ の射の合成として得られることをいう。
特に、これは $\mathcal{F}$ が有限個の対象しかもたないことを含意する。
まず、$\mathcal{I}$ の任意の有限生成部分圏が、一意な終対象をもつ有限生成部分圏へ
拡張できるという性質を $\mathcal{I}$ がもつ場合へ、証明を帰着する。

\medskip\noindent
$\omega$ を有限順序数からなる有向集合とし、これをフィルター付き圏とみなす。
積圏 $\mathcal{I}\times \omega$ もフィルター付きであり、射影
$\Pi : \mathcal{I} \times \omega \to \mathcal{I}$ は共終であることは、
容易に確かめられる。

\medskip\noindent
ここで $\mathcal{F}$ を $\mathcal{I}\times \omega$ の任意の有限生成部分圏とする。
フィルター付き圏の公理と、$\mathcal{F}$ の有限な生成元集合に関する単純な帰納法を用いて、
図式 $\mathcal{F} \to \mathcal{I} \times \omega$ に対する
$\mathcal{I} \times \omega$ の余錐 $(\{f_i\}, i_\infty)$ を構成できる。
すなわち、$\mathcal{F}$ の各対象 $i$ に対して射 $f_i : i \to i_\infty$ があり、
$\mathcal{F}$ の各射 $f : i \to i'$ に対して
$f_i = f_{i'} \circ f$ が成り立つ。さらに、$i_\infty$ から
$\mathcal{F}$ の対象への射が存在しないように $i_\infty$ を選べる。
実際、射 $f_i$ を、$\mathcal{I}$ 成分では恒等射であり
$\omega$ 成分では狭義に増加する射と常に後合成できる。
$\mathcal{F}^+$ を、$\mathcal{F}$ のすべての対象と射に加えて、対象
$i_\infty$、恒等射 $\text{id}_{i_\infty}$、および射 $f_i$ からなる圏とする。
$\mathcal{F}^+$ において $i_\infty$ から $\mathcal{F}$ の対象への射はないので、
$\mathcal{F}^+$ の射集合は合成について閉じており、したがって
$\mathcal{F}^+$ は確かに圏である。構成により、これは $i_\infty$ を一意な終対象にもつ
$\mathcal{I}$ の有限生成部分圏である。Lemma \ref{categories-lemma-cofinal} により、
任意の図式 $M : \mathcal{I} \to \mathcal{C}$ の余極限は
$M\circ\Pi$ の余極限と一致するので、望む帰着が得られる。

\medskip\noindent
一意な終対象をもつ $\mathcal{I}$ のすべての有限生成部分圏からなる集合には、
包含関係による自然な順序が入る。$\mathcal{I}_0$ をこの集合に対応する圏とする。
さらに、関手 $M_0 : \mathcal{I}_0 \to \mathcal{I}$ があり、これは
$\mathcal{I}_0$ の射 $\mathcal{F} \subset \mathcal{F'}$ を、
$\mathcal{F}$ の終対象から $\mathcal{F}'$ の終対象への一意な写像へ送る。
$\mathcal{I}$ の任意の二つの有限生成部分圏が与えられると、
この二圏から生成される圏も有限生成である。$\mathcal{I}$ に関する仮定により、
これはさらに、一意な終対象をもつ $\mathcal{I}$ の有限生成部分圏に含まれる。
したがって $\mathcal{I}_0$ は有向である。

\medskip\noindent
最後に $M_0$ が共終であることを確かめる。$\mathcal{I}$ の任意の対象は、
その対象と恒等射だけからなる部分圏の終対象なので、関手 $M_0$ は対象上全射である。
特に Definition \ref{categories-definition-cofinal} の条件 (1) が満たされる。
$\mathcal{I}$ の対象 $i$、$\mathcal{I}_0$ の対象
$\mathcal{F}_1, \mathcal{F}_2$、および $\mathcal{I}$ の写像
$\varphi_1 : i \to M_0(\mathcal{F}_1)$ と
$\varphi_2 : i \to M_0(\mathcal{F}_2)$ が与えられたとする。
$\mathcal{F}_{12}$ として、$\mathcal{F}_1$, $\mathcal{F}_2$ および射
$\varphi_1, \varphi_2$ を含み、一意な終対象をもつ有限生成圏を取れる。
得られる図式
$$
\xymatrix{
& M_0(\mathcal{F}_{12}) & \\
M_0(\mathcal{F}_{1}) \ar[ru] & & M_0(\mathcal{F}_{2}) \ar[lu] \\
& i \ar[lu] \ar[ru]
}
$$
は、圏 $\mathcal{F}_{12}$ の中にあり、$M_0(\mathcal{F}_{12})$ が
この圏の終対象なので可換である。したがって条件 (2) も満たされ、証明が完了する。
\end{proof}

\begin{remark}
\label{categories-remark-trick-needed}
有限有向集合 $(I, \geq)$ は常に最大対象 $i_\infty$ をもつことに注意されたい。
したがって、このような集合上の系 $(M_i, f_{ii'})$ の余極限は、
余極限が $M_{i_\infty}$ に等しいという意味で常に自明である。
これに対して、有限フィルター付き圏を添字とする余極限は自明とは限らない。
例えば、$\mathcal{I}$ を、対象が一つ $i$ で、$e = e \circ e$ を満たす
非自明な射 $e$ が一つだけある圏とする。図式
$M : \mathcal{I} \to Sets$ の余極限は冪等射 $M(e)$ の像である。
これは、Lemma \ref{categories-lemma-directed-category-system} の証明で
$\mathcal{I}\times \omega$ へ移る仕掛けのようなものが本質的であることを示している。
\end{remark}

\begin{lemma}
\label{categories-lemma-nonempty-limit}
$S : \mathcal{I} \to \textit{Sets}$ が集合の余フィルター付き図式で、
すべての $S_i$ が有限非空ならば、$\lim_i S_i$ は非空である。
言い換えると、有限非空集合からなる有向逆系の極限は非空である。
\end{lemma}

\begin{proof}
二つの主張は Lemma \ref{categories-lemma-directed-category-system} により同値である。
$I$ を有向集合とし、$(S_i)_{i \in I}$ を $I$ 上の有限非空集合の逆系とする。
{\it 部分系} $T$ とは、非空な部分集合 $T_i \subset S_i$ からなる族
$T = (T_i)_{i \in I}$ であって、すべての $i' \geq i$ に対して
$T_{i'}$ が遷移写像 $S_{i'} \to S_i$ により $T_i$ の中へ写るものをいう。
部分系の集合を $\mathcal{T}$ と書き、$\mathcal{T}$ を包含関係で順序づける。
$T_\alpha$, $\alpha \in A$ が $\mathcal{T}$ の元からなる全順序族であるとする。
$T_\alpha = (T_{\alpha, i})_{i \in I}$ と書く。このとき
$T_i = \bigcap_{\alpha \in A} T_{\alpha, i}$ と置くことにより
下界 $T = (T_i)_{i \in I}$ を得る。これは、すべての
$T_{\alpha, i}$ が非空で、$T_\alpha$ が全順序族をなすので、明らかに
$S_i$ の有限非空部分集合である。したがってツォルンの補題を適用し、
$\mathcal{T}$ が極小元をもつことが分かる。

\medskip\noindent
極小元 $T \in \mathcal{T}$ の形を調べよう。まず、写像
$T_{i'} \to T_i$ はすべて全射であることに注意する。
実際、$I$ は有向集合で $T_i$ は有限だから、共通部分
$T'_i = \bigcap_{i' \geq i} \Im(T_{i'} \to T_i)$ は非空である。
したがって $T' = (T'_i)$ は $T$ に含まれる部分系であり、極小性により
$T' = T$ である。最後に、各 $i$ に対して $T_i$ は一元集合であると主張する。
実際、$x \in T_i$ ならば、$i' \geq i$ に対して
$T'_{i'} = (T_{i'} \to T_i)^{-1}(\{x\})$ と置き、
$j \not \geq i$ ならば $T'_j = T_j$ と置ける。
上で部分系 $T$ の遷移写像が全射であることを見たので、これも部分系である。
極小性により $T = T'$ となり、これは確かに $T_i$ が一元集合であることを含意する。
これはすべての $i \in I$ について成り立つ。したがって、ある $x_i \in S_i$ に対して
$T_i = \{x_i\}$ であり、すべての $i' \geq i$ に対し、写像
$S_{i'} \to S_i$ のもとで $x_{i'} \mapsto x_i$ となる。
言い換えると $(x_i) \in \lim S_i$ であり、補題が証明された。
\end{proof}


\section{本質的に定値な系}
\label{categories-section-essentially-constant}

\noindent
$M : \mathcal{I} \to \mathcal{C}$ を圏 $\mathcal{C}$ における図式とする。
添字圏 $\mathcal{I}$ はフィルター付きであると仮定する。この場合、
$\mathcal{C}$ の対象 $X$ を指定する、順に強くなる三つの概念がある。
第一のものは単に
$$
X = \colim_{i \in \mathcal{I}} M_i.
$$
である。このとき $X$ には余射影 $M_i \to X$ が備わる。
さらに強い条件は、$X$ が余極限であり、ある $i \in \mathcal{I}$ と射
$X \to M_i$ が存在して、合成 $X \to M_i \to X$ が $\text{id}_X$ となることを
要請するものである。さらに強い条件が次である。

\begin{definition}
\label{categories-definition-essentially-constant-diagram}
$M : \mathcal{I} \to \mathcal{C}$ を圏 $\mathcal{C}$ における図式とする。
\begin{enumerate}
\item 添字圏 $\mathcal{I}$ はフィルター付きであると仮定し、
$(X, \{M_i \to X\}_i)$ を $M$ の余錐とする。Remark
\ref{categories-remark-cones-and-cocones} を参照されたい。ある $i \in \mathcal{I}$ と
射 $X \to M_i$ が存在して次を満たすとき、$M$ は {\it 値} $X$ をもつ
{\it 本質的に定値} な図式であるという。
\begin{enumerate}
\item $X \to M_i \to X$ は $\text{id}_X$ である。
\item すべての $j$ に対して、ある $k$ と射 $i \to k$, $j \to k$ が存在し、
射 $M_j \to M_k$ は合成 $M_j \to X \to M_i \to M_k$ に等しい。
\end{enumerate}
\item 添字圏 $\mathcal{I}$ は余フィルター付きであると仮定し、
$(X, \{X \to M_i\}_i)$ を $M$ の錐とする。Remark
\ref{categories-remark-cones-and-cocones} を参照されたい。ある $i \in \mathcal{I}$ と
射 $M_i \to X$ が存在して次を満たすとき、$M$ は {\it 値} $X$ をもつ
{\it 本質的に定値} な図式であるという。
\begin{enumerate}
\item $X \to M_i \to X$ は $\text{id}_X$ である。
\item すべての $j$ に対して、ある $k$ と射 $k \to i$, $k \to j$ が存在し、
射 $M_k \to M_j$ は合成 $M_k \to M_i \to X \to M_j$ に等しい。
\end{enumerate}
\end{enumerate}
この定義を用いる際には Lemma
\ref{categories-lemma-essentially-constant-is-limit-colimit} を念頭に置かれたい。
\end{definition}

\noindent
二つのどちらを意味するかは文脈から明らかになる。混同のおそれがある場合には、
第一の場合を $M$ が {\it ind-対象} として本質的に定値であるといい、
第二の場合を pro-対象として本質的に定値であるという。
この用語は Remarks \ref{categories-remark-ind-category} および
\ref{categories-remark-pro-category} でさらに説明する。
実際には「本質的に定値な系」という用語をしばしば用いるが、
厳密には有向集合上の系についてだけ定義される。

\begin{definition}
\label{categories-definition-essentially-constant-system}
$\mathcal{C}$ を圏とする。有向系 $(M_i, f_{ii'})$ が
{\it 本質的に定値な系} であるとは、$M$ を関手 $I \to \mathcal{C}$ とみたとき、
それが本質的に定値な図式を定めることをいう。有向逆系 $(M_i, f_{ii'})$ が
{\it 本質的に定値な逆系} であるとは、$M$ を関手
$I^{opp} \to \mathcal{C}$ とみたとき、それが本質的に定値な逆図式を定めることをいう。
\end{definition}

\noindent
$(M_i, f_{ii'})$ が本質的に定値な系で、射 $f_{ii'}$ がモノ射ならば、
十分大きいすべての $i \leq i'$ に対して射 $f_{ii'}$ は同型である。
一方、写像 $(a, b) \mapsto (a + b, 0)$ による系
$$
\mathbf{Z}^2 \to \mathbf{Z}^2 \to \mathbf{Z}^2 \to \ldots
$$
を考える。この系は値 $\mathbf{Z}$ をもつ本質的に定値な系だが、
すべての遷移写像は核をもつ。

\medskip\noindent
本質的に定値でない系の例を挙げる。
$M = \bigoplus_{n \geq 0} \mathbf{Z}$ とし、$S : M \to M$ をシフト作用素
$(a_0, a_1, \ldots) \mapsto (a_1, a_2, \ldots)$ とする。
この場合、遷移写像が $S$ である系 $M \to M \to M \to \ldots$ の余極限は
$0$ であり、合成 $0 \to M \to 0$ は恒等射であるが、この系は本質的に定値ではない。

\medskip\noindent
次の補題は定義の健全性を確かめるものである。

\begin{lemma}
\label{categories-lemma-essentially-constant-is-limit-colimit}
$M : \mathcal{I} \to \mathcal{C}$ を図式とする。
$\mathcal{I}$ がフィルター付きで、$M$ が ind-対象として本質的に定値ならば、
$X = \colim M_i$ が存在し、$M$ は値 $X$ をもつ本質的に定値な図式である。
$\mathcal{I}$ が余フィルター付きで、$M$ が pro-対象として本質的に定値ならば、
$X = \lim M_i$ が存在し、$M$ は値 $X$ をもつ本質的に定値な図式である。
\end{lemma}

\begin{proof}
省略する。定義に関するよい演習である。
\end{proof}

\begin{remark}
\label{categories-remark-ind-category}
$\mathcal{C}$ を圏とする。$\mathcal{C}$ の {\it ind-対象} からなる大きな圏
$\text{Ind-}\mathcal{C}$ が存在する。すなわち、
$F : \mathcal{I} \to \mathcal{C}$ と $G : \mathcal{J} \to \mathcal{C}$ が
$\mathcal{C}$ におけるフィルター付き図式ならば、
$$
\Mor_{\text{Ind-}\mathcal{C}}(F, G) =
\lim_i \colim_j \Mor_\mathcal{C}(F(i), G(j)).
$$
と定められる。$X$ を $X$ 上の {\it 定値系} へ送る標準的関手
$\mathcal{C} \to \text{Ind-}\mathcal{C}$ がある。これは充満忠実な埋め込みである。
この言葉では、図式 $F$ が本質的に定値であるための必要十分条件は、
$F$ が定値系と同型であることである。この内容が必要になれば補題として定式化し、
ここで証明することにする。
\end{remark}


\begin{remark}
\label{categories-remark-pro-category}
$\mathcal{C}$ を圏とする。$\mathcal{C}$ の {\it pro-対象} からなる大きな圏
$\text{Pro-}\mathcal{C}$ が存在する。すなわち、
$F : \mathcal{I} \to \mathcal{C}$ と $G : \mathcal{J} \to \mathcal{C}$ が
$\mathcal{C}$ における余フィルター付き図式ならば、
$$
\Mor_{\text{Pro-}\mathcal{C}}(F, G) =
\lim_j \colim_i \Mor_\mathcal{C}(F(i), G(j)).
$$
と定められる。$X$ を $X$ 上の {\it 定値系} へ送る標準的関手
$\mathcal{C} \to \text{Pro-}\mathcal{C}$ がある。これは充満忠実な埋め込みである。
この言葉では、図式 $F$ が本質的に定値であるための必要十分条件は、
$F$ が定値系と同型であることである。この内容が必要になれば補題として定式化し、
ここで証明することにする。
\end{remark}

\begin{example}
\label{categories-example-pro-morphism-inverse-systems}
$\mathcal{C}$ を圏とする。$(X_n)$ と $(Y_n)$ を、通常の順序をもつ
$\mathbf{N}$ 上の $\mathcal{C}$ における逆系とする。図示すると
$$
\ldots \to X_3 \to X_2 \to X_1
\quad\text{and}\quad
\ldots \to Y_3 \to Y_2 \to Y_1
$$
である。$a : (X_n) \to (Y_n)$ を $\mathcal{C}$ の pro-対象の射とする。
$a$ は具体的にどのようなものだろうか。各 $n \in \mathbf{N}$ に対して
$m(n)$ と射 $a_n : X_{m(n)} \to Y_n$ が存在するはずである。
これらの射は次の意味で整合しなければならない。すべての $n' \geq n$ に対して、
図式
$$
\xymatrix{
X_{m(n, n')} \ar[rr] \ar[d] & & X_{m(n)} \ar[d]^{a_n} \\
X_{m(n')} \ar[r]^{a_{n'}} & Y_{n'} \ar[r] & Y_n
}
$$
が可換となるような $m(n', n) \geq m(n'), m(n)$ が存在する。
$m(n)$ を $\max_{k, l \leq n}\{m(n, k), m(k, l)\}$ で置き換えると、
$\ldots \geq m(3) \geq m(2) \geq m(1)$ と可換図式
$$
\xymatrix{
\ldots \ar[r] &
X_{m(3)} \ar[d]^{a_3} \ar[r] &
X_{m(2)} \ar[d]^{a_2} \ar[r] &
X_{m(1)} \ar[d]^{a_1} \\
\ldots \ar[r] &
Y_3 \ar[r] &
Y_2 \ar[r] &
Y_1
}
$$
を得る。$m' \geq m$ を満たす増加写像
$m' : \mathbf{N} \to \mathbf{N}$ が与えられ、
$a'_i : X_{m'(i)} \to X_{m(i)} \to Y_i$ と置くと、組
$(m', a')$ は同じ pro-系の射を定める。逆に、上のような二つの組
$(m_1, a_1)$ と $(m_1, a_2)$ が与えられたとき、これらが同じ pro-対象の射を
定めるための必要十分条件は、$a'_1 = a'_2$ となる
$m' \geq m_1, m_2$ を見いだせることである。
\end{example}

\begin{remark}
\label{categories-remark-pro-category-copresheaves}
$\mathcal{C}$ を圏とする。$F : \mathcal{I} \to \mathcal{C}$ と
$G : \mathcal{J} \to \mathcal{C}$ を $\mathcal{C}$ における
余フィルター付き図式とする。関手 $A, B : \mathcal{C} \to \textit{Sets}$ を
$$
A(X) = \colim_i \Mor_\mathcal{C}(F(i), X)
\quad\text{and}\quad
B(X) = \colim_j \Mor_\mathcal{C}(G(j), X)
$$
によって定める。$F$ から $G$ への pro-系の射は、関手の変換
$t : B \to A$ と同じものであると主張する。実際、$t$ が与えられたとき、
$B(G(j))$ における $\text{id}_{G(j)}$ の類へ $t$ を適用して、整合する元の系
$\xi_j \in A(G(j)) = \colim_i \Mor_\mathcal{C}(F(i), G(j))$
を得られる。これはまさに Remark \ref{categories-remark-pro-category} における
$\text{Pro-}\mathcal{C}$ の射の定義である。
$F$ から $G$ への pro-対象の射から変換 $B \to A$ を構成する過程は省略する。
\end{remark}

\begin{lemma}
\label{categories-lemma-image-essentially-constant}
$\mathcal{C}$ を圏とする。$M : \mathcal{I} \to \mathcal{C}$ を、
添字圏 $\mathcal{I}$ がフィルター付き（それぞれ余フィルター付き）である図式とする。
$F : \mathcal{C} \to \mathcal{D}$ を関手とする。$M$ が ind-対象
（それぞれ pro-対象）として本質的に定値ならば、
$F \circ M : \mathcal{I} \to \mathcal{D}$ も同様である。
\end{lemma}

\begin{proof}
$X$ が $M$ の値ならば、定義から直ちに $F(X)$ は $F \circ M$ の値である。
\end{proof}

\begin{lemma}
\label{categories-lemma-characterize-essentially-constant-ind}
$\mathcal{C}$ を圏とする。$M : \mathcal{I} \to \mathcal{C}$ を、
添字圏 $\mathcal{I}$ がフィルター付きである図式とする。次は同値である。
\begin{enumerate}
\item $M$ は本質的に定値な ind-対象である。
\item 余錐 $(X, \{M_i \to X\}_i)$ が存在して、$\mathcal{C}$ の任意の
$W$ に対し写像
$\colim_i \Mor_\mathcal{C}(W, M_i) \to \Mor_\mathcal{C}(W, X)$
は全単射である。
\item $X = \colim_i M_i$ が存在して、$\mathcal{C}$ の任意の $W$ に対し写像
$\colim_i \Mor_\mathcal{C}(W, M_i) \to \Mor_\mathcal{C}(W, X)$
は全単射である。
\item $\mathcal{I}$ のある $i$ と射 $X \to M_i$ が存在して、
$\mathcal{C}$ の任意の $W$ に対し写像
$\Mor_\mathcal{C}(W, X) \to
\colim_{j \in \mathcal{I}} \Mor_\mathcal{C}(W, M_j)$
は全単射である。
\end{enumerate}
(2), (3), (4) の場合、本質的に定値な系の値は $X$ である。
\end{lemma}

\begin{proof}
(3) が (2) を含意することは明らかである。(2) を仮定する。
このとき $\text{id}_X \in \Mor_\mathcal{C}(X, X)$ は、ある $i$ に対する射
$X \to M_i$ から生じる。すなわち $X \to M_i \to X$ は恒等射である。
すると、すべての $W$ に対して二つの写像
$$
\Mor_\mathcal{C}(W, X)
\longrightarrow
\colim_i \Mor_\mathcal{C}(W, M_i)
\longrightarrow
\Mor_\mathcal{C}(W, X)
$$
はいずれも全単射である。ここで第一の写像は上で見いだした射 $X \to M_i$ から誘導され、
合成は恒等写像である。したがって合成
$$
\colim_i \Mor_\mathcal{C}(W, M_i)
\longrightarrow
\Mor_\mathcal{C}(W, X)
\longrightarrow
\colim_i \Mor_\mathcal{C}(W, M_i)
$$
も恒等写像である。$W = M_j$ と置き、余極限における
$\text{id}_{M_j}$ から始めると、十分大きいある $k$ に対して
$M_j \to X \to M_i \to M_k$ が $M_j \to M_k$ に等しいことが分かる。
よって (1) が成り立つ。

\medskip\noindent
(4) を仮定する。$k$ を $\mathcal{I}$ の対象とする。$W = M_k$ と置くと、
一意な射 $M_k \to X$ が存在し、ある $j$ と $\mathcal{I}$ の射
$k \to j$, $i \to j$ が存在して、
$M_k \to X \to M_i \to M_j$ が $M_k \to M_j$ に等しくなることが従う。
一意性により余錐 $(X, \{M_k \to X\})$ が得られる。
このようにして (4) が (2) を含意することが分かる。若干の詳細は省略する。

\medskip\noindent
(1) が他の条件を含意することの証明は省略する。
\end{proof}

\begin{lemma}
\label{categories-lemma-characterize-essentially-constant-pro}
$\mathcal{C}$ を圏とする。$M : \mathcal{I} \to \mathcal{C}$ を、
添字圏 $\mathcal{I}$ が余フィルター付きである図式とする。次は同値である。
\begin{enumerate}
\item $M$ は本質的に定値な pro-対象である。
\item 錐 $(X, \{X \to M_i\})$ が存在して、$\mathcal{C}$ の任意の
$W$ に対し写像
$\colim_{i \in \mathcal{I}^{opp}} \Mor_\mathcal{C}(M_i, W) \to
\Mor_\mathcal{C}(X, W)$ は全単射である。
\item $X = \lim_i M_i$ が存在して、$\mathcal{C}$ の任意の $W$ に対し写像
$\colim_{i \in \mathcal{I}^{opp}} \Mor_\mathcal{C}(M_i, W) \to
\Mor_\mathcal{C}(X, W)$ は全単射である。
\item $\mathcal{I}$ のある $i$ と射 $M_i \to X$ が存在して、
$\mathcal{C}$ の任意の $W$ に対し写像 $\Mor_\mathcal{C}(X, W) \to
\colim_{j \in \mathcal{I}^{opp}} \Mor_\mathcal{C}(M_j, W)$
は全単射である。
\end{enumerate}
(2), (3), (4) の場合、本質的に定値な系の値は $X$ である。
\end{lemma}

\begin{proof}
この補題は Lemma \ref{categories-lemma-characterize-essentially-constant-ind} の双対である。
\end{proof}

\begin{lemma}
\label{categories-lemma-cofinal-essentially-constant}
$\mathcal{C}$ を圏とする。$H : \mathcal{I} \to \mathcal{J}$ を
フィルター付き添字圏の間の関手とする。$H$ が共終ならば、
任意の図式 $M : \mathcal{J} \to \mathcal{C}$ について、それが
本質的に定値であることと $M \circ H$ が本質的に定値であることは同値である。
\end{lemma}

\begin{proof}
これは Lemmas \ref{categories-lemma-characterize-essentially-constant-ind} and
\ref{categories-lemma-cofinal} から形式的に従う。
\end{proof}

\begin{lemma}
\label{categories-lemma-essentially-constant-over-product}
$\mathcal{I}$ と $\mathcal{J}$ をフィルター付き圏とし、射影を
$p : \mathcal{I} \times \mathcal{J} \to \mathcal{J}$ と表す。
このとき $\mathcal{I} \times \mathcal{J}$ はフィルター付きであり、
図式 $M : \mathcal{J} \to \mathcal{C}$ が本質的に定値であることと
$M \circ p : \mathcal{I} \times \mathcal{J} \to \mathcal{C}$ が
本質的に定値であることは同値である。
\end{lemma}

\begin{proof}
$\mathcal{I} \times \mathcal{J}$ がフィルター付きであることの確認は省略する。
この同値は $p$ が共終であること（確認は省略する）と
Lemma \ref{categories-lemma-cofinal-essentially-constant} から従う。
\end{proof}

\begin{lemma}
\label{categories-lemma-initial-essentially-constant}
$\mathcal{C}$ を圏とする。$H : \mathcal{I} \to \mathcal{J}$ を
余フィルター付き添字圏の間の関手とする。$H$ が共始ならば、
任意の図式 $M : \mathcal{J} \to \mathcal{C}$ について、それが
本質的に定値であることと $M \circ H$ が本質的に定値であることは同値である。
\end{lemma}

\begin{proof}
これは Lemmas \ref{categories-lemma-characterize-essentially-constant-pro},
\ref{categories-lemma-initial}, \ref{categories-lemma-cofinal} と、$\mathcal{I}$ が
$\mathcal{J}$ において共始ならば $\mathcal{I}^{opp}$ は
$\mathcal{J}^{opp}$ において共終であるという事実から形式的に従う。
\end{proof}





\section{完全関手}
\label{categories-section-exact-functor}

\noindent
この節では完全関手を定義する。

\begin{definition}
\label{categories-definition-exact}
$F : \mathcal{A} \to \mathcal{B}$ を関手とする。
\begin{enumerate}
\item $\mathcal{A}$ にすべての有限極限が存在すると仮定する。
すべての有限極限と可換であるとき、$F$ を\textit{左完全}という。
\item $\mathcal{A}$ にすべての有限余極限が存在すると仮定する。
すべての有限余極限と可換であるとき、$F$ を\textit{右完全}という。
\item 左完全かつ右完全であるとき、$F$ を\textit{完全}という。
\end{enumerate}
\end{definition}

\begin{lemma}
\label{categories-lemma-characterize-left-exact}
$F : \mathcal{A} \to \mathcal{B}$ を関手とする。
$\mathcal{A}$ にすべての有限極限が存在すると仮定する。
Lemma \ref{categories-lemma-finite-limits-exist} を参照せよ。次は同値である。
\begin{enumerate}
\item $F$ は左完全である。
\item $F$ は有限積および等化子と可換である。
\item $F$ は $\mathcal{A}$ の終対象を $\mathcal{B}$ の終対象へ移し、
ファイバー積と可換である。
\end{enumerate}
\end{lemma}

\begin{proof}
Lemma \ref{categories-lemma-limits-products-equalizers} により (2) は (1) を含意する。
(3) を仮定する。終対象上のファイバー積は積である。
$a, b : A \to B$ を $\mathcal{A}$ の射とすると、$a, b$ の等化子は
$$
(A \times_{a, B, b} A)\times_{(\text{pr}_1, \text{pr}_2), A \times A, \Delta} A.
$$
したがって (3) は (2) を含意する。最後に、空極限は終対象であり、
ファイバー積は極限であるから、(1) は (3) を含意する。
\end{proof}

\begin{lemma}
\label{categories-lemma-characterize-right-exact}
$F : \mathcal{A} \to \mathcal{B}$ を関手とする。
$\mathcal{A}$ にすべての有限余極限が存在すると仮定する。
Lemma \ref{categories-lemma-colimits-exist} を参照せよ。次は同値である。
\begin{enumerate}
\item $F$ は右完全である。
\item $F$ は有限余積および余等化子と可換である。
\item $F$ は $\mathcal{A}$ の始対象を $\mathcal{B}$ の始対象へ移し、
押し出しと可換である。
\end{enumerate}
\end{lemma}

\begin{proof}
Lemma \ref{categories-lemma-characterize-left-exact} の双対である。
\end{proof}





\section{随伴関手}
\label{categories-section-adjoint}

\begin{definition}
\label{categories-definition-adjoint}
$\mathcal{C}$, $\mathcal{D}$ を圏とする。
$u : \mathcal{C} \to \mathcal{D}$ および
$v : \mathcal{D} \to \mathcal{C}$ を関手とする。
$X \in \Ob(\mathcal{C})$ および $Y \in \Ob(\mathcal{D})$ に関して関手的な全単射
$$
\Mor_\mathcal{D}(u(X), Y)
\longrightarrow
\Mor_\mathcal{C}(X, v(Y))
$$
が存在するとき、$u$ を $v$ の\textit{左随伴}、または
$v$ を $u$ の\textit{右随伴}という。
\end{definition}

\noindent
言い換えると、これは $\mathcal{C}^{opp} \times \mathcal{D} \to \textit{Sets}$
という関手の間に、$\Mor_\mathcal{D}(u(-), -)$ から
$\Mor_\mathcal{C}(-, v(-))$ への\textit{指定された}同型があることを意味する。
$\mathcal{C}$ の任意の対象 $X$ に対して、$\text{id}_{u(X)}$ に対応する射
$X \to v(u(X))$ を得る。同様に、$\mathcal{D}$ の任意の対象 $Y$ に対して、
$\text{id}_{v(Y)}$ に対応する射 $u(v(Y)) \to Y$ を得る。
これらの写像を\textit{随伴写像}という。随伴写像は $X$ と $Y$ に関して関手的なので、
関手の射
$$
\eta : \text{id}_\mathcal{C} \to v \circ u\quad (\text{unit})
\quad\text{and}\quad
\epsilon : u \circ v \to \text{id}_\mathcal{D}\quad (\text{counit}).
$$
を得る。さらに、$\alpha : u(X) \to Y$ と $\beta : X \to v(Y)$ を射とすると、
次は同値である。
\begin{enumerate}
\item $\alpha$ と $\beta$ は定義の全単射によって互いに対応する。
\item $\beta$ は合成 $X \to v(u(X)) \xrightarrow{v(\alpha)} v(Y)$ である。
\item $\alpha$ は合成 $u(X) \xrightarrow{u(\beta)} u(v(Y)) \to Y$ である。
\end{enumerate}
このようにして、随伴関手の概念を随伴写像によって言い換えることができる。

\begin{lemma}
\label{categories-lemma-adjoint-exists}
$u : \mathcal{C} \to \mathcal{D}$ を圏の間の関手とする。
各 $y \in \Ob(\mathcal{D})$ に対して関手
$x \mapsto \Mor_\mathcal{D}(u(x), y)$ が表現可能ならば、
$u$ は右随伴をもつ。
\end{lemma}

\begin{proof}
各 $y$ に対して対象 $v(y)$ と関手の同型
$\Mor_\mathcal{C}(-, v(y)) \to \Mor_\mathcal{D}(u(-), y)$ を選ぶ。
米田の補題 (Lemma \ref{categories-lemma-yoneda}) により、任意の射 $g : y \to y'$ に対して、
関手の変換
$$
\Mor_\mathcal{C}(-, v(y)) \to \Mor_\mathcal{D}(u(-), y) \to
\Mor_\mathcal{D}(u(-), y') \to \Mor_\mathcal{C}(-, v(y'))
$$
は一意な射 $v(g) : v(y) \to v(y')$ に対応する。
$v$ が関手であり、$u$ の右随伴であることの確認は省略する。
\end{proof}

\begin{lemma}
\label{categories-lemma-left-adjoint-composed-fully-faithful}
\begin{reference}
Bhargav Bhatt, private communication.
\end{reference}
Definition \ref{categories-definition-adjoint} のように、$u$ を $v$ の左随伴とする。
\begin{enumerate}
\item $v \circ u$ が充満忠実ならば、$u$ は充満忠実である。
\item $u \circ v$ が充満忠実ならば、$v$ は充満忠実である。
\end{enumerate}
\end{lemma}

\begin{proof}
(2) を証明する。$u \circ v$ が充満忠実であると仮定する。
$X$, $Y$ を $\mathcal{D}$ の対象とする。このとき自然な合成写像
$$
\Mor(X,Y) \to \Mor(v(X),v(Y)) \to \Mor(u(v(X)), u(v(Y)))
$$
は全単射なので、少なくとも $v$ は忠実である。充満忠実性を示すには、
上の第二の写像が単射であることを示せばよい。
一方、$u$ と $v$ の随伴性により
$$
\Mor(v(X), v(Y)) \to \Mor(u(v(X)), u(v(Y))) \to \Mor(u(v(X)), Y)
$$
は全単射である。ここで第一の写像は自然な写像であり、第二の写像は随伴の余単位
$u(v(Y)) \to Y$ から生じる。したがって
$\Mor(v(X), v(Y)) \to \Mor(u(v(X)), u(v(Y)))$ も単射であり、これが望む結論である。
(1) の証明はこれと双対である。
\end{proof}

\begin{lemma}
\label{categories-lemma-adjoint-fully-faithful}
Definition \ref{categories-definition-adjoint} のように、$u$ を $v$ の左随伴とする。このとき
\begin{enumerate}
\item $u$ が充満忠実 $\Leftrightarrow$ $\text{id} \cong v \circ u$
$\Leftrightarrow$ $\eta : \text{id} \to v \circ u$ が同型である。
\item $v$ が充満忠実 $\Leftrightarrow$
$u \circ v \cong \text{id}$ $\Leftrightarrow$
$\epsilon : u \circ v \to \text{id}$ が同型である。
\end{enumerate}
\end{lemma}

\begin{proof}
(1) を証明する。$u$ が充満忠実であると仮定する。
$\eta_X : X \to v(u(X))$ が同型であることを示す。
$X' \to v(u(X))$ を任意の射とする。随伴性により、これは射
$u(X') \to u(X)$ に対応する。$u$ の充満忠実性により、これは一意な射
$X' \to X$ に対応する。したがって、$\eta_X$ との後合成は全単射
$\Mor(X', X) \to \Mor(X', v(u(X)))$ を定める。よって $\eta_X$ は同型である。
関手の同型 $\text{id} \cong v \circ u$ が存在するならば、$v \circ u$ は充満忠実である。
Lemma \ref{categories-lemma-left-adjoint-composed-fully-faithful} により、$u$ は充満忠実である。
上で示したことから、これは $\eta$ が同型であることを含意する。
したがってすべての $3$ 条件は同値である（さらに、これらの条件は
$v \circ u$ が充満忠実であることとも同値である）。

\medskip\noindent
(2) は (1) と双対である。
\end{proof}

\begin{lemma}
\label{categories-lemma-adjoint-exact}
Definition \ref{categories-definition-adjoint} のように、$u$ を $v$ の左随伴とする。
\begin{enumerate}
\item $M : \mathcal{I} \to \mathcal{C}$ を図式とし、
$\colim_\mathcal{I} M$ が $\mathcal{C}$ に存在すると仮定する。このとき
$u(\colim_\mathcal{I} M) = \colim_\mathcal{I} u \circ M$ である。
言い換えると、$u$ は（表現可能な）余極限と可換である。
\item $M : \mathcal{I} \to \mathcal{D}$ を図式とし、
$\lim_\mathcal{I} M$ が $\mathcal{D}$ に存在すると仮定する。このとき
$v(\lim_\mathcal{I} M) = \lim_\mathcal{I} v \circ M$ である。
言い換えると、$v$ は表現可能な極限と可換である。
\end{enumerate}
\end{lemma}

\begin{proof}
余極限から対象への射は、極限の各構成要素からその対象への整合する射の系と同じものである。
Remark \ref{categories-remark-limit-colim} を参照せよ。したがって
$$
\begin{matrix}
\Mor_\mathcal{D}(u(\colim_{i \in \mathcal{I}} M_i), Y) &
= & \Mor_\mathcal{C}(\colim_{i \in \mathcal{I}} M_i, v(Y)) \\
& = &
\lim_{i \in \mathcal{I}^{opp}}
\Mor_\mathcal{C}(M_i, v(Y)) \\
& = &
\lim_{i \in \mathcal{I}^{opp}}
\Mor_\mathcal{D}(u(M_i), Y)
\end{matrix}
$$
により、$u(\colim_{i \in \mathcal{I}} M_i)$ は求める余極限である。
もう一方の主張も同様に証明できる。
\end{proof}

\begin{lemma}
\label{categories-lemma-exact-adjoint}
Definition \ref{categories-definition-adjoint} のように、$u$ を $v$ の左随伴とする。
\begin{enumerate}
\item $\mathcal{C}$ が有限余極限をもつならば、$u$ は右完全である。
\item $\mathcal{D}$ が有限極限をもつならば、$v$ は左完全である。
\end{enumerate}
\end{lemma}

\begin{proof}
定義および Lemma \ref{categories-lemma-adjoint-exact} から明らかである。
\end{proof}

\begin{lemma}
\label{categories-lemma-unit-counit-relations}
$u : \mathcal{C} \to \mathcal{D}$ を関手
$v : \mathcal{D} \to \mathcal{C}$ の左随伴とする。
$\eta_X : X \to v(u(X))$ を単位、$\epsilon_Y : u(v(Y)) \to Y$ を余単位とする。このとき
$$
u(X) \xrightarrow{u(\eta_X)} u(v(u(X))
\xrightarrow{\epsilon_{u(X)}} u(X)
\quad\text{and}\quad
v(Y) \xrightarrow{\eta_{v(Y)}} v(u(v(Y))) \xrightarrow{v(\epsilon_Y)}
v(Y)
$$
は恒等射である。
\end{lemma}

\begin{proof}
省略する。
\end{proof}

\begin{lemma}
\label{categories-lemma-transformation-between-functors-and-adjoints}
$u_1, u_2 : \mathcal{C} \to \mathcal{D}$ を、右随伴
$v_1, v_2 : \mathcal{D} \to \mathcal{C}$ をもつ関手とする。
$\beta : u_2 \to u_1$ を関手の変換とし、
$\beta^\vee : v_1 \to v_2$ を対応する随伴関手の変換とする。このとき
$$
\xymatrix{
u_2 \circ v_1 \ar[r]_\beta \ar[d]_{\beta^\vee} &
u_1 \circ v_1 \ar[d] \\
u_2 \circ v_2 \ar[r] & \text{id}
}
$$
は可換である。ここでラベルのない矢印は余単位変換である。
\end{lemma}

\begin{proof}
これは、$\beta^\vee_D : v_1D \to v_2D$ が次の性質をもつ一意な射だからである。
すなわち、誘導される写像 $\Mor(C, v_1D) \to \Mor(C, v_2D)$ は、
$\beta_C : u_2C \to u_1C$ によって誘導される写像
$\Mor(u_1C, D) \to \Mor(u_2C, D)$ である。これは、
$\beta_{v_1 D}$ によって誘導される写像
$$
\Mor(u_1 v_1 D, D') \to \Mor(u_2 v_1 D, D')
$$
が、$\beta^\vee_{D'}$ によって誘導される写像
$$
\Mor(v_1 D, v_1 D') \to \Mor(v_1 D, v_2 D')
$$
と同じであることを意味する。$D' = D$ と置くと、余単位
$u_1 v_1 D \to D$ と $\beta_{v_1D}$ の前合成は、随伴によって
$\beta^\vee_D$ に対応することが分かる。これはまさに、$D$ 上で評価したときに
図式が可換であることを意味する。
\end{proof}

\begin{lemma}
\label{categories-lemma-compose-counits}
$\mathcal{A}$, $\mathcal{B}$, $\mathcal{C}$ を圏とする。
$v : \mathcal{A} \to \mathcal{B}$ および
$v' : \mathcal{B} \to \mathcal{C}$ を、それぞれ左随伴 $u$ および $u'$ をもつ関手とする。
このとき
\begin{enumerate}
\item 関手 $v'' = v' \circ v$ は $u'' = u \circ u'$ に等しい左随伴をもつ。
\item $\mathcal{A}$ の $X$ に対して
\begin{equation}
\label{categories-equation-compose-counits}
\epsilon_X^v \circ u(\epsilon^{v'}_{v(X)}) = \epsilon^{v''}_X :
u''(v''(X)) \to X
\end{equation}
が成り立つ。ここで $\epsilon$ は各随伴の余単位である。
\end{enumerate}
\end{lemma}

\begin{proof}
(2) の公式を展開する。これにより (1) も直ちに証明される。
まず、組 $(u, v)$ および $(u', v')$ の随伴の余単位は写像
$\epsilon_X^v : u(v(X)) \to X$ および
$\epsilon_Y^{v'} : u'(v'(Y)) \to Y$ である。
Definition \ref{categories-definition-adjoint} に続く議論を参照せよ。
(1) の $u''$ と $v''$ を用いてすべてを展開すると
$$
u''(v''(X)) = u(u'(v'(v(X)))) \xrightarrow{u(\epsilon_{v(X)}^{v'})}
u(v(X)) \xrightarrow{\epsilon_X^v} X
$$
となり、(\ref{categories-equation-compose-counits}) の左辺の写像を得る。
当面これを $\epsilon_X^{v''}$ と表す。
これが随伴対 $(u'', v'')$ の余単位であることを示すには、$\mathcal{C}$ の
$Z$ が与えられたとき、射 $\beta : Z \to v''(X)$ を
$\alpha = \epsilon_X^{v''} \circ u''(\beta) : u''(Z) \to X$ へ送る規則が
射の集合の間の全単射であることを示せばよい。
実際、この規則は、$\beta$ を $\epsilon_{v(X)}^{v'} \circ u'(\beta)$ へ送る規則と、
次にこれを $\epsilon_X^v \circ u(\epsilon_{v(X)}^{v'} \circ u'(\beta))$ へ送る規則との
合成である。前者は $(u', v')$ についての仮定により全単射であり、
後者は $(u, v)$ についての仮定により全単射である。
\end{proof}





\section{表現可能性の判定条件}
\label{categories-section-representable}

\noindent
次の補題は、大きな圏における普遍対象の存在を証明する際にしばしば有用である。
Remark \ref{categories-remark-how-to-use-it} の議論を参照せよ。

\begin{lemma}
\label{categories-lemma-a-version-of-brown}
$\mathcal{C}$ を極限をもつ大きな\footnote{Remark \ref{categories-remark-big-categories} を参照せよ。}
圏とする。$F : \mathcal{C} \to \textit{Sets}$ を関手とし、次を仮定する。
\begin{enumerate}
\item $F$ は極限と可換である。
\item $\mathcal{C}$ の対象の族 $\{x_i\}_{i \in I}$ と、各 $i \in I$ に対する
元 $f_i \in F(x_i)$ が存在して、$y \in \Ob(\mathcal{C})$ および
$g \in F(y)$ に対し、$F(\varphi)(f_i) = g$ を満たす $i$ と射
$\varphi : x_i \to y$ が存在する。
\end{enumerate}
このとき $F$ は表現可能である。すなわち、$y$ に関して関手的に
$$
F(y) = \Mor_\mathcal{C}(x, y)
$$
となる $\mathcal{C}$ の対象 $x$ が存在する。
\end{lemma}

\begin{proof}
$\mathcal{I}$ を、対象が組 $(x_i, f_i)$ であり、射
$(x_i, f_i) \to (x_{i'}, f_{i'})$ が
$F(\varphi)(f_i) = f_{i'}$ を満たす $\mathcal{C}$ の写像
$\varphi : x_i \to x_{i'}$ である圏とする。
$$
x = \lim_{(x_i, f_i) \in \mathcal{I}} x_i
$$
と置く（これは求める $x$ ではない。以下を参照せよ）。
仮定によりこの極限は存在する。$F$ は極限と可換なので
$$
F(x) = \lim_{(x_i, f_i) \in \mathcal{I}} F(x_i).
$$
したがって、射影写像 $x \to x_i$ に $F$ を適用した写像の下で
$f_i \in F(x_i)$ へ移る普遍元 $f \in F(x)$ が存在する。
$f$ を用いて関手の変換
$$
\xi : \Mor_\mathcal{C}(x, - ) \longrightarrow F(-)
$$
を得る。Section \ref{categories-section-opposite} を参照せよ。
$y$ を $\mathcal{C}$ の任意の対象、$g \in F(y)$ とする。
仮定により可能なので、$f_i$ を $g$ へ移す $x_i \to y$ を選ぶ。
すると、写像
$$
x \longrightarrow x_i \longrightarrow y
$$
（第一の写像は $x$ を定める極限の射影写像）に $F$ を適用すると、
$f$ は $g$ へ移る。したがって変換 $\xi$ は全射である。

\medskip\noindent
$F$ を表現する対象を得るため、$F(\varphi)(f) = f$ を満たすすべての自己写像
$\varphi : x \to x$ の等化子を $e : x' \to x$ とする。
$F$ は極限と可換するので等化子とも可換し、$F(x)$ において $f$ へ移る
$f' \in F(x')$ が存在することが分かる。$\xi$ は全射であり、$f'$ は $f$ へ移るので、
$\xi' : \Mor_\mathcal{C}(x', -) \to F(-)$ も全射である。
最後に、$a, b : x' \to y$ を $F(a)(f') = F(b)(f')$ を満たす二つの写像とする。
$a = b$ を示さなければならない。等化子 $e' : x'' \to x'$ を考える。
再び、$f'$ へ移る $f'' \in F(x'')$ を得る。
$F(\psi)(f) = f''$ を満たす写像 $\psi : x \to x''$ を選ぶ。
すると $e \circ e' \circ \psi : x \to x$ は
$F(e \circ e' \circ \psi)(f) = f$ を満たす射である。したがって
$e \circ e' \circ \psi \circ e = e$ である。$e$ は単射なので、
$e'$ は全射であり、望むように $a = b$ となる。
\end{proof}

\begin{remark}
\label{categories-remark-how-to-use-it}
上の補題は、何かの上の自由な何かを構成するためにしばしば用いられる。
例えば、集合上の自由アーベル群、集合上の自由群などである。
集合 $E$ 上の自由群の場合には、関手
$$
F : \textit{Groups} \to \textit{Sets},\quad
G \longmapsto \text{Map}(E, G)
$$
を考える。この関手は極限と可換である。対象の族として、濃度が高々
$\max(\aleph_0, |E|)$ である群 $G_i$ と集合写像 $E \to G_i$ からなる族
$E \to G_i$ で、このような構造の各同型類が少なくとも一度現れるものを取れる。
実際、$E$ から群 $G$ への写像 $E \to G$ が与えられたとき、その像が生成する
部分群 $G'$ の濃度は高々 $\max(\aleph_0, |E|)$ である。
補題によりこの関手は表現可能なので、
$\Mor_{\textit{Groups}}(F_E, G) = \text{Map}(E, G)$ を満たす群 $F_E$ が存在する。
特に、$F_E$ の恒等射は写像 $E \to F_E$ に対応し、関係式を一切課さずに
$F_E$ がその像によって生成されることを示せる。

\medskip\noindent
もう一つの典型的な応用として、極限の存在が分かっているとき、補題を用いて
余極限を構成できる。Topology, Lemma \ref{topology-lemma-limits} により極限をもつ
位相空間の圏を用いて説明する。$\mathcal{I} \to \textit{Top}$,
$i \mapsto X_i$ を関手とする。このとき
$$
F : \textit{Top} \longrightarrow \textit{Sets},\quad
Y \longmapsto \lim_\mathcal{I} \Mor_{\textit{Top}}(X_i, Y)
$$
を考えられる。この関手は極限と可換である。さらに、任意の位相空間 $Y$ と
$F(Y)$ の元 $(\varphi_i : X_i \to Y)$ が与えられたとき、各射 $\varphi_i$ が
$Y'$ の中へ写る、濃度が高々 $|\coprod X_i|$ の部分空間 $Y' \subset Y$ が存在する。
実際、各 $\varphi_i$ の像の合併に誘導位相を入れればよい。
したがって補題の仮定が満たされることは明らかであり、関手 $F$ を表現する
位相空間 $X$ を得る。これはまさに $X$ が図式 $i \mapsto X_i$ の
余極限であることを意味する。
\end{remark}

\begin{theorem}[随伴関手定理]
\label{categories-theorem-adjoint-functor}
$G : \mathcal{C} \to \mathcal{D}$ を大きな圏の間の関手とする。
$\mathcal{C}$ は極限をもち、$G$ はそれらと可換すると仮定する。
また、$\mathcal{D}$ のすべての対象 $y$ に対して、
$x_i \in \Ob(\mathcal{C})$, $f_i \in \Mor_\mathcal{D}(y, G(x_i))$ を満たす
組の集合 $(x_i, f_i)_{i \in I}$ が存在し、
$x \in \Ob(\mathcal{C})$, $f \in \Mor_\mathcal{D}(y, G(x))$ を満たす任意の組
$(x, f)$ に対して、$f = G(h) \circ f_i$ を満たす $i$ と射
$h : x_i \to x$ が存在すると仮定する。このとき $G$ は左随伴 $F$ をもつ。
\end{theorem}

\begin{proof}
仮定により、$\mathcal{D}$ の各対象 $y$ に対して関手
$x \mapsto \Mor_\mathcal{D}(y, G(x))$ は
Lemma \ref{categories-lemma-a-version-of-brown} の仮定を満たす。
したがってこれは、$F(y)$ と呼ぶことにする対象によって表現される。
米田の補題 (Lemma \ref{categories-lemma-yoneda}) を適用すると、規則
$y \mapsto F(y)$ は関手となり、構成により $G$ の随伴である。詳細は省略する。
\end{proof}






\section{圏論的コンパクト対象}
\label{categories-section-compact}

\noindent
圏の「小さい」対象について少し述べる。

\begin{definition}
\label{categories-definition-compact-object}
$\mathcal{C}$ を大きな\footnote{Remark \ref{categories-remark-big-categories} を参照せよ。}
圏とする。$\mathcal{C}$ の対象 $X$ が\textit{圏論的コンパクト}であるとは、
$\colim_i M_i$ が存在するすべてのフィルター付き図式
$M : \mathcal{I} \to \mathcal{C}$ に対して
$$
\Mor_\mathcal{C}(X, \colim_i M_i) =
\colim_i \Mor_\mathcal{C}(X, M_i)
$$
が成り立つことをいう。
\end{definition}

\noindent
この定義は、$\mathcal{C}$ がすべてのフィルター付き余極限をもつという仮定の下でのみ
なされることも多い。

\begin{lemma}
\label{categories-lemma-extend-functor-by-colim}
$\mathcal{C}$ と $\mathcal{D}$ をフィルター付き余極限をもつ大きな圏とする。
$\mathcal{C}' \subset \mathcal{C}$ を $\mathcal{C}$ の圏論的コンパクト対象からなる
小さい充満部分圏とし、$\mathcal{C}$ のすべての対象が $\mathcal{C}'$ の対象の
フィルター付き余極限であると仮定する。このとき、すべての関手
$F' : \mathcal{C}' \to \mathcal{D}$ は、フィルター付き余極限と可換する一意な延長
$F : \mathcal{C} \to \mathcal{D}$ をもつ。
\end{lemma}

\begin{proof}
$\mathcal{C}$ の各対象 $X$ に対して、$X_i \in \Ob(\mathcal{C}')$ を用いた
フィルター付き余極限 $X = \colim X_i$ として $X$ を書ける。このとき
$$
F(X) = \colim F'(X_i)
$$
と $\mathcal{D}$ において置く。以下で、この構成が $X$ の余極限表示の選択に依存しないことを示す。

\medskip\noindent
$\mathcal{C}$ の射 $\alpha : X \to Y$ が与えられ、
$X = \colim_{i \in I} X_i$ および $Y = \colim_{j \in J} Y_i$ が
$\mathcal{C}'$ の対象のフィルター付き余極限として書かれているとする。
各 $i \in I$ に対して、$X_i$ は $\mathcal{C}$ の圏論的コンパクト対象なので、
$j \in J$ と可換図式
$$
\xymatrix{
X_i \ar[r] \ar[d] & X \ar[d]^\alpha \\
Y_j \ar[r] & Y
}
$$
を見いだせる。このとき射 $F'(X_i) \to F'(Y_j) \to F(Y)$ を得る。
第二の射は $F(Y) = \colim F'(Y_j)$ への余射影である。
矢印 $\beta_i : F'(X_i) \to F(Y)$ は $j$ の選択に依存しない。
$i \leq i'$ に対して合成
$$
F'(X_i) \to F'(X_{i'}) \xrightarrow{\beta_{i'}} F(Y)
$$
は $\beta_i$ に等しい。したがって余極限の普遍性により、良定義な矢印
$$
F(\alpha) : F(X) = \colim F(X_i) \to F(Y)
$$
を得る。$\alpha' : Y \to Z$ を $\mathcal{C}$ の第二の射とし、
$Z = \colim Z_k$ も $\mathcal{C}'$ の対象のフィルター付き余極限として書かれているとする。
このとき、誘導される射 $F(\alpha) : F(X) \to F(Y)$ と
$F(\alpha') : F(Y) \to F(Z)$ の合成が射
$F(\alpha' \circ \alpha)$ に等しいことを示すのはよい演習である。詳細は省略する。

\medskip\noindent
特に、$\mathcal{C}'$ の系のフィルター付き余極限として二つの表示
$X = \colim X_i$ と $X = \colim X'_{i'}$ が与えられると、互いに逆な矢印
$\colim F'(X_i) \to \colim F'(X'_{i'})$ および
$\colim F'(X'_{i'}) \to \colim F'(X_i)$ を得る。
言い換えると、値 $F(X)$ は、$X$ を $\mathcal{C}'$ の対象のフィルター付き余極限として
表示する方法に依存せず良定義である。前段落で論じた $F$ の関手性と合わせると、
$F$ は関手である。また、$X \in \Ob(\mathcal{C}')$ ならば
$F(X) = F'(X)$ であることも明らかである。

\medskip\noindent
補題の一意性の主張は、$F$ がフィルター付き余極限と可換することを示せば明らかである
（そうでなければこの主張は意味をなさない）。これを示すため、
$X = \colim_{\lambda \in \Lambda} X_\lambda$ が $\mathcal{C}$ の
フィルター付き余極限であると仮定する。$F$ は関手なので、写像
$$
\colim_\lambda F(X_\lambda) \longrightarrow F(X)
$$
を得る。一方、$X = \colim X_i$ を $\mathcal{C}'$ の対象の
フィルター付き余極限として書く。上と同様に、各 $i \in I$ に対して
$\lambda \in \Lambda$ と可換図式
$$
\xymatrix{
X_i \ar[rr] \ar[rd] & & X_\lambda \ar[ld] \\
& X
}
$$
を選べる。上と同様に、これは遷移射と整合する良定義な射
$F'(X_i) \to \colim_\lambda F(X_\lambda)$、したがって射
$$
F(X) = \colim_i F'(X_i) \longrightarrow \colim_\lambda F(X_\lambda)
$$
を定める。この射は上の射の逆であり（詳細は省略する）、望むように
$F(X) = \colim_\lambda F(X_\lambda)$ であることを示す。
\end{proof}







\section{圏の局所化}
\label{categories-section-localization}

\noindent
この節の基本的な考えは、圏 $\mathcal{C}$ と矢印の集合 $S$ が与えられたとき、
関手
$F : \mathcal{C} \to S^{-1}\mathcal{C}$
であって、$S$ のすべての元が $S^{-1}\mathcal{C}$ において可逆となり、
かつこの性質をもつすべての関手の中で $F$ が普遍的となるものを構成することである。
この節の参考文献は \cite[Chapter I, Section 2]{GZ} および
\cite[Chapter II, Section 2]{Verdier} である。

\begin{definition}
\label{categories-definition-multiplicative-system}
$\mathcal{C}$ を圏とする。$\mathcal{C}$ の矢印の集合 $S$ が
\textit{左乗法系}であるとは、次の性質をもつことをいう。
\begin{enumerate}
\item[LMS1] $\mathcal{C}$ の各対象の恒等射は $S$ に属し、合成可能な
$S$ の二元の合成も $S$ に属する。
\item[LMS2] $t \in S$ をもつ任意の実線図式
$$
\xymatrix{
X \ar[d]_t \ar[r]_g & Y \ar@{..>}[d]^s \\
Z \ar@{..>}[r]^f & W
}
$$
は、$s \in S$ である可換な点線の正方形へ拡張できる。
\item[LMS3] 射の任意の組 $f, g : X \to Y$ と、終点が $X$ である
$t \in S$ が $f \circ t = g \circ t$ を満たすとき、始点が $Y$ である
$s \in S$ で $s \circ f = s \circ g$ を満たすものが存在する。
\end{enumerate}
$\mathcal{C}$ の矢印の集合 $S$ が\textit{右乗法系}であるとは、
次の性質をもつことをいう。
\begin{enumerate}
\item[RMS1] $\mathcal{C}$ の各対象の恒等射は $S$ に属し、合成可能な
$S$ の二元の合成も $S$ に属する。
\item[RMS2] $s \in S$ をもつ任意の実線図式
$$
\xymatrix{
X \ar@{..>}[d]_t \ar@{..>}[r]_g & Y \ar[d]^s \\
Z \ar[r]^f & W
}
$$
は、$t \in S$ である可換な点線の正方形へ拡張できる。
\item[RMS3] 射の任意の組 $f, g : X \to Y$ と、始点が $Y$ である
$s \in S$ が $s \circ f = s \circ g$ を満たすとき、終点が $X$ である
$t \in S$ で $f \circ t = g \circ t$ を満たすものが存在する。
\end{enumerate}
$\mathcal{C}$ の矢印の集合 $S$ が\textit{乗法系}であるとは、
左乗法系かつ右乗法系であることをいう。言い換えると、これは MS1, MS2, MS3 が
成り立つことを意味する。ここで MS1 $=$ LMS1 $+$ RMS1,
MS2 $=$ LMS2 $+$ RMS2, MS3 $=$ LMS3 $+$ RMS3 である。
（もちろん LMS1 $=$ RMS1 $=$ MS1 である。）
\end{definition}

\noindent
これらの条件は、次のように圏 $S^{-1}\mathcal{C}$ を構成するために有用である。

\medskip\noindent
{\bf 左分数計算。}
$\mathcal{C}$ を圏、$S$ を左乗法系とする。新しい圏
$S^{-1}\mathcal{C}$ を次のように定義する（これが機能することは
Lemma \ref{categories-lemma-left-localization} の証明で確認する）。
\begin{enumerate}
\item $\Ob(S^{-1}\mathcal{C}) = \Ob(\mathcal{C})$ と置く。
\item $S^{-1}\mathcal{C}$ の射 $X \to Y$ は、同値関係を除いて
$s \in S$ を満たす組 $(f : X \to Y', s : Y \to Y')$ によって与えられる。
（同値関係は以下で定義する。組 $(f, s)$ の同値類を
$s^{-1}f : X \to Y$ と考えよ。）
\item 二つの組 $(f_1 : X \to Y_1, s_1 : Y \to Y_1)$ および
$(f_2 : X \to Y_2, s_2 : Y \to Y_2)$ が同値であるとは、第三の組
$(f_3 : X \to Y_3, s_3 : Y \to Y_3)$ と、可換図式
$$
\xymatrix{
 & Y_1 \ar[d]^u & \\
X \ar[ru]^{f_1} \ar[r]^{f_3} \ar[rd]_{f_2} &
Y_3 &
Y \ar[lu]_{s_1} \ar[l]_{s_3} \ar[ld]^{s_2} \\
& Y_2 \ar[u]_v &
}
$$
に入る $\mathcal{C}$ の射 $u : Y_1 \to Y_3$ および
$v : Y_2 \to Y_3$ が存在することをいう。
\item 組 $(f : X \to Y', s : Y \to Y')$ と
$(g : Y \to Z', t : Z \to Z')$ の同値類の合成を、組
$(h \circ f : X \to Z'', u \circ t : Z \to Z'')$ の同値類として定義する。
ここで $h$ および $u \in S$ は、仮定により存在する可換図式
$$
\xymatrix{
Y \ar[d]_s \ar[r]_g & Z' \ar[d]^u \\
Y' \ar[r]^h & Z''
}
$$
に入るように選ぶ。
\item $S^{-1} \mathcal{C}$ における恒等射 $X \to X$ は、組
$(\text{id} : X \to X, \text{id} : X \to X)$ の同値類である。
\end{enumerate}

\begin{lemma}
\label{categories-lemma-left-localization}
$\mathcal{C}$ を圏、$S$ を左乗法系とする。
\begin{enumerate}
\item 上で定義した組の間の関係は同値関係である。
\item 上で与えた合成規則は同値類上で良定義である。
\item 合成は結合的であり（恒等射は恒等律を満たす）、したがって
$S^{-1}\mathcal{C}$ は圏である。
\end{enumerate}
\end{lemma}

\begin{proof}
(1) を証明する。二つの組
$p_1 = (f_1 : X \to Y_1, s_1 : Y \to Y_1)$ および
$p_2 = (f_2 : X \to Y_2, s_2 : Y \to Y_2)$ が\textit{初等同値}であるとは、
$a \circ f_1 = f_2$ および $a \circ s_1 = s_2$ を満たす
$\mathcal{C}$ の射 $a : Y_1 \to Y_2$ が存在することをいう。図式で書けば
$$
\xymatrix{
X \ar@{=}[d] \ar[r]_{f_1} & Y_1 \ar[d]^a & Y \ar[l]^{s_1} \ar@{=}[d] \\
X \ar[r]^{f_2} & Y_2 & Y \ar[l]_{s_2}
}
$$
である。この性質を $p_1Ep_2$ と表すことにする。
$pEp$ および $aEb, bEc \Rightarrow aEc$ であることに注意する。
（名称とは異なり、$E$ は同値関係ではない。）
(1) の主張は、関係
$p \sim p' \Leftrightarrow \exists q: pEq \wedge p'Eq$
（ここで $q$ は $p$ および $p'$ と同じ条件を満たす組とする）が
同値関係であるということである。上の $E$ の性質を用いる簡単な形式的議論により、
$p_3Ep_1, p_3Ep_2 \Rightarrow p_1 \sim p_2$ を証明すれば十分である。
そこで、$s_i \in S$ をもち、可換な図式
$$
\xymatrix{
 & Y_1 & \\
X \ar[ru]^{f_1} \ar[r]^{f_3} \ar[rd]_{f_2} &
Y_3 \ar[u]_{a_{31}} \ar[d]^{a_{32}} &
Y \ar[lu]_{s_1} \ar[l]_{s_3} \ar[ld]^{s_2} \\
& Y_2 &
}
$$
が与えられているとする。まず LMS2 を適用して、$a_{24} \in S$ をもち可換な図式
$$
\xymatrix{
Y \ar[d]_{s_1} \ar[r]_{s_2} & Y_2 \ar@{..>}[d]^{a_{24}} \\
Y_1 \ar@{..>}[r]^{a_{14}} & Y_4
}
$$
を得る。このとき
$$
a_{14} \circ a_{31} \circ s_3 =
a_{14} \circ s_1 =
a_{24} \circ s_2 =
a_{24} \circ a_{32} \circ s_3.
$$
したがって LMS3 により、$s_{44} \in S$ であり、
$s_{44} \circ a_{14} \circ a_{31}
= s_{44} \circ a_{24} \circ a_{32}$ を満たす射
$s_{44} : Y_4 \to Y'_4$ が存在する。
よって $Y_4$, $a_{14}$, $a_{24}$ をそれぞれ
$Y'_4$, $s_{44} \circ a_{14}$, $s_{44} \circ a_{24}$ で置き換えると、
$a_{14} \circ a_{31} = a_{24} \circ a_{32}$ と仮定できる
（なお $a_{24} \in S$ および $a_{14} \circ s_1 = a_{24} \circ s_2$ は成り立つ）。
$$
f_4 =
a_{14} \circ f_1 =
a_{14} \circ a_{31} \circ f_3 =
a_{24} \circ a_{32} \circ f_3 =
a_{24} \circ f_2
$$
および $s_4 = a_{14} \circ s_1 = a_{24} \circ s_2$ と置く。このとき図式
$$
\xymatrix{
X \ar@{=}[d] \ar[r]_{f_1} & Y_1 \ar[d]^{a_{14}} & Y \ar[l]^{s_1} \ar@{=}[d] \\
X \ar[r]^{f_4} & Y_4 & Y \ar[l]_{s_4}
}
$$
は可換であり、（LMS1 により）$s_4 \in S$ である。
したがって、$p_4 = (f_4, s_4)$ と置けば $p_1 E p_4$ である。
同様に $p_2 E p_4$ である。これらを合わせると $p_1 \sim p_2$ を得る。

\medskip\noindent
(2) を証明する。$p = (f : X \to Y', s : Y \to Y')$ および
$q = (g : Y \to Z', t : Z \to Z')$ を上の合成の定義に現れる組とする。
合成するため、$u_2 \in S$ をもつ図式
$$
\xymatrix{
Y \ar[d]_s \ar[r]_g & Z' \ar[d]^{u_2} \\
Y' \ar[r]^{h_2} & Z_2
}
$$
を選ぶ。まず、組
$r_2 = (h_2 \circ f : X \to Z_2, u_2 \circ t : Z \to Z_2)$ の同値類が
$(Z_2, h_2, u_2)$ の選択に依存しないことを示す。
実際、$(Z_3, h_3, u_3)$ をもう一つの選択とし、対応する合成を
$r_3 = (h_3 \circ f : X \to Z_3, u_3 \circ t : Z \to Z_3)$ とする。
LMS2 により、$u_{34} \in S$ をもつ図式
$$
\xymatrix{
Z' \ar[d]_{u_2} \ar[r]_{u_3} & Z_3 \ar[d]^{u_{34}} \\
Z_2 \ar[r]^{h_{24}} & Z_4
}
$$
を選べる。$h_2 \circ s = u_2 \circ g$ であり、同様に
$h_3 \circ s = u_3 \circ g$ である。ここで
$$
u_{34} \circ h_3 \circ s
= u_{34} \circ u_3 \circ g
= h_{24} \circ u_2 \circ g
= h_{24} \circ h_2 \circ s.
$$
したがって LMS3 により、
$s_{44} \circ u_{34} \circ h_3 = s_{44} \circ h_{24} \circ h_2$ を満たす
$Z'_4$ および $s_{44} : Z_4 \to Z'_4$ が存在する。
$Z_4$, $h_{24}$, $u_{34}$ をそれぞれ $Z'_4$,
$s_{44} \circ h_{24}$, $s_{44} \circ u_{34}$ で置き換えると、
$u_{34} \circ h_3 = h_{24} \circ h_2$ と仮定できる。
一方、関係 $u_{34} \circ u_3 = h_{24} \circ u_2$ および
$u_{34} \in S$ は引き続き成り立つ。そこで
$h_4 = u_{34} \circ h_3 = h_{24} \circ h_2$ および
$u_4 = u_{34} \circ u_3 = h_{24} \circ u_2$ と置ける。このとき可換図式
$$
\xymatrix{
X \ar@{=}[d] \ar[r]_{h_2\circ f} &
Z_2 \ar[d]^{h_{24}} &
Z \ar[l]^{u_2 \circ t} \ar@{=}[d] \\
X \ar@{=}[d] \ar[r]^{h_4\circ f} &
Z_4 &
Z \ar@{=}[d] \ar[l]_{u_4 \circ t} \\
X \ar[r]^{h_3 \circ f} &
Z_3 \ar[u]^{u_{34}} &
Z \ar[l]_{u_3 \circ t}
}
$$
を得る。したがって組
$r_4 =
(h_4 \circ f : X \to Z_4, u_4 \circ t : Z \to Z_4)$ を得て、上の図式は
$r_2Er_4$ および $r_3Er_4$ を示す。よって望むように $r_2 \sim r_3$ である。
したがって、上のように得られるすべての組 $r$ の同値類として
$p \circ q$ を定義することができる。

\medskip\noindent
(2) の証明を完了するには、組 $p_1, p_2, q$ が $p_1Ep_2$ を満たすとき、
合成が意味をもつ限り $p_1 \circ q = p_2 \circ q$ および
$q \circ p_1 = q \circ p_2$ であることを示さなければならない。
このため、$p_1 = (f_1 : X \to Y_1, s_1 : Y \to Y_1)$ および
$p_2 = (f_2 : X \to Y_2, s_2 : Y \to Y_2)$ と書き、
$f_2 = a \circ f_1$ および $s_2 = a \circ s_1$ を満たす
$\mathcal{C}$ の射 $a : Y_1 \to Y_2$ を取る。
まず $q = (g : Y \to Z', t : Z \to Z')$ と仮定する。
この場合、左側のような可換図式
$$
\vcenter{
\xymatrix{
Y \ar[d]_{s_2} \ar[r]^g & Z' \ar[d]^u \\
Y_2 \ar[r]^h & Z''
}
}
\quad
\Rightarrow
\quad
\vcenter{
\xymatrix{
Y \ar[d]_{s_1} \ar[r]^g & Z' \ar[d]^u \\
Y_1 \ar[r]^{h \circ a} & Z''
}
}
$$
（$u \in S$）を選ぶ。すると右側の図式も可換となる。
これらの図式を用いると、二つの合成 $q \circ p_1$ および
$q \circ p_2$ はいずれも
$(h \circ a \circ f_1 : X \to Z'', u \circ t : Z \to Z'')$ の同値類である。
したがって $q \circ p_1 = q \circ p_2$ である。
もう一方の場合、すなわち $p_1 \circ q = p_2 \circ q$ を示す証明は省略する。
（これは上で扱った場合と同様である。）

\medskip\noindent
(3) を証明する。合成の結合律を証明しなければならない。実線図式
$$
\xymatrix{
& & & Z \ar[d] \\
& & Y \ar[d] \ar[r] & Z' \ar@{..>}[d] \\
& X \ar[d] \ar[r] & Y' \ar@{..>}[d] \ar@{..>}[r] & Z'' \ar@{..>}[d] \\
W \ar[r] & X' \ar@{..>}[r] & Y'' \ar@{..>}[r] & Z'''
}
$$
を考える（その垂直矢印は $S$ に属する）。これは三つの合成可能な組を与える。
LMS2 を用いると、各正方形を可換にし、かつ垂直矢印が $S$ に属するように
点線の矢印を選べる。このとき三つの組の合成が、最下段の水平矢印と
右端列の垂直矢印をそれぞれ合成して得られる組
$(W \to Z''', Z \to Z''')$ の同値類であることは明らかである。

\medskip\noindent
恒等律の確認は読者に委ねる。
\end{proof}

\begin{remark}
\label{categories-remark-motivation-localization}
$S^{-1} \mathcal{C}$ の構成の動機は、$S$ の射に逆射を人工的に作ることで、
それらを可逆に「強制」することにある（その代償として、既存のいくつかの射が
互いに同一視されることがある）。これは可換環を乗法的部分集合で局所化すること、
さらに一般には非可換環を右分母集合で局所化することと類似している
（\cite[Section 10A]{Lam} を参照せよ）。これは単なる類似以上のものである。
$S^{-1} \mathcal{C}$ の構成（より正確には、加法圏 $\mathcal{C}$ に対する版）は、
実際に後者の局所化を一般化する。すなわち、非可換環は一つの対象をもつ前加法圏
（射が環の元である）とみなせる。この環の乗法的部分集合は LMS1（別名 RMS1）を
満たす射の集合 $S$ となる。このとき、この圏と部分集合 $S$ に対する条件 RMS2 と
RMS3 は、右分母集合の二条件（「右可換」と「右可逆」）に移り
（LMS と左分母集合についても同様）、$S^{-1} \mathcal{C}$ は
（加法構造を適切に定義すれば）環の局所化に対応する一対象圏となる。
\end{remark}

\begin{definition}
\label{categories-definition-left-localization-as-fraction}
$\mathcal{C}$ を圏、$S$ を $\mathcal{C}$ の射の左乗法系とする。
$\mathcal{C}$ の任意の射 $f : X \to Y'$ と、$S$ の任意の射
$s : Y \to Y'$ が与えられたとき、組
$(f : X \to Y', s : Y \to Y')$ の同値類を\textit{$s^{-1} f$} と表す。
これは $S^{-1} \mathcal{C}$ における $X$ から $Y$ への射である。
\end{definition}

\noindent
この記法が示唆する性質は実際に成り立つ。
$\mathcal{C}$ の任意の射 $f : X \to Y'$ と、$S$ の任意の二射
$s : Y \to Y'$ および $t : Y' \to Y''$ に対して
$\left(t \circ s\right)^{-1} \left(t \circ f\right) = s^{-1} f$ である。
また、$\mathcal{C}$ の任意の $f : X \to Y'$ と $g : Y' \to Z'$、および
$S$ のすべての $s : Z \to Z'$ に対して
$s^{-1} \left(g \circ f\right) = \left(s^{-1} g\right) \circ
\left(\text{id}_{Y'}^{-1} f\right)$ である。
最後に、$\mathcal{C}$ の任意の $f : X \to Y'$、$S$ のすべての
$s : Y \to Y'$ および $S$ のすべての $t : Z \to Y$ に対して
$\left(s \circ t\right)^{-1} f
= \left(t^{-1} \text{id}_Y\right)
\circ \left(s^{-1} f\right)$ である。
これらはすべて定義から明らかである。
「同じ終点をもつ有限個の射を、共通分母をもつ分数として書く」ことができる。

\begin{lemma}
\label{categories-lemma-morphisms-left-localization}
$\mathcal{C}$ を圏、$S$ を $\mathcal{C}$ の射の左乗法系とする。
$S^{-1}\mathcal{C}$ の射の任意の有限族 $g_i : X_i \to Y$
（$i$ で添字付ける）が与えられたとき、$S$ の元 $s : Y \to Y'$ と
$\mathcal{C}$ の射の族 $f_i : X_i \to Y'$ で、各 $g_i$ が組
$(f_i : X_i \to Y', s : Y \to Y')$ の同値類となるものを見いだせる。
\end{lemma}

\begin{proof}
各 $i$ に対して $g_i$ の代表
$(X_i \to Y_i, s_i : Y \to Y_i)$ を選ぶ。
$S$ に属する射 $s : Y \to Y'$ で、各 $i$ に対して
$a_i \circ s_i = s$ を満たす射 $a_i : Y_i \to Y'$ が存在するものを
見いだせれば、補題が従う。二つの添字 $i = 1, 2$ がある場合には、
Definition \ref{categories-definition-multiplicative-system} により可能なように、
$t_2 \in S$ をもつ正方形
$$
\xymatrix{
Y \ar[d]_{s_1} \ar[r]_{s_2} & Y_2 \ar[d]^{t_2} \\
Y_1 \ar[r]^{a_1} & Y'
}
$$
を完成させればよい。このとき $s = t_2 \circ s_2 \in S$ が使える。
$n > 2$ 個の射がある場合には、上の方法で $n - 1$ 個の場合へ帰着し、
帰納法によって結論を得る。
\end{proof}

\noindent
同じ分母をもつ場合には、射の相等性を容易に特徴づけられる。

\begin{lemma}
\label{categories-lemma-equality-morphisms-left-localization}
$\mathcal{C}$ を圏、$S$ を $\mathcal{C}$ の射の左乗法系とする。
$A, B : X \to Y$ を $S^{-1}\mathcal{C}$ の射で、それぞれ
$(f : X \to Y', s : Y \to Y')$ および
$(g : X \to Y', s : Y \to Y')$ の同値類であるものとする。次は同値である。
\begin{enumerate}
\item $A = B$
\item $t \circ f = t \circ g$ を満たす $S$ の射
$t : Y' \to Y''$ が存在する。
\item $a \circ f = a \circ g$ および $a \circ s \in S$ を満たす射
$a : Y' \to Y''$ が存在する。
\end{enumerate}
\end{lemma}

\begin{proof}
$S^{-1}\mathcal{C}$ が圏であること
（Lemma \ref{categories-lemma-left-localization}）を用いる。また、
Definition \ref{categories-definition-left-localization-as-fraction} の記法と、
その定義に続く議論を用いて $S^{-1}\mathcal{C}$ のいくつかの射を同一視する。
したがって $A = s^{-1}f$ および $B = s^{-1}g$ と書く。

\medskip\noindent
$A = B$ ならば
$(\text{id}_{Y'}^{-1}s) \circ A = (\text{id}_{Y'}^{-1}s) \circ B$ である。
さらに $(\text{id}_{Y'}^{-1}s) \circ A = \text{id}_{Y'}^{-1}f$ および
$(\text{id}_{Y'}^{-1}s) \circ B = \text{id}_{Y'}^{-1}g$ である。
$\text{id}_{Y'}^{-1}f$ と $\text{id}_{Y'}^{-1}g$ の相等性は、定義により、
$t \in S$ をもつ可換図式
$$
\xymatrix{
 & Y' \ar[d]^u & \\
X \ar[ru]^f \ar[r]^h \ar[rd]_g &
Z &
Y' \ar[lu]_{\text{id}_{Y'}} \ar[l]_t \ar[ld]^{\text{id}_{Y'}} \\
& Y' \ar[u]_v &
}
$$
が存在することを意味する。特に $u = v = t \in S$ および
$t \circ f = t\circ g$ である。したがって (1) は (2) を含意する。

\medskip\noindent
(2) $\Rightarrow$ (3) は直ちに従う。$a$ が (3) のとおりであると仮定する。
$s' = a \circ s \in S$ と表す。このとき $\text{id}_{Y''}^{-1}s'$ は圏
$S^{-1}\mathcal{C}$ における同型である
（その逆は $(s')^{-1}\text{id}_{Y''}$ である）。
したがって $A = B$ を確認するには
$\text{id}_{Y''}^{-1}s' \circ A = \text{id}_{Y''}^{-1}s' \circ B$ を
確認すれば十分である。Definition
\ref{categories-definition-left-localization-as-fraction} に続く本文で述べた規則を用いて計算すると
$\text{id}_{Y''}^{-1}s' \circ A =
\text{id}_{Y''}^{-1}(a \circ s) \circ s^{-1}f =
\text{id}_{Y''}^{-1}(a \circ f) =
\text{id}_{Y''}^{-1}(a \circ g) =
\text{id}_{Y''}^{-1}(a \circ s) \circ s^{-1}g =
\text{id}_{Y''}^{-1}s' \circ B$ となり、(1) が成り立つことが分かる。
\end{proof}

\begin{remark}
\label{categories-remark-left-localization-morphisms-colimit}
$\mathcal{C}$ を圏、$S$ を左乗法系とする。$\mathcal{C}$ の対象 $Y$ が与えられたとき、
対象が $s \in S$ を満たす $s : Y \to Y'$ であり、射が可換図式
$$
\xymatrix{
& Y \ar[ld]_s \ar[rd]^t & \\
Y' \ar[rr]^a & & Y''
}
$$
（ここで $a : Y' \to Y''$ は任意）である圏を $Y/S$ と表す。
圏 $Y/S$ はフィルター付きであると主張する
（Definition \ref{categories-definition-directed} を参照せよ）。
実際、LMS1 により $\text{id}_Y : Y \to Y$ は $Y/S$ に属するので、
$Y/S$ は空でない。LMS2 により、$s_1 : Y \to Y_1$ および
$s_2 : Y \to Y_2$ が与えられたとき、$t \in S$ をもつ図式
$$
\xymatrix{
Y \ar[d]_{s_1} \ar[r]_{s_2} & Y_2 \ar[d]^t \\
Y_1 \ar[r]^a & Y_3
}
$$
を見いだせる。したがって $s_1 : Y \to Y_1$ と $s_2 : Y \to Y_2$ はいずれも
$Y/S$ において $t \circ s_2 : Y \to Y_3$ への写像をもつ。
最後に、$Y/S$ において $s_1 : Y \to Y_1$ から
$s_2 : Y \to Y_2$ への二つの射 $a, b$ が与えられると、
$a \circ s_1 = b \circ s_1$ である。よって LMS3 により、
$t \circ a = t \circ b$ を満たす $S$ の射 $t : Y_2 \to Y_3$ が存在する。
Lemmas \ref{categories-lemma-morphisms-left-localization} and
\ref{categories-lemma-equality-morphisms-left-localization} の結果を合わせると
\begin{equation}
\label{categories-equation-left-localization-morphisms-colimit}
\Mor_{S^{-1}\mathcal{C}}(X, Y) =
\colim_{(s : Y \to Y') \in Y/S} \Mor_\mathcal{C}(X, Y')
\end{equation}
を得る。$S^{-1}\mathcal{C}$ の射集合を $\mathcal{C}$ の射集合の
フィルター付き余極限として表すこの公式は、ときに有用である。
\end{remark}

\begin{lemma}
\label{categories-lemma-properties-left-localization}
$\mathcal{C}$ を圏、$S$ を $\mathcal{C}$ の射の左乗法系とする。
\begin{enumerate}
\item 規則 $X \mapsto X$ および
$(f : X \to Y) \mapsto (f : X \to Y, \text{id}_Y : Y \to Y)$ は関手
$Q : \mathcal{C} \to S^{-1}\mathcal{C}$ を定める。
\item 任意の $s \in S$ に対して、射 $Q(s)$ は
$S^{-1}\mathcal{C}$ における同型である。
\item $G : \mathcal{C} \to \mathcal{D}$ を、すべての $s \in S$ に対して
$G(s)$ が可逆である任意の関手とする。このとき $H \circ Q = G$ を満たす
一意な関手 $H : S^{-1}\mathcal{C} \to \mathcal{D}$ が存在する。
\end{enumerate}
\end{lemma}

\begin{proof}
(1) と (2) は明らかである。（(2) において、$Q(s)$ の逆は組
$(\text{id}_Y, s)$ の同値類である。）(3) を示すには、
$H(X) = G(X)$ と置き、
$H((f : X \to Y', s : Y \to Y')) = G(s)^{-1} \circ G(f)$ と置けばよい。
詳細は省略する。
\end{proof}

\begin{lemma}
\label{categories-lemma-left-localization-limits}
$\mathcal{C}$ を圏、$S$ を $\mathcal{C}$ の射の左乗法系とする。
局所化関手 $Q : \mathcal{C} \to S^{-1}\mathcal{C}$ は有限余極限と可換である。
\end{lemma}

\begin{proof}
$\mathcal{I}$ を有限圏とし、$\mathcal{I} \to \mathcal{C}$,
$i \mapsto X_i$ を余極限が存在する関手とする。
(\ref{categories-equation-left-localization-morphisms-colimit})、$Y/S$ がフィルター付きであること、
および Lemma \ref{categories-lemma-directed-commutes} を用いると
\begin{align*}
\Mor_{S^{-1}\mathcal{C}}(Q(\colim X_i), Q(Y))
& =
\colim_{(s : Y \to Y') \in Y/S} \Mor_\mathcal{C}(\colim X_i, Y') \\
& =
\colim_{(s : Y \to Y') \in Y/S} \lim_i \Mor_\mathcal{C}(X_i, Y') \\
& =
\lim_i \colim_{(s : Y \to Y') \in Y/S} \Mor_\mathcal{C}(X_i, Y') \\
& =
\lim_i \Mor_{S^{-1}\mathcal{C}}(Q(X_i), Q(Y))
\end{align*}
を得る。この同型は、各 $j \in \Ob(\mathcal{I})$ に対する集合
$\Mor_{S^{-1}\mathcal{C}}(Q(X_j), Q(Y))$ への両辺からの射影と可換である。
したがって $Q(\colim X_i)$ は関手 $i \mapsto Q(X_i)$ の余極限の普遍性を満たす。
よって望むように、これはその余極限である。
\end{proof}

\begin{lemma}
\label{categories-lemma-left-localization-diagram}
$\mathcal{C}$ を圏、$S$ を左乗法系とする。
$f : X \to Y$, $f' : X' \to Y'$ を $\mathcal{C}$ の二射とし、
$$
\xymatrix{
Q(X) \ar[d]_{Q(f)} \ar[r]_a & Q(X') \ar[d]^{Q(f')} \\
Q(Y) \ar[r]^b & Q(Y')
}
$$
を $S^{-1}\mathcal{C}$ における可換図式とする。このとき、
$\mathcal{C}$ の射 $f'' : X'' \to Y''$ と、$s, t \in S$ および
$a = s^{-1}g$, $b = t^{-1}h$ を満たす $\mathcal{C}$ の可換図式
$$
\xymatrix{
X \ar[d]_f \ar[r]_g & X'' \ar[d]^{f''} & X' \ar[d]^{f'} \ar[l]^s \\
Y \ar[r]^h & Y'' & Y' \ar[l]_t
}
$$
が存在する。
\end{lemma}

\begin{proof}
写像と対象を次のように選ぶ。まず、$S$ のある $s : X' \to X''$ と
$g : X \to X''$ によって $a = s^{-1}g$ と書く。
LMS2 により、$S$ の $t : Y' \to Y''$ と $f'' : X'' \to Y''$ で、
$$
\xymatrix{
X \ar[d]_f \ar[r]_g & X'' \ar[d]^{f''} & X' \ar[d]^{f'} \ar[l]^s \\
Y & Y'' & Y' \ar[l]_t
}
$$
を可換にするものを見いだせる。これから、この図式における
$$
X'' \xrightarrow{f''} Y'' \xleftarrow{t} Y'
$$
の選択を、$d \circ t \in S$ を満たす射 $d : Y'' \to Y'''$ により
$t$ と $f''$ の両方を後合成することで繰り返し変更する。
Remark \ref{categories-remark-left-localization-morphisms-colimit} によれば、
このような置換の後、$b = t^{-1}h$ を満たす射 $h : Y \to Y''$ が存在すると
仮定できる\footnote{これをより直接に見る方法は次のとおりである。
$S$ のある $q : Y' \to Z$ と $i : Y \to Z$ によって
$b = q^{-1}i$ と書く。LMS2 により、$j \circ q = r \circ t$ を満たす
$S$ の $r : Y'' \to Y'''$ と $j : Z \to Y'''$ を見いだせる。
ここで $d = r$ および $h = j \circ i$ と置く。}。
この時点で、図式の左側の正方形が可換であるか分からないことを除き、補題のすべてを得ている。
しかし $S^{-1} \mathcal{C}$ における合成の定義によれば、
$b \circ Q\left(f\right)$ は組
$(h \circ f : X \to Y'', t : Y' \to Y'')$ の同値類である
（$b$ は組 $(h : Y \to Y'', t : Y' \to Y'')$ の同値類であり、
$Q\left(f\right)$ は組 $(f : X \to Y, \text{id} : Y \to Y)$ の同値類である）。
一方、$Q\left(f'\right) \circ a$ は組
$(f'' \circ g : X \to Y'', t : Y' \to Y'')$ の同値類である
（$a$ は組 $(g : X \to X'', s : X' \to X'')$ の同値類であり、
$Q\left(f'\right)$ は組
$(f' : X' \to Y', \text{id} : Y' \to Y')$ の同値類である）。
$b \circ Q\left(f\right) = Q\left(f'\right) \circ a$ であることが分かっているので、
組 $(h \circ f : X \to Y'', t : Y' \to Y'')$ と
$(f'' \circ g : X \to Y'', t : Y' \to Y'')$ の同値類は等しい。
したがって Lemma \ref{categories-lemma-equality-morphisms-left-localization} により、
$d \circ t \in S$ および
$d \circ h \circ f = d \circ f'' \circ g$ を満たす射
$d : Y'' \to Y'''$ を見いだせる。よって上で述べた種類の置換をもう一度行えば、
結論を得る。
\end{proof}

\noindent
{\bf 右分数計算。}
$\mathcal{C}$ を圏、$S$ を右乗法系とする。新しい圏
$S^{-1}\mathcal{C}$ を次のように定義する（これが機能することは
Lemma \ref{categories-lemma-right-localization} の証明で確認する）。
\begin{enumerate}
\item $\Ob(S^{-1}\mathcal{C}) = \Ob(\mathcal{C})$ と置く。
\item $S^{-1}\mathcal{C}$ の射 $X \to Y$ は、同値関係を除いて
$s \in S$ を満たす組 $(f : X' \to Y, s : X' \to X)$ によって与えられる。
（同値関係は以下で定義する。組 $(f, s)$ の同値類を
$fs^{-1} : X \to Y$ と考えよ。）
\item 二つの組 $(f_1 : X_1 \to Y, s_1 : X_1 \to X)$ および
$(f_2 : X_2 \to Y, s_2 : X_2 \to X)$ が同値であるとは、第三の組
$(f_3 : X_3 \to Y, s_3 : X_3 \to X)$ と、可換図式
$$
\xymatrix{
 & X_1 \ar[ld]_{s_1} \ar[rd]^{f_1} & \\
X &
X_3 \ar[l]_{s_3} \ar[u]_u \ar[d]^v \ar[r]^{f_3} &
Y \\
& X_2 \ar[lu]^{s_2} \ar[ru]_{f_2} &
}
$$
に入る $\mathcal{C}$ の射 $u : X_3 \to X_1$ および
$v : X_3 \to X_2$ が存在することをいう。
\item 組 $(f : X' \to Y, s : X' \to X)$ と
$(g : Y' \to Z, t : Y' \to Y)$ の同値類の合成を、組
$(g \circ h : X'' \to Z, s \circ u : X'' \to X)$ の同値類として定義する。
ここで $h$ および $u \in S$ は、仮定により存在する可換図式
$$
\xymatrix{
X'' \ar[d]_u \ar[r]^h & Y' \ar[d]^t \\
X' \ar[r]^f & Y
}
$$
に入るように選ぶ。
\item $S^{-1} \mathcal{C}$ における恒等射 $X \to X$ は、組
$(\text{id} : X \to X, \text{id} : X \to X)$ の同値類である。
\end{enumerate}

\begin{lemma}
\label{categories-lemma-right-localization}
$\mathcal{C}$ を圏、$S$ を右乗法系とする。
\begin{enumerate}
\item 上で定義した組の間の関係は同値関係である。
\item 上で与えた合成規則は同値類上で良定義である。
\item 合成は結合的であり（恒等射は恒等律を満たす）、したがって
$S^{-1}\mathcal{C}$ は圏である。
\end{enumerate}
\end{lemma}

\begin{proof}
この補題は Lemma \ref{categories-lemma-left-localization} と双対である。
$\mathcal{C}$ を、その中で $S$ が左乗法系となる反対圏で置き換えることにより、
同補題から形式的に従う。
\end{proof}

\begin{definition}
\label{categories-definition-right-localization-as-fraction}
$\mathcal{C}$ を圏、$S$ を $\mathcal{C}$ の射からなる右乗法系とする。
$\mathcal{C}$ の任意の射 $f : X' \to Y$ と、$S$ の任意の射
$s : X' \to X$ に対し、組 $(f : X' \to Y, s : X' \to X)$ の同値類を
{\it $f s^{-1}$} と記す。これは $S^{-1} \mathcal{C}$ における
$X$ から $Y$ への射である。
\end{definition}

\noindent
Definition \ref{categories-definition-left-localization-as-fraction} におけるものと同様の
（実際には双対な）恒等式が成り立つ。
「同じ始域をもつ任意の有限個の射を、共通分母をもつ分数として書く」ことができる。

\begin{lemma}
\label{categories-lemma-morphisms-right-localization}
$\mathcal{C}$ を圏、$S$ を $\mathcal{C}$ の射からなる右乗法系とする。
$S^{-1}\mathcal{C}$ の任意の有限個の射 $g_i : X \to Y_i$
（$i$ で添字付けられている）が与えられたとする。このとき、
$S$ の元 $s : X' \to X$ と、$\mathcal{C}$ の射の族
$f_i : X' \to Y_i$ で、$g_i$ が組
$(f_i : X' \to Y_i, s : X' \to X)$ の同値類となるものを見いだせる。
\end{lemma}

\begin{proof}
この補題は Lemma \ref{categories-lemma-morphisms-left-localization} と双対であり、
現れるすべての圏をその反対圏で置き換えることにより、同補題から形式的に従う。
\end{proof}

\noindent
二つの射が同じ分母をもつ場合、その等しさには簡単な特徴付けがある。

\begin{lemma}
\label{categories-lemma-equality-morphisms-right-localization}
$\mathcal{C}$ を圏、$S$ を $\mathcal{C}$ の射からなる右乗法系とする。
$A, B : X \to Y$ を $S^{-1}\mathcal{C}$ の射で、それぞれ
$(f : X' \to Y, s : X' \to X)$ および
$(g : X' \to Y, s : X' \to X)$ の同値類であるものとする。
次は同値である。
\begin{enumerate}
\item $A = B$、
\item $f \circ t = g \circ t$ を満たす $S$ の射
$t : X'' \to X'$ が存在すること、
\item $f \circ a = g \circ a$ かつ $s \circ a \in S$ を満たす射
$a : X'' \to X'$ が存在すること。
\end{enumerate}
\end{lemma}

\begin{proof}
これは Lemma \ref{categories-lemma-equality-morphisms-left-localization} と双対である。
\end{proof}

\begin{remark}
\label{categories-remark-right-localization-morphisms-colimit}
$\mathcal{C}$ を圏、$S$ を右乗法系とする。
$\mathcal{C}$ の対象 $X$ が与えられたとき、$S/X$ により次の圏を表す。
その対象は $s \in S$ を満たす $s : X' \to X$ であり、その射は可換図式
$$
\xymatrix{
X' \ar[rd]_s \ar[rr]_a & & X'' \ar[ld]^t \\
& X
}
$$
である。ここで $a : X' \to X''$ は任意である。圏 $S/X$ は余フィルター付きである
（Definition \ref{categories-definition-codirected} を参照）。
（これは Remark \ref{categories-remark-left-localization-morphisms-colimit} の対応する主張と双対である。）
さて、Lemmas \ref{categories-lemma-morphisms-right-localization} および
\ref{categories-lemma-equality-morphisms-right-localization} の結果を合わせると、
\begin{equation}
\label{categories-equation-right-localization-morphisms-colimit}
\Mor_{S^{-1}\mathcal{C}}(X, Y) =
\colim_{(s : X' \to X) \in (S/X)^{opp}} \Mor_\mathcal{C}(X', Y)
\end{equation}
を得る。$S^{-1}\mathcal{C}$ の射を $\mathcal{C}$ の射のフィルター付き余極限として
表すこの公式は、ときに有用である。
\end{remark}

\begin{lemma}
\label{categories-lemma-properties-right-localization}
$\mathcal{C}$ を圏、$S$ を $\mathcal{C}$ の射からなる右乗法系とする。
\begin{enumerate}
\item 規則 $X \mapsto X$ および
$(f : X \to Y) \mapsto (f : X \to Y, \text{id}_X : X \to X)$ は
関手 $Q : \mathcal{C} \to S^{-1}\mathcal{C}$ を定める。
\item 任意の $s \in S$ に対し、射 $Q(s)$ は
$S^{-1}\mathcal{C}$ における同型である。
\item $G : \mathcal{C} \to \mathcal{D}$ を、すべての $s \in S$ に対して
$G(s)$ が可逆となる任意の関手とする。このとき、$H \circ Q = G$ を満たす
一意的な関手 $H : S^{-1}\mathcal{C} \to \mathcal{D}$ が存在する。
\end{enumerate}
\end{lemma}

\begin{proof}
この補題は Lemma \ref{categories-lemma-properties-left-localization} と双対であり、
現れるすべての圏をその反対圏で置き換えることにより、同補題から形式的に従う。
\end{proof}

\begin{lemma}
\label{categories-lemma-right-localization-limits}
$\mathcal{C}$ を圏、$S$ を $\mathcal{C}$ の射からなる右乗法系とする。
局所化関手 $Q : \mathcal{C} \to S^{-1}\mathcal{C}$ は有限極限と可換である。
\end{lemma}

\begin{proof}
これは Lemma \ref{categories-lemma-left-localization-limits} と双対である。
\end{proof}

\begin{lemma}
\label{categories-lemma-right-localization-diagram}
$\mathcal{C}$ を圏、$S$ を右乗法系とする。
$f : X \to Y$, $f' : X' \to Y'$ を $\mathcal{C}$ の二射とし、
$$
\xymatrix{
Q(X) \ar[d]_{Q(f)} \ar[r]_a & Q(X') \ar[d]^{Q(f')} \\
Q(Y) \ar[r]^b & Q(Y')
}
$$
を $S^{-1}\mathcal{C}$ における可換図式とする。このとき、
$\mathcal{C}$ の射 $f'' : X'' \to Y''$ と、$s, t \in S$ および
$a = gs^{-1}$, $b = ht^{-1}$ を満たす $\mathcal{C}$ の可換図式
$$
\xymatrix{
X \ar[d]_f & X'' \ar[l]^s \ar[d]^{f''} \ar[r]_g & X' \ar[d]^{f'} \\
Y & Y'' \ar[l]_t \ar[r]^h & Y'
}
$$
が存在する。
\end{lemma}

\begin{proof}
この補題は Lemma \ref{categories-lemma-left-localization-diagram} と双対である。
\end{proof}

\noindent
{\bf 乗法系と両側分数計算。}
$S$ が乗法系であるとき、左分数計算と右分数計算は標準的に同型な圏を与える。

\begin{lemma}
\label{categories-lemma-multiplicative-system}
$\mathcal{C}$ を圏、$S$ を乗法系とする。左分数の圏と右分数の圏
$S^{-1}\mathcal{C}$ は標準的に同型である。
\end{lemma}

\begin{proof}
二つの分数圏を $\mathcal{C}_{left}$, $\mathcal{C}_{right}$ と記す。
Lemmas \ref{categories-lemma-properties-left-localization} および
\ref{categories-lemma-properties-right-localization} の普遍性により、関手
$\mathcal{C}_{left} \to \mathcal{C}_{right}$ と
$\mathcal{C}_{right} \to \mathcal{C}_{left}$ を得る。
普遍性における一意性の主張により、これらの関手は互いに逆である。
\end{proof}

\begin{definition}
\label{categories-definition-saturated-multiplicative-system}
$\mathcal{C}$ を圏、$S$ を乗法系とする。MS1, MS2, MS3 に加えて次も成り立つとき、
$S$ は {\it 飽和している} という。
\begin{enumerate}
\item[MS4] 合成可能な三射 $f, g, h$ が与えられ、
$fg, gh \in S$ ならば $g \in S$ である。
\end{enumerate}
\end{definition}

\noindent
飽和乗法系はすべての同型を含むことに注意せよ。
さらに、$f, g, h$ が圏における合成可能な射で、$fg, gh$ が同型ならば、
$g$ は同型である（このとき $g$ は左逆射と右逆射をともにもつため、可逆である）。

\begin{lemma}
\label{categories-lemma-what-gets-inverted}
$\mathcal{C}$ を圏、$S$ を乗法系とする。局所化関手を
$Q : \mathcal{C} \to S^{-1}\mathcal{C}$ と記す。集合
$$
\hat S = \{f \in \text{Arrows}(\mathcal{C}) \mid
Q(f) \text{ is an isomorphism}\}
$$
は
$$
S' = \{f \in \text{Arrows}(\mathcal{C}) \mid
\text{there exist }g, h\text{ such that }gf, fh \in S\}
$$
に等しく、$S$ を含む最小の飽和乗法系である。特に、$S$ が飽和していれば
$\hat S = S$ である。
\end{lemma}

\begin{proof}
$S'$ の元は $S^{-1}\mathcal{C}$ において左逆射と右逆射をともにもつ射へ写るので、
$S \subset S' \subset \hat S$ は明らかである。$S'$ は MS4 を満たし、
$\hat S$ は MS1 を満たすことに注意せよ。次に $S' = \hat S$ を示す。

\medskip\noindent
$f \in \hat S$ とする。$S^{-1}\mathcal{C}$ における逆射を
$s^{-1}g = ht^{-1}$ とする。（$S^{-1}\mathcal{C}$ の射を記述するために
左分数と右分数の両方を用いてよい。Lemma \ref{categories-lemma-multiplicative-system} を参照。）
$S^{-1}\mathcal{C}$ における関係 $\text{id}_X = s^{-1}gf$ は、ある射
$f', u, v$ と $s' \in S$ に対して可換図式
$$
\xymatrix{
 & X' \ar[d]^u & \\
X \ar[ru]^{gf} \ar[r]^{f'} \ar[rd]_{\text{id}_X} &
X'' &
X \ar[lu]_s \ar[l]_{s'} \ar[ld]^{\text{id}_X} \\
& X \ar[u]_v &
}
$$
が存在することを意味する。したがって $ugf = s' \in S$ である。
同様に、$\text{id}_Y = fht^{-1}$ を用いると、ある $w$ に対して
$fhw \in S$ であることが分かる。ゆえに $f \in S'$ である。
したがって $S' = \hat S$ である。$S' = \hat S$ が乗法系であることを示せば、
$S' = \hat S$ が $S$ を含む最小の飽和系であることは明らかである。

\medskip\noindent
上の注意により MS1 と MS4 は処理済みなので、補題の証明を終えるには、
LMS2, RMS2, LMS3, RMS3 が $\hat S$ に対して成り立つことを示せばよい。
$\hat S$ に対する LMS2 を確認しよう。$t \in \hat S$ をもつ実線の図式
$$
\xymatrix{
X \ar[d]_t \ar[r]_g & Y \ar@{..>}[d]^s \\
Z \ar@{..>}[r]^f & W
}
$$
が与えられたとする。$at \in S$ を満たす射 $a : Z \to Z'$ を選ぶ。
すると $S$ に対する LMS2 を用いて可換図式
$$
\xymatrix{
X \ar[d]_t \ar[r]_g & Y \ar[dd]^s \\
Z \ar[d]_a \\
Z' \ar[r]^{f'} & W
}
$$
を見いだせ、$f = f' \circ a$ と置けばよい。RMS2 の証明はこれと双対である。
最後に、射の対 $f, g : X \to Y$ と、終域が $X$ である
$t \in \hat S$ が与えられ、$ft = gt$ であるとする。
$tb \in S$ を満たす射 $b$ を選ぶ。このとき $ftb = gtb$ であり、
$S$ に対する LMS3 により、始域が $Y$ で $sf = sg$ を満たす
$s \in S$ が存在する。これが望むものである。RMS3 の証明はこれと双対である。
\end{proof}

\section{形式的性質}
\label{categories-section-formal-cat-cat}

\noindent
この節では、圏からなる $2$-圏のいくつかの形式的性質を論じる。
これは後で（狭義の）$2$-圏を定義することにつながる。

\medskip\noindent
すべての圏からなる類を $\Ob(\textit{Cat})$ と記そう。
圏の各対 $\mathcal{A}, \mathcal{B} \in \Ob(\textit{Cat})$ に対し、
関手からなる「小さい」圏 $\text{Fun}(\mathcal{A}, \mathcal{B})$ がある。
次のような関手の変換の合成
$$
\xymatrix{
\mathcal{A}
\rruppertwocell^{F''}{t'}
\ar[rr]_(.3){F'}
\rrlowertwocell_F{t}
& &
\mathcal{B}
}
\text{ composes to }
\xymatrix{
\mathcal{A}
\rrtwocell^{F''}_F{\ \ t \circ t'}
& &
\mathcal{B}
}
$$
を {\it 垂直}合成という。これには通常の記号 $\circ$ を用いる。
次に {\it 水平}合成を定義する。そのため、手元の構造をもう少し説明する。

\medskip\noindent
すなわち、圏の各三つ組 $\mathcal{A}$, $\mathcal{B}$, $\mathcal{C}$ に対して、
関手の合成から得られる合成則
$$
\circ : \Ob(\text{Fun}(\mathcal{B}, \mathcal{C}))
\times
\Ob(\text{Fun}(\mathcal{A}, \mathcal{B}))
\longrightarrow
\Ob(\text{Fun}(\mathcal{A}, \mathcal{C}))
$$
がある。この合成則は結合的であり、恒等関手は単位元として作用する。
言い換えると、関手の変換を忘れれば、$\textit{Cat}$ は圏をなすことが分かる。
この構造は関手間の射とどのように相互作用するだろうか。

\medskip\noindent
$t : F \to F'$ を関手 $F, F' : \mathcal{A} \to \mathcal{B}$ の変換とし、
$G : \mathcal{B} \to \mathcal{C}$ を関手とする。このとき関手の変換
$G\circ F \to G \circ F'$ を定義できる。この変換を ${}_Gt$ と記す。
これはすべての $x \in \mathcal{A}$ に対し、公式
$({}_Gt)_x = G(t_x) : G(F(x)) \to G(F'(x))$ によって与えられる。
このようにして、$G$ との合成は関手
$$
\text{Fun}(\mathcal{A}, \mathcal{B})
\longrightarrow
\text{Fun}(\mathcal{A}, \mathcal{C}).
$$
となる。これを見るには ${}_G(\text{id}_F) = \text{id}_{G \circ F}$ および
${}_G(t_1 \circ t_2) = {}_Gt_1 \circ {}_Gt_2$ を確認すればよい。
もちろん ${}_{\text{id}_\mathcal{B}}t = t$ も成り立つ。

\medskip\noindent
同様に、$s : G \to G'$ を関手 $G, G' : \mathcal{B} \to \mathcal{C}$ の変換、
$F : \mathcal{A} \to \mathcal{B}$ を関手とする。このとき $s_F$ を、
すべての $x \in \mathcal{A}$ に対し
$(s_F)_x = s_{F(x)} : G(F(x)) \to G'(F(x))$ で与えられる関手の変換
$G\circ F \to G' \circ F$ として定義できる。このようにして、
$F$ との合成は関手
$$
\text{Fun}(\mathcal{B}, \mathcal{C})
\longrightarrow
\text{Fun}(\mathcal{A}, \mathcal{C}).
$$
となる。これを見るには $(\text{id}_G)_F = \text{id}_{G\circ F}$ および
$(s_1 \circ s_2)_F = s_{1, F} \circ s_{2, F}$ を確認すればよい。
もちろん $s_{\text{id}_\mathcal{B}} = s$ も成り立つ。

\medskip\noindent
これらの構成は、意味をなすときにはいつでも、さらに次の性質を満たす。
$$
{}_{G_1}({}_{G_2}t) = {}_{G_1\circ G_2}t,
\ (s_{F_1})_{F_2} = s_{F_1 \circ F_2},
\text{ and }{}_H(s_F) = ({}_Hs)_F
$$
最後に、関手 $F, F' : \mathcal{A} \to \mathcal{B}$ および
$G, G' : \mathcal{B} \to \mathcal{C}$ と、変換
$t : F \to F'$, $s : G \to G'$ が与えられたとき、図式
$$
\xymatrix{
G \circ F \ar[r]^{{}_Gt} \ar[d]_{s_F}
&
G \circ F' \ar[d]^{s_{F'}} \\
G' \circ F \ar[r]_{{}_{G'}t}
&
G' \circ F'
}
$$
は可換である。言い換えると ${}_{G'}t \circ s_F = s_{F'}\circ {}_Gt$ である。
これを証明するには、任意の対象 $x \in \Ob(\mathcal{A})$ 上で何が起こるかを考えればよい。
$$
\xymatrix{
G(F(x)) \ar[r]^{G(t_x)} \ar[d]_{s_{F(x)}}
&
G(F'(x)) \ar[d]^{s_{F'(x)}} \\
G'(F(x)) \ar[r]_{G'(t_x)}
&
G'(F'(x))
}
$$
これは $s$ が関手の変換であるため可換である。この両立関係により、水平合成を定義できる。

\begin{definition}
\label{categories-definition-horizontal-composition}
次の左辺にあるような図式が与えられたとする。
$$
\xymatrix{
\mathcal{A}
\rtwocell^F_{F'}{t}
&
\mathcal{B}
\rtwocell^G_{G'}{s}
&
\mathcal{C}
}
\text{ gives }
\xymatrix{
\mathcal{A}
\rrtwocell^{G \circ F} _{G' \circ F'}{\ \ s \star t}
& &
\mathcal{C}
}
$$
{\it 水平}合成 $s \star t$ を、関手の変換
${}_{G'}t \circ s_F = s_{F'}\circ {}_Gt$ として定義する。
\end{definition}

\noindent
すると、先に構成した変換 ${}_Gt$ と $s_F$ は
$ {}_Gt = \text{id}_G \star t $ および $ s_F = s \star \text{id}_F $ として復元できる。
さらに、上で得たすべての規則は、次の補題に述べる性質から従う。

\begin{lemma}
\label{categories-lemma-properties-2-cat-cats}
水平合成と垂直合成は次の性質をもつ。
\begin{enumerate}
\item $\circ$ と $\star$ は結合的である。
\item 恒等変換 $\text{id}_F$ は $\circ$ の単位元である。
\item 恒等関手の恒等変換 $\text{id}_{\text{id}_\mathcal{A}}$ は
$\star$ と $\circ$ の単位元である。
\item 図式
$$
\xymatrix{
\mathcal{A}
\rruppertwocell^F{t}
\ar[rr]_(.3){F'}
\rrlowertwocell_{F''}{t'}
& &
\mathcal{B}
\rruppertwocell^G{s}
\ar[rr]_(.3){G'}
\rrlowertwocell_{G''}{s'}
& &
\mathcal{C}
}
$$
が与えられるとき、$ (s' \circ s) \star (t' \circ t) = (s' \star t') \circ (s \star t)$ である。
\end{enumerate}
\end{lemma}

\begin{proof}
最後の主張は、先の記法を用いると次の等式となる。
$$
s'_{F''}
\circ
{}_{G'}t'
\circ
s_{F'}
\circ
{}_Gt
=
(s' \circ s)_{F''}
\circ
{}_G(t' \circ t).
$$
上の結果を中央の合成に適用すると、左辺を
$
s'_{F''}
\circ
s_{F''}
\circ
{}_Gt'
\circ
{}_Gt
$
と書き直せる。これが右辺に等しいことは容易に分かる。
\end{proof}

\noindent
補題の条件 (4) を別の形で述べると、関手の合成と関手の変換の水平合成が、
積圏を始域とする関手
$$
(\circ, \star) :
\text{Fun}(\mathcal{B}, \mathcal{C})
\times
\text{Fun}(\mathcal{A}, \mathcal{B})
\longrightarrow
\text{Fun}(\mathcal{A}, \mathcal{C})
$$
を生じさせるということである。Definition \ref{categories-definition-product-category} を参照。

\section{2-圏}
\label{categories-section-2-categories}

\noindent
スタックの文脈に現れる（狭義の）$2$-圏を定義する。
これを読む前に Section \ref{categories-section-formal-cat-cat} および
Example \ref{categories-example-2-1-category-of-categories} を参照されたい。
基本的には、この例を取り、そこにおける対象、$1$-射、$2$-射が満たす
すべての規則を書き下す。

\begin{definition}
\label{categories-definition-2-category}
（狭義の）{\it $2$-圏} $\mathcal{C}$ は次のデータからなる。
\begin{enumerate}
\item 対象の集合 $\Ob(\mathcal{C})$。
\item 各対 $x, y \in \Ob(\mathcal{C})$ に対する圏
$\Mor_\mathcal{C}(x, y)$。$\Mor_\mathcal{C}(x, y)$ の対象を
{\it $1$-射}と呼び、$F : x \to y$ と記す。これらの $1$-射の間の射を
{\it $2$-射}と呼び、$t : F' \to F$ と記す。
$\Mor_\mathcal{C}(x, y)$ における $2$-射の合成を {\it 垂直}合成と呼び、
$t : F' \to F$ および $t' : F'' \to F'$ に対して $t \circ t'$ と記す。
\item 各三つ組 $x, y, z\in \Ob(\mathcal{C})$ に対する関手
$$
(\circ, \star) :
\Mor_\mathcal{C}(y, z) \times \Mor_\mathcal{C}(x, y)
\longrightarrow
\Mor_\mathcal{C}(x, z).
$$
左辺の $1$-射の対 $(F, G)$ の像を $F$ と $G$ の {\it 合成}と呼び、
$F\circ G$ と記す。$2$-射の対 $(t, s)$ の像を {\it 水平}合成と呼び、
$t \star s$ と記す。
\end{enumerate}
これらのデータは次の規則を満たすものとする。
\begin{enumerate}
\item 対象の集合と $1$-射の集合は、$1$-射の合成とともに圏をなす。
\item $2$-射の水平合成は結合的である。
\item 恒等 $1$-射 $\text{id}_x$ の恒等 $2$-射
$\text{id}_{\text{id}_x}$ は水平合成の単位元である。
\end{enumerate}
\end{definition}

\noindent
これは明らかに、扱ううえであまり快適な種類の対象ではない。
一方、どのように扱うかがかなり明瞭な例は多数ある。
ここまでにある唯一の例は、対象が与えられた圏の集まり、$1$-射がそれらの圏の間の関手、
$2$-射が関手の自然変換である $2$-圏である。
Section \ref{categories-section-formal-cat-cat} を参照。この本文で扱うすべての $2$-圏は、
この例の部分 $2$-圏となる。部分 $2$-圏の意味は次のとおりである。

\begin{definition}
\label{categories-definition-sub-2-category}
$\mathcal{C}$ を $2$-圏とする。$\mathcal{C}$ の {\it 部分 $2$-圏}
$\mathcal{C}'$ は、$\Ob(\mathcal{C})$ の部分集合 $\Ob(\mathcal{C}')$ と、
すべての $x, y \in \Ob(\mathcal{C}')$ に対する圏
$\Mor_\mathcal{C}(x, y)$ の部分圏 $\Mor_{\mathcal{C}'}(x, y)$ であって、
これらが演算 $\circ$（$1$-射の合成）、$\circ$（$2$-射の垂直合成）、
$\star$（水平合成）とともに $2$-圏をなすものによって与えられる。
\end{definition}

\begin{remark}
\label{categories-remark-big-2-categories}
大きな $2$-圏。
多くの文献では $2$-圏が対象の類をもつことを許す
（ただし「類からなる類」は許されないことが望ましい）。
ここでも、次に列挙する場合に限って、これらの「大きな」$2$-圏を許す
（この一覧は進行に応じて更新する）。
\begin{enumerate}
\item 圏の $2$-圏 $\textit{Cat}$。
\item 圏の $(2, 1)$-圏 $\textit{Cat}$。
\item グルーポイドの $2$-圏 $\textit{Groupoids}$。これは $(2, 1)$-圏である。
\item 固定した圏上のファイバー圏の $2$-圏。
\item 固定した圏上のファイバー圏の $(2, 1)$-圏。
\item 固定した圏上でグルーポイドにファイバー化された圏の $2$-圏。
これは $(2, 1)$-圏である。
\item 固定したサイト上のスタックの $2$-圏。
\item 固定したサイト上のスタックの $(2, 1)$-圏。
\item 固定したサイト上のグルーポイドのスタックの $2$-圏。
これは $(2, 1)$-圏である。
\item 固定したサイト上の setoid のスタックの $2$-圏。
これは $(2, 1)$-圏である。
\item 固定した概形上の代数スタックの $2$-圏。
これは $(2, 1)$-圏である。
\end{enumerate}
Definition \ref{categories-definition-2-1-category} を参照。
各場合に $2$-圏 $\mathcal{C}$ の対象の類は真の類であるが、
すべての対象 $x, y \in \Ob(C)$ に対して圏 $\Mor_\mathcal{C}(x, y)$ は
（本書の約束の意味で）「小さい」ことに注意せよ。
\end{remark}

\noindent
Section \ref{categories-section-definition-categories} で定義した圏同値の概念は、
次のように $2$-圏というより一般の状況へ拡張される。

\begin{definition}
\label{categories-definition-equivalence}
$2$-圏の二対象 $x, y$ が {\it 同値}であるとは、$1$-射
$F : x \to y$ および $G : y \to x$ が存在し、$F \circ G$ が
$\text{id}_y$ と $2$-同型であり、$G \circ F$ が $\text{id}_x$ と
$2$-同型であることをいう。
\end{definition}

\noindent
圏から $2$-圏への関手が何を意味するか述べる必要があることもある。

\begin{definition}
\label{categories-definition-functor-into-2-category}
$\mathcal{A}$ を圏、$\mathcal{C}$ を $2$-圏とする。
\begin{enumerate}
\item 通常の圏から $2$-圏への {\it 関手}は、特に断らない限り
$2$-射を無視する。言い換えると、すべての 2-射を忘れて 2-圏から作られる圏への
「通常の」関手とする。
\item $\mathcal{A}$ から 2-圏 $\mathcal{C}$ への {\it 弱関手}、または
{\it 擬関手} $\varphi$ は、次のデータによって与えられる。
\begin{enumerate}
\item 写像 $\varphi : \Ob(\mathcal{A}) \to \Ob(\mathcal{C})$、
\item 各対 $x, y\in \Ob(\mathcal{A})$ と各射 $f : x \to y$ に対する
$1$-射 $\varphi(f) : \varphi(x) \to \varphi(y)$、
\item 各 $x\in \Ob(A)$ に対する $2$-射
$\alpha_x : \text{id}_{\varphi(x)} \to \varphi(\text{id}_x)$、
\item $\mathcal{A}$ の合成可能な射の各対 $f : x \to y$,
$g : y \to z$ に対する $2$-射
$\alpha_{g, f} : \varphi(g \circ f) \to \varphi(g) \circ \varphi(f)$。
\end{enumerate}
これらのデータは次の条件を満たすものとする。
\begin{enumerate}
\item $2$-射 $\alpha_x$ および $\alpha_{g, f}$ はすべて同型である。
\item $\mathcal{A}$ の任意の射 $f : x \to y$ に対して
$\alpha_{\text{id}_y, f} = \alpha_y \star \text{id}_{\varphi(f)}$ である。
$$
\xymatrix{
\varphi(x)
\rrtwocell^{\varphi(f)}_{\varphi(f)}{\ \ \ \ \text{id}_{\varphi(f)}}
& &
\varphi(y)
\rrtwocell^{\text{id}_{\varphi(y)}}_{\varphi(\text{id}_y)}{\alpha_y}
& &
\varphi(y)
}
=
\xymatrix{
\varphi(x)
\rrtwocell^{\varphi(f)}_{\varphi(\text{id}_y) \circ \varphi(f)}{\ \ \ \ \alpha_{\text{id}_y, f}}
& &
\varphi(y)
}
$$
\item $\mathcal{A}$ の任意の射 $f : x \to y$ に対して
$\alpha_{f, \text{id}_x} = \text{id}_{\varphi(f)} \star \alpha_x$ である。
\item $\mathcal{A}$ の合成可能な射の任意の三つ組
$f : w \to x$, $g : x \to y$, $h : y \to z$ に対して
$$
(\text{id}_{\varphi(h)} \star \alpha_{g, f})
\circ
\alpha_{h, g \circ f}
=
(\alpha_{h, g} \star \text{id}_{\varphi(f)})
\circ
\alpha_{h \circ g, f}
$$
である。言い換えると、対象が $1$-射、矢印が $2$-射である次の図式が可換である。
$$
\xymatrix{
\varphi(h \circ g \circ f)
\ar[d]_{\alpha_{h, g \circ f}}
\ar[rr]_{\alpha_{h \circ g, f}}
& &
\varphi(h \circ g) \circ \varphi(f)
\ar[d]^{\alpha_{h, g} \star \text{id}_{\varphi(f)}} \\
\varphi(h) \circ \varphi(g \circ f)
\ar[rr]^{\text{id}_{\varphi(h)} \star \alpha_{g, f}}
& &
\varphi(h) \circ \varphi(g) \circ \varphi(f)
}
$$
\end{enumerate}
\end{enumerate}
\end{definition}

\noindent
これもやはりあまり扱いやすい概念ではないが、ときには現れる。
任意の擬関手は関手と同型である、という定理がある。
最後に、{\it $2$-圏の間の関手}および {\it $2$-圏の間の擬関手}という概念がある。
後者の概念は $3$-圏の領域へ導く。ほとんどどんな代償を払ってでも、
これを定義する必要は避けたい。

\section{(2, 1)-圏}
\label{categories-section-2-1-categories}

\noindent
すべての $2$-射が同型であるという性質をもつ $2$-圏がある。
これらは以下で重要な役割を果たし、また扱いやすい。

\begin{definition}
\label{categories-definition-2-1-category}
（狭義の）{\it $(2, 1)$-圏}とは、すべての $2$-射が同型である $2$-圏をいう。
\end{definition}

\begin{example}
\label{categories-example-2-1-category-of-categories}
$2$-圏 $\textit{Cat}$（Remark \ref{categories-remark-big-2-categories} を参照）は、
関手の同型だけを $2$-射として許すことにより $(2, 1)$-圏にできる。
\end{example}

\noindent
実際、より一般に任意の $2$-圏 $\mathcal{C}$ から、対象と $1$-射は同じで、
$2$-射を $\mathcal{C}$ の可逆な $2$-射とする部分 $2$-圏
$\mathcal{C}'$ を考えることにより、$(2, 1)$-圏が得られる。
この状況では「{\it $\mathcal{C}'$ を $\mathcal{C}$ に付随する
$(2, 1)$-圏とする}」などという。例えば、グルーポイドの $(2, 1)$-圏とは、
対象がグルーポイド、$1$-射が関手、$2$-射が関手の同型である $2$-圏をいう。
ただし、グルーポイド間の関手の変換は自動的に同型なので、これは良い例ではない。

\begin{remark}
\label{categories-remark-other-2-categories}
したがって、上の Example \ref{categories-example-2-1-category-of-categories} の構成には、
グルーポイドの $2$-圏、固定した圏上でグルーポイドにファイバー化された圏、
またはスタックを考える変種がある。以下同様である。
\end{remark}

\section{2-ファイバー積}
\label{categories-section-2-fibre-products}

\noindent
この節では $2$-ファイバー積を導入する。$\mathcal{C}$ を 2-圏とする。図式
$$
\xymatrix{
w \ar[r] \ar[d] & y \ar[d] \\
x \ar[r] & z }
$$
が 2-可換であるとは、二つの 1-射 $w \to y \to z$ と $w \to x \to z$ が
2-同型であることをいう。2-圏では、実際に可換する図式よりも図式の 2-可換性を
求める方が自然である。（実際、狭義の 2-圏をそもそも扱うべきでないと言う人もおり、
「弱い」2-圏では 1-射からなる可換図式という概念さえ意味をなさない。）
それに応じて、ファイバー積の概念を調整しなければならない。

\medskip\noindent
$\mathcal{C}$ を $2$-圏とする。$x, y, z\in \Ob(\mathcal{C})$ とし、
$f\in \Mor_\mathcal{C}(x, z)$ および $g\in \Mor_{\mathcal C}(y, z)$ とする。
$f$ と $g$ の 2-ファイバー積を定義するため、2-可換図式
$$
\xymatrix{
& w \ar[r]_a \ar[d]_b & x \ar[d]^{f} \\
& y \ar[r]^{g} & z. }
$$
を考える。圏の場合、ファイバー積はこのような図式からなる圏の終対象である。
それに対応して、2-ファイバー積は 2-圏における終対象である
（下の定義を参照）。{\it $f$ と $g$ 上の $2$-可換図式の $2$-圏}を
次の $2$-圏として定義する。
\begin{enumerate}
\item 対象は上のような四つ組 $(w, a, b, \phi)$ であり、$\phi$ は可逆な 2-射
$\phi : f \circ a \to g \circ b$ である。
\item $(w', a', b', \phi')$ から $(w, a, b, \phi)$ への 1-射は、
可換図式
$$
\xymatrix{
f \circ a'
\ar[rr]_{\text{id}_f \star \alpha}
\ar[d]_{\phi'}
& &
f \circ a \circ k
\ar[d]^{\phi \star \text{id}_k}
\\
g \circ b'
\ar[rr]^{\text{id}_g \star \beta}
& &
g \circ b \circ k
}
$$
を満たす
$(k : w' \to w, \alpha : a' \to a \circ k, \beta : b' \to b \circ k)$
によって与えられる。
\item 第二の $1$-射
$(k', \alpha', \beta') : (w'', a'', b'', \phi'') \to
(w', \alpha', \beta', \phi')$ が与えられたとき、$1$-射の合成は規則
$$
(k, \alpha, \beta) \circ (k', \alpha', \beta') =
(k \circ k',
(\alpha \star \text{id}_{k'}) \circ \alpha',
(\beta \star \text{id}_{k'}) \circ \beta'),
$$
によって与えられる。
\item 始域と終域が同じ $1$-射
$(k_i, \alpha_i, \beta_i)$, $i = 1, 2$ の間の 2-射は、図式
$$
\xymatrix{
a'
\ar[rd]_{\alpha_2}
\ar[r]_{\alpha_1} &
a \circ k_1
\ar[d]^{\text{id}_a \star \delta} &
&
b \circ k_1
\ar[d]_{\text{id}_b \star \delta} &
b'
\ar[l]^{\beta_1}
\ar[ld]^{\beta_2}
\\
&
a \circ k_2
&
&
b \circ k_2
&
}
$$
を可換にする 2-射 $\delta : k_1 \to k_2$ によって与えられる。
\item $2$-射の垂直合成は、$\mathcal{C}$ における射 $\delta$ の垂直合成によって与えられる。
\item 図式
$$
\xymatrix{
(w'', a'', b'', \phi'')
\rrtwocell^{(k'_1, \alpha'_1, \beta'_1)}_{(k'_2, \alpha'_2, \beta'_2)}{\delta'}
& &
(w', a', b', \phi')
\rrtwocell^{(k_1, \alpha_1, \beta_1)}_{(k_2, \alpha_2, \beta_2)}{\delta}
& &
(w, a, b, \phi)
}
$$
の水平合成は図式
$$
\xymatrix@C=12pc{
(w'', a'', b'', \phi'')
\rtwocell^{(k_1 \circ k'_1, (\alpha_1 \star \text{id}_{k'_1}) \circ \alpha'_1, (\beta_1 \star \text{id}_{k'_1}) \circ \beta'_1)}_{(k_2 \circ k'_2, (\alpha_2 \star \text{id}_{k'_2}) \circ \alpha'_2, (\beta_2 \star \text{id}_{k'_2}) \circ \beta'_2)}{\ \ \ \delta \star \delta'}
&
(w, a, b, \phi)
}
$$
によって与えられる。
\end{enumerate}
$\mathcal{C}$ が実際に $(2, 1)$-圏であれば、上の (2) における射
$\alpha$ と $\beta$ も自動的に同型となることに注意せよ
\footnote{実際、$2$-圏の場合には、矢印 $\alpha$, $\beta$ の向きを逆にした、
あるいはその一方だけの向きを逆にした別の 2-可換図式の 2-圏を定義できるように思われる。
このため、後では $(2, 1)$-圏に制限する。}。
さらに、$\mathcal{C}$ が $(2, 1)$-圏ならば、$2$-可換図式の $2$-圏も
$(2, 1)$-圏である。

\begin{definition}
\label{categories-definition-final-object-2-category}
$(2, 1)$-圏 $\mathcal{C}$ の {\it 終対象}とは、次を満たす対象 $x$ をいう。
\begin{enumerate}
\item すべての $y \in \Ob(\mathcal{C})$ に対して射 $y \to x$ が存在する。
\item 任意の二つの射 $y \to x$ は、一意な 2-射によって同型である。
\end{enumerate}
\end{definition}

\noindent
より一般の $2$-圏の場合には、終対象にも異なる種類があると思われる。
ここではそれに立ち入りたくないので、$(2, 1)$ の場合にのみ $2$-ファイバー積を定義する。

\begin{definition}
\label{categories-definition-2-fibre-products}
$\mathcal{C}$ を $(2, 1)$-圏とする。$x, y, z\in \Ob(\mathcal{C})$ とし、
$f\in \Mor_\mathcal{C}(x, z)$ および $g\in \Mor_{\mathcal C}(y, z)$ とする。
{\it $f$ と $g$ の 2-ファイバー積}とは、上で述べた 2-可換図式の圏における
終対象である。2-ファイバー積が存在するとき、それを
$x \times_z y\in \Ob(\mathcal{C})$ と記し、必要な射を
$p\in \Mor_{\mathcal C}(x \times_z y, x)$ および
$q\in \Mor_{\mathcal C}(x \times_z y, y)$ と記す。これらは図式
$$
\xymatrix{
& x \times_z y \ar[r]^{p} \ar[d]_q & x \ar[d]^{f} \\
& y \ar[r]^{g} & z }
$$
を 2-可換にする。このことを示す所与の可逆 2-射を
$\psi : f \circ p \to g \circ q$ と記す。
\end{definition}

\noindent
したがって次の普遍性が成り立つ。任意の $w\in \Ob(\mathcal{C})$ と射
$a \in \Mor_{\mathcal C}(w, x)$, $b \in \Mor_\mathcal{C}(w, y)$、および
所与の 2-同型 $\phi : f \circ a \to g\circ b$ に対し、図式
$$
\xymatrix{
w\ar[rrrd]^a \ar@{-->}[rrd]_\gamma \ar[rrdd]_b & & \\
& & x \times_z y \ar[r]_p \ar[d]_q & x \ar[d]^{f} \\
& & y \ar[r]^{g} & z }
$$
を 2-可換にする $\gamma \in \Mor_{\mathcal C}(w, x \times_z y)$ が存在し、
$a \to p \circ \gamma$ および $b \to q \circ \gamma$ を適切に選べば図式
$$
\xymatrix{
f \circ a \ar[r] \ar[d]_\phi &
f \circ p \circ \gamma
\ar[d]^{\psi \star \text{id}_\gamma}
\\
g\circ b
\ar[r]
&
g \circ q \circ \gamma
}
$$
が可換となる。さらに $\gamma$ は同型を除いて一意である。
もちろん、正確な性質はこれより精密である。後で必要となる 2-ファイバー積はすべて、
次に述べる圏の 2-圏における 2-ファイバー積の例から得られる。

\begin{example}
\label{categories-example-2-fibre-product-categories}
$\mathcal{A}$, $\mathcal{B}$, $\mathcal{C}$ を圏とする。
$F : \mathcal{A} \to \mathcal{C}$ および
$G : \mathcal{B} \to \mathcal{C}$ を関手とする。圏
$\mathcal{A} \times_\mathcal{C} \mathcal{B}$ を次のように定義する。
\begin{enumerate}
\item $\mathcal{A} \times_\mathcal{C} \mathcal{B}$ の対象は三つ組
$(A, B, f)$ である。ここで $A\in \Ob(\mathcal{A})$,
$B\in \Ob(\mathcal{B})$ であり、$f : F(A) \to G(B)$ は
$\mathcal{C}$ における同型である。
\item 射 $(A, B, f) \to (A', B', f')$ は対 $(a, b)$ によって与えられる。
ここで $a : A \to A'$ は $\mathcal{A}$ の射、$b : B \to B'$ は
$\mathcal{B}$ の射であり、図式
$$
\xymatrix{
F(A) \ar[r]^f \ar[d]^{F(a)} & G(B) \ar[d]^{G(b)} \\
F(A') \ar[r]^{f'} & G(B')
}
$$
が可換である。
\end{enumerate}
さらに、
$p : \mathcal{A} \times_\mathcal{C}\mathcal{B} \to \mathcal{A}$ および
$q : \mathcal{A} \times_\mathcal{C}\mathcal{B} \to \mathcal{B}$ を
$$
p(A, B, f) = A, \quad q(A, B, f) = B,
$$
と置くことによって関手を定義する。言い換えれば、これらは忘却関手である。
関手の変換 $\psi : F \circ p \to G \circ q$ を定義する。
対象 $\xi = (A, B, f)$ 上では
$\psi_\xi = f : F(p(\xi)) = F(A) \to G(B) = G(q(\xi))$ によって与えられる。
\end{example}

\begin{lemma}
\label{categories-lemma-2-fibre-product-categories}
圏の $(2, 1)$-圏では $2$-ファイバー積が存在し、
Example \ref{categories-example-2-fibre-product-categories} の構成によって与えられる。
\end{lemma}

\begin{proof}
普遍性を確認しよう。$\mathcal{W}$ を圏とし、
$a : \mathcal{W} \to \mathcal{A}$ および
$b : \mathcal{W} \to \mathcal{B}$ を関手とし、
$t : F \circ a \to G \circ b$ を関手の同型とする。

\medskip\noindent
$W \mapsto (a(W), b(W), t_W)$ によって与えられる関手
$\gamma : \mathcal{W} \to \mathcal{A} \times_\mathcal{C}\mathcal{B}$ を考える。
（これが関手であることの確認は省略されている。）
さらに、恒等式 $p \circ \gamma = a$ および $q \circ \gamma = b$ から得られる
$\alpha : a \to p \circ \gamma$ と $\beta : b \to q \circ \gamma$ を考える。
すると $(\gamma, \alpha, \beta)$ が $(W, a, b, t)$ から
$(\mathcal{A} \times_\mathcal{C} \mathcal{B}, p, q, \psi)$ への射であることは明らかである。

\medskip\noindent
第二のそのような射
$(k, \alpha', \beta') :
(W, a, b, t) \to (\mathcal{A} \times_\mathcal{C} \mathcal{B}, p, q, \psi)$
をとる。$\mathcal{W}$ の対象 $W$ に対して
$k(W) = (a_k(W), b_k(W), t_{k, W})$ と書く。したがって
$p(k(W)) = a_k(W)$ などである。写像 $\alpha'$ は関手的な写像
$\alpha' : a(W) \to a_k(W)$ に対応する。圏の $(2, 1)$-圏で考えているので、
実際、各写像 $a(W) \to a_k(W)$ は同型である。これら（および対応する
$b(W) \to b_k(W)$）を用いて同型
$$
\delta_W :
\gamma(W) = (a(W), b(W), t_W)
\longrightarrow
(a_k(W), b_k(W), t_{k, W}) = k(W).
$$
を得る。$\delta$ が、望むように $2$-可換図式の $2$-圏における
$\gamma$ と $k$ の間の $2$-同型を定めることは、直接に示せる。
\end{proof}

\begin{remark}
\label{categories-remark-other-description-2-fibre-product}
$\mathcal{A}$, $\mathcal{B}$, $\mathcal{C}$ を圏とし、
$F : \mathcal{A} \to \mathcal{C}$ および $G : \mathcal{B} \to \mathcal{C}$ を
関手とする。$2$-ファイバー積 $\mathcal{A} \times_\mathcal{C} \mathcal{B}$ の
別の、やや対称的な構成は次のとおりである。対象は五つ組
$(A, B, C, a, b)$ であり、$A, B, C$ は
$\mathcal{A}, \mathcal{B}, \mathcal{C}$ の対象、
$a : F(A) \to C$ と $b : G(B) \to C$ は同型である。
射 $(A, B, C, a, b) \to (A', B', C', a', b')$ は、射の三つ組
$A \to A', B \to B', C \to C'$ で、射 $a, b, a', b'$ と両立するものによって
与えられる。これが $2$-ファイバー積を与えることを直接証明できる。
しかし、写像 $(A, B, C, a, b) \mapsto (A, B, b^{-1} \circ a)$ が
五つ組の圏から Example \ref{categories-example-2-fibre-product-categories} で構成した圏への
同値を与えることに注意する方が容易である。
\end{remark}

\begin{lemma}
\label{categories-lemma-functoriality-2-fibre-product}
圏の $2$-可換図式
$$
\xymatrix{
& \mathcal{Y} \ar[d]_I \ar[rd]^K & \\
\mathcal{X} \ar[r]^H \ar[rd]^L &
\mathcal{Z} \ar[rd]^M & \mathcal{B} \ar[d]^G \\
& \mathcal{A} \ar[r]^F & \mathcal{C}
}
$$
が与えられたとする。同型
$\alpha : G \circ K \to M \circ I$ および
$\beta : M \circ H \to F \circ L$ の選択は、この状況に付随する
$2$-ファイバー積の射
$$
\mathcal{X} \times_\mathcal{Z} \mathcal{Y}
\longrightarrow
\mathcal{A} \times_\mathcal{C} \mathcal{B}
$$
を定める。
\end{lemma}

\begin{proof}
対象上では関手
$$
(X, Y, \phi) \longmapsto (L(X), K(Y),
\alpha^{-1}_Y \circ M(\phi) \circ \beta^{-1}_X)
$$
を、射上では
$$
(a, b) \longmapsto (L(a), K(b))
$$
を用いればよい。
\end{proof}

\begin{lemma}
\label{categories-lemma-equivalence-2-fibre-product}
仮定を Lemma \ref{categories-lemma-functoriality-2-fibre-product} と同じとする。
\begin{enumerate}
\item $K$ と $L$ が忠実ならば、射
$\mathcal{X} \times_\mathcal{Z} \mathcal{Y} \to
\mathcal{A} \times_\mathcal{C} \mathcal{B}$ は忠実である。
\item $K$ と $L$ が充満忠実で、$M$ が忠実ならば、射
$\mathcal{X} \times_\mathcal{Z} \mathcal{Y} \to
\mathcal{A} \times_\mathcal{C} \mathcal{B}$ は充満忠実である。
\item $K$ と $L$ が圏同値で、$M$ が充満忠実ならば、射
$\mathcal{X} \times_\mathcal{Z} \mathcal{Y} \to
\mathcal{A} \times_\mathcal{C} \mathcal{B}$ は圏同値である。
\end{enumerate}
\end{lemma}

\begin{proof}
$(X, Y, \phi)$ と $(X', Y', \phi')$ を
$\mathcal{X} \times_\mathcal{Z} \mathcal{Y}$ の対象とする。
$Z = H(X)$ と置き、$\phi$ によってこれを $I(Y)$ と同一視する。
また、$\alpha_X$ によって $M(Z)$ を $F(L(X))$ と同一視し、
$\beta_Y$ によって $M(Z)$ を $G(K(Y))$ と同一視する。
$Z' = H(X')$ および $M(Z')$ についても同様とする。
射上の写像は
$$
\xymatrix{
\Mor_\mathcal{X}(X, X')
\times_{\Mor_\mathcal{Z}(Z, Z')}
\Mor_\mathcal{Y}(Y, Y')
\ar[d] \\
\Mor_\mathcal{A}(L(X), L(X'))
\times_{\Mor_\mathcal{C}(M(Z), M(Z'))}
\Mor_\mathcal{B}(K(Y), K(Y'))
}
$$
である。したがって (1) と (2) が従う。さらに、$K$ と $L$ が圏同値で、
$M$ が充満忠実ならば、任意の対象 $(A, B, \phi)$ は次の理由で本質像に属する。
$L(X) \cong A$ および $K(Y) \cong B$ を満たす $X$, $Y$ を選ぶ。
すると $M$ の充満忠実性により、同型 $H(X) \cong I(Y)$ を見いだせる。
いくつかの詳細は省略する。
\end{proof}

\begin{lemma}
\label{categories-lemma-associativity-2-fibre-product}
圏と関手の図式
$$
\xymatrix{
\mathcal{A} \ar[rd] & & \mathcal{C} \ar[ld] \ar[rd] & & \mathcal{E} \ar[ld] \\
& \mathcal{B} & & \mathcal{D}
}
$$
が与えられたとする。このとき圏の標準的な同型
$$
(\mathcal{A} \times_\mathcal{B} \mathcal{C}) \times_\mathcal{D} \mathcal{E}
\cong
\mathcal{A} \times_\mathcal{B} (\mathcal{C} \times_\mathcal{D} \mathcal{E})
$$
が存在する。
\end{lemma}

\begin{proof}
意味が明らかであれば、関手
$$
((A, C, \phi), E, \psi)
\longmapsto
(A, (C, E, \psi), \phi)
$$
を用いればよい。
\end{proof}

\noindent
以後、三つ以上の圏のファイバー積を扱うときには括弧を書かない。

\begin{lemma}
\label{categories-lemma-triple-2-fibre-product-pr02}
圏と関手の可換図式
$$
\xymatrix{
\mathcal{A} \ar[rd] & & \mathcal{C} \ar[ld] \ar[rd] & & \mathcal{E} \ar[ld] \\
& \mathcal{B} \ar[rd]_F & & \mathcal{D} \ar[ld]^G \\
& & \mathcal{F} &
}
$$
が与えられたとする。このとき圏の標準的な関手
$$
\text{pr}_{02} :
\mathcal{A} \times_\mathcal{B} \mathcal{C} \times_\mathcal{D} \mathcal{E}
\longrightarrow
\mathcal{A} \times_\mathcal{F} \mathcal{E}
$$
が存在する。
\end{lemma}

\begin{proof}
$\mathcal{A} \times_\mathcal{B} \mathcal{C}
\times_\mathcal{D} \mathcal{E}$ を
$(\mathcal{A} \times_\mathcal{B} \mathcal{C})
\times_\mathcal{D} \mathcal{E}$ と書けば、意味が明らかである限り、関手
$$
((A, C, \phi), E, \psi)
\longmapsto
(A, E, G(\psi) \circ F(\phi))
$$
を用いればよい。
\end{proof}

\begin{lemma}
\label{categories-lemma-2-fibre-product-erase-factor}
圏と関手の図式
$$
\mathcal{A} \to
\mathcal{B} \leftarrow \mathcal{C} \leftarrow \mathcal{D}
$$
が与えられたとする。このとき圏の標準的な同型
$$
\mathcal{A} \times_\mathcal{B} \mathcal{C} \times_\mathcal{C} \mathcal{D}
\cong
\mathcal{A} \times_\mathcal{B} \mathcal{D}
$$
が存在する。
\end{lemma}

\begin{proof}
省略する。
\end{proof}

\noindent
これは、これらの $2$-ファイバー積を通常のファイバー積と同じように扱えることを
意味すると主張する。後で実際に現れる補題をさらにいくつか挙げる。

\begin{lemma}
\label{categories-lemma-diagonal-1}
圏の $2$-ファイバー積
$$
\xymatrix{
\mathcal{C}_3 \ar[r] \ar[d] & \mathcal{S} \ar[d]^\Delta \\
\mathcal{C}_1 \times \mathcal{C}_2 \ar[r]^{G_1 \times G_2} &
\mathcal{S} \times \mathcal{S}
}
$$
が与えられたとする。このとき標準的な同型
$\mathcal{C}_3 \cong
\mathcal{C}_1 \times_{G_1, \mathcal{S}, G_2} \mathcal{C}_2$ が存在する。
\end{lemma}

\begin{proof}
$\mathcal{C}_3$ は Example \ref{categories-example-2-fibre-product-categories} で構成した圏
$(\mathcal{C}_1 \times \mathcal{C}_2)\times_{\mathcal{S} \times \mathcal{S}}
\mathcal{S}$ であると仮定してよい。したがって対象は三つ組
$((X_1, X_2), S, \phi)$ であり、
$\phi = (\phi_1, \phi_2) : (G_1(X_1), G_2(X_2)) \to (S, S)$ は同型である。
そこで、これに三つ組 $(X_1, X_2, \phi_2^{-1} \circ \phi_1)$ を対応させられる。
逆に、$(X_1, X_2, \psi)$ が
$\mathcal{C}_1 \times_{G_1, \mathcal{S}, G_2} \mathcal{C}_2$ の対象ならば、
これに三つ組
$((X_1, X_2), G_2(X_2), (\psi, \text{id}_{G_2(X_2)}))$ を対応させられる。
これらの構成が互いに逆な関手を与えると主張する。
射をどのように扱うかの記述と、これらが互いに逆であることの証明は省略する。
\end{proof}

\begin{lemma}
\label{categories-lemma-diagonal-2}
圏の $2$-ファイバー積
$$
\xymatrix{
\mathcal{C}' \ar[r] \ar[d] & \mathcal{S} \ar[d]^\Delta \\
\mathcal{C} \ar[r]^{(G_1, G_2)} &
\mathcal{S} \times \mathcal{S}
}
$$
が与えられたとする。このとき標準的な同型
$$
\mathcal{C}' \cong
(\mathcal{C} \times_{G_1, \mathcal{S}, G_2} \mathcal{C})
\times_{(p, q), \mathcal{C} \times \mathcal{C}, \Delta}
\mathcal{C}.
$$
が存在する。
\end{lemma}

\begin{proof}
右辺の対象は $((C_1, C_2, \phi), C_3, \psi)$ によって与えられる。
ここで $\phi : G_1(C_1) \to G_2(C_2)$ は同型であり、
$\psi = (\psi_1, \psi_2) : (C_1, C_2) \to (C_3, C_3)$ も同型である。
したがって、これに三つ組
$(C_3, G_1(C_1), (G_1(\psi_1^{-1}), \phi^{-1} \circ G_2(\psi_2^{-1})))$
を対応させられる。これは $\mathcal{C}'$ の対象である。詳細は省略する。
\end{proof}

\begin{lemma}
\label{categories-lemma-fibre-product-after-map}
$\mathcal{A} \to \mathcal{C}$, $\mathcal{B} \to \mathcal{C}$ および
$\mathcal{C} \to \mathcal{D}$ を圏の間の関手とする。このとき図式
$$
\xymatrix{
\mathcal{A} \times_\mathcal{C} \mathcal{B} \ar[d] \ar[r] &
\mathcal{A} \times_\mathcal{D} \mathcal{B} \ar[d] \\
\mathcal{C} \ar[r]^-{\Delta_{\mathcal{C}/\mathcal{D}}} \ar[r] &
\mathcal{C} \times_\mathcal{D} \mathcal{C}
}
$$
は $2$-ファイバー積図式である。
\end{lemma}

\begin{proof}
省略する。
\end{proof}

\begin{lemma}
\label{categories-lemma-base-change-diagonal}
圏の $2$-ファイバー積
$$
\xymatrix{
\mathcal{U} \ar[d] \ar[r] & \mathcal{V} \ar[d] \\
\mathcal{X} \ar[r] & \mathcal{Y}
}
$$
が与えられたとする。このとき図式
$$
\xymatrix{
\mathcal{U} \ar[d] \ar[r] &
\mathcal{U} \times_\mathcal{V} \mathcal{U} \ar[d] \\
\mathcal{X} \ar[r] &
\mathcal{X} \times_\mathcal{Y} \mathcal{X}
}
$$
は $2$-カルテシアンである。
\end{lemma}

\begin{proof}
これは純粋に $2$-圏論的な主張であり、$2$-ファイバー積をもつ任意の
$(2, 1)$-圏で成り立つ。具体的には、次の同値の連鎖から従う。
\begin{align*}
\mathcal{X} \times_{(\mathcal{X} \times_\mathcal{Y} \mathcal{X})}
(\mathcal{U} \times_\mathcal{V} \mathcal{U})
& =
\mathcal{X} \times_{(\mathcal{X} \times_\mathcal{Y} \mathcal{X})}
((\mathcal{X} \times_\mathcal{Y} \mathcal{V})
\times_\mathcal{V} (\mathcal{X} \times_\mathcal{Y} \mathcal{V})) \\
& =
\mathcal{X} \times_{(\mathcal{X} \times_\mathcal{Y} \mathcal{X})}
(\mathcal{X} \times_\mathcal{Y} \mathcal{X}
\times_\mathcal{Y} \mathcal{V}) \\
& =
\mathcal{X} \times_\mathcal{Y} \mathcal{V} = \mathcal{U}
\end{align*}
Lemmas \ref{categories-lemma-associativity-2-fibre-product} および
\ref{categories-lemma-2-fibre-product-erase-factor} を参照。
\end{proof}

\section{圏上の圏}
\label{categories-section-categories-over-categories}

\noindent
本節では関手 $p : \mathcal{S} \to \mathcal{C}$ を考える。
$\mathcal{S}$ を上側に、$\mathcal{C}$ を下側にあるものと考える。
何を論じているのかを明確にするため、$\mathcal{C}$ 上の圏からなる
$2$-圏を定義する。

\begin{definition}
\label{categories-definition-categories-over-C}
$\mathcal{C}$ を圏とする。{\it $\mathcal{C}$ 上の圏の $2$-圏}とは、
次のように定義される $2$-圏である。
\begin{enumerate}
\item その対象は関手 $p : \mathcal{S} \to \mathcal{C}$ である。
\item その $1$-射 $(\mathcal{S}, p) \to (\mathcal{S}', p')$ は、
$p' \circ G = p$ を満たす関手 $G : \mathcal{S} \to \mathcal{S}'$ である。
\item
$G, H : (\mathcal{S}, p) \to (\mathcal{S}', p')$ に対するその $2$-射
$t : G \to H$ は、すべての $x \in \Ob(\mathcal{S})$ について
$p'(t_x) = \text{id}_{p(x)}$ を満たす関手の射である。
\end{enumerate}
この状況では
$$
\Mor_{\textit{Cat}/\mathcal{C}}(\mathcal{S}, \mathcal{S}')
$$
によって $(\mathcal{S}, p)$ と $(\mathcal{S}', p')$ の間の
$1$-射からなる圏を表す。
\end{definition}

\noindent
この $2$-圏における水平合成と垂直合成は、
Section \ref{categories-section-formal-cat-cat} で $\textit{Cat}$ に対して行ったのと
まったく同様に定義する。$2$-圏の公理は、$\textit{Cat}$ におけるのと
同じ理由で満たされる。これを確認するには、$\textit{Cat}$ で公理が
成り立つことを用い、「$\mathcal{C}$ 上の $2$-射の垂直合成が再び
$\mathcal{C}$ 上の $2$-射を与える」といった事柄を検証してもよい。
これは明らかである。

\medskip\noindent
空間の写像のファイバーと同様に、ファイバー圏という概念、およびこの状況に
付随するいくつかの持ち上げの概念がある。

\begin{definition}
\label{categories-definition-fibre-category}
$\mathcal{C}$ を圏とし、$p : \mathcal{S} \to \mathcal{C}$ を
$\mathcal{C}$ 上の圏とする。
\begin{enumerate}
\item 対象 $U\in \Ob(\mathcal{C})$ 上の {\it ファイバー圏}とは、対象が
$$
\Ob(\mathcal{S}_U) = \{x\in \Ob(\mathcal{S}) :
p(x) = U\}
$$
であり、射が
$$
\Mor_{\mathcal{S}_U}(x, y) = \{ \phi \in \Mor_\mathcal{S}(x, y) :
p(\phi) = \text{id}_U\}.
$$
である圏 $\mathcal{S}_U$ をいう。
\item 対象 $U \in \Ob(\mathcal{C})$ の {\it 持ち上げ}とは、
$p(x) = U$ を満たす対象 $x\in \Ob(\mathcal{S})$、すなわち
$x\in \Ob(\mathcal{S}_U)$ をいう。また、{\it $x$ は $U$ 上にある}ともいう。
\item 同様に、$\mathcal{C}$ における射 $f : V \to U$ の
{\it 持ち上げ}とは、$p(\phi) = f$ を満たす $\mathcal{S}$ の射
$\phi : y \to x$ をいう。{\it $\phi$ は $f$ 上にある}ともいう。
\end{enumerate}
\end{definition}

\noindent
ここでいくつか注意できる。例えば、
$F : (\mathcal{S}, p) \to (\mathcal{S}', p')$ が $\mathcal{C}$ 上の圏の
$1$-射ならば、$F$ はファイバー圏の関手
$F : \mathcal{S}_U \to \mathcal{S}'_U$ を誘導する。$2$-射についても同様である。

\medskip\noindent
次は、$\mathcal{C}$ 上の圏からなる $(2, 1)$-圏における
$2$-ファイバー積を記述する標準的な補題である。

\begin{lemma}
\label{categories-lemma-2-product-categories-over-C}
$\mathcal{C}$ を圏とする。$\mathcal{C}$ 上の圏からなる $(2, 1)$-圏は
2-ファイバー積をもつ。$F : \mathcal{X} \to \mathcal{S}$ および
$G : \mathcal{Y} \to \mathcal{S}$ を $\mathcal{C}$ 上の圏の射とする。
具体的な 2-ファイバー積 $\mathcal{X} \times_\mathcal{S}\mathcal{Y}$ は
次のように記述される。
\begin{enumerate}
\item $\mathcal{X} \times_\mathcal{S} \mathcal{Y}$ の対象は四つ組
$(U, x, y, f)$ である。ここで $U \in \Ob(\mathcal{C})$,
$x\in \Ob(\mathcal{X}_U)$, $y\in \Ob(\mathcal{Y}_U)$ であり、
$f : F(x) \to G(y)$ は $\mathcal{S}_U$ における同型である。
\item 射 $(U, x, y, f) \to (U', x', y', f')$ は対 $(a, b)$ によって
与えられる。ここで $a : x \to x'$ は $\mathcal{X}$ の射、
$b : y \to y'$ は $\mathcal{Y}$ の射であり、次を満たす。
\begin{enumerate}
\item $a$ と $b$ は同じ射 $U \to U'$ を誘導する。
\item 図式
$$
\xymatrix{
F(x) \ar[r]^f \ar[d]^{F(a)} & G(y) \ar[d]^{G(b)} \\
F(x') \ar[r]^{f'} & G(y')
}
$$
は可換である。
\end{enumerate}
\end{enumerate}
この場合、関手 $p : \mathcal{X} \times_\mathcal{S}\mathcal{Y} \to \mathcal{X}$
および $q : \mathcal{X} \times_\mathcal{S}\mathcal{Y} \to \mathcal{Y}$ は
忘却関手である。変換 $\psi : F \circ p \to G \circ q$ は、対象
$\xi = (U, x, y, f)$ 上で
$\psi_\xi = f : F(p(\xi)) = F(x) \to G(y) = G(q(\xi))$ によって与えられる。
\end{lemma}

\begin{proof}
普遍性を確認しよう。$p_\mathcal{W} : \mathcal{W}\to \mathcal{C}$ を
$\mathcal{C}$ 上の圏とし、$X : \mathcal{W} \to \mathcal{X}$ および
$Y : \mathcal{W} \to \mathcal{Y}$ を $\mathcal{C}$ 上の関手とし、
$t : F \circ X \to G \circ Y$ を $\mathcal{C}$ 上の関手の同型とする。
求める関手
$\gamma : \mathcal{W} \to \mathcal{X} \times_\mathcal{S} \mathcal{Y}$ は
$W \mapsto (p_\mathcal{W}(W), X(W), Y(W), t_W)$ によって与えられる。
詳細は省略する。Lemma \ref{categories-lemma-2-fibre-product-categories} と比較せよ。
\end{proof}

\begin{example}
\label{categories-example-different-2-fibre-products}
Lemma \ref{categories-lemma-2-product-categories-over-C} で与えた圏上の圏の
$2$-ファイバー積の構成と、Lemma \ref{categories-lemma-2-fibre-product-categories}
（Example \ref{categories-example-2-fibre-product-categories} の構成）で与えた圏の
2-ファイバー積の構成とは、一般には同値でない出力を生じる。
実際、$\mathcal{S}$ を一つの対象と二つの矢をもつグルーポイド圏とし、
$\mathcal{X}$ を一つの対象をもつ離散圏とする。圏として
$2$-ファイバー積 $\mathcal{X} \times_\mathcal{S} \mathcal{X}$ をとると、
二つの対象をもつ離散圏が得られる。しかし、これらをすべて
$\mathcal{S}$ 上の圏とみなすと、$\mathcal{S}$ 上の圏としての
$2$-ファイバー積 $\mathcal{X} \times_\mathcal{S} \mathcal{X}$ は、
一つの対象をもつ離散圏となる。相違は、
Lemma \ref{categories-lemma-2-product-categories-over-C} の記法で、最初の状況では
比較同型 $f$ を任意に選べたのに対し、第二の状況では恒等射しか
選べなかったことにある。
\end{example}

\begin{lemma}
\label{categories-lemma-fibre-2-fibre-product-categories-over-C}
$\mathcal{C}$ を圏とする。$f : \mathcal{X} \to \mathcal{S}$ および
$g : \mathcal{Y} \to \mathcal{S}$ を $\mathcal{C}$ 上の圏の射とする。
$\mathcal{C}$ の任意の対象 $U$ に対して、ファイバー圏の等式
$$
\left(\mathcal{X} \times_\mathcal{S}\mathcal{Y}\right)_U
=
\mathcal{X}_U \times_{\mathcal{S}_U} \mathcal{Y}_U
$$
が成り立つ。
\end{lemma}

\begin{proof}
省略する。
\end{proof}

\section{ファイバー圏}
\label{categories-section-fibred-categories}

\noindent
ファイバー圏についてごく簡単に論じておく必要がある。

\medskip\noindent
$p : \mathcal{S} \to \mathcal{C}$ を $\mathcal{C}$ 上の圏とする。
$p(x) = U$ を満たす対象 $x \in \mathcal{S}$ と射 $f : V \to U$ が
与えられたとき、ある種の「ファイバー積 $V \times_U x$」
（または $V \to U$ による $x$ の {\it 基底変換}）をとろうとすることができる。
すなわち、対象 $z \in \mathcal{S}$ から「$V \times_U x$」への射は、
$p(\varphi) = f \circ g$ を満たす対 $(\varphi, g)$ によって与えられるべきである。
ここで $\varphi : z \to x$, $g : p(z) \to V$ である。図示すると
$$
\xymatrix{
z \ar@{~>}[d]^p \ar@{-}[r] &
? \ar[r] \ar@{~>}[d]^p &
x \ar@{~>}[d]^p \\
p(z) \ar[r] & V \ar[r]^f & U
}
$$
となる。このような射 $V \times_U x \to x$ が存在するとき、
それを強カルテシアン射と呼ぶ。

\begin{definition}
\label{categories-definition-cartesian-over-C}
$\mathcal{C}$ を圏とし、$p : \mathcal{S} \to \mathcal{C}$ を
$\mathcal{C}$ 上の圏とする。{\it 強カルテシアン射}、より正確には
{\it 強 $\mathcal{C}$-カルテシアン射}とは、$\mathcal{S}$ の射
$\varphi : y \to x$ であって、すべての $z \in \Ob(\mathcal{S})$ に対して写像
$$
\Mor_\mathcal{S}(z, y)
\longrightarrow
\Mor_\mathcal{S}(z, x)
\times_{\Mor_\mathcal{C}(p(z), p(x))}
\Mor_\mathcal{C}(p(z), p(y)),
$$
が全単射となるものをいう。この写像は
$\psi \longmapsto (\varphi \circ \psi, p(\psi))$ によって与えられる。
\end{definition}

\noindent
Yoneda Lemma \ref{categories-lemma-yoneda} により、$U \in \Ob(\mathcal{C})$ 上にある
$x \in \Ob(\mathcal{S})$ と $\mathcal{C}$ の射 $f : V \to U$ が与えられ、
$p(\varphi) = f$ を満たす強カルテシアン射 $\varphi : y \to x$ が
存在するならば、$(y, \varphi)$ は一意な同型を除いて一意である。
これは上の定義から明らかである。実際、関手
$$
z
\longmapsto
\Mor_\mathcal{S}(z, x)
\times_{\Mor_\mathcal{C}(p(z), U)}
\Mor_\mathcal{C}(p(z), V)
$$
はデータ $(x, U, f : V \to U)$ のみに依存する。したがって、$f$ の
持ち上げである強カルテシアン射を表すために、しばしば
$V \times_U x \to x$ または $f^*x \to x$ を用いる。

\begin{lemma}
\label{categories-lemma-composition-cartesian}
$\mathcal{C}$ を圏とし、$p : \mathcal{S} \to \mathcal{C}$ を
$\mathcal{C}$ 上の圏とする。
\begin{enumerate}
\item 二つの強カルテシアン射の合成は強カルテシアンである。
\item $\mathcal{S}$ の任意の同型は強カルテシアンである。
\item $p(\varphi)$ が同型であるような任意の強カルテシアン射
$\varphi$ は同型である。
\end{enumerate}
\end{lemma}

\begin{proof}
(1) の証明。$\varphi : y \to x$ と $\psi : z \to y$ を
強カルテシアンとする。$t$ を $\mathcal{S}$ の任意の対象とする。このとき
\begin{align*}
& \Mor_\mathcal{S}(t, z) \\
& =
\Mor_\mathcal{S}(t, y)
\times_{\Mor_\mathcal{C}(p(t), p(y))}
\Mor_\mathcal{C}(p(t), p(z)) \\
& =
\Mor_\mathcal{S}(t, x)
\times_{\Mor_\mathcal{C}(p(t), p(x))}
\Mor_\mathcal{C}(p(t), p(y))
\times_{\Mor_\mathcal{C}(p(t), p(y))}
\Mor_\mathcal{C}(p(t), p(z)) \\
& =
\Mor_\mathcal{S}(t, x)
\times_{\Mor_\mathcal{C}(p(t), p(x))}
\Mor_\mathcal{C}(p(t), p(z))
\end{align*}
を得るので、$z \to x$ は強カルテシアンである。

\medskip\noindent
(2) の証明。$y \to x$ を同型とする。このとき $p(y) \to p(x)$ も同型である。
したがって
$$
\Mor_\mathcal{C}(p(z), p(y)) \to
\Mor_\mathcal{C}(p(z), p(x))
$$
は全単射である。ゆえに
$\Mor_\mathcal{S}(z, x)
\allowbreak\times_{\Mor_\mathcal{C}(p(z), p(x))}
\Mor_\mathcal{C}(p(z), p(y))$ は
$\Mor_\mathcal{S}(z, x)$ と全単射である。
また $y \to x$ が同型なので、Definition \ref{categories-definition-cartesian-over-C}
の表示された写像は全単射であり、$y \to x$ は強カルテシアンである。

\medskip\noindent
(3) の証明。$\varphi : y \to x$ を強カルテシアンとし、
$p(\varphi) : p(y) \to p(x)$ が同型であると仮定する。
$z = x$ として定義を適用すると、$(\text{id}_x, p(\varphi)^{-1})$ は
一意な射 $\chi : x \to y$ から得られる。$\chi$ が $\varphi$ の逆射であることの
確認は省略する。
\end{proof}

\begin{lemma}
\label{categories-lemma-cartesian-over-cartesian}
$F : \mathcal{A} \to \mathcal{B}$ および $G : \mathcal{B} \to \mathcal{C}$ を
圏の間の合成可能な関手とする。$x \to y$ を $\mathcal{A}$ の射とする。
$x \to y$ が強 $\mathcal{B}$-カルテシアンで、$F(x) \to F(y)$ が
強 $\mathcal{C}$-カルテシアンならば、$x \to y$ は
強 $\mathcal{C}$-カルテシアンである。
\end{lemma}

\begin{proof}
これは定義から直接従う。
\end{proof}

\begin{lemma}
\label{categories-lemma-strongly-cartesian-fibre-product}
$\mathcal{C}$ を圏とし、$p : \mathcal{S} \to \mathcal{C}$ を
$\mathcal{C}$ 上の圏とする。$x \to y$ と $z \to y$ を
$\mathcal{S}$ の射とする。次を仮定する。
\begin{enumerate}
\item $x \to y$ は強カルテシアンである。
\item $p(x) \times_{p(y)} p(z)$ が存在する。
\item $p(w) = p(x) \times_{p(y)} p(z)$ および
$p(a) = \text{pr}_2 : p(x) \times_{p(y)} p(z) \to p(z)$ を満たす
$\mathcal{S}$ の強カルテシアン射 $a : w \to z$ が存在する。
\end{enumerate}
このときファイバー積 $x \times_y z$ が存在し、$w$ と同型である。
\end{lemma}

\begin{proof}
$x \to y$ が強カルテシアンなので、$p(b) = \text{pr}_1$ を満たす
一意な射 $b : w \to x$ が存在する。$w$ がファイバー積であることを見るため、
次を計算する。
\begin{align*}
& \Mor_\mathcal{S}(t, w) \\
& = \Mor_\mathcal{S}(t, z)
\times_{\Mor_\mathcal{C}(p(t), p(z))}
\Mor_\mathcal{C}(p(t), p(w)) \\
& = \Mor_\mathcal{S}(t, z)
\times_{\Mor_\mathcal{C}(p(t), p(z))}
(\Mor_\mathcal{C}(p(t), p(x))
\times_{\Mor_\mathcal{C}(p(t), p(y))}
\Mor_\mathcal{C}(p(t), p(z))) \\
& = \Mor_\mathcal{S}(t, z)
\times_{\Mor_\mathcal{C}(p(t), p(y))}
\Mor_\mathcal{C}(p(t), p(x)) \\
& = \Mor_\mathcal{S}(t, z)
\times_{\Mor_\mathcal{S}(t, y)}
\Mor_\mathcal{S}(t, y)
\times_{\Mor_\mathcal{C}(p(t), p(y))}
\Mor_\mathcal{C}(p(t), p(x)) \\
& = \Mor_\mathcal{S}(t, z)
\times_{\Mor_\mathcal{S}(t, y)}
\Mor_\mathcal{S}(t, x)
\end{align*}
となり、望む結論を得る。最初の等式は $a : w \to z$ が強カルテシアンで
あることから、最後の等式は $x \to y$ が強カルテシアンであることから従う。
\end{proof}

\begin{definition}
\label{categories-definition-fibred-category}
$\mathcal{C}$ を圏とする。
$p : \mathcal{S} \to \mathcal{C}$ を $\mathcal{C}$ 上の圏とする。
任意の $x \in \Ob(\mathcal{S})$ が $U \in \Ob(\mathcal{C})$ 上にあり、
$\mathcal{C}$ の任意の射 $f : V \to U$ が与えられたとき、$f$ 上にある
強カルテシアン射 $f^*x \to x$ が存在するならば、$\mathcal{S}$ を
{\it $\mathcal{C}$ 上のファイバー圏}という。
\end{definition}

\noindent
$p : \mathcal{S} \to \mathcal{C}$ をファイバー圏と仮定する。
定義における各 $f : V \to U$ と $x\in \Ob(\mathcal{S}_U)$ に対し、
$f$ 上にある強カルテシアン射 $f^\ast x \to x$ を一つ選ぶことができる。
選択公理により、すべての $f: V \to U = p(x)$ に対して
$f^*x \to x$ を同時に選ぶことができる。
$\mathcal{S}_U$ の各射 $\phi : x \to x'$ と $f : V \to U$ に対し、
$\mathcal{S}_V$ の一意な射 $f^\ast \phi : f^\ast x \to f^\ast x'$ であって
図式
$$
\xymatrix{
f^\ast x \ar[r]_{f^\ast \phi} \ar[d] & f^\ast x' \ar[d] \\
x \ar[r]^{\phi} & x' }
$$
を可換にするものが存在すると主張する。実際、射 $f^*x' \to x'$ が
強カルテシアンであることから、この矢は存在して一意である。この矢の一意性により、
（射についても定義された）$f^\ast$ は関手
$ f^\ast : \mathcal{S}_U \to \mathcal{S}_V$ となる。

\begin{definition}
\label{categories-definition-pullback-functor-fibred-category}
$p : \mathcal{S} \to \mathcal{C}$ をファイバー圏と仮定する。
\begin{enumerate}
\item $p : \mathcal{S} \to \mathcal{C}$ に対する{\it 引き戻しの選択}
\footnote{この用語はおそらく標準的ではない。文献によってはこれを
``cleavage'' と呼ぶが、誤った連想を招く。おそらく ``cleaving'' の方が
よい語であろう。関連する概念に ``splitting'' があるが、多くの文献では
``splitting'' は、任意の合成可能な射の対に対して
$g^*f^* = (f \circ g)^*$ となるような引き戻しの選択を意味する。
Definition \ref{categories-definition-split-fibred-category} も比較せよ。}
とは、$\mathcal{C}$ の任意の射 $f: V \to U$ と任意の
$x \in \Ob(\mathcal{S}_U)$ に対し、$f$ 上にある強カルテシアン射
$f^\ast x \to x$ を選ぶことである。
\item 引き戻しの選択が与えられたとき、$\mathcal{C}$ の任意の射
$f : V \to U$ に対し、上で述べた関手
$f^* : \mathcal{S}_U \to \mathcal{S}_V$ を（上で行った選択
$f^*x \to x$ に付随する）{\it 引き戻し関手}という。
\end{enumerate}
\end{definition}

\noindent
もちろん、引き戻しの選択が $\text{id}_U^*x = x$ を満たすと常に仮定できる。
もっとも、引き戻しにさらに仮定を課さなければ、実際上これは役に立たない性質である。

\begin{lemma}
\label{categories-lemma-fibred}
$p : \mathcal{S} \to \mathcal{C}$ をファイバー圏と仮定し、
$p : \mathcal{S} \to \mathcal{C}$ に対する引き戻しの選択が与えられているとする。
\begin{enumerate}
\item 任意の合成可能な射 $f : V \to U$, $g : W \to V$ の対に対し、
関手 $\mathcal{S}_U \to \mathcal{S}_W$ の一意な同型
$$
\alpha_{g, f} :
(f \circ g)^\ast
\longrightarrow
g^\ast \circ f^\ast
$$
であって、すべての $y\in \Ob(\mathcal{S}_U)$ に対し次の図式を
可換にするものが存在する。
$$
\xymatrix{
g^\ast f^\ast y \ar[r]
&
f^\ast y \ar[d] \\
(f \circ g)^\ast y \ar[r]
\ar[u]^{(\alpha_{g, f})_y}
&
y
}
$$
\item $f = \text{id}_U$ ならば、関手 $\mathcal{S}_U \to \mathcal{S}_U$
の標準的な同型 $\alpha_U : \text{id} \to (\text{id}_U)^*$ が存在する。
\item 四つ組
$(U \mapsto \mathcal{S}_U, f \mapsto f^*, \alpha_{g, f}, \alpha_U)$
は $\mathcal{C}^{opp}$ から圏の $(2, 1)$-圏への擬関手を定める。
Definition \ref{categories-definition-functor-into-2-category} を参照せよ。
\end{enumerate}
\end{lemma}

\begin{proof}
実際、(1) の可換図式がファイバー圏 $\mathcal{S}_W$ の射
$(\alpha_{g, f})_y$ を一意に定めることは明らかである。
射 $(f \circ g)^*y \to y$ と合成
$g^*f^*y \to f^*y \to y$ はいずれも $f \circ g$ を持ち上げる
強カルテシアン射なので、これは同型である（Definition
\ref{categories-definition-cartesian-over-C} の後の議論および
Lemma \ref{categories-lemma-composition-cartesian} を参照）。同様に、
$\text{id}_x : x \to x$ は明らかに $\text{id}_U$ 上で強カルテシアン
（ここで $U = p(x)$）なので、同型
$(\alpha_U)_x : x \to (\text{id}_U)^*x$ が存在する。
（もちろん、$f$ が恒等射ならば常に $f^*x = x$ であるとあらかじめ
仮定することもできたが、一般性のためには仮定しない方がよい。）
このようにして得られた $\alpha_{g, f}$ と $\alpha_U$ が関手の変換であることの
確認は省略する。また (3) の確認も省略する。
\end{proof}

\begin{lemma}
\label{categories-lemma-fibred-equivalent}
$\mathcal{C}$ を圏とする。$\mathcal{S}_1$, $\mathcal{S}_2$ を
$\mathcal{C}$ 上の圏とする。$\mathcal{S}_1$ と $\mathcal{S}_2$ が
$\mathcal{C}$ 上の圏として同値であると仮定する。このとき、
$\mathcal{S}_1$ が $\mathcal{C}$ 上でファイバー化されていることと、
$\mathcal{S}_2$ が $\mathcal{C}$ 上でファイバー化されていることは同値である。
\end{lemma}

\begin{proof}
与えられた関手を $p_i : \mathcal{S}_i \to \mathcal{C}$ と記す。
$F : \mathcal{S}_1 \to \mathcal{S}_2$,
$G : \mathcal{S}_2 \to \mathcal{S}_1$ を $\mathcal{C}$ 上の関手とし、
$i : F \circ G \to \text{id}_{\mathcal{S}_2}$,
$j : G \circ F \to \text{id}_{\mathcal{S}_1}$ を
$\mathcal{C}$ 上の関手の同型とする。このとき $F$ は強カルテシアン射を
強カルテシアン射に写すと主張する。実際、$\varphi : y \to x$ を
$\mathcal{S}_1$ の強カルテシアン射とする。
$f : V \to U$ を $p_1(\varphi)$ とおく。
$z' \in \Ob(\mathcal{S}_2)$ とし、$W = p_2(z')$ とおく。また、
$p_2(\psi') = f \circ g$ を満たす $g : W \to V$ と
$\psi' : z' \to F(x)$ が与えられているとする。このとき
$$
\psi = j \circ G(\psi') : G(z') \to G(F(x)) \to x
$$
は $p_1(\psi) = f \circ g$ を満たす $\mathcal{S}_1$ の射である。
したがって仮定により、$g$ 上にある一意な射 $\xi : G(z') \to y$ であって
$\psi = \varphi \circ \xi$ を満たすものが存在する。これはさらに、$g$ 上にあり
$\psi' = F(\varphi) \circ \xi'$ を満たす射
$$
\xi' = F(\xi) \circ i^{-1} : z' \to F(G(z')) \to F(y)
$$
を与える。$\xi'$ が一意であることの確認は省略する。
\end{proof}

\noindent
Lemma \ref{categories-lemma-fibred-equivalent} から、同値は強カルテシアン射を
強カルテシアン射に写すことが分かる。しかし、$\mathcal{C}$ 上の
ファイバー圏の間の任意の関手については、これは成り立つとは限らない。
そこでファイバー圏の $2$-圏を次のように定義する。

\begin{definition}
\label{categories-definition-fibred-categories-over-C}
$\mathcal{C}$ を圏とする。{\it $\mathcal{C}$ 上のファイバー圏の $2$-圏}
とは、$\mathcal{C}$ 上の圏の $2$-圏（Definition
\ref{categories-definition-categories-over-C} を参照）の部分 $2$-圏で、
次のように定められるものである。
\begin{enumerate}
\item その対象はファイバー圏 $p : \mathcal{S} \to \mathcal{C}$ である。
\item その $1$-射 $(\mathcal{S}, p) \to (\mathcal{S}', p')$ は、
$p' \circ G = p$ を満たし、$G$ が強カルテシアン射を強カルテシアン射に写す
関手 $G : \mathcal{S} \to \mathcal{S}'$ である。
\item $G, H : (\mathcal{S}, p) \to (\mathcal{S}', p')$ に対する
$2$-射 $t : G \to H$ は、すべての $x \in \Ob(\mathcal{S})$ に対して
$p'(t_x) = \text{id}_{p(x)}$ を満たす関手の射である。
\end{enumerate}
この状況では、$(\mathcal{S}, p)$ と $(\mathcal{S}', p')$ の間の
$1$-射の圏を
$$
\Mor_{\textit{Fib}/\mathcal{C}}(\mathcal{S}, \mathcal{S}')
$$
と記す。
\end{definition}

\noindent
$1$-射に課された条件に注意せよ。また、これは真の $2$-圏であって
$(2, 1)$-圏ではないことにも注意せよ。したがって $2$-ファイバー積を
とるときは、まず付随する $(2, 1)$-圏に移る。

\begin{lemma}
\label{categories-lemma-2-product-fibred-categories-over-C}
$\mathcal{C}$ を圏とする。$\mathcal{C}$ 上のファイバー圏の
$(2, 1)$-圏は 2-ファイバー積をもち、それらは
Lemma \ref{categories-lemma-2-product-categories-over-C} のように記述される。
\end{lemma}

\begin{proof}
ここで本質的に示すべきことは、$\mathcal{C}$ 上のファイバー圏の射
$F : \mathcal{X} \to \mathcal{S}$ と $G : \mathcal{Y} \to \mathcal{S}$
が与えられたとき、Lemma \ref{categories-lemma-2-product-categories-over-C} で述べた圏
$\mathcal{X} \times_\mathcal{S} \mathcal{Y}$ がファイバー化されていることである。
$\mathcal{X} \times_\mathcal{S} \mathcal{Y}$ が十分多くの強カルテシアン射を
もつことを示そう。実際、$(U, x, y, \phi)$ を
$\mathcal{X} \times_\mathcal{S} \mathcal{Y}$ の対象とし、
$f : V \to U$ を $\mathcal{C}$ の射とする。
$\mathcal{X}$ において $f$ 上にある強カルテシアン射
$a : f^*x \to x$ と、$\mathcal{Y}$ において $f$ 上にある
強カルテシアン射 $b : f^*y \to y$ を選ぶ。
仮定により $F(a)$ と $G(b)$ は強カルテシアンである。
$\phi : F(x) \to G(y)$ は同型なので、強カルテシアン射の一意性により、
$G(b) \circ f^*\phi = \phi \circ F(a)$ を満たす一意な同型
$f^*\phi : F(f^*x) \to G(f^*y)$ が得られる。言い換えれば、
$(a, b) : (V, f^*x, f^*y, f^*\phi) \to (U, x, y, \phi)$ は
$\mathcal{X} \times_\mathcal{S} \mathcal{Y}$ の射である。
これが強カルテシアン射であること（そして実際にはこれらが強カルテシアン射の
すべてであること）の確認は省略する。
\end{proof}

\begin{lemma}
\label{categories-lemma-cute}
$\mathcal{C}$ を圏とし、$U \in \Ob(\mathcal{C})$ とする。
$p : \mathcal{S} \to \mathcal{C}$ がファイバー圏であり、$p$ が
$p' : \mathcal{S} \to \mathcal{C}/U$ を経由して分解するならば、
$p' : \mathcal{S} \to \mathcal{C}/U$ はファイバー圏である。
\end{lemma}

\begin{proof}
$\varphi : x' \to x$ が $p$ に関して強カルテシアンであると仮定する。
$\varphi$ は $p'$ に関しても強カルテシアンであると主張する。
$g = p'(\varphi)$ とおく。このとき、ある射 $f : V \to U$ と
$f' : V' \to U$ に対し $g : V'/U \to V/U$ である。
$z \in \Ob(\mathcal{S})$ とし、$p'(z) = (W \to U)$ とおく。
$\varphi$ が $p'$ に関して強カルテシアンであることを示すには、
$\psi' \longmapsto (\varphi \circ \psi', p'(\psi'))$ によって与えられる写像
$$
\Mor_\mathcal{S}(z, x')
\longrightarrow
\Mor_\mathcal{S}(z, x)
\times_{\Mor_{\mathcal{C}/U}(W/U, V/U)}
\Mor_{\mathcal{C}/U}(W/U, V'/U),
$$
が全単射であることを示せばよい。右辺の元 $(\psi, h)$ が与えられたとする。
すると特に $g \circ h = p(\psi)$ であり、$\varphi$ が強カルテシアンである
という条件により、$\psi = \varphi \circ \psi'$ および
$p(\psi') = h$ を満たす一意な射 $\psi' : z \to x'$ を得る。
さて $p'(\psi') : W/U \to V/U$ は、対応する写像 $W \to V$ が $h$ である
射なので、$\mathcal{C}/U$ の射として $h$ に等しい。したがって $\psi'$ は、
与えられた対 $(\psi, h)$ に写る一意な射 $z \to x'$ である。これで主張が示された。

\medskip\noindent
最後に、$g : V'/U \to V/U$ と $p'(x) = V/U$ を満たす $x$ が
与えられたとする。$p : \mathcal{S} \to \mathcal{C}$ はファイバー圏なので、
$p(\varphi) = g$ を満たす強カルテシアン射 $\varphi : x' \to x$ が存在する。
上と同じ議論により $p'(\varphi) = g : V'/U \to V/U$ が従う。
また上で見たように、射 $\varphi$ は強カルテシアンである。
したがって Definition \ref{categories-definition-fibred-category} の条件が満たされ、
証明が完了する。
\end{proof}

\begin{lemma}
\label{categories-lemma-fibred-over-fibred}
$\mathcal{A} \to \mathcal{B} \to \mathcal{C}$ を圏の間の関手とする。
$\mathcal{A}$ が $\mathcal{B}$ 上でファイバー化され、$\mathcal{B}$ が
$\mathcal{C}$ 上でファイバー化されているならば、$\mathcal{A}$ は
$\mathcal{C}$ 上でファイバー化されている。
\end{lemma}

\begin{proof}
これは定義と Lemma \ref{categories-lemma-cartesian-over-cartesian} から従う。
\end{proof}

\begin{lemma}
\label{categories-lemma-fibred-category-representable-goes-up}
$p : \mathcal{S} \to \mathcal{C}$ をファイバー圏とする。
$x \to y$ と $z \to y$ を $\mathcal{S}$ の射とし、$x \to y$ は
強カルテシアンであるとする。$p(x) \times_{p(y)} p(z)$ が存在するならば、
$x \times_y z$ が存在し、$p(x \times_y z) = p(x) \times_{p(y)} p(z)$ であり、
$x \times_y z \to z$ は強カルテシアンである。
\end{lemma}

\begin{proof}
$\text{pr}_2 : p(x) \times_{p(y)} p(z) \to p(z)$ 上にある
強カルテシアン射 $\text{pr}_2^*z \to z$ を選ぶ。このとき
Lemma \ref{categories-lemma-strongly-cartesian-fibre-product} により
$\text{pr}_2^*z = x \times_y z$ である。
\end{proof}

\begin{lemma}
\label{categories-lemma-ameliorate-morphism-fibred-categories}
$\mathcal{C}$ を圏とし、$F : \mathcal{X} \to \mathcal{Y}$ を
$\mathcal{C}$ 上のファイバー圏の $1$-射とする。
$\mathcal{C}$ 上のファイバー圏の $1$-射
$$
\xymatrix{
\mathcal{X} \ar@<1ex>[r]^u &
\mathcal{X}' \ar[r]^v \ar@<1ex>[l]^w & \mathcal{Y}
}
$$
であって $F = v \circ u$ を満たし、さらに次を満たすものが存在する。
\begin{enumerate}
\item $u : \mathcal{X} \to \mathcal{X}'$ は充満忠実である。
\item $w$ は $u$ の左随伴である。
\item $v : \mathcal{X}' \to \mathcal{Y}$ はファイバー圏である。
\end{enumerate}
\end{lemma}

\begin{proof}
構造関手を $p : \mathcal{X} \to \mathcal{C}$ および
$q : \mathcal{Y} \to \mathcal{C}$ と記す。$\mathcal{X}'$ を次のように
具体的に構成する。$\mathcal{X}'$ の対象は四つ組 $(U, x, y, f)$ であり、
$x \in \Ob(\mathcal{X}_U)$, $y \in \Ob(\mathcal{Y}_U)$ で、
$f : y \to F(x)$ は $\mathcal{Y}_U$ の射である。
射 $(a, b) : (U, x, y, f) \to (U', x', y', f')$ は、
$p(a) = q(b) : U \to U'$ および $f' \circ b = F(a) \circ f$ を満たす
$a : x \to x'$ と $b : y \to y'$ によって与えられる。

\medskip\noindent
$p$ と $q$ の双方に対する引き戻しの選択を行い、それらを同じ記法で表す。
$(U, x, y, f)$ を対象とし、$h : V \to U$ を射とする。
与えられた強カルテシアン射 $h^*x \to x$ と $h^*y \to y$ から得られる射
$c : (V, h^*x, h^*y, h^*f) \to (U, x, y, f)$ を考える。
$c$ は $\mathcal{C}$ 上の $\mathcal{X}'$ において強カルテシアンであると主張する。
実際、$(W, x', y', f')$ を $\mathcal{X}'$ の対象とし、
$(a, b) : (W, x', y', f') \to (U, x, y, f)$ を $W \to U$ 上にある射とし、
$W \to U$ の $h$ を経由する分解 $W \to V \to U$ が与えられたとする。
$h^*x \to x$ と $h^*y \to y$ は強カルテシアンなので、与えられた射
$W \to V$ 上にある射 $a' : x' \to h^*x$ と $b' : y' \to h^*y$ を得る。
図式
$$
\xymatrix{
y' \ar[d]_{f'} \ar[r] & h^*y \ar[r] \ar[d]_{h^*f} & y \ar[d]_f \\
F(x') \ar[r] & F(h^*x) \ar[r] & F(x)
}
$$
を考える。外側の長方形と右の正方形は可換である。
$F$ はファイバー圏の $1$-射なので、射 $F(h^*x) \to F(x)$ は
強カルテシアンである。したがって強カルテシアン射の普遍性により、
左の正方形も可換である。これにより $\mathcal{X}'$ が
$\mathcal{C}$ 上でファイバー化されていることが示された。

\medskip\noindent
関手 $u : \mathcal{X} \to \mathcal{X}'$ は
$x \mapsto (p(x), x, F(x), \text{id})$ によって与えられ、充満忠実である。
関手 $\mathcal{X}' \to \mathcal{Y}$ は $(U, x, y, f) \mapsto y$ によって与えられる。
関手 $w : \mathcal{X}' \to \mathcal{X}$ は $(U, x, y, f) \mapsto x$ によって
与えられる。$\mathcal{C}$ 上の $\mathcal{X}'$ の強カルテシアン射についての
上の記述により、これらの関手はいずれも $\mathcal{C}$ 上のファイバー圏の
$1$-射である。$w$ と $u$ の随伴性は
$$
\Mor_\mathcal{X}(x, x') =
\Mor_{\mathcal{X}'}((U, x, y, f), (p(x'), x', F(x'), \text{id})),
$$
を意味し、これは定義から直ちに従う。

\medskip\noindent
最後に、$\mathcal{X}' \to \mathcal{Y}$ がファイバー圏であることを示す必要がある。
$c : y' \to y$ を $\mathcal{Y}$ の射とし、$(U, x, y, f)$ を $y$ 上にある
$\mathcal{X}'$ の対象とする。$V = q(y')$ とおき、
$h = q(c) : V \to U$ とおく。$a : h^*x \to x$ と
$b : h^*y \to y$ を $h$ を被覆する強カルテシアン射とする。
$F$ はファイバー圏の $1$-射なので、強カルテシアン射
$F(a) : F(h^*x) \to F(x)$ とともに $h^*F(x) = F(h^*x)$ と同一視できる。
$b : h^*y \to y$ の普遍性により、$c = b \circ c'$ を満たす
$\mathcal{Y}_V$ の射 $c' : y' \to h^*y$ が存在する。
$$
(a, c) : (V, h^*x, y', h^*f \circ c') \longrightarrow (U, x, y, f)
$$
が $\mathcal{Y}$ 上の $\mathcal{X}'$ において強カルテシアンであると主張する。
これを見るため、$(W, x_1, y_1, f_1)$ を $\mathcal{X}'$ の対象とし、
$(a_1, b_1) : (W, x_1, y_1, f_1) \to (U, x, y, f)$ を射とし、
ある射 $b_1' : y_1 \to y'$ に対して $b_1 = c \circ b_1'$ とする。このとき
$$
(a_1', b_1') : (W, x_1, y_1, f_1) \longrightarrow (V, h^*x, y', h^*f \circ c')
$$
が求める射である。ここで $a_1' : x_1 \to h^*x$ は、
$a_1 = a \circ a_1'$ を満たし、与えられた射 $q(b_1') : W \to V$ 上にある
一意な射である。
\end{proof}

\section{惰性}
\label{categories-section-inertia}

\noindent
圏 $\mathcal{C}$ 上のファイバー圏 $p : \mathcal{S} \to \mathcal{C}$ と
$p' : \mathcal{S}' \to \mathcal{C}$、および $1$-射
$F : \mathcal{S} \to \mathcal{S}'$ が与えられたとき、
$\mathcal{C}$ 上のファイバー圏の $(2, 1)$-圏には対角射
$$
\Delta = \Delta_{\mathcal{S}/\mathcal{S}'} :
\mathcal{S} \longrightarrow \mathcal{S} \times_{\mathcal{S}'} \mathcal{S}
$$
がある。

\begin{lemma}
\label{categories-lemma-inertia-fibred-category}
$\mathcal{C}$ を圏とする。$p : \mathcal{S} \to \mathcal{C}$ と
$p' : \mathcal{S}' \to \mathcal{C}$ をファイバー圏とする。
$F : \mathcal{S} \to \mathcal{S}'$ を $\mathcal{C}$ 上のファイバー圏の
$1$-射とする。$\mathcal{C}$ 上の圏
$\mathcal{I}_{\mathcal{S}/\mathcal{S}'}$ を次のように定める。
\begin{enumerate}
\item 対象は対 $(x, \alpha)$ である。ここで $x \in \Ob(\mathcal{S})$ であり、
$\alpha : x \to x$ は $F(\alpha) = \text{id}$ を満たす自己同型である。
\item 射 $(x, \alpha) \to (y, \beta)$ は、図式
$$
\xymatrix{
x\ar[r]_\phi\ar[d]_\alpha &
y\ar[d]^{\beta} \\
x\ar[r]^\phi &
y \\
}
$$
を可換にする射 $\phi : x \to y$ によって与えられる。
\item 関手 $\mathcal{I}_{\mathcal{S}/\mathcal{S}'} \to \mathcal{C}$ は
$(x, \alpha) \mapsto p(x)$ によって与えられる。
\end{enumerate}
このとき次が成り立つ。
\begin{enumerate}
\item $\mathcal{C}$ 上の圏の $(2, 1)$-圏に同値
$$
\mathcal{I}_{\mathcal{S}/\mathcal{S}'} \longrightarrow
\mathcal{S}
\times_{\Delta, (\mathcal{S} \times_{\mathcal{S}'} \mathcal{S}), \Delta}
\mathcal{S}
$$
が存在する。
\item $\mathcal{I}_{\mathcal{S}/\mathcal{S}'}$ は $\mathcal{C}$ 上の
ファイバー圏である。
\end{enumerate}
\end{lemma}

\begin{proof}
(2) は Lemma \ref{categories-lemma-2-product-fibred-categories-over-C} と
Lemma \ref{categories-lemma-fibred-equivalent} により (1) から従うことに注意せよ。
したがって (1) を証明すれば十分である。以下では
Lemma \ref{categories-lemma-2-product-fibred-categories-over-C} の $2$-ファイバー積の構成を
断りなく用いる。特に
$\mathcal{S}
\times_{\Delta, (\mathcal{S} \times_{\mathcal{S}'} \mathcal{S}), \Delta}
\mathcal{S}$ の対象は三つ組 $(x, y, (\iota, \kappa))$ である。
ここで $x$ と $y$ は $\mathcal{S}$ の対象であり、
$(\iota, \kappa) : (x, x, \text{id}_{F(x)}) \to (y, y, \text{id}_{F(y)})$ は
$\mathcal{S} \times_{\mathcal{S}'} \mathcal{S}$ の同型である。
これは単に、$\iota, \kappa : x \to y$ が同型で、
$F(\iota) = F(\kappa)$ であることを意味する。関手
$$
I_{\mathcal{S}/\mathcal{S}'}
\longrightarrow
\mathcal{S}
\times_{\Delta, (\mathcal{S} \times_{\mathcal{S}'} \mathcal{S}), \Delta}
\mathcal{S}
$$
を考える。この関手は左辺の対象 $(x, \alpha)$ に右辺の対象
$(x, x, (\alpha, \text{id}_x))$ を対応させ、左辺の射 $\phi$ に
右辺の射 $(\phi, \phi)$ を対応させる。この射の擬逆は関手
$$
\mathcal{S}
\times_{\Delta, (\mathcal{S} \times_{\mathcal{S}'} \mathcal{S}), \Delta}
\mathcal{S}
\longrightarrow
I_{\mathcal{S}/\mathcal{S}'}
$$
によって与えられると主張する。この関手は左辺の対象
$(x, y, (\iota, \kappa))$ に右辺の対象
$(x, \kappa^{-1} \circ \iota)$ を対応させ、左辺の射
$(\phi, \phi') : (x, y, (\iota, \kappa)) \to (z, w, (\lambda, \mu))$ に
射 $\phi$ を対応させる。実際、上の二つの関手を合成して得られる
$I_{\mathcal{S}/\mathcal{S}'}$ の自己関手は厳密に恒等関手であり、
$\mathcal{S}
\times_{\Delta, (\mathcal{S} \times_{\mathcal{S}'} \mathcal{S}), \Delta}
\mathcal{S}$ 上に誘導される自己関手は、自然同型
$$
(\text{id}_x, \kappa) :
(x, x, (\kappa^{-1} \circ \iota, \text{id}_x))
\longrightarrow
(x, y, (\iota, \kappa)).
$$
によって恒等関手と同型である。いくつかの詳細は省略する。
\end{proof}

\begin{definition}
\label{categories-definition-inertia-fibred-category}
$\mathcal{C}$ を圏とする。
\begin{enumerate}
\item $F : \mathcal{S} \to \mathcal{S}'$ を $\mathcal{C}$ 上のファイバー圏の
$1$-射とする。$\mathcal{S}$ の $\mathcal{S}'$ 上の{\it 相対惰性}とは、
Lemma \ref{categories-lemma-inertia-fibred-category} のファイバー圏
$\mathcal{I}_{\mathcal{S}/\mathcal{S}'} \to \mathcal{C}$ のことである。
\item $\mathcal{S}$ の{\it 惰性ファイバー圏 $\mathcal{I}_\mathcal{S}$}とは
$\mathcal{I}_\mathcal{S} = \mathcal{I}_{\mathcal{S}/\mathcal{C}}$ のことである。
\end{enumerate}
\end{definition}

\noindent
$\mathcal{C}$ 上のファイバー圏の標準的な $1$-射
\begin{equation}
\label{categories-equation-inertia-structure-map}
\mathcal{I}_{\mathcal{S}/\mathcal{S}'} \longrightarrow \mathcal{S}
\quad\text{and}\quad
\mathcal{I}_\mathcal{S} \longrightarrow \mathcal{S}
\end{equation}
があることに注意せよ。Lemma \ref{categories-lemma-inertia-fibred-category} の記述では、
これらは対象 $(x, \alpha)$ を対象 $x$ に、射
$\phi : (x, \alpha) \to (y, \beta)$ を射 $\phi : x \to y$ に写す。
さらに、規則 $x \mapsto (x, \text{id}_x)$ および
$(\phi : x \to y) \mapsto \phi$ により定まる{\it 中立切断}
\begin{equation}
\label{categories-equation-neutral-section}
e : \mathcal{S} \to \mathcal{I}_{\mathcal{S}/\mathcal{S}'}
\quad\text{and}\quad
e : \mathcal{S} \to \mathcal{I}_\mathcal{S}
\end{equation}
がある。これは (\ref{categories-equation-inertia-structure-map}) の右逆である。
$2$-可換図式
$$
\xymatrix{
\mathcal{S}_1 \ar[d]_{F_1} \ar[r]_G & \mathcal{S}_2 \ar[d]^{F_2} \\
\mathcal{S}'_1 \ar[r]^{G'} & \mathcal{S}'_2
}
$$
が与えられたとき、規則 $(x, \alpha) \mapsto (G(x), G(\alpha))$ および
$\phi \mapsto G(\phi)$ により定まる{\it 関手性写像}
\begin{equation}
\label{categories-equation-functorial}
\mathcal{I}_{\mathcal{S}_1/\mathcal{S}'_1}
\longrightarrow
\mathcal{I}_{\mathcal{S}_2/\mathcal{S}'_2}
\quad\text{and}\quad
\mathcal{I}_{\mathcal{S}_1}
\longrightarrow
\mathcal{I}_{\mathcal{S}_2}
\end{equation}
がある。特に常に比較写像
\begin{equation}
\label{categories-equation-comparison}
\mathcal{I}_{\mathcal{S}/\mathcal{S}'}
\longrightarrow
\mathcal{I}_\mathcal{S}
\end{equation}
があり、上のすべての写像はこれと両立する。

\begin{lemma}
\label{categories-lemma-relative-inertia-as-fibre-product}
$F : \mathcal{S} \to \mathcal{S}'$ を圏 $\mathcal{C}$ 上でファイバー化された
圏の $1$-射とする。このとき図式
$$
\xymatrix{
\mathcal{I}_{\mathcal{S}/\mathcal{S}'}
\ar[d]_{F \circ (\ref{categories-equation-inertia-structure-map})}
\ar[rr]_{(\ref{categories-equation-comparison})} & &
\mathcal{I}_\mathcal{S} \ar[d]^{(\ref{categories-equation-functorial})} \\
\mathcal{S}' \ar[rr]^e & &
\mathcal{I}_{\mathcal{S}'}
}
$$
は $2$-ファイバー積である。
\end{lemma}

\begin{proof}
省略する。
\end{proof}

\section{グルーポイドにファイバー化された圏}
\label{categories-section-fibred-groupoids}

\noindent
この節では、グルーポイドにファイバー化された圏をどのように捉えるかを説明し、
それらが本質的にはグルーポイドの $(2, 1)$-圏に値をとる関手と同じであることを見る。

\begin{definition}
\label{categories-definition-fibred-groupoids}
$p : \mathcal{S} \to \mathcal{C}$ を関手とする。次の二条件が成り立つとき、
$\mathcal{S}$ は $\mathcal{C}$ 上{\it グルーポイドにファイバー化されている}という。
\begin{enumerate}
\item $\mathcal{C}$ の任意の射 $f : V \to U$ と $U$ の任意の持ち上げ
$x$ に対し、終域が $x$ である $f$ の持ち上げ $\phi : y \to x$ が存在する。
\item 任意の二射 $\phi : y \to x$ と $ \psi : z \to x$、および
$p(\phi) \circ f = p(\psi)$ を満たす任意の射 $f : p(z) \to p(y)$ に対し、
$\phi \circ \chi = \psi$ を満たす $f$ の一意な持ち上げ
$\chi : z \to y$ が存在する。
\end{enumerate}
\end{definition}

\noindent
条件 (2) を別の言い方で述べると、関手 $p$ を適用することにより、
次の可換図式における破線矢印の集合の間に全単射が得られる。
$$
\xymatrix{
y \ar[r] & x & p(y) \ar[r] & p(x) \\
z \ar@{-->}[u] \ar[ru] & & p(z) \ar@{-->}[u]\ar[ru] & \\
}
$$
第二の条件は次のように考えることもできる。$g : W \to V$ と
$f : V \to U$ を $\mathcal{C}$ の射とし、
$x \in \Ob(\mathcal{S}_U)$ とする。第一の条件により、まず $f$ を
$ \phi : y \to x$ に持ち上げ、次に $g$ を $\psi : z \to y$ に
持ち上げることができる。この二段階の過程の代わりに、$g \circ f$ を
$\gamma : z' \to x$ に直接持ち上げることができる。これにより、図式
\begin{equation}
\label{categories-equation-fibred-groupoids}
\vcenter{
\xymatrix{
z' \ar@{-->}[d]\ar[rrd]^\gamma & & \\
z \ar@{-->}[u] \ar[r]^\psi \ar@{~>}[d]^p &
y \ar[r]^\phi \ar@{~>}[d]^p &
x \ar@{~>}[d]^p
\\
W \ar[r]^g & V \ar[r]^f & U \\
}
}
\end{equation}
の実線矢印が得られる。ここで波形矢印は射ではなく関手 $p$ を表す。
第二の条件を射 $\phi \circ \psi$, $\gamma$, $\text{id}_W$ に適用すると、
$\gamma \circ \chi = \phi \circ \psi$ を満たす $\mathcal{S}_W$ の
一意な射 $\chi : z \to z'$ が存在することが分かる。同様に、一意な射
$z' \to z$ が存在する。一意性により、射 $z' \to z$ と $z\to z'$ は
互いに逆、すなわち同型である。

\medskip\noindent
この議論から、グルーポイドにファイバー化された圏がファイバー圏と
極めて密接に関係することは明らかであろう。結果は次のとおりである。

\begin{lemma}
\label{categories-lemma-fibred-groupoids}
$p : \mathcal{S} \to \mathcal{C}$ を関手とする。次は同値である。
\begin{enumerate}
\item $p : \mathcal{S} \to \mathcal{C}$ はグルーポイドに
ファイバー化された圏である。
\item すべてのファイバー圏がグルーポイドであり、$\mathcal{S}$ は
$\mathcal{C}$ 上のファイバー圏である。
\end{enumerate}
さらに、この場合 $\mathcal{S}$ のすべての射は強カルテシアンである。
加えて、すべての $f: V \to U = p(x)$ に対し $f$ 上にある
$f^\ast x \to x$ が与えられると、Lemma \ref{categories-lemma-fibred} で構成したデータ
$(U \mapsto \mathcal{S}_U, f \mapsto f^*, \alpha_{f, g}, \alpha_U)$ は、
$\mathcal{C}^{opp}$ からグルーポイドの $(2, 1)$-圏への擬関手を定める。
\end{lemma}

\begin{proof}
$p : \mathcal{S} \to \mathcal{C}$ がグルーポイドにファイバー化されていると
仮定する。$U \in \Ob(\mathcal{C})$ に対するすべてのファイバー圏
$\mathcal{S}_U$ がグルーポイドであることを示すには、$\mathcal{S}_U$ の
各 $f : y \to x$ に対して逆射を構成すればよい。
左の図式（$\mathcal{S}_U$ 内）は $p$ によって右の図式に写される。
$$
\xymatrix{
y \ar[r]^f & x & U \ar[r]^{\text{id}_U} & U \\
x \ar@{-->}[u] \ar[ru]_{\text{id}_x} & &
U \ar@{-->}[u]\ar[ru]_{\text{id}_U} & \\
}
$$
右の図式を可換にするのは $\text{id}_U$ だけなので、左の図式を可換にする
一意な $g : x \to y$ が存在し、$fg = \text{id}_x$ となる。同様の議論により、
$gh = \text{id}_y$ を満たす一意な $h : y \to x$ が存在する。
このとき $fgh = f : y \to x$ である。$fg = \text{id}_x$ なので $h = f$ である。
Definition \ref{categories-definition-fibred-groupoids} の条件 (2) は、$\mathcal{S}$ の
すべての射が強カルテシアンであることをちょうど意味する。したがって
Definition \ref{categories-definition-fibred-groupoids} の条件 (1) から、
$\mathcal{S}$ が $\mathcal{C}$ 上のファイバー圏であることが従う。

\medskip\noindent
逆に、すべてのファイバー圏がグルーポイドであり、$\mathcal{S}$ が
$\mathcal{C}$ 上のファイバー圏であると仮定する。
Definition \ref{categories-definition-fibred-groupoids} の条件 (1) と (2) を確認する。
第一の条件は自明に従う。$\phi : y \to x$, $\psi : z \to x$ および
$p(\phi) \circ f = p(\psi)$ を満たす $f : p(z) \to p(y)$ を、
Definition \ref{categories-definition-fibred-groupoids} の条件 (2) のとおりにとる。
$U = p(x)$, $V = p(y)$, $W = p(z)$,
$p(\phi) = g : V \to U$, $p(\psi) = h : W \to U$ と書く。
$g$ 上にある強カルテシアン射 $g^*x \to x$ を選ぶ。
すると $\mathcal{S}_V$ の射 $i : y \to g^*x$ が得られ、これは同型である。
また、$g^*x \to x$ が強カルテシアンであることから、対 $(\psi, f)$ に
対応する射 $j : z \to g^*x$ も得られる。このとき
$\chi = i^{-1} \circ j$ が解であることを確認できる。

\medskip\noindent
(1) $\Rightarrow$ (2) の証明で、$\mathcal{S}$ のすべての射が
強カルテシアンであることを見た。最後の主張は Lemma \ref{categories-lemma-fibred} から
直接従う。
\end{proof}

\begin{lemma}
\label{categories-lemma-fibred-gives-fibred-groupoids}
$\mathcal{C}$ を圏とし、$p : \mathcal{S} \to \mathcal{C}$ を
ファイバー圏とする。$\mathcal{S}'$ を $\mathcal{S}$ の部分圏として
次のように定める。
\begin{enumerate}
\item $\Ob(\mathcal{S}') = \Ob(\mathcal{S})$ である。
\item $x, y \in \Ob(\mathcal{S}')$ に対し、$\mathcal{S}'$ における
$x$ と $y$ の間の射の集合は、$\mathcal{S}$ における $x$ と $y$ の間の
強カルテシアン射の集合である。
\end{enumerate}
$p' : \mathcal{S}' \to \mathcal{C}$ を $p$ の $\mathcal{S}'$ への制限とする。
このとき $p' : \mathcal{S}' \to \mathcal{C}$ はグルーポイドに
ファイバー化されている。
\end{lemma}

\begin{proof}
この構成が意味をもつことに注意せよ。実際、
Lemma \ref{categories-lemma-composition-cartesian} により、$\mathcal{S}$ の任意の対象の
恒等射は強カルテシアンであり、強カルテシアン射の合成も強カルテシアンである。
Definition \ref{categories-definition-fibred-groupoids} の第一の持ち上げ性質は、
ファイバー圏において、任意の射 $f : V \to U$ と $U$ 上にある $x$ に対し、
$f$ 上にある強カルテシアン射 $\varphi : y \to x$ が存在することから従う。
$\mathcal{C}$ 上の圏 $p' : \mathcal{S}' \to \mathcal{C}$ について、
Definition \ref{categories-definition-fibred-groupoids} の第二の持ち上げ性質を確認しよう。
Definition \ref{categories-definition-fibred-groupoids} の後の議論と同様に論じる。
したがって Diagram \ref{categories-equation-fibred-groupoids} における射
$\phi$, $\psi$, $\gamma$ は $\mathcal{S}$ の強カルテシアン射である。
ゆえに $\gamma$ と $\phi \circ \psi$ は $\mathcal{S}$ の強カルテシアン射で、
$\mathcal{C}$ の同じ矢上にあり、
$\mathcal{S}$ において同じ終域をもつ強カルテシアン射である。
Definition \ref{categories-definition-cartesian-over-C} の後の議論により、これは望むように
二つの矢が同型であることを意味する（ここで再び
Lemma \ref{categories-lemma-composition-cartesian} により $\mathcal{S}$ の任意の同型が
強カルテシアンであることも用いる）。
\end{proof}

\begin{example}
\label{categories-example-group-homomorphism-fibreedingroupoids}
群準同型 $p : G \to H$ は、Example
\ref{categories-example-group-homomorphism-functor} のように関手
$p : \mathcal{S}\to\mathcal{C}$ を与える。この関手
$p : \mathcal{S}\to\mathcal{C}$ がグルーポイドにファイバー化されている
ことと、$p$ が全射であることは同値である。（一意な）対象
$U\in \Ob(\mathcal{C})$ 上のファイバー圏 $\mathcal{S}_U$ は、
Example \ref{categories-example-group-groupoid} のように $p$ の核に付随する圏である。
\end{example}

\noindent
$p : \mathcal{S} \to \mathcal{C}$ が与えられたとき、すべての
$U \in \Ob(\mathcal{C})$ に対してファイバー圏 $\mathcal{S}_U$ が
グルーポイドならば、$\mathcal{S}$ は $\mathcal{C}$ 上グルーポイドに
ファイバー化されているか、と問うことができる。答えが否であることは
次のように分かる。グルーポイドにファイバー化された圏
$p : \mathcal{S} \to \mathcal{C}$ から始める。ある対象上の恒等射に写らない
$\mathcal{S}$ の射を変えても、圏 $\mathcal{S}_U$ は変わらない。
したがって持ち上げの存在条件と一意性条件を破ることができる。
一例は、$G \to H$ が全射でない場合の Example
\ref{categories-example-group-homomorphism-fibreedingroupoids} の関手である。
別の例を次に挙げる。

\begin{example}
\label{categories-example-not-fibred-in-groupoids-but-fibre-cats-are}
$\Ob(\mathcal{C}) = \{A, B, T\}$ とし、
$\Mor_\mathcal{C}(A, B) = \{f\}$, $\Mor_\mathcal{C}(B, T) = \{g\}$,
$\Mor_\mathcal{C}(A, T) = \{h\} = \{gf\}, $ と、各対象の恒等射をとる。
この圏の図を下に示す。次に $\Ob(\mathcal{S}) = \{A', B', T'\}$ とし、
$\Mor_\mathcal{S}(A', B') = \emptyset$,
$\Mor_\mathcal{S}(B', T') = \{g'\}$,
$\Mor_\mathcal{S}(A', T') = \{h'\}, $ と、恒等射をとる。
関手 $p : \mathcal{S} \to \mathcal{C}$ は明らかである。このとき、すべての
$U \in \Ob(\mathcal{C})$ に対し、$\mathcal{S}_U$ は一つの対象とその対象の
恒等射からなる圏、したがってグルーポイドであるが、射 $f: A \to B$ は
持ち上げられない。同様に、
$\Mor_\mathcal{S}(A', B') = \{f'_1, f'_2\}$ および
$ \Mor_\mathcal{S}(A', T') = \{h'\} = \{g'f'_1 \} = \{g'f'_2\}$
と定めれば、ファイバー圏は同じであり、下図の $f: A \to B$ は二つの
持ち上げをもつ。
$$
\xymatrix{
B' \ar[r]^{g'} & T' & & B \ar[r]^g & T & \\
A' \ar@{-->}[u]^{??} \ar[ru]_{h'} & & \ar@{}[u]^{above} &
A \ar[u]^f \ar[ru]_{gf = h} & \\
}
$$
\end{example}

\noindent
後で「$\mathcal{C}$ 上グルーポイドにファイバー化された任意の圏は、
分裂したものと同値である」や、「ファイバー圏が集合的であるような、
グルーポイドにファイバー化された任意の圏は、集合にファイバー化された圏と
同値である」といった主張をしたい。同値の概念は、どの $2$-圏で
作業しているかに依存する。

\begin{definition}
\label{categories-definition-categories-fibred-in-groupoids-over-C}
$\mathcal{C}$ を圏とする。{\it $\mathcal{C}$ 上グルーポイドに
ファイバー化された圏の $2$-圏}とは、$\mathcal{C}$ 上のファイバー圏の
$2$-圏（Definition \ref{categories-definition-fibred-categories-over-C} を参照）の
部分 $2$-圏で、次のように定められるものである。
\begin{enumerate}
\item その対象は、グルーポイドにファイバー化された圏
$p : \mathcal{S} \to \mathcal{C}$ である。
\item その $1$-射 $(\mathcal{S}, p) \to (\mathcal{S}', p')$ は、
$p' \circ G = p$ を満たす関手 $G : \mathcal{S} \to \mathcal{S}'$ である
（すべての射が強カルテシアンなので、$G$ はそれらを自動的に保つ）。
\item $G, H : (\mathcal{S}, p) \to (\mathcal{S}', p')$ に対する
$2$-射 $t : G \to H$ は、すべての $x \in \Ob(\mathcal{S})$ に対し
$p'(t_x) = \text{id}_{p(x)}$ を満たす関手の射である。
\end{enumerate}
\end{definition}

\noindent
すべての $2$-射は自動的に同型であることに注意せよ。したがってこれは実際には
単なる $2$-圏ではなく $(2, 1)$-圏である。次は $2$-ファイバー積に関する
当然の補題である。

\begin{lemma}
\label{categories-lemma-2-product-fibred-categories}
$\mathcal{C}$ を圏とする。$\mathcal{C}$ 上グルーポイドにファイバー化された
圏の $2$-圏は 2-ファイバー積をもち、それらは
Lemma \ref{categories-lemma-2-product-categories-over-C} のように記述される。
\end{lemma}

\begin{proof}
Lemma \ref{categories-lemma-2-product-fibred-categories-over-C} により、
Lemma \ref{categories-lemma-2-product-categories-over-C} で記述したファイバー積は
ファイバー圏である。したがって Lemma \ref{categories-lemma-fibred-groupoids} により、
ファイバー圏がグルーポイドであることを証明すれば十分である。
Lemma \ref{categories-lemma-fibre-2-fibre-product-categories-over-C} により、
グルーポイドの $2$-ファイバー積がグルーポイドであることを示せば十分であり、
これは明らかである（例えば Lemma \ref{categories-lemma-2-fibre-product-categories} の
構成から従う）。
\end{proof}

\begin{remark}
\label{categories-remark-2-product-fibred-categories-same}
$\mathcal{C}$ を圏とする。$f : \mathcal{X} \to \mathcal{S}$ と
$g : \mathcal{Y} \to \mathcal{S}$ を、$\mathcal{C}$ 上グルーポイドに
ファイバー化された圏の $1$-射とする。
$p : \mathcal{S} \to \mathcal{C}$ を与えられた関手とする。
Lemma \ref{categories-lemma-2-product-fibred-categories} の $2$-ファイバー積は、
Example \ref{categories-example-2-fibre-product-categories} のものと（圏として）
標準的に同値であると主張する。前者の対象は $p(\alpha) = \text{id}_U$ を
満たす四つ組 $(U, x, y, \alpha)$ であり
（Lemma \ref{categories-lemma-2-product-categories-over-C} を参照）、後者の対象は三つ組
$(x, y, \alpha)$ である（Example \ref{categories-example-2-fibre-product-categories} を参照）。
二つの圏の間の同値は、規則
$(U, x, y, \alpha) \mapsto (x, y, \alpha)$ および
$(x, y, \alpha) \mapsto (p(f(x)), x, y', \alpha')$ によって与えられる。
ここで $\alpha' = g(\gamma)^{-1} \circ \alpha$ であり、
$\gamma : y' \to y$ は矢 $p(\alpha) : p(f(x)) \to p(g(y))$ の持ち上げである。
詳細は省略する。
\end{remark}

\begin{lemma}
\label{categories-lemma-equivalence-fibred-categories}
$p : \mathcal{S}\to \mathcal{C}$ と
$p' : \mathcal{S'}\to \mathcal{C}$ をグルーポイドにファイバー化された圏とし、
$G : \mathcal{S}\to \mathcal {S}'$ を $\mathcal{C}$ 上の関手とする。
\begin{enumerate}
\item $G$ が忠実（それぞれ充満忠実、同値）であることと、各
$U\in\Ob(\mathcal{C})$ に対し誘導される関手
$G_U : \mathcal{S}_U\to \mathcal{S}'_U$ が忠実
（それぞれ充満忠実、同値）であることは同値である。
\item $G$ が同値ならば、$G$ は $\mathcal{C}$ 上グルーポイドに
ファイバー化された圏の $2$-圏における同値である。
\end{enumerate}
\end{lemma}

\begin{proof}
$x, y$ を、同じ対象 $U$ 上にある $\mathcal{S}$ の対象とする。
可換図式
$$
\xymatrix{
\Mor_\mathcal{S}(x, y) \ar[rd]_p \ar[rr]_G & &
\Mor_{\mathcal{S}'}(G(x), G(y)) \ar[ld]^{p'} \\
& \Mor_\mathcal{C}(U, U) &
}
$$
を考える。この図式から、$G$ が忠実（それぞれ充満忠実）ならば、各
$G_U$ もそうであることは明らかである。

\medskip\noindent
$G$ が同値であると仮定する。$\mathcal{S}'$ の各対象 $x'$ に対し、
$G(x)$ が $x'$ と同型になる $\mathcal{S}$ の対象 $x$ が存在する。
$x'$ が $U'$ 上にあり、$x$ が $U$ 上にあるとする。このとき
$\mathcal{C}$ の同型 $f : U' \to U$ が存在する。すなわち、同型
$x' \to G(x)$ に $p'$ を適用したものである。
グルーポイドにファイバー化された圏の公理により、$f$ 上にある
$\mathcal{S}$ の矢 $f^*x \to x$ が存在する。したがって、
$p'(\alpha) = \text{id}_{U'}$ を満たす同型
$\alpha : x' \to G(f^*x)$ が存在する（今度は $\mathcal{S}'$ の公理による）。
結局、$\mathcal{S}'$ の各対象 $x'$ に対して、$\mathcal{S}$ の対象
$o_{x'}$ と、$p'(\alpha_{x'}) = \text{id}_{p'(x')}$ を満たす同型
$\alpha_{x'} : x' \to G(o_{x'})$ からなる対 $(o_{x'}, \alpha_{x'})$ を
選ぶことができる。ここからは通常どおりに進める
（Lemma \ref{categories-lemma-equivalence-categories} の証明を参照）。
逆関手 $F : \mathcal{S}' \to \mathcal{S}$ を、$x' \mapsto o_{x'}$ とし、
$\varphi' : x' \to y'$ を
$\alpha_{y'}^{-1} \circ G(\varphi_{\varphi'}) \circ \alpha_{x'} = \varphi'$ を
満たす一意な矢 $\varphi_{\varphi'} : o_{x'} \to o_{y'}$ に写すことで構成する。
この選択のもとで $F$ は $\mathcal{C}$ 上の関手である。
$G \circ F$ と $F \circ G$ がそれぞれの恒等関手と $2$-同型であること
（$\mathcal{C}$ 上グルーポイドにファイバー化された圏の $2$-圏において）の
確認は省略する。

\medskip\noindent
すべての $U\in\Ob(\mathcal C)$ に対し $G_U$ が忠実（それぞれ充満忠実）
であると仮定する。$G$ が忠実（それぞれ充満忠実）であることを示すには、
任意の対象 $x, y\in\Ob(\mathcal{S})$ に対し、$G$ が
$\Mor_\mathcal{S}(x, y)$ と $\Mor_{\mathcal{S}'}(G(x), G(y))$ の間に
単射（それぞれ全単射）を誘導することを示せばよい。
$U = p(x)$, $V = p(y)$ とおく。任意の射 $f : U \to V$ に対し、
$G$ が $f$ 上にある射 $x \to y$ から $f$ 上にある射
$G(x) \to G(y)$ への単射（それぞれ全単射）を誘導することを示せば十分である。
$f : U \to V$ を固定し、$f^*y \to y$ を引き戻しと記す。
すると $G(f^*y) \to G(y)$ も引き戻しである。$f$ 上にある
$x$ から $y$ への射の集合は、$\text{id}_U$ 上にある $x$ と $f^*y$ の間の
射の集合と全単射である（グルーポイドにファイバー化された圏の第二公理による）。
同様に、$f$ 上にある $G(x)$ から $G(y)$ への射の集合は、
$\text{id}_U$ 上にある $G(x)$ と $G(f^*y)$ の間の射の集合と全単射である。
したがって、$G_U$ が忠実（それぞれ充満忠実）であることから望む結果が従う。

\medskip\noindent
最後に、すべての $U$ に対しすべての $G_U$ が同値であると仮定する。
するとそれは充満忠実かつ本質的全射である。これから $G$ が充満忠実であることは
既に見たので、それが同値であることを証明するには本質的全射であることを
証明すればよい。これは明らかである。実際、$z'\in \Ob(\mathcal{S}')$ ならば、
$U = p'(z')$ とおくと $z'\in \Ob(\mathcal{S}'_U)$ である。
$G_U$ は本質的全射なので、ある $z\in \Ob(\mathcal{S}_U)$ に対し、
$z'$ は $\mathcal{S}'_U$ において $G_U(z)$ の形の対象と同型である。
しかし $\mathcal{S}'_U$ の射は $\mathcal{S}'$ の射なので、$z'$ は
$\mathcal{S}'$ において $G(z)$ と同型である。
\end{proof}

\begin{lemma}
\label{categories-lemma-fully-faithful-diagonal-equivalence}
$\mathcal{C}$ を圏とする。$p : \mathcal{S}\to \mathcal{C}$ と
$p' : \mathcal{S'}\to \mathcal{C}$ をグルーポイドにファイバー化された圏とする。
$G : \mathcal{S}\to \mathcal {S}'$ を $\mathcal{C}$ 上の関手とする。
このとき $G$ が充満忠実であることと、対角射
$$
\Delta_G :
\mathcal{S}
\longrightarrow
\mathcal{S} \times_{G, \mathcal{S}', G} \mathcal{S}
$$
が同値であることは同値である。
\end{lemma}

\begin{proof}
Lemma \ref{categories-lemma-equivalence-fibred-categories} により、$\mathcal{C}$ の対象
$U$ 上のファイバー圏だけを見れば十分である。右辺の対象は三つ組
$(x, x', \alpha)$ であり、$\alpha : G(x) \to G(x')$ は
$\mathcal{S}'_U$ の射である。関手 $\Delta_G$ は $\mathcal{S}_U$ の対象
$x$ を三つ組 $(x, x, \text{id}_{G(x)})$ に写す。
$(x, x', \alpha)$ が $\Delta_G$ の本質像に入ることと、ある
$\mathcal{S}_U$ の射 $\beta : x \to x'$ に対して
$\alpha = G(\beta)$ となることは同値である（詳細は省略する）。
したがって $\Delta_G$ が同値であるためには、すべての $\alpha$ がある射
$\beta : x \to x'$ の像でなければならず、さらに相異なる任意の二射
$\beta, \beta' : x \to x'$ が相異なる射 $G(\beta), G(\beta')$ を
与えなければならない。これで補題が示された。
\end{proof}

\begin{lemma}
\label{categories-lemma-morphisms-equivalent-fibred-groupoids}
$\mathcal{C}$ を圏とする。$\mathcal{S}_i$, $i = 1, 2, 3, 4$ を
$\mathcal{C}$ 上グルーポイドにファイバー化された圏とする。
$\varphi : \mathcal{S}_1 \to \mathcal{S}_2$ と
$\psi : \mathcal{S}_3 \to \mathcal{S}_4$ を $\mathcal{C}$ 上の同値とする。
このとき
$$
\Mor_{\textit{Cat}/\mathcal{C}}(\mathcal{S}_2, \mathcal{S}_3)
\longrightarrow
\Mor_{\textit{Cat}/\mathcal{C}}(\mathcal{S}_1, \mathcal{S}_4),
\quad \alpha \longmapsto \psi \circ \alpha \circ \varphi
$$
は圏同値である。
\end{lemma}

\begin{proof}
これは一般論であり、任意の $2$-圏で成り立つ。
\end{proof}

\begin{lemma}
\label{categories-lemma-inertia-fibred-groupoids}
$\mathcal{C}$ を圏とする。$p : \mathcal{S} \to \mathcal{C}$ が
グルーポイドにファイバー化されているならば、惰性ファイバー圏
$\mathcal{I}_\mathcal{S} \to \mathcal{C}$ もそうである。
\end{lemma}

\begin{proof}
Lemma \ref{categories-lemma-inertia-fibred-category} の構成から明らかである。または、
同じ補題による同値
$I_\mathcal{S} \to \mathcal{S}
\times_{\Delta, \mathcal{S} \times_\mathcal{C} \mathcal{S}, \Delta}\mathcal{S}$
と Lemma \ref{categories-lemma-2-product-fibred-categories} を用いてもよい。
\end{proof}

\begin{lemma}
\label{categories-lemma-cute-groupoids}
$\mathcal{C}$ を圏とし、$U \in \Ob(\mathcal{C})$ とする。
$p : \mathcal{S} \to \mathcal{C}$ がグルーポイドにファイバー化された圏であり、
$p$ が $p' : \mathcal{S} \to \mathcal{C}/U$ を経由して分解するならば、
$p' : \mathcal{S} \to \mathcal{C}/U$ はグルーポイドにファイバー化されている。
\end{lemma}

\begin{proof}
Lemma \ref{categories-lemma-cute} で、$p'$ がファイバー圏であることは既に見た。
したがって Lemma \ref{categories-lemma-fibred-groupoids} により、ファイバー圏が
グルーポイドであることを証明すれば十分である。$V \in \Ob(\mathcal{C})$ に対し
$$
\mathcal{S}_V = \coprod\nolimits_{f : V \to U} \mathcal{S}_{(f : V \to U)}
$$
が成り立つ。ここで左辺は $p$ のファイバー圏であり、右辺は $p'$ の
ファイバー圏の直和である。したがって結果が従う。
\end{proof}

\begin{lemma}
\label{categories-lemma-fibred-in-groupoids-over-fibred-in-groupoids}
$\mathcal{A} \to \mathcal{B} \to \mathcal{C}$ を圏の間の関手とする。
$\mathcal{A}$ が $\mathcal{B}$ 上グルーポイドにファイバー化され、
$\mathcal{B}$ が $\mathcal{C}$ 上グルーポイドにファイバー化されているならば、
$\mathcal{A}$ は $\mathcal{C}$ 上グルーポイドにファイバー化されている。
\end{lemma}

\begin{proof}
これは定義から直接証明できるが、Lemma \ref{categories-lemma-fibred-groupoids} の判定法を
用いて論じる。Lemma \ref{categories-lemma-fibred-over-fibred} により、$\mathcal{A}$ は
$\mathcal{C}$ 上でファイバー化されている。証明を終えるため、$\mathcal{C}$ の
$U$ に対してファイバー圏 $\mathcal{A}_U$ がグルーポイドであることを示す。
実際、$x \to y$ を $\mathcal{A}_U$ の射とすると、その $\mathcal{B}$ における
像は、$\mathcal{B}_U$ がグルーポイドなので同型である。しかしこのとき
$x \to y$ も同型である。例えば Lemma \ref{categories-lemma-composition-cartesian} と、
$\mathcal{A}$ のすべての射が強 $\mathcal{B}$-カルテシアンであること
（Lemma \ref{categories-lemma-fibred-groupoids} を参照）から従う。
\end{proof}

\begin{lemma}
\label{categories-lemma-fibred-groupoids-fibre-product-goes-up}
$p : \mathcal{S} \to \mathcal{C}$ をグルーポイドにファイバー化された圏とする。
$x \to y$ と $z \to y$ を $\mathcal{S}$ の射とする。
$p(x) \times_{p(y)} p(z)$ が存在するならば、$x \times_y z$ が存在し、
$p(x \times_y z) = p(x) \times_{p(y)} p(z)$ である。
\end{lemma}

\begin{proof}
Lemma \ref{categories-lemma-fibred-category-representable-goes-up} から従う。
\end{proof}

\begin{lemma}
\label{categories-lemma-ameliorate-morphism-categories-fibred-groupoids}
$\mathcal{C}$ を圏とし、$F : \mathcal{X} \to \mathcal{Y}$ を
$\mathcal{C}$ 上グルーポイドにファイバー化された圏の $1$-射とする。
$\mathcal{C}$ 上グルーポイドにファイバー化された圏の $1$-射による分解
$\mathcal{X} \to \mathcal{X}' \to \mathcal{Y}$ であって、
$\mathcal{X} \to \mathcal{X}'$ が $\mathcal{C}$ 上の同値であり、かつ
$\mathcal{X}'$ が $\mathcal{Y}$ 上グルーポイドにファイバー化された圏と
なるものが存在する。
\end{lemma}

\begin{proof}
構造関手を $p : \mathcal{X} \to \mathcal{C}$ および
$q : \mathcal{Y} \to \mathcal{C}$ と記す。$\mathcal{X}'$ を次のように
具体的に構成する。$\mathcal{X}'$ の対象は四つ組 $(U, x, y, f)$ であり、
$x \in \Ob(\mathcal{X}_U)$, $y \in \Ob(\mathcal{Y}_U)$ で、
$f : F(x) \to y$ は $\mathcal{Y}_U$ の同型である。
射 $(a, b) : (U, x, y, f) \to (U', x', y', f')$ は、
$p(a) = q(b)$ および $f' \circ F(a) = b \circ f$ を満たす
$a : x \to x'$ と $b : y \to y'$ によって与えられる。言い換えれば、
$\mathcal{X}' = \mathcal{X} \times_{F, \mathcal{Y}, \text{id}} \mathcal{Y}$ であり、
ここでは Lemma \ref{categories-lemma-2-product-categories-over-C} の $2$-ファイバー積の
構成を用いる。Lemma \ref{categories-lemma-2-product-fibred-categories} により、
$\mathcal{X}'$ は $\mathcal{C}$ 上グルーポイドにファイバー化された圏であり、
$\mathcal{X}' \to \mathcal{Y}$ は $\mathcal{C}$ 上の圏の射である。
関手 $\mathcal{X} \to \mathcal{X}'$ として、対象上では
$x \mapsto (p(x), x, F(x), \text{id}_{F(x)})$、射上では
$(a : x \to x') \mapsto (a, F(a))$ をとる。合成
$\mathcal{X} \to \mathcal{X}' \to \mathcal{Y}$ が $F$ に等しいことは明らかである。
$\mathcal{X} \to \mathcal{X}'$ が $\mathcal{C}$ 上のファイバー圏の同値である
ことの確認は省略する。

\medskip\noindent
最後に、$\mathcal{X}' \to \mathcal{Y}$ がグルーポイドにファイバー化された圏で
あることを示す必要がある。$b : y' \to y$ を $\mathcal{Y}$ の射とし、
$(U, x, y, f)$ を $y$ 上にある $\mathcal{X}'$ の対象とする。
$\mathcal{X}$ は $\mathcal{C}$ 上グルーポイドにファイバー化されているので、
$U' = q(y') \to q(y) = U$ 上にある射 $a : x' \to x$ を見つけられる。
$\mathcal{Y}$ は $\mathcal{C}$ 上グルーポイドにファイバー化されており、
$F(x') \to F(x)$ と $y' \to y$ はいずれも同じ射 $U' \to U$ 上にあるので、
$f \circ F(a) = b \circ f'$ を満たし $\text{id}_{U'}$ 上にある
$f' : F(x') \to y'$ を見つけられる。したがって
$(a, b) : (U', x', y', f') \to (U, x, y, f)$ を得る。
これにより Definition \ref{categories-definition-fibred-groupoids} の第一条件 (1) が確認された。
(2) を見るため、$\mathcal{X}'$ の射
$(a, b) : (U', x', y', f') \to (U, x, y, f)$ と
$(a', b') : (U'', x'', y'', f'') \to (U, x, y, f)$、および
$b' \circ b'' = b$ を満たす $\mathcal{Y}$ の射 $b'' : y' \to y''$ をとる。
$f'' \circ F(a'') = b'' \circ f'$ および
$(a', b') \circ (a'', b'') = (a, b)$ を満たす一意な射
$a'' : x' \to x''$ が存在することを示す必要がある。
$\mathcal{X}$ はグルーポイドにファイバー化されているので、
$a' \circ a'' = a$ と $p(a'') = q(b'')$ を満たす一意な射
$a'' : x' \to x''$ が存在する。$\mathcal{Y}$ はグルーポイドに
ファイバー化されているので、$F(a'')$ は
$F(a') \circ F(a'') = F(a)$ と $q(F(a'')) = q(b'')$ を満たす一意な射
$F(x') \to F(x'')$ である。このことと、与えられた関係
$f \circ F(a) = b \circ f'$ および $f \circ F(a') = b' \circ f''$ から
$f'' \circ F(a'') = b'' \circ f'$ が従う。
\end{proof}

\begin{lemma}
\label{categories-lemma-amelioration-unique}
$\mathcal{C}$ を圏とし、$F : \mathcal{X} \to \mathcal{Y}$ を
$\mathcal{C}$ 上グルーポイドにファイバー化された圏の $1$-射とする。
$2$-可換図式
$$
\xymatrix{
\mathcal{X}' \ar[rd]_f &
\mathcal{X} \ar[l]^a \ar[d]^F \ar[r]_b &
\mathcal{X}'' \ar[ld]^g \\
& \mathcal{Y}
}
$$
が与えられていると仮定する。ここで $a$ と $b$ は $\mathcal{C}$ 上の
圏同値であり、$f$ と $g$ はグルーポイドにファイバー化された圏である。
このとき、$\mathcal{Y}$ 上の圏同値 $h : \mathcal{X}'' \to \mathcal{X}'$ であって、
$h \circ b$ が $\mathcal{C}$ 上の圏の $1$-射として $a$ と $2$-同型に
なるものが存在する。上の図式が実際に可換ならば、$h \circ b$ が
$\mathcal{Y}$ 上の圏の $1$-射として $a$ と $2$-同型になるようにできる。
\end{lemma}

\begin{proof}
$\mathcal{Y}$ 上の $\mathcal{X}'$ と $\mathcal{X}''$ はいずれも、
$\mathcal{Y}$ 上グルーポイドにファイバー化された圏
$\mathcal{X} \times_{F, \mathcal{Y}, \text{id}} \mathcal{Y}$ と同値であることを
示す。Lemma \ref{categories-lemma-ameliorate-morphism-categories-fibred-groupoids} の証明を
参照せよ。$\mathcal{C}$ 上の圏の $2$-圏における擬逆
$b^{-1} : \mathcal{X}'' \to \mathcal{X}$ を選ぶ。図式の右の三角形は
$2$-可換なので、図式
$$
\xymatrix{
\mathcal{X} \ar[d]_F & \mathcal{X}'' \ar[l]^{b^{-1}} \ar[d]^g \\
\mathcal{Y} & \mathcal{Y} \ar[l]
}
$$
は $2$-可換である。したがって $2$-ファイバー積の普遍性により $1$-射
$c : \mathcal{X}'' \to
\mathcal{X} \times_{F, \mathcal{Y}, \text{id}} \mathcal{Y}$ を得る。
さらに $c$ は $\mathcal{Y}$ 上の圏の射であり（！）、かつ同値である
（$b$ が同値であるという仮定と Lemma
\ref{categories-lemma-equivalence-2-fibre-product} による）。したがって
Lemma \ref{categories-lemma-equivalence-fibred-categories} により、$c$ は
$\mathcal{Y}$ 上グルーポイドにファイバー化された圏の $2$-圏における同値である。

\medskip\noindent
なお、$c \circ b$ と、Lemma
\ref{categories-lemma-ameliorate-morphism-categories-fibred-groupoids} の証明で構成した関手
$d : \mathcal{X} \to
\mathcal{X} \times_{F, \mathcal{Y}, \text{id}} \mathcal{Y}$,
$x \mapsto (p(x), x, F(x), \text{id}_{F(x)})$ の間の $2$-同型を構成する必要がある。
$\alpha : F \to g \circ b$ と $\beta : b^{-1} \circ b \to \text{id}$ を、
$\mathcal{C}$ 上の圏の $1$-射の間の $2$-同型とする。
$c \circ b$ は対象上で規則
$$
x \mapsto (p(x), b^{-1}(b(x)), g(b(x)), \alpha_x \circ F(\beta_x))
$$
によって与えられることに注意せよ。このとき
$$
(\beta_x, \alpha_x) :
(p(x), x, F(x), \text{id}_{F(x)})
\longrightarrow
(p(x), b^{-1}(b(x)), g(b(x)), \alpha_x \circ F(\beta_x))
$$
は、求める $2$-射 $d \to b \circ c$ を与える関手的な同型である。
最後に、図式が可換ならば、すべての $x$ に対して $\alpha_x$ は恒等射であり、
この $2$-射が $\mathcal{Y}$ 上の圏の $2$-圏における $2$-射であることが分かる。
\end{proof}

\section{圏の前層}
\label{categories-section-presheaves-categories}

\noindent
この節では、ファイバー圏の概念と、密接に関連する「圏の前層」の概念を比較する。
基本的な構成を次の例で説明する。

\begin{example}
\label{categories-example-functor-categories}
$\mathcal{C}$ を圏とする。$F : \mathcal{C}^{opp} \to \textit{Cat}$ を
圏の $2$-圏への関手とする。Definition
\ref{categories-definition-functor-into-2-category} を参照せよ。
$\mathcal{C}$ の $f : V \to U$ に対し、$F(U)$ から $F(V)$ への関手を
示唆的に $F(f) = f^\ast$ と書く。これから $\mathcal{C}$ 上のファイバー圏
$\mathcal{S}_F$ を次のように構成できる。まず
$$
\Ob(\mathcal{S}_F) =
\{(U, x) \mid U\in \Ob(\mathcal{C}), x\in \Ob(F(U))\}.
$$
と定める。$(U, x), (V, y) \in \Ob(\mathcal{S}_F)$ に対して
\begin{align*}
\Mor_{\mathcal{S}_F}((V, y), (U, x)) & =
\{ (f, \phi) \mid f \in \Mor_\mathcal{C}(V, U),
\phi \in \Mor_{F(V)}(y, f^\ast x)\} \\
& =
\coprod\nolimits_{f \in \Mor_\mathcal{C}(V, U)}
\Mor_{F(V)}(y, f^\ast x)
\end{align*}
と定める。合成を定めるため、$\mathcal{C}$ の合成可能な射の対に対して
$g^\ast \circ f^\ast = (f \circ g)^\ast$ であることを用いる
（$2$-圏への関手の定義による）。すなわち、
$\psi : z \to g^\ast y$ と $ \phi : y \to f^\ast x$ の合成を
$ g^\ast(\phi) \circ \psi$ と定める。関手
$p_F : \mathcal{S}_F \to \mathcal{C}$ は規則 $(U, x) \mapsto U$ によって
与えられる。これが実際にファイバー圏であることを確認しよう。
$\mathcal{C}$ の $f: V \to U$ と、$U$ の持ち上げ $(U, x)$ が与えられたとき、
$(f, \text{id}_{f^\ast x}): (V, {f^\ast x}) \to (U, x)$ は $f$ の
強カルテシアンな持ち上げであると主張する。左図の $h$ が右図の
$(h, \nu)$ を定めることを示す必要がある。
$$
\xymatrix{
V \ar[r]^f &
U &
(V, f^*x) \ar[r]^{(f, \text{id}_{f^*x})} &
(U, x) \\
W \ar@{-->}[u]^h \ar[ru]_g & &
(W, z) \ar@{-->}[u]^{(h, \nu)} \ar[ru]_{(g, \psi)} &
}
$$
$\nu = \psi$ ととればよい。これは $f \circ h = g$ であり、したがって
$g^*x = h^*f^*x$ なので機能する。さらに、右図を可換にする持ち上げは
これだけである。
\end{example}

\begin{definition}
\label{categories-definition-split-fibred-category}
$\mathcal{C}$ を圏とする。$F : \mathcal{C}^{opp} \to \textit{Cat}$ を
圏の $2$-圏への関手とする。Example \ref{categories-example-functor-categories} で構成した
ファイバー圏を $p_F : \mathcal{S}_F \to \mathcal{C}$ と書く。
{\it 分裂ファイバー圏}とは、$\mathcal{C}$ 上でこれらの圏
{\it $\mathcal{S}_F$} の一つと同型（！）なファイバー圏である。
\end{definition}

\begin{lemma}
\label{categories-lemma-when-split}
$\mathcal{C}$ を圏とし、$\mathcal{S}$ を $\mathcal{C}$ 上のファイバー圏とする。
$\mathcal{S}$ が分裂していることと、ある引き戻しの選択
（Definition \ref{categories-definition-pullback-functor-fibred-category} を参照）に対して
引き戻し関手 $(f \circ g)^*$ と $g^* \circ f^*$ が等しいことは同値である。
\end{lemma}

\begin{proof}
これは定義から直ちに従う。
\end{proof}

\begin{lemma}
\label{categories-lemma-fibred-strict}
$ p : \mathcal{S} \to \mathcal{C}$ をファイバー圏とする。
反変関手 $F : \mathcal{C} \to \textit{Cat}$ であって、$\mathcal{S}$ が
$\mathcal{C}$ 上のファイバー圏の $2$-圏において $\mathcal{S}_F$ と
同値になるものが存在する。言い換えれば、すべてのファイバー圏は
分裂したものと同値である。
\end{lemma}

\begin{proof}
引き戻しの選択を行う（Definition
\ref{categories-definition-pullback-functor-fibred-category} を参照）。
Lemma \ref{categories-lemma-fibred} により、$\mathcal{C}$ の各射 $f$ に対する
引き戻し関手 $f^*$ を得る。

\medskip\noindent
新しい圏 $\mathcal{S}'$ を次のように構成する。$\mathcal{S}'$ の対象は、
$\mathcal{C}$ の射 $f : V \to U$ と $U$ 上にある $\mathcal{S}$ の対象 $x$、
すなわち $x\in \Ob(\mathcal{S}_U)$ からなる対 $(x, f)$ である。
関手 $p' : \mathcal{S}' \to \mathcal{C}$ は対 $(x, f)$ を射 $f$ の始域に写す。
言い換えれば $p'(x, f : V\to U) = V$ である。
射 $\varphi : (x_1, f_1: V_1 \to U_1) \to (x_2, f_2 : V_2 \to U_2)$ は、
射 $g : V_1 \to V_2$ と、$p(\varphi) = g$ を満たす射
$\varphi : f_1^\ast x_1 \to f_2^\ast x_2$ からなる対 $(\varphi, g)$ によって
与えられる。合成則を定めることに問題はない。任意の合成可能な射の対に対し、
$(\varphi, g) \circ (\psi, h) =
(\varphi \circ \psi, g\circ h)$ とする。自然な関手
$\mathcal{S} \to \mathcal{S}'$ は、$U$ 上の $x$ を単に対
$(x, \text{id}_U)$ に写す。

\medskip\noindent
ここで、$p'$ が $\mathcal{S}'$ を $\mathcal{C}$ 上のファイバー圏にすること、
および $\mathcal{S} \to \mathcal{S}'$ が強カルテシアン射を
強カルテシアン射に写す $\mathcal{C}$ 上の圏同値であることを確認する必要がある。
これらの確認は省略する。

\medskip\noindent
最後に、$\mathcal{S}'$ 上の引き戻し関手を次のように定められる。
$g : V' \to V$ と $f : V \to U$ に対し、対象上では
$g^\ast(x, f) = (x, f \circ g)$ とおく。$\mathcal{S}'_V$ の射の間の射
$(\varphi, \text{id}_V) : (x_1, f_1) \to (x_2, f_2)$ に対し、射上では
$g^\ast(\varphi, \text{id}_V) =
(g^\ast\varphi, \text{id}_{V'})$ とおく。ここで Lemma \ref{categories-lemma-fibred} の
一意な同一視 $g^\ast f_i^\ast x_i = (f_i \circ g)^\ast x_i$ を用い、
$g^\ast\varphi$ を $(f_1 \circ g)^\ast x_1$ から
$(f_2 \circ g)^\ast x_2$ への射とみなす。明らかに、これらの引き戻し関手
$g^\ast$ は $g_1^\ast \circ g_2^\ast = (g_2\circ g_1)^\ast$ を満たす。
言い換えれば、$\mathcal{S}'$ は望むように分裂している。
\end{proof}

\section{グルーポイドの前層}
\label{categories-section-presheaves-groupoids}

\noindent
この節では、グルーポイドにファイバー化された圏の概念と、密接に関連する
「グルーポイドの前層」の概念を比較する。基本的な構成を次の例で説明する。

\begin{example}
\label{categories-example-functor-groupoids}
これは Example \ref{categories-example-functor-categories} において
「圏の前層」を「グルーポイドの前層」に置き換えた類似例である。
出力はファイバー圏ではなく、グルーポイドにファイバー化された圏となる。
$F : \mathcal{C}^{opp} \to \textit{Groupoids}$ をグルーポイドの圏への関手と
仮定する。Definition \ref{categories-definition-functor-into-2-category} を参照せよ。
$\mathcal{C}$ の $f : V \to U$ に対し、$F(U)$ から $F(V)$ への関手を
示唆的に $F(f) = f^\ast$ と書く。$\mathcal{C}$ 上グルーポイドに
ファイバー化された圏 $\mathcal{S}_F$ を次のように構成する。まず
$$
\Ob(\mathcal{S}_F) =
\{(U, x) \mid U\in \Ob(\mathcal{C}), x\in \Ob(F(U))\}.
$$
と定める。$(U, x), (V, y) \in \Ob(\mathcal{S}_F)$ に対して
\begin{align*}
\Mor_{\mathcal{S}_F}((V, y), (U, x))
& =
\{ (f, \phi) \mid f \in \Mor_\mathcal{C}(V, U),
\phi \in \Mor_{F(V)}(y, f^\ast x)\} \\
& =
\coprod\nolimits_{f \in \Mor_\mathcal{C}(V, U)}
\Mor_{F(V)}(y, f^\ast x)
\end{align*}
と定める。合成を定めるため、$\mathcal{C}$ の合成可能な射の対に対して
$g^\ast \circ f^\ast = (f \circ g)^\ast$ であることを用いる
（$2$-圏への関手の定義による）。すなわち、
$\psi : z \to g^\ast y$ と $ \phi : y \to f^\ast x$ の合成を
$ g^\ast(\phi) \circ \psi$ と定める。関手
$p_F : \mathcal{S}_F \to \mathcal{C}$ は規則 $(U, x) \mapsto U$ によって
与えられる。すべての $U$ に対して $F(U)$ がグルーポイドであるという条件により、
$\mathcal{S}_F$ は $\mathcal{C}$ 上グルーポイドにファイバー化されている。
実際、Example \ref{categories-example-functor-categories} ですでに
$\mathcal{S}_F$ がファイバー圏であることを見たので、Lemma
\ref{categories-lemma-fibred-groupoids} を参照すればよい。しかし、Definition
\ref{categories-definition-fibred-groupoids} の条件 (1), (2) を次のように直接証明することも
できる。(1) 射の持ち上げは存在する。実際、$\mathcal{C}$ の
$f: V \to U$ と、$U$ 上にある $\mathcal{S}_F$ の対象 $(U, x)$ が
与えられれば、
$(f, \text{id}_{f^\ast x}): (V, {f^\ast x}) \to (U, x)$ は $f$ の持ち上げである。
(2) 次の実線の図式が与えられたと仮定する。
$$
\xymatrix{
V \ar[r]^f & U & (V, y) \ar[r]^{(f, \phi)} & (U, x) \\
W \ar@{-->}[u]^h \ar[ru]_g & &
(W, z) \ar@{-->}[u]^{(h, \nu)} \ar[ru]_{(g, \psi)} & \\
}
$$
このとき点線の矢印に対して $\nu = (h^\ast \phi)^{-1} \circ \psi$ である。
したがって $h$ が与えられれば $\nu$ が存在し、逆射の一意性により一意である。
\end{example}

\begin{definition}
\label{categories-definition-split-category-fibred-in-groupoids}
$\mathcal{C}$ を圏とする。
$F : \mathcal{C}^{opp} \to \textit{Groupoids}$ をグルーポイドの $2$-圏への
関手とする。Example \ref{categories-example-functor-groupoids} で構成した、
グルーポイドにファイバー化された圏を
$p_F : \mathcal{S}_F \to \mathcal{C}$ と書く。
{it 分裂したグルーポイドにファイバー化された圏}とは、
$\mathcal{C}$ 上でこれらの圏 {it $\mathcal{S}_F$} の一つと同型（！）な、
グルーポイドにファイバー化された圏である。
\end{definition}

\begin{lemma}
\label{categories-lemma-fibred-groupoids-strict}
$ p : \mathcal{S} \to \mathcal{C}$ をグルーポイドにファイバー化された圏とする。
反変関手 $F : \mathcal{C} \to \textit{Groupoids}$ であって、
$\mathcal{S}$ が $\mathcal{C}$ 上で $\mathcal{S}_F$ と同値になるものが存在する。
言い換えれば、グルーポイドにファイバー化されたすべての圏は、分裂したものと
同値である。
\end{lemma}

\begin{proof}
引き戻しの選択を行う（Definition
\ref{categories-definition-pullback-functor-fibred-category} を参照）。Lemmas
\ref{categories-lemma-fibred} および \ref{categories-lemma-fibred-groupoids} により、
$\mathcal{C}$ の各射 $f$ に対する引き戻し関手 $f^*$ を得る。

\medskip\noindent
新しい圏 $\mathcal{S}'$ を次のように構成する。$\mathcal{S}'$ の対象は、
$\mathcal{C}$ の射 $f : V \to U$ と $U$ 上にある $\mathcal{S}$ の対象 $x$、
すなわち $x\in \Ob(\mathcal{S}_U)$ からなる対 $(x, f)$ である。
関手 $p' : \mathcal{S}' \to \mathcal{C}$ は対 $(x, f)$ を射 $f$ の始域に写す。
言い換えれば $p'(x, f : V\to U) = V$ である。
射 $\varphi : (x_1, f_1: V_1 \to U_1) \to (x_2, f_2 : V_2 \to U_2)$ は、
射 $g : V_1 \to V_2$ と、$p(\varphi) = g$ を満たす射
$\varphi : f_1^\ast x_1 \to f_2^\ast x_2$ からなる対 $(\varphi, g)$ によって
与えられる。合成則を定めることに問題はない。任意の合成可能な射の対に対し、
$(\varphi, g) \circ (\psi, h) =
(\varphi \circ \psi, g\circ h)$ とする。自然な関手
$\mathcal{S} \to \mathcal{S}'$ は、$U$ 上の $x$ を単に対
$(x, \text{id}_U)$ に写す。

\medskip\noindent
ここで、$p'$ が $\mathcal{S}'$ を $\mathcal{C}$ 上グルーポイドに
ファイバー化された圏にすること、および $\mathcal{S} \to \mathcal{S}'$ が
$\mathcal{C}$ 上の圏同値であることを確認する必要がある。これらの確認は省略する。

\medskip\noindent
最後に、$\mathcal{S}'$ 上の引き戻し関手を次のように定められる。
$g : V' \to V$ と $f : V \to U$ に対し、対象上では
$g^\ast(x, f) = (x, f \circ g)$ とおく。$\mathcal{S}'_V$ の射の間の射
$(\varphi, \text{id}_V) : (x_1, f_1) \to (x_2, f_2)$ に対し、射上では
$g^\ast(\varphi, \text{id}_V) =
(g^\ast\varphi, \text{id}_{V'})$ とおく。ここで Lemma
\ref{categories-lemma-fibred-groupoids} の一意な同一視
$g^\ast f_i^\ast x_i = (f_i \circ g)^\ast x_i$ を用い、
$g^\ast\varphi$ を $(f_1 \circ g)^\ast x_1$ から
$(f_2 \circ g)^\ast x_2$ への射とみなす。明らかに、これらの引き戻し関手
$g^\ast$ は $g_1^\ast \circ g_2^\ast = (g_2\circ g_1)^\ast$ を満たす。
言い換えれば、$\mathcal{S}'$ は望むように分裂している。
\end{proof}

\noindent
この補題の別証明は Section \ref{categories-section-representable-1-morphisms} で見る。

\section{集合にファイバー化された圏}
\label{categories-section-fibred-in-sets}

\begin{definition}
\label{categories-definition-discrete}
圏の射が恒等射だけであるとき、その圏を {\it 離散的} であるという。
\end{definition}

\noindent
離散圏がもつ興味深い情報は、その対象の集合だけである。したがって、
離散圏と集合とを同一視することがある。

\begin{definition}
\label{categories-definition-category-fibred-sets}
$\mathcal{C}$ を圏とする。{\it 集合にファイバー化された圏}、または
{\it 離散圏にファイバー化された圏}とは、すべてのファイバー圏が離散的である、
グルーポイドにファイバー化された圏をいう。
\end{definition}

\noindent
集合にファイバー化された圏と前層（Definition
\ref{categories-definition-presheaf} を参照）の関係を明確にしたい。そのためには、
まず次の定義をするのが自然である。

\begin{definition}
\label{categories-definition-categories-fibred-in-sets-over-C}
$\mathcal{C}$ を圏とする。{\it $\mathcal{C}$ 上集合にファイバー化された圏の
$2$-圏}とは、$\mathcal{C}$ 上グルーポイドにファイバー化された圏の圏
（Definition \ref{categories-definition-categories-fibred-in-groupoids-over-C} を参照）の
部分 $2$-圏であって、次のように定められるものをいう。
\begin{enumerate}
\item その対象は、集合にファイバー化された圏
$p : \mathcal{S} \to \mathcal{C}$ とする。
\item その $1$-射 $(\mathcal{S}, p) \to (\mathcal{S}', p')$ は、
$p' \circ G = p$ を満たす関手 $G : \mathcal{S} \to \mathcal{S}'$ とする
（すべての射が強カルテシアンなので、$G$ はそれらを自動的に保つ）。
\item $G, H : (\mathcal{S}, p) \to (\mathcal{S}', p')$ に対するその $2$-射
$t : G \to H$ は、すべての $x \in \Ob(\mathcal{S})$ に対して
$p'(t_x) = \text{id}_{p(x)}$ を満たす関手の射とする。
\end{enumerate}
\end{definition}

\noindent
すべての $2$-射が自動的に同型であることに注意せよ。したがって、この
$2$-圏は実際には $(2, 1)$-圏である。次は $2$-ファイバー積の存在に関する
当然の補題である。

\begin{lemma}
\label{categories-lemma-2-product-categories-fibred-sets}
$\mathcal{C}$ を圏とする。$\mathcal{C}$ 上集合にファイバー化された圏の
2-圏は 2-ファイバー積をもつ。より正確には、Lemma
\ref{categories-lemma-2-product-categories-over-C} で記述した 2-ファイバー積は、
集合にファイバー化された圏から始めれば、集合にファイバー化された圏を返す。
\end{lemma}

\begin{proof}
省略する。
\end{proof}

\begin{example}
\label{categories-example-presheaf}
これは Examples \ref{categories-example-functor-categories} および
\ref{categories-example-functor-groupoids} において、「圏の前層」の代わりに前層を
扱う類似例である。出力はファイバー圏ではなく、集合にファイバー化された圏となる。
$F : \mathcal{C}^{opp} \to \textit{Sets}$ を前層とする。
$\mathcal{C}$ の $f : V \to U$ に対し、示唆的に
$F(f) = f^\ast : F(U) \to F(V)$ と書く。$\mathcal{C}$ 上集合に
ファイバー化された圏 $\mathcal{S}_F$ を次のように構成する。まず
$$
\Ob(\mathcal{S}_F) =
\{(U, x) \mid U \in \Ob(\mathcal{C}), x \in \Ob(F(U))\}.
$$
と定める。$(U, x), (V, y) \in \Ob(\mathcal{S}_F)$ に対して
\begin{align*}
\Mor_{\mathcal{S}_F}((V, y), (U, x))
& =
\{f \in \Mor_\mathcal{C}(V, U) \mid f^*x = y\}
\end{align*}
と定める。合成は $\mathcal{C}$ の合成から受け継ぐ。これは
$\mathcal{C}$ の合成可能な射の対に対して
$g^\ast \circ f^\ast = (f \circ g)^\ast$ だから機能する。
関手 $p_F : \mathcal{S}_F \to \mathcal{C}$ は規則
$(U, x) \mapsto U$ によって与えられる。各ファイバー圏
$\mathcal{S}_{F, U}$ は台集合 $F(U)$ をもつ離散圏であり、Example
\ref{categories-example-functor-groupoids} ですでに $\mathcal{S}_F$ がグルーポイドに
ファイバー化された圏であることを見たので、$\mathcal{S}_F$ は集合に
ファイバー化されていると結論できる。
\end{example}

\begin{lemma}
\label{categories-lemma-2-category-fibred-sets}
\begin{slogan}
集合にファイバー化された圏は、まさに前層である。
\end{slogan}
$\mathcal{C}$ を圏とする。集合にファイバー化された圏の間の $2$-射は
恒等射だけである。言い換えれば、集合にファイバー化された圏の $2$-圏は圏である。
さらに、圏同値
$$
\left\{
\begin{matrix}
\text{the category of presheaves}\\
\text{of sets over }\mathcal{C}
\end{matrix}
\right\}
\leftrightarrow
\left\{
\begin{matrix}
\text{the category of categories}\\
\text{fibred in sets over }\mathcal{C}
\end{matrix}
\right\}
$$
が存在する。左から右への関手は Example \ref{categories-example-presheaf} で論じた構成
$F \to \mathcal{S}_F$ である。右から左への関手は
$p : \mathcal{S} \to \mathcal{C}$ に対象の前層
$U \mapsto \Ob(\mathcal{S}_U)$ を対応させる。
\end{lemma}

\begin{proof}
第一の主張は明らかである。ファイバー圏の射は恒等射だけだからである。

\medskip\noindent
$p : \mathcal{S} \to \mathcal{C}$ が集合にファイバー化されていると仮定する。
$f : V \to U$ を $\mathcal{C}$ の射とし、
$x \in \Ob(\mathcal{S}_U)$ とする。このとき対象 $f^\ast x$ の選択は
ちょうど一つだけである。したがって、Lemma \ref{categories-lemma-fibred-groupoids} の
$f, g$ に対して $(f \circ g)^\ast x = g^\ast(f^\ast x)$ である。
ゆえに、対応 $U \mapsto \Ob(\mathcal{S}_U)$ および $f \mapsto f^\ast$ を
$\mathcal{C}$ 上の前層とみなせる。
\end{proof}

\noindent
集合にファイバー化された圏の重要な例を挙げる。

\begin{example}
\label{categories-example-fibred-category-from-functor-of-points}
$\mathcal{C}$ を圏とし、$X \in \Ob(\mathcal{C})$ とする。表現可能前層
$h_X = \Mor_\mathcal{C}(-, X)$ を考える
（Example \ref{categories-example-hom-functor} を参照）。一方、Example
\ref{categories-example-category-over-X} の圏
$p : \mathcal{C}/X \to \mathcal{C}$ を考える。ファイバー圏
$(\mathcal{C}/X)_U$ の対象は射 $h : U \to X$ であり、その射は恒等射だけである。
したがって、Lemma \ref{categories-lemma-2-category-fibred-sets} の対応のもとで
$$
h_X \longleftrightarrow \mathcal{C}/X.
$$
である。言い換えれば、圏 $\mathcal{C}/X$ は、Example
\ref{categories-example-presheaf} において $h_X$ に対応する圏
$\mathcal{S}_{h_X}$ と標準的に同値である。
\end{example}

\noindent
このため、グルーポイドにファイバー化された圏の 2-圏における
「表現可能」対象を、対応する前層が表現可能であるような、集合に
ファイバー化された圏として定義したくなる。しかし、同値のもとで不変な概念を
望むので、これは用いるのに適した定義ではない。これを明確にするため、
グルーポイドにファイバー化された圏のうち、どれが集合にファイバー化された圏と
同値であるかを正確に調べる。

\section{セトイドにファイバー化された圏}
\label{categories-section-fibred-in-setoids}

\begin{definition}
\label{categories-definition-setoid}
圏が、すべての対象がちょうど一つの自己同型、すなわち恒等射をもつ
グルーポイドであるとき、その圏を {\it セトイド} と呼ぶ
\footnote{ステロイドを投与された集合！？}。
\end{definition}

\noindent
$C$ が同値関係 $\sim$ を備えた集合ならば、次のようにセトイド
$\mathcal{C}$ を作ることができる：
$\Ob(\mathcal{C}) = C$ とし、$x \sim y$ でない限り
$\Mor_\mathcal{C}(x, y) = \emptyset$ とし、そうであれば
$\Mor_\mathcal{C}(x, y) = \{1\}$ とする。$\sim$ の推移性により射を合成できる。
逆に、任意のセトイド圏はその対象上に同値関係（同型）を定め、この手順で
記述された圏を（一意な同型まで、同値までではなく）復元できる。

\medskip\noindent
離散圏はセトイドである。任意のセトイド $\mathcal{C}$ に対して、それと同値な
離散圏を作る標準的な手順がある。すなわち、$\Ob(\mathcal{C})$ を同型類の集合で
置き換え、恒等射を加える。同値関係を備えた集合の言葉では、これは同値関係による
商を取ることに対応する。

\begin{definition}
\label{categories-definition-category-fibred-setoids}
$\mathcal{C}$ を圏とする。{\it セトイドにファイバー化された圏}とは、すべての
ファイバー圏がセトイドである、グルーポイドにファイバー化された圏をいう。
\end{definition}

\noindent
以下では、セトイドにファイバー化された圏と集合にファイバー化された圏の関係を
明確にする。

\begin{definition}
\label{categories-definition-categories-fibred-in-setoids-over-C}
$\mathcal{C}$ を圏とする。{\it $\mathcal{C}$ 上セトイドにファイバー化された圏の
$2$-圏}とは、$\mathcal{C}$ 上グルーポイドにファイバー化された圏の圏
（Definition \ref{categories-definition-categories-fibred-in-groupoids-over-C} を参照）の
部分 $2$-圏であって、次のように定められるものをいう。
\begin{enumerate}
\item その対象は、セトイドにファイバー化された圏
$p : \mathcal{S} \to \mathcal{C}$ とする。
\item その $1$-射 $(\mathcal{S}, p) \to (\mathcal{S}', p')$ は、
$p' \circ G = p$ を満たす関手 $G : \mathcal{S} \to \mathcal{S}'$ とする
（すべての射が強カルテシアンなので、$G$ はそれらを自動的に保つ）。
\item $G, H : (\mathcal{S}, p) \to (\mathcal{S}', p')$ に対するその $2$-射
$t : G \to H$ は、すべての $x \in \Ob(\mathcal{S})$ に対して
$p'(t_x) = \text{id}_{p(x)}$ を満たす関手の射とする。
\end{enumerate}
\end{definition}

\noindent
すべての $2$-射が自動的に同型であることに注意せよ。したがって、この
$2$-圏は実際には $(2, 1)$-圏である。

\noindent
次は $2$-ファイバー積の存在に関する当然の補題である。

\begin{lemma}
\label{categories-lemma-2-product-categories-fibred-setoids}
$\mathcal{C}$ を圏とする。$\mathcal{C}$ 上セトイドにファイバー化された圏の
2-圏は 2-ファイバー積をもつ。より正確には、Lemma
\ref{categories-lemma-2-product-categories-over-C} で記述した 2-ファイバー積は、
セトイドにファイバー化された圏から始めれば、セトイドにファイバー化された圏を返す。
\end{lemma}

\begin{proof}
省略する。
\end{proof}

\begin{lemma}
\label{categories-lemma-setoid-fibres}
$\mathcal{C}$ を圏とし、$\mathcal{S}$ を $\mathcal{C}$ 上の圏とする。
\begin{enumerate}
\item $\mathcal{S} \to \mathcal{S}'$ が $\mathcal{C}$ 上の同値であり、
$\mathcal{S}'$ が $\mathcal{C}$ 上集合にファイバー化されているならば、
\begin{enumerate}
\item $\mathcal{S}$ は $\mathcal{C}$ 上セトイドにファイバー化され、かつ
\item 各 $U \in \Ob(\mathcal{C})$ に対し、写像
$\Ob(\mathcal{S}_U) \to \Ob(\mathcal{S}'_U)$ は、始域の同型類の集合として
終域を同定する。
\end{enumerate}
\item $p : \mathcal{S} \to \mathcal{C}$ がセトイドにファイバー化された圏ならば、
集合にファイバー化された圏
$p' : \mathcal{S}' \to \mathcal{C}$ と、$\mathcal{C}$ 上の同値
$\text{can} : \mathcal{S} \to \mathcal{S}'$ が存在する。
\end{enumerate}
\end{lemma}

\begin{proof}
(2) を証明する。圏 $\mathcal{S}'$ の対象を、$U \in \Ob(\mathcal{C})$ と
$\mathcal{S}_U$ の対象の同型類 $\xi$ からなる対 $(U, \xi)$ とする。
射 $(U, \xi) \to (V , \psi)$ は、$x \in \xi$ および $y \in \psi$ を満たす
対象の射 $x \to y$ で与えられる。ここで、二つの射 $x \to y$ と
$x' \to y'$ が、同じ射 $U \to V$ を誘導し、かつ
$\mathcal{S}_U$ における $x \to x'$ と $\mathcal{S}_V$ における
$y \to y'$ という同型の選択に対して、合成
$x \to x' \to y'$ と $x \to y \to y'$ が一致するとき、これらを同一視する。
構成により、対象と射について $\mathcal{S} \to \mathcal{S}'$ は全射写像となる。
$\mathcal{S}'$ における射の合成を、$\mathcal{S} \to \mathcal{S}'$ を関手にする
唯一の法則として定める。詳細は省略する。
\end{proof}

\noindent
したがって、セトイドにファイバー化された圏とは、グルーポイドにファイバー化
された圏のうち、集合にファイバー化された圏と同値なものにほかならない。
さらに、集合にファイバー化された圏の同値は、Lemma
\ref{categories-lemma-2-category-fibred-sets} により同型である。

\begin{lemma}
\label{categories-lemma-2-category-fibred-setoids}
$\mathcal{C}$ を圏とする。Lemma \ref{categories-lemma-setoid-fibres} の構成 (2) は、
セトイドにファイバー化された圏に前層を対応させる関手
$$
F :
\left\{
\begin{matrix}
\text{the 2-category of categories}\\
\text{fibred in setoids over }\mathcal{C}
\end{matrix}
\right\}
\longrightarrow
\left\{
\begin{matrix}
\text{the category of categories}\\
\text{fibred in sets over }\mathcal{C}
\end{matrix}
\right\}
$$
を与える（Definition \ref{categories-definition-functor-into-2-category} を参照）。
この関手は、次の意味で同値である。
\begin{enumerate}
\item $F(f) = F(g)$ を満たす任意の二つの 1-射
$f, g : \mathcal{S}_1 \to \mathcal{S}_2$ に対して、一意な 2-同型 $f \to g$ が存在する。
\item 任意の射 $h : F(\mathcal{S}_1) \to F(\mathcal{S}_2)$ に対して、
$F(f) = h$ を満たす 1-射 $f : \mathcal{S}_1 \to \mathcal{S}_2$ が存在し、
\item 任意の集合にファイバー化された圏 $\mathcal{S}$ は $F(\mathcal{S})$ に等しい。
\end{enumerate}
特に、$F_i(U) = \Ob(\mathcal{S}_{i, U})/\cong$ という規則で
$F_i \in \textit{PSh}(\mathcal{C})$ を定めると、
$$
\Mor_{\textit{Cat}/\mathcal{C}}(\mathcal{S}_1, \mathcal{S}_2)
\Big/
2\text{-isomorphism}
=
\Mor_{\textit{PSh}(\mathcal{C})}(F_1, F_2)
$$
となる。より正確には、任意の写像 $\phi : F_1 \to F_2$ に対して、
対象の同型類上で $\phi$ を誘導する $1$-射
$f : \mathcal{S}_1 \to \mathcal{S}_2$ が存在し、それは一意な
$2$-同型を除いて一意である。
\end{lemma}

\begin{proof}
Lemma \ref{categories-lemma-2-category-fibred-sets} により、$F$ の終域は圏であり、
したがって主張は意味をもつ。Lemma \ref{categories-lemma-setoid-fibres} の構成 (2) は、
$\mathcal{S}$ に、$U$ 上の値が $\mathcal{S}_U$ の同型類の集合である
集合にファイバー化された圏を対応させる。したがって、これは記述された関手を
定めることが明らかである。
$F(f) = F(g)$ を満たす $f, g : \mathcal{S}_1 \to \mathcal{S}_2$ を (1) のように
とる。$\mathcal{C}$ の各対象 $U$ と $\mathcal{S}_{1, U}$ の各対象 $x$ に対して、
仮定により $f(x) \cong g(x)$ であることが分かる。$\mathcal{S}_2$ はセトイドに
ファイバー化されているので、$\mathcal{S}_{2, U}$ に一意な同型
$t_x : f(x) \to g(x)$ が存在する。明らかに $x \mapsto t_x$ は求める
$2$-同型 $f \to g$ を与える。(2) と (3) の証明は省略する。
最後の主張を見るには、Lemma \ref{categories-lemma-2-category-fibred-sets} により右辺が
$\Mor_{\textit{Cat}/\mathcal{C}}(F(\mathcal{S}_1), F(\mathcal{S}_2))$ に等しいことを
確認し、上の (1) と (2) を適用すればよい。
\end{proof}

\noindent
ここで、すべてのグルーポイドにファイバー化された圏の中から
セトイドにファイバー化された圏を特徴づける別の方法を述べる。

\begin{lemma}
\label{categories-lemma-characterize-fibred-setoids-inertia}
$\mathcal{C}$ を圏とする。$p : \mathcal{S} \to \mathcal{C}$ を
グルーポイドにファイバー化された圏とする。次は同値である。
\begin{enumerate}
\item $p : \mathcal{S} \to \mathcal{C}$ はセトイドにファイバー化された圏であり、
\item 標準的な $1$-射 $\mathcal{I}_\mathcal{S} \to \mathcal{S}$（
\ref{categories-equation-inertia-structure-map} を参照）が
（$\mathcal{C}$ 上の圏の）同値である。
\end{enumerate}
\end{lemma}

\begin{proof}
(2) を仮定する。圏 $\mathcal{I}_\mathcal{S}$ の対象は $(x, \alpha)$ であり、
$x \in \mathcal{S}$、$p(x) = U$ とすると、$\alpha : x \to x$ は
$\mathcal{S}_U$ における射である。
したがって $\mathcal{I}_\mathcal{S} \to \mathcal{S}$ が $\mathcal{C}$ 上の同値ならば、
任意の対象対 $(x, \alpha)$ と $(x, \alpha')$ は
$\mathcal{I}_\mathcal{S}$ の $U$ 上のファイバー圏で同型である。
$\mathcal{I}_\mathcal{S}$ における射の定義を見ると、$\alpha$ と $\alpha'$ は
$x$ の自己同型群で共役である。ここで $\alpha' = \text{id}_x$ とすれば、
$x$ のすべての自己同型は恒等射に等しいことが分かる。
$\mathcal{S} \to \mathcal{C}$ はグルーポイドにファイバー化されているので、
これは $\mathcal{S} \to \mathcal{C}$ がセトイドにファイバー化されていることを意味する。
(1) $\Rightarrow$ (2) の証明は省略する。
\end{proof}

\begin{lemma}
\label{categories-lemma-category-fibred-setoids-presheaves-products}
$\mathcal{C}$ を圏とする。セトイドにファイバー化された圏に前層を対応させる、
Lemma \ref{categories-lemma-2-category-fibred-setoids} の構成は、積と両立する。すなわち、
$2$-ファイバー積 $\mathcal{X} \times_\mathcal{Y} \mathcal{Z}$ に対応する前層は、
$\mathcal{X}, \mathcal{Y}, \mathcal{Z}$ に対応する前層のファイバー積である。
\end{lemma}

\begin{proof}
$U \in \Ob(\mathcal{C})$ とする。この補題は
$$
\Ob((\mathcal{X} \times_\mathcal{Y} \mathcal{Z})_U)/\!\cong
\quad \text{equals} \quad
\Ob(\mathcal{X}_U)/\!\cong
\ \times_{\Ob(\mathcal{Y}_U)/\!\cong}
\ \Ob(\mathcal{Z}_U)/\!\cong
$$
が成り立つというだけであり、その証明は省略する。（ただし、
$\mathcal{Y}_U$ がセトイドでない場合には、一般にはこれは成り立たないことに注意せよ。）
\end{proof}
\section{表現可能なグルーポイドにファイバー化された圏}
\label{categories-section-representable-fibred-groupoids}

\noindent
ここで、グルーポイドにファイバー化された表現可能な圏を定義する。
予告したとおり、これは同値のもとで不変である。

\begin{definition}
\label{categories-definition-representable-fibred-category}
$\mathcal{C}$ を圏とする。
グルーポイドにファイバー化された圏
$p : \mathcal{S} \to \mathcal{C}$ が表現可能であるとは、
ある対象 $X$ が $\mathcal{C}$ に存在し、かつ
同値 $j : \mathcal{S} \to \mathcal{C}/X$ が存在することをいう
（$\mathcal{C}$ 上のグルーポイドにファイバー化された圏の
$2$-圏において）。
\end{definition}

\noindent
通常の記法の濫用として、同値 $j$ に言及せずに
「$X$ が $\mathcal{S}$ を表現する」と言う。
これが意味することを詳しく述べよう。

\begin{lemma}
\label{categories-lemma-characterize-representable-fibred-category}
$\mathcal{C}$ を圏とする。
$p : \mathcal{S} \to \mathcal{C}$ をグルーポイドにファイバー化された圏とする。
\begin{enumerate}
\item $\mathcal{S}$ が表現可能であることと、次の条件が成り立つことは同値である。
\begin{enumerate}
\item $\mathcal{S}$ はセトイドにファイバー化されており、
\item 前層 $U \mapsto \Ob(\mathcal{S}_U)/\cong$ は表現可能である。
\end{enumerate}
\item $\mathcal{S}$ が表現可能ならば、$j$ が
同値 $j : \mathcal{S} \to \mathcal{C}/X$ であるような組 $(X, j)$ は、
同型を除いて一意に定まる。
\end{enumerate}
\end{lemma}

\begin{proof}
第一の主張は補題 \ref{categories-lemma-setoid-fibres} から直ちに従う。
第二の主張について、$j' : \mathcal{S} \to \mathcal{C}/X'$ が
もう一つのそのような組であるとする。
$j' \circ t' \cong \text{id}_{\mathcal{C}/X'}$ および
$t' \circ j' \cong \text{id}_\mathcal{S}$ を満たす
$1$-射 $t' : \mathcal{C}/X' \to \mathcal{S}$ を選ぶ。
すると
$j \circ t' : \mathcal{C}/X' \to \mathcal{C}/X$ は同値である。
したがって、補題 \ref{categories-lemma-2-category-fibred-sets} によれば、
これは同型である。
よって、例えば補題 \ref{categories-lemma-yoneda}（例
\ref{categories-example-fibred-category-from-functor-of-points} を介して）により、
これは同型 $X' \to X$ によって与えられる。
\end{proof}

\begin{lemma}
\label{categories-lemma-morphisms-representable-fibred-categories}
$\mathcal{C}$ を圏とする。
$\mathcal{X}$ と $\mathcal{Y}$ を $\mathcal{C}$ 上の
グルーポイドにファイバー化された圏とする。
$\mathcal{X}$ と $\mathcal{Y}$ はそれぞれ
$\mathcal{C}$ の対象 $X$ と $Y$ によって表現されると仮定する。
このとき
$$
\Mor_{\textit{Cat}/\mathcal{C}}(\mathcal{X}, \mathcal{Y})
\Big/
2\text{-isomorphism}
=
\Mor_\mathcal{C}(X, Y)
$$
が成り立つ。
より正確には、$\phi : X \to Y$ が与えられると、
同型類の対象上で $\phi$ を誘導する
$1$-射 $f : \mathcal{X} \to \mathcal{Y}$ が存在し、
これは一意な $2$-同型を除いて一意である。
\end{lemma}

\begin{proof}
例 \ref{categories-example-fibred-category-from-functor-of-points} により、
$\mathcal{C}/X = \mathcal{S}_{h_X}$ および
$\mathcal{C}/Y = \mathcal{S}_{h_Y}$ である。
補題 \ref{categories-lemma-2-category-fibred-setoids} により、
$$
\Mor_{\textit{Cat}/\mathcal{C}}(\mathcal{X}, \mathcal{Y})
\Big/
2\text{-isomorphism}
=
\Mor_{\textit{PSh}(\mathcal{C})}(h_X, h_Y)
$$
である。
ヨネダ補題 \ref{categories-lemma-yoneda} により、
$\Mor_{\textit{PSh}(\mathcal{C})}(h_X, h_Y)
= \Mor_\mathcal{C}(X, Y)$ である。
\end{proof}
\section{2-ヨネダ補題}
\label{categories-section-2-yoneda}

\noindent
$\mathcal{C}$ を圏とする。$\mathcal{C}$ 上のファイバー化された圏の
$2$-圏は、定義
\ref{categories-definition-fibred-categories-over-C} で構成・定義した。
$\mathcal{S}$ と $\mathcal{S}'$ が $\mathcal{C}$ 上の
ファイバー化された圏ならば、
$$
\Mor_{\textit{Fib}/\mathcal{C}}(\mathcal{S}, \mathcal{S}')
$$
はこの $2$-圏における $1$-射の圏を表す。
ここで、ファイバー化された圏の設定における
ヨネダ補題の $2$-圏類似を述べる。

\begin{lemma}[ファイバー化された圏に対する2-ヨネダ補題]
\label{categories-lemma-yoneda-2category-fibred}
$\mathcal{C}$ を圏とする。
$\mathcal{S} \to \mathcal{C}$ を $\mathcal{C}$ 上のファイバー化された圏とする。
$U \in \Ob(\mathcal{C})$ とする。
関手
$$
\Mor_{\textit{Fib}/\mathcal{C}}(\mathcal{C}/U, \mathcal{S})
\longrightarrow
\mathcal{S}_U
$$
を $G \mapsto G(\text{id}_U)$ によって定めると、これは同値である。
\end{lemma}

\begin{proof}
$\mathcal{S}$ に対する引き戻しの選択を一つ固定する
（定義 \ref{categories-definition-pullback-functor-fibred-category} を参照）。
次のように関手
$$
\mathcal{S}_U
\longrightarrow
\Mor_{\textit{Fib}/\mathcal{C}}(\mathcal{C}/U, \mathcal{S})
$$
を定義する。$x \in \Ob(\mathcal{S}_U)$ が与えられたとき、
対応する関手は次で定まる。
\begin{enumerate}
\item 対象上では、$(f : V \to U) \mapsto f^*x$ とし、
\item 射上では、射 $(g : V'/U \to V/U)$ を
次の合成に写す。
$$
(f \circ g)^*x \xrightarrow{(\alpha_{g, f})_x} g^*f^*x \rightarrow f^*x
$$
ここで $\alpha_{g, f}$ は補題 \ref{categories-lemma-fibred} におけるものである。
\end{enumerate}
これが補題の関手の逆であることの検証は省略する。
\end{proof}

\noindent
$\mathcal{C}$ を圏とする。$\mathcal{C}$ 上のグルーポイドに
ファイバー化された圏の $2$-圏は、$\mathcal{C}$ 上の圏の
$2$-圏の「充満」部分 $2$-圏である
（定義 \ref{categories-definition-categories-fibred-in-groupoids-over-C} を参照）。
したがって、$\mathcal{S}$ と $\mathcal{S}'$ が $\mathcal{C}$ 上で
グルーポイドにファイバー化されているならば、
$$
\Mor_{\textit{Cat}/\mathcal{C}}(\mathcal{S}, \mathcal{S}')
$$
はこの $2$-圏における $1$-射の圏を表す
（定義 \ref{categories-definition-categories-over-C} を参照）。
これらはすべてグルーポイドである。これは定義
\ref{categories-definition-categories-fibred-in-groupoids-over-C} に続く注意を参照。
ここでヨネダ補題の $2$-圏類似を述べる。

\begin{lemma}[2-ヨネダ補題]
\label{categories-lemma-yoneda-2category}
$\mathcal{S}\to \mathcal{C}$ をグルーポイドにファイバー化された圏とする。
$U \in \Ob(\mathcal{C})$ とする。
関手
$$
\Mor_{\textit{Cat}/\mathcal{C}}(\mathcal{C}/U, \mathcal{S})
\longrightarrow
\mathcal{S}_U
$$
を $G \mapsto G(\text{id}_U)$ によって定めると、これは同値である。
\end{lemma}

\begin{proof}
$\mathcal{S}$ に対する引き戻しの選択を一つ固定する
（定義 \ref{categories-definition-pullback-functor-fibred-category} を参照）。
次のように関手
$$
\mathcal{S}_U
\longrightarrow
\Mor_{\textit{Cat}/\mathcal{C}}(\mathcal{C}/U, \mathcal{S})
$$
を定義する。$x \in \Ob(\mathcal{S}_U)$ が与えられたとき、
対応する関手は次で定まる。
\begin{enumerate}
\item 対象上では、$(f : V \to U) \mapsto f^*x$ とし、
\item 射上では、射 $(g : V'/U \to V/U)$ を
次の合成に写す。
$$
(f \circ g)^*x \xrightarrow{(\alpha_{g, f})_x} g^*f^*x \rightarrow f^*x
$$
ここで $\alpha_{g, f}$ は補題 \ref{categories-lemma-fibred-groupoids} におけるものである。
\end{enumerate}
これが補題の関手の逆であることの検証は省略する。
\end{proof}

\begin{remark}
\label{categories-remark-alternative-fibred-groupoids-strict}
$2$-ヨネダ補題を用いると、補題
\ref{categories-lemma-fibred-groupoids-strict} の別証明を与えられる。
$p : \mathcal{S} \to \mathcal{C}$ をグルーポイドにファイバー化された圏とする。
次のように、$\mathcal{C}$ からグルーポイドの圏への反変関手
$F$ を定義する。$U\in \Ob(\mathcal{C})$ に対して
$$
F(U) = \Mor_{\textit{Cat}/\mathcal{C}}(\mathcal{C}/U, \mathcal{S}).
$$
$f : U \to V$ ならば、誘導された関手
$\mathcal{C}/U \to \mathcal{C}/V$ は射
$F(f) : F(V) \to F(U)$ を誘導する。明らかに $F$ は関手である。
例 \ref{categories-example-functor-groupoids} から得られる対応する
グルーポイドにファイバー化された圏を $\mathcal{S}'$ とする。
$\mathcal{C}$ 上の明らかな関手 $G : \mathcal{S}' \to \mathcal{S}$ があり、
$U \in \Ob(\mathcal{C})$ かつ $x \in F(U)$ である組 $(U, x)$ を
$x(\text{id}_U) \in \mathcal{S}$ に写す。
ここで補題 \ref{categories-lemma-yoneda-2category} により、各 $U$ について
$$
G_U : \mathcal{S}'_U = F(U)=
\Mor_{\textit{Cat}/\mathcal{C}}(\mathcal{C}/U, \mathcal{S})
\to
\mathcal{S}_U
$$
は同値であり、したがって補題
\ref{categories-lemma-equivalence-fibred-categories} により、$G$ は
$\mathcal{S}$ と $\mathcal{S}'$ の間の同値である。
\end{remark}
\section{表現可能な1-射}
\label{categories-section-representable-1-morphisms}

\noindent
$\mathcal{C}$ を圏とする。
この節では、$\mathcal{C}$ 上のグルーポイドにファイバー化された圏の間の
$1$-射が表現可能であるとは何を意味するかを説明する。

\medskip\noindent
$\mathcal{C}$ を圏とする。
$\mathcal{X}$ と $\mathcal{Y}$ を $\mathcal{C}$ 上の
グルーポイドにファイバー化された圏とする。
$U \in \Ob(\mathcal{C})$ とする。
$F : \mathcal{X} \to \mathcal{Y}$ および
$G : \mathcal{C}/U \to \mathcal{Y}$ を、
$\mathcal{C}$ 上のグルーポイドにファイバー化された圏の
$1$-射とする。
次の $2$-ファイバー積を記述したい。
$$
\xymatrix{
(\mathcal{C}/U) \times_\mathcal{Y} \mathcal{X} \ar[r] \ar[d] &
\mathcal{X} \ar[d]^F \\
\mathcal{C}/U \ar[r]^G &
\mathcal{Y}
}
$$
$y = G(\text{id}_U) \in \mathcal{Y}_U$ とおく。
$\mathcal{Y}$ に対する引き戻しを一つ選ぶ
（定義 \ref{categories-definition-pullback-functor-fibred-category} を参照）。
すると $G$ は関手 $(f : V \to U) \mapsto f^*y$ と同型である。
補題 \ref{categories-lemma-yoneda-2category} とその証明を参照せよ。
補題 \ref{categories-lemma-2-product-categories-over-C} により、
$(\mathcal{C}/U) \times_\mathcal{Y} \mathcal{X}$ の対象を
四つ組 $(V, f : V \to U, x, \phi)$ とみなせる。
上の $G$ の記述を用いると、$\phi$ は
$\mathcal{Y}_V$ における同型
$\phi : f^*y \to F(x)$ とみなせる。

\begin{lemma}
\label{categories-lemma-identify-fibre-product}
上の状況において、$\mathcal{C}/U$ の対象
$f : V \to U$ 上の
$(\mathcal{C}/U) \times_\mathcal{Y} \mathcal{X}$ のファイバー圏は、
次のように記述される圏である。
\begin{enumerate}
\item 対象は組 $(x, \phi)$ であり、
$x \in \Ob(\mathcal{X}_V)$ および
$\phi : f^*y \to F(x)$ は $\mathcal{Y}_V$ における射である。
\item $(x, \phi)$ と $(x', \phi')$ の間の射の集合は、
$\mathcal{X}_V$ における射 $\psi : x \to x'$ であって
$F(\psi) = \phi' \circ \phi^{-1}$ を満たすものの集合である。
\end{enumerate}
\end{lemma}

\begin{proof}
上の議論を参照せよ。
\end{proof}

\begin{lemma}
\label{categories-lemma-prepare-representable-map-stack-in-groupoids}
$\mathcal{C}$ を圏とする。
$\mathcal{X}$ と $\mathcal{Y}$ を $\mathcal{C}$ 上の
グルーポイドにファイバー化された圏とする。
$F : \mathcal{X} \to \mathcal{Y}$ を $1$-射とする。
$G : \mathcal{C}/U \to \mathcal{Y}$ を $1$-射とする。
このとき
$$
(\mathcal{C}/U) \times_\mathcal{Y} \mathcal{X}
\longrightarrow
\mathcal{C}/U
$$
はグルーポイドにファイバー化された圏である。
\end{lemma}

\begin{proof}
補題 \ref{categories-lemma-2-product-fibred-categories} で、
合成
$$
(\mathcal{C}/U) \times_\mathcal{Y} \mathcal{X}
\longrightarrow
\mathcal{C}/U
\longrightarrow
\mathcal{C}
$$
がグルーポイドにファイバー化された圏であることをすでに見た。
したがって、この補題は補題 \ref{categories-lemma-cute-groupoids} から従う。
\end{proof}

\begin{definition}
\label{categories-definition-representable-map-categories-fibred-in-groupoids}
$\mathcal{C}$ を圏とする。
$\mathcal{X}$ と $\mathcal{Y}$ を $\mathcal{C}$ 上の
グルーポイドにファイバー化された圏とする。
$F : \mathcal{X} \to \mathcal{Y}$ を $1$-射とする。
$F$ が{\it 表現可能}である、または
{\it $\mathcal{X}$ が $\mathcal{Y}$ 上で相対的に表現可能}であるとは、
任意の $U \in \Ob(\mathcal{C})$ および任意の
$G : \mathcal{C}/U \to \mathcal{Y}$ に対して、
グルーポイドにファイバー化された圏
$$
(\mathcal{C}/U) \times_\mathcal{Y} \mathcal{X}
\longrightarrow
\mathcal{C}/U
$$
が表現可能であることをいう。
\end{definition}

\begin{lemma}
\label{categories-lemma-spell-out-representable-map-stack-in-groupoids}
$\mathcal{C}$ を圏とする。
$\mathcal{X}$ と $\mathcal{Y}$ を $\mathcal{C}$ 上の
グルーポイドにファイバー化された圏とする。
$F : \mathcal{X} \to \mathcal{Y}$ を $1$-射とする。
$F$ が表現可能ならば、各ファイバー圏の間の関手
$$
F_U : \mathcal{X}_U \longrightarrow \mathcal{Y}_U
$$
はすべて忠実である。
\end{lemma}

\begin{proof}
これはファイバー圏の記述、すなわち補題
\ref{categories-lemma-identify-fibre-product} と、表現可能な
ファイバー化された圏の特徴付け、すなわち補題
\ref{categories-lemma-characterize-representable-fibred-category} から明らかである。
\end{proof}

\begin{lemma}
\label{categories-lemma-criterion-representable-map-stack-in-groupoids}
$\mathcal{C}$ を圏とする。
$\mathcal{X}$ と $\mathcal{Y}$ を $\mathcal{C}$ 上の
グルーポイドにファイバー化された圏とする。
$F : \mathcal{X} \to \mathcal{Y}$ を $1$-射とする。
$\mathcal{Y}$ に対する引き戻しを一つ選ぶ。
\begin{enumerate}
\item 各ファイバー圏の間の関手
$F_U : \mathcal{X}_U \longrightarrow \mathcal{Y}_U$ が忠実であり、かつ
\item 各 $U$ および各 $y \in \mathcal{Y}_U$ に対して、前層
$$
(f : V \to U)
\longmapsto
\{(x, \phi) \mid x \in \mathcal{X}_V, \phi : f^*y \to F(x)\}/\cong
$$
が $\mathcal{C}/U$ 上の表現可能な前層である。
\end{enumerate}
ならば $F$ は表現可能である。
\end{lemma}

\begin{proof}
これはファイバー圏の記述、すなわち補題
\ref{categories-lemma-identify-fibre-product} と、表現可能な
ファイバー化された圏の特徴付け、すなわち補題
\ref{categories-lemma-characterize-representable-fibred-category} から明らかである。
\end{proof}

\noindent
次の補題を述べる前に、グルーポイドにファイバー化された圏の
$2$-圏が $(2, 1)$-圏であり、したがってそれが終対象をもつという
意味が分かっていることを指摘しておく
（定義 \ref{categories-definition-final-object-2-category} を参照）。
そして実際、終対象は
$\text{id} : \mathcal{C} \to \mathcal{C}$ である。
そこで、$\mathcal{C}$ 上のグルーポイドにファイバー化された圏の
{\it $2$-積}を、次の $2$-ファイバー積として定義する。
$$
\mathcal{X} \times \mathcal{Y} :=
\mathcal{X} \times_\mathcal{C} \mathcal{Y}.
$$
この定義により、次の補題は意味を持つ。

\begin{lemma}
\label{categories-lemma-representable-diagonal-groupoids}
$\mathcal{C}$ を圏とする。
$\mathcal{S} \to \mathcal{C}$ をグルーポイドにファイバー化された圏とする。
$\mathcal{C}$ は対象の任意の二つ組の積とファイバー積をもつと仮定する。
次の条件は同値である。
\begin{enumerate}
\item 対角 $\mathcal{S} \to \mathcal{S} \times \mathcal{S}$ は表現可能である。
\item $\mathcal{C}$ の任意の $U$ に対して、任意の
$G : \mathcal{C}/U \to \mathcal{S}$ は表現可能である。
\end{enumerate}
\end{lemma}

\begin{proof}
対角が表現可能であると仮定し、$U, G$ が与えられたとする。
$V \in \Ob(\mathcal{C})$ と任意の
$G' : \mathcal{C}/V \to \mathcal{S}$ を考える。
$\mathcal{C}/U \times \mathcal{C}/V = \mathcal{C}/U \times V$ は
表現可能であることに注意する。したがって、ファイバー積
$$
\xymatrix{
(\mathcal{C}/U \times V)
\times_{(\mathcal{S} \times \mathcal{S})}
\mathcal{S}
\ar[r] \ar[d] &
\mathcal{S} \ar[d] \\
\mathcal{C}/U \times V \ar[r]^{(G, G')} &
\mathcal{S} \times \mathcal{S}
}
$$
は仮定により表現可能である。
これは、$\mathcal{C}$ に $W \to U \times V$ が存在して、
$$
\xymatrix{
\mathcal{C}/W \ar[d] \ar[r] & \mathcal{S} \ar[d] \\
\mathcal{C}/U \times \mathcal{C}/V \ar[r] & \mathcal{S} \times \mathcal{S}
}
$$
がカルテシアンであることを意味する。したがって
$\mathcal{C}/W \cong \mathcal{C}/U \times_\mathcal{S} \mathcal{C}/V$
である（補題 \ref{categories-lemma-diagonal-1} と注意
\ref{categories-remark-2-product-fibred-categories-same} を参照）。
これが望んだ結論である。

\medskip\noindent
(2) を仮定する。$V \in \Ob(\mathcal{C})$ と、任意の
$(G, G') : \mathcal{C}/V \to \mathcal{S} \times \mathcal{S}$ を考える。
$\mathcal{C}/V \times_{\mathcal{S} \times \mathcal{S}} \mathcal{S}$
が表現可能であることを示さなければならない。既知なのは
$\mathcal{C}/V \times_{G, \mathcal{S}, G'} \mathcal{C}/V$
が、$a : W \to V$（$\mathcal{C}/V$ 内）によって表現されるということである。
同値
$$
\mathcal{C}/W \to \mathcal{C}/V \times_{G, \mathcal{S}, G'} \mathcal{C}/V
$$
第二射影を $\mathcal{C}/V$ へ続けると、第二の射 $a' : W \to V$ が得られる。ここで
$W' = W \times_{(a, a'), V \times V} V$ とおく。
次の同値が存在する。
$$
\mathcal{C}/W' \cong
\mathcal{C}/V \times_{\mathcal{S} \times \mathcal{S}} \mathcal{S}
$$
実際、
\begin{eqnarray*}
\mathcal{C}/W' & \cong &
\mathcal{C}/W \times_{(\mathcal{C}/V \times \mathcal{C}/V)} \mathcal{C}/V \\
& \cong &
\left(\mathcal{C}/V \times_{(G, \mathcal{S}, G')} \mathcal{C}/V\right)
\times_{(\mathcal{C}/V \times \mathcal{C}/V)} \mathcal{C}/V \\
& \cong &
\mathcal{C}/V \times_{(\mathcal{S} \times \mathcal{S})} \mathcal{S}
\end{eqnarray*}
となる（最後の同型については補題 \ref{categories-lemma-diagonal-2} と注意
\ref{categories-remark-2-product-fibred-categories-same} を参照）。これで補題が示された。
\end{proof}

\noindent
{\bf 文献注：}
この一部は Vistoli のノート \cite{Vis2} から取った。
\section{モノイド圏}
\label{categories-section-monoidal}

\noindent
$\mathcal{C}$ を圏とする。関手
$$
\otimes : \mathcal{C} \times \mathcal{C} \longrightarrow \mathcal{C}
$$
が与えられているとする。
$\otimes$ が結合則を満たすか、また $\otimes$ に単位元があるかを
知りたいことが多い。

\medskip\noindent
$(\mathcal{C}, \otimes)$ に対する{\it 結合性制約}とは、次を満たす
関手的同型
$$
\phi_{X, Y, Z} : X \otimes (Y \otimes Z) \to (X \otimes Y) \otimes Z
$$
である。すべての対象 $X, Y, Z, W$ に対して、図式
\begin{equation}
\label{categories-equation-pentagon}
\vcenter{
\xymatrix{
X \otimes (Y \otimes ( Z \otimes W)) \ar[r] \ar[d] &
X \otimes ((Y \otimes Z) \otimes W) \ar[d] \\
(X \otimes Y) \otimes (Z \otimes W) \ar[d] &
(X \otimes (Y \otimes Z)) \otimes W \ar[ld] \\
((X \otimes Y) \otimes Z) \otimes W
}
}
\end{equation}
が可換となる。ただし、各射は $\phi$ の適切な適用と
$\otimes$ の関手性によって定まる。

\medskip\noindent
\cite[定理3.1]{associativity} では、上のような三つ組
$(\mathcal{C}, \otimes, \phi)$ がコヒーレントであることが説明されている。
言い換えると、このような三つ組から、すべての $n \geq 1$ に対して
よく定義された関手の体系
$$
\mathcal{C} \times \ldots \times \mathcal{C} \longrightarrow \mathcal{C},
\quad
(X_1, \ldots, X_n) \longmapsto X_1 \otimes \ldots \otimes X_n
$$
および $n, m \geq 1$ に対する関手的同型
$$
(X_1 \otimes \ldots \otimes X_n) \otimes (Y_1 \otimes \ldots \otimes Y_m)
\to
X_1 \otimes \ldots \otimes X_n \otimes Y_1 \otimes \ldots \otimes Y_m
$$
が得られ、これらの同型から作られる可能なすべての図式が可換となる
（上の五角形図式は一例である）。
$n = 5$ のとき、これは次の式が
$$
\begin{matrix}
(((X_1 \otimes X_2) \otimes X_3) \otimes X_4) \otimes X_5 &
((X_1 \otimes (X_2 \otimes X_3)) \otimes X_4) \otimes X_5 \\
((X_1 \otimes X_2) \otimes (X_3 \otimes X_4)) \otimes X_5 &
(X_1 \otimes ((X_2 \otimes X_3) \otimes X_4)) \otimes X_5 \\
(X_1 \otimes (X_2 \otimes (X_3 \otimes X_4))) \otimes X_5 &
((X_1 \otimes X_2) \otimes X_3) \otimes (X_4 \otimes X_5) \\
(X_1 \otimes (X_2 \otimes X_3)) \otimes (X_4 \otimes X_5) &
(X_1 \otimes X_2) \otimes ((X_3 \otimes X_4) \otimes X_5) \\
(X_1 \otimes X_2) \otimes (X_3 \otimes (X_4 \otimes X_5)) &
X_1 \otimes (((X_2 \otimes X_3) \otimes X_4) \otimes X_5) \\
X_1 \otimes ((X_2 \otimes (X_3 \otimes X_4)) \otimes X_5) &
X_1 \otimes ((X_2 \otimes X_3) \otimes (X_4 \otimes X_5)) \\
X_1 \otimes (X_2 \otimes ((X_3 \otimes X_4) \otimes X_5)) &
X_1 \otimes (X_2 \otimes (X_3 \otimes (X_4 \otimes X_5)))
\end{matrix}
$$
すべて $X_1 \otimes \ldots \otimes X_5$ と関手的に同型であり、
これらの同型が、$\phi$ を可能な三つの連続する場所に用いて得られる
これらの関手間のすべての同型と整合することを意味する。
以下ではこれを断りなく用い、$\otimes$ の反復を書くときに括弧を
もはや記さない。

\medskip\noindent
上のような三つ組 $(\mathcal{C}, \otimes, \phi)$ に対する{\it 単位}とは、
$\mathcal{C}$ の対象 $\mathbf{1}$ と、関手的同型
$$
l : \mathbf{1} \otimes X \to X
\quad\text{and}\quad
r : X \otimes \mathbf{1} \to X
$$
の組であって、すべての対象 $X, Y$ に対して図式
\begin{equation}
\label{categories-equation-triangle}
\vcenter{
\xymatrix{
X \otimes (\mathbf{1} \otimes Y) \ar[rr]_\phi \ar[rd]_{\text{id} \otimes l} & &
(X \otimes \mathbf{1}) \otimes Y \ar[ld]^{r \otimes \text{id}} \\
& X \otimes Y
}
}
\end{equation}
が可換となるものである。以下の補題のように、単位を組
$(\mathbf{1}, 1)$ と考えることが多い。

\begin{lemma}
\label{categories-lemma-monoidal-unit}
$(\mathcal{C}, \otimes, \phi)$ を上記のものとする。
$\mathcal{C}$ における単位 $(\mathbf{1}, l, r)$ と、
$\mathbf{1}$ が $\mathcal{C}$ の対象であり
$1 : \mathbf{1} \otimes \mathbf{1} \to \mathbf{1}$ が同型で、
関手 $L : X \mapsto \mathbf{1} \otimes X$ と
$R : X \mapsto X \otimes \mathbf{1}$ が同値であるような組
$(\mathbf{1}, 1)$ の間には、一対一対応がある。
\end{lemma}

\begin{proof}
単位 $(\mathbf{1}, l, r)$ が与えられると、同型
$r : \mathbf{1} \otimes \mathbf{1} \to \mathbf{1}$ が得られ、
$L$ と $R$ は恒等関手と同型なので同値である。
逆に、$L$ と $R$ が同値となる組 $(\mathbf{1}, 1)$ が与えられたとする。
関手的同型 $l_X : \mathbf{1} \otimes X \to X$ と
$r_X : X \otimes \mathbf{1} \to X$ を、
$L(l_X) = 1 \otimes \text{id}_X$ および
$R(r_X) = \text{id}_X \otimes 1$ によって特徴付ける。
次に、すべての $X$ と $Y$ に対し、
$X \otimes \mathbf{1} \otimes Y$ から $X \otimes Y$ への二つの射
$\text{id}_X \otimes l_Y$ と $r_X \otimes \text{id}_Y$ が
同じであることを示さなければならない。この性質は $X$ と $Y$ の
同型類だけに依存する。
$R$ と $L$ は同値なので、ある対象 $Z$ と $W$ に対して
$X = Z \otimes \mathbf{1}$ および $Y = \mathbf{1} \otimes W$ の場合だけ
示せば十分である。
言い換えると、
$$
\text{id}_Z \otimes \text{id}_\mathbf{1} \otimes l_{\mathbf{1} \otimes W}
=
r_{Z \otimes \mathbf{1}} \otimes \text{id}_\mathbf{1} \otimes \text{id}_Y
$$
を示す必要がある。構成により、これらの写像は
$\text{id}_Z \otimes 1 \otimes \text{id}_\mathbf{1} \otimes \text{id}_W$ と
$\text{id}_Z \otimes \text{id}_\mathbf{1} \otimes 1 \otimes \text{id}_W$
にそれぞれ等しい。したがって
$1 \otimes \text{id}_\mathbf{1} = \text{id}_\mathbf{1} \otimes 1$ を
示せば十分である。

\medskip\noindent
次が成り立つ。
$\text{id}_\mathbf{1} \otimes 1 = r_\mathbf{1} \otimes \text{id}_\mathbf{1}$
および
$l_\mathbf{1} \otimes \text{id}_\mathbf{1} =
\text{id}_\mathbf{1} \otimes 1$ である。
ある $\mathbf{1}$ の自己同型 $a, b$ に対して
$r_\mathbf{1} = a \circ 1$ および $l_\mathbf{1} = b \circ 1$ と書ける。
したがって
$\text{id}_\mathbf{1} \otimes 1 =
(a \otimes \text{id}_\mathbf{1}) \circ
(1 \otimes \text{id}_\mathbf{1})$
および
$1 \otimes \text{id}_\mathbf{1} =
(\text{id}_\mathbf{1} \otimes b) \circ
(\text{id}_\mathbf{1} \otimes 1)$
である。
すると
\begin{align*}
(1 \otimes \text{id}_\mathbf{1}) \circ
(\text{id}_\mathbf{1} \otimes 1 \otimes \text{id}_\mathbf{1})
& =
(1 \otimes \text{id}_\mathbf{1}) \circ
(\text{id}_\mathbf{1} \otimes \text{id}_\mathbf{1} \otimes b)
\circ
(\text{id}_\mathbf{1} \otimes \text{id}_\mathbf{1} \otimes 1) \\
& =
(\text{id}_\mathbf{1} \otimes b) \circ
(1 \otimes \text{id}_\mathbf{1}) \circ
(\text{id}_\mathbf{1} \otimes \text{id}_\mathbf{1} \otimes 1) \\
& =
(\text{id}_\mathbf{1} \otimes b) \circ
(1 \otimes 1)
\end{align*}
であり、また
\begin{align*}
(1 \otimes \text{id}_\mathbf{1}) \circ
(\text{id}_\mathbf{1} \otimes 1 \otimes \text{id}_\mathbf{1})
& =
(\text{id}_\mathbf{1} \otimes b) \circ
(\text{id}_\mathbf{1} \otimes 1) \circ
(a \otimes \text{id}_\mathbf{1} \otimes \text{id}_\mathbf{1}) \circ
(1 \otimes \text{id}_\mathbf{1} \otimes \text{id}_\mathbf{1}) \\
& =
(\text{id}_\mathbf{1} \otimes b) \circ
(a \otimes \text{id}_\mathbf{1}) \circ
(1 \otimes 1)
\end{align*}
である。
これは $a \otimes \text{id}_\mathbf{1}$ が恒等であることを示し、
したがって所望どおり $a$ は恒等である。

\medskip\noindent
証明を終えるため、上の規則が、（適切に定義した）単位の圏と
組 $(\mathbf{1}, 1)$ の圏との間の逆向きの圏同値を定めることに注意する。
\end{proof}

\begin{lemma}
\label{categories-lemma-monoidal-unit-unique}
$(\mathcal{C}, \otimes, \phi)$ を上記のものとする。
$(\mathbf{1}, 1)$ を単位とする（補題 \ref{categories-lemma-monoidal-unit} を参照）。
このとき
\begin{enumerate}
\item $1 \otimes \text{id}_\mathbf{1} = \text{id}_\mathbf{1} \otimes 1$
\item $\Gamma = \text{Mor}(\mathbf{1}, \mathbf{1})$ は可換モノイドである。
\item $a = 1 \circ (a \otimes \text{id}_\mathbf{1}) \circ 1^{-1} =
1 \circ (\text{id}_\mathbf{1} \otimes a) \circ 1^{-1}$ はすべての
$a \in \Gamma$ に対して成り立つ。
\item 他の任意の単位は $(\mathbf{1}, 1)$ と一意な同型によって同型である。
\end{enumerate}
\end{lemma}

\begin{proof}
(1) は補題 \ref{categories-lemma-monoidal-unit} の証明で示した。
$a \in \Gamma$ に対して
\begin{align*}
(1 \circ (a \otimes \text{id}_\mathbf{1})) \otimes \text{id}_\mathbf{1}
& =
(1 \otimes \text{id}_\mathbf{1}) \circ
(a \otimes \text{id}_\mathbf{1} \otimes \text{id}_\mathbf{1}) \\
& =
(\text{id}_\mathbf{1} \otimes 1) \circ
(a \otimes \text{id}_\mathbf{1} \otimes \text{id}_\mathbf{1}) \\
& =
(a \otimes \text{id}_\mathbf{1}) \circ
(\text{id}_\mathbf{1} \otimes 1) \\
& =
(a \otimes \text{id}_\mathbf{1}) \circ
(1 \otimes \text{id}_\mathbf{1}) \\
& =
(a \circ 1) \otimes \text{id}_\mathbf{1}
\end{align*}
である。したがって
$1 \circ (a \otimes \text{id}_\mathbf{1}) = a \circ 1$ であり、
これは (3) の半分を意味する。残りの半分も同様に従う。
(2) を見るため、$a, b \in \Gamma$ に対して
$$
a \circ b =
1 \circ (a \otimes \text{id}_\mathbf{1}) \circ 1^{-1}
\circ
1 \circ (\text{id}_\mathbf{1} \otimes b) \circ 1^{-1}
=
1 \circ (a \otimes b) \circ \circ 1^{-1}
$$
であり、$b \circ a$ も同じ式に評価されることに注意する。
したがって (2) が成り立つ。
(4) を見るため、$(\mathbf{1}', 1')$ が第二の単位であるとする。
$r$ と $l'$ を用いると、
$\mathbf{1}' \otimes \mathbf{1} \to \mathbf{1}'$ と
$\mathbf{1}' \otimes \mathbf{1} \to \mathbf{1}$ という同型がある。
したがって同型 $t : \mathbf{1}' \to \mathbf{1}$ が存在する。
すると図式
$$
\xymatrix{
\mathbf{1}' \otimes \mathbf{1}' \ar[r]_{t \otimes t} \ar[d]_{1'} &
\mathbf{1} \otimes \mathbf{1} \ar[d]^1 \\
\mathbf{1}' \ar[r]^t & \mathbf{1}
}
$$
は $\mathbf{1}$ の自己同型 $a$ を除いて可換である。
$t$ を $a \circ t$ で置き換えれば、図式は可換となる
（ヒント： (3) を用いて $1 \circ (a \otimes a) = a^2 \circ 1$ を
見るとよい）。
\end{proof}
\begin{definition}
\label{categories-definition-monoidal-category}
$\mathcal{C}$ が圏であり、$\otimes : \mathcal{C} \times \mathcal{C} \to \mathcal{C}$ が関手で、
$\phi$ が結合性制約であるような三つ組
$(\mathcal{C}, \otimes, \phi)$ を、単位 $\mathbf{1}$ が存在するとき
{\it モノイド圏}と呼ぶ。
\end{definition}

\noindent
モノイド圏の単位を表すために常に $\mathbf{1}$ と書き、
$1 : \mathbf{1} \otimes \mathbf{1} \to \mathbf{1}$ を選んだ同型とする。
組 $(\mathbf{1}, 1)$ は一意な同型を除いて定まるので
（補題 \ref{categories-lemma-monoidal-unit-unique}）、一つを選んでも差し支えない。

\begin{lemma}
\label{categories-lemma-monoidal-coherence}
モノイド圏 $\mathcal{C}, \otimes, \phi, \mathbf{1}, 1$ において、
補題 \ref{categories-lemma-monoidal-unit} の証明における記法を用いると、
\begin{enumerate}
\item 射 $1, r_\mathbf{1}, l_\mathbf{1} :
\mathbf{1} \otimes \mathbf{1} \to \mathbf{1}$ は一致し、
\item 射
$l_X \otimes \text{id}_Y, l_{X \otimes Y} :
\mathbf{1} \otimes X \otimes Y \to X \otimes Y$ は一致し、
\item 射
$\text{id}_X \otimes r_Y , r_{X \otimes Y} :
X \otimes Y \otimes \mathbf{1} \to X \otimes Y$ は一致する。
\end{enumerate}
モノイド圏は \cite[定理5.2]{associativity} の仮定を満たす。
\end{lemma}

\begin{proof}
(1) は補題 \ref{categories-lemma-monoidal-unit} の証明で見た。
補題 \ref{categories-lemma-monoidal-unit} の証明において、
$l_X$ と $l_{X \otimes Y}$ は
$\text{id}_\mathbf{1} \otimes l_{X \otimes Y} =
1 \otimes \text{id}_{X \otimes Y}$ および
$\text{id}_\mathbf{1} \otimes l_X = 1 \otimes \text{id}_X$ を満たす
一意な射であることを見た。(2) は直ちに従う。(3) も同様に証明される。
これと (\ref{categories-equation-pentagon}) および
(\ref{categories-equation-triangle}) の可換性を合わせると、補題の最後の主張が成り立つ。
\end{proof}

\noindent
\cite[定理5.2]{associativity} では、上の補題のような五つ組
$(\mathcal{C}, \otimes, \phi, \mathbf{1}, 1)$ がコヒーレントであることが説明されている。
言い換えると、モノイド圏（ここでの意味）では、すべての $n \geq i \geq 0$ に対して
関手的同型
$$
X_1 \otimes \ldots \otimes X_i \otimes \mathbf{1}
\otimes X_{i + 1} \otimes \ldots \otimes X_n
\to
X_1 \otimes \ldots \otimes X_n
$$
があり、これらの同型から作られるすべての図式が可換となる。
例えば $\mathbf{1}$ を二度挿入して異なる順序で同型を使っても、
同じ射が得られる。$\otimes$ の反復を書く際に括弧を省くという
規約に加え、今後は $X \otimes \mathbf{1}$ と
$\mathbf{1} \otimes X$ を、断りなく $X$ と同一視する。
さらに、$\otimes, \phi, \mathbf{1}$ を理解したうえで
「$\mathcal{C}$ をモノイド圏とする」と言う。

\begin{definition}
\label{categories-definition-functor-monoidal-categories}
$\mathcal{C}$ と $\mathcal{C}'$ をモノイド圏とする。
{\it モノイド圏の関手} $F : \mathcal{C} \to \mathcal{C}'$ とは、
指定された関手 $F$ と、$X$ と $Y$ に関して関手的な同型
$$
F(X) \otimes F(Y) \to F(X \otimes Y)
$$
の組であって、すべての対象 $X$, $Y$, $Z$ に対して図式
$$
\xymatrix{
F(X) \otimes (F(Y) \otimes F(Z)) \ar[r] \ar[d] &
F(X) \otimes F(Y \otimes Z) \ar[r] &
F(X \otimes (Y \otimes Z)) \ar[d] \\
(F(X) \otimes F(Y)) \otimes F(Z) \ar[r] &
F(X \otimes Y) \otimes F(Z) \ar[r] &
F((X \otimes Y) \otimes Z)
}
$$
が可換であり、かつ $F(\mathbf{1})$ が $\mathcal{C}'$ の単位となるものをいう。
\end{definition}

\noindent
単位に関する規約により、$F$ がモノイド圏の関手ならば
常に $F(\mathbf{1}) = \mathbf{1}$ と仮定できる。
例えば、$A \to B$ が環準同型ならば、関手
$M \mapsto M \otimes_A B$ は $\text{Mod}_A$ から $\text{Mod}_B$ への
モノイド圏の関手である。

\begin{lemma}
\label{categories-lemma-invertible}
$\mathcal{C}$ をモノイド圏、$X$ を $\mathcal{C}$ の対象とする。
次の条件は同値である。
\begin{enumerate}
\item 関手 $L : Y \mapsto X \otimes Y$ は同値である。
\item 関手 $R : Y \mapsto Y \otimes X$ は同値である。
\item $X \otimes X' \cong X' \otimes X \cong \mathbf{1}$ を満たす対象
$X'$ が存在する。
\end{enumerate}
\end{lemma}

\begin{proof}
(1) を仮定する。$L(X') = \mathbf{1}$、すなわち
$X \otimes X' \cong \mathbf{1}$ となるような $X'$ を選ぶ。
$X'$ に対応する関手を $L'$ および $R'$ と書く。
$X \otimes X' \cong \mathbf{1}$ という式から
$L \circ L' \cong \text{id}$ が従う。したがって、$L'$ は
(仮定により存在する) $L$ の準逆でなければならない。
よって $L' \circ L \cong \text{id}$ である。
したがって $X' \otimes X \cong \mathbf{1}$ であり、(3) が成り立つ。

\medskip\noindent
(2) $\Rightarrow$ (3) の証明は、いま述べたものの双対である。

\medskip\noindent
(3) を仮定する。このとき $L'$ と $L$ は互いに準逆であり、
$R'$ と $R$ も互いに準逆であることは明らかである。
したがって (1) と (2) が成り立つ。
\end{proof}

\begin{definition}
\label{categories-definition-invertible}
$\mathcal{C}$ をモノイド圏とする。$\mathcal{C}$ の対象 $X$ が{\it 可逆}であるとは、
補題 \ref{categories-lemma-invertible} の同値な条件のいずれか（したがってすべて）が
成り立つことをいう。
\end{definition}

\noindent
$F : \mathcal{C} \to \mathcal{C}'$ がモノイド圏の関手ならば、
$F$ は可逆対象を可逆対象へ写すことに注意する。
\begin{definition}
\label{categories-definition-dual}
モノイド圏 $(\mathcal{C}, \otimes, \phi)$ と対象 $X$ が与えられたとする。
{\it 左双対}とは、射
$\eta : \mathbf{1} \to X \otimes Y$ および
$\epsilon : Y \otimes X \to \mathbf{1}$ を備えた対象 $Y$ であって、
図式
$$
\vcenter{
\xymatrix{
X \ar[rd]_1 \ar[r]_-{\eta \otimes 1} &
X \otimes Y \otimes X \ar[d]^{1 \otimes \epsilon} \\
& X
}
}
\quad\text{and}\quad
\vcenter{
\xymatrix{
Y \ar[rd]_1 \ar[r]_-{1 \otimes \eta} &
Y \otimes X \otimes Y \ar[d]^{\epsilon \otimes 1} \\
& Y
}
}
$$
が可換となるものをいう。この状況で、$X$ は $Y$ の{\it 右双対}であるという。
\end{definition}

\noindent
$F : \mathcal{C} \to \mathcal{C}'$ がモノイド圏の関手ならば、
$Y$ が $X$ の左双対であるとき $F(Y)$ は $F(X)$ の左双対であることに注意する。

\begin{lemma}
\label{categories-lemma-left-dual}
$\mathcal{C}$ をモノイド圏とする。$Y$ が $X$ の左双対ならば、
$$
\Mor(Z' \otimes X, Z) = \Mor(Z', Z \otimes Y)
\quad\text{and}\quad
\Mor(Y \otimes Z', Z) = \Mor(Z', X \otimes Z)
$$
は $Z$ と $Z'$ に関して関手的である。
\end{lemma}

\begin{proof}
写像
$$
\Mor(Z' \otimes X, Z) \to
\Mor(Z' \otimes X \otimes Y, Z \otimes Y) \to
\Mor(Z', Z \otimes Y)
$$
を考える。ここで第二の射には $\eta$ を用いる。
また、写像の列
$$
\Mor(Z', Z \otimes Y) \to
\Mor(Z' \otimes X, Z \otimes Y \otimes X) \to
\Mor(Z' \otimes X, Z)
$$
を考える。第二の射には $\epsilon$ を用いる。
これらの射が互いに逆であることを示すため、
$a : Z' \to Z \otimes Y$ という写像を考える。次を示す必要がある。
$$
Z' \xrightarrow{\text{id}_{Z'} \otimes \eta}
Z' \otimes X \otimes Y \xrightarrow{a \otimes \text{id}_{X \otimes Y}}
Z \otimes Y \otimes X \otimes Y
\xrightarrow{\text{id}_Z \otimes \epsilon \otimes \text{id}_Y}
Z \otimes Y
$$
これは $a$ に等しい。
最初の二つの射の合成は
$(\text{id}_{Z \otimes Y} \otimes \eta) \circ a :
Z' \to Z \otimes Y \to Z \otimes Y \otimes X \otimes Y$ に等しい。
次に、$\text{id}_{Z \otimes Y} \otimes \eta$ と
$\text{id}_Z \otimes \epsilon \otimes \text{id}_Y$ の合成は、
双対の定義により恒等である。同様にもう一方の合成についても成り立つ。
第二の等式の証明は省略する。
\end{proof}

\begin{remark}
\label{categories-remark-left-dual-adjoint}
補題 \ref{categories-lemma-left-dual} は特に、
$Z \mapsto Z \otimes Y$ が $Z' \mapsto Z' \otimes X$ の右随伴であると述べている。
特に、随伴関手の一意性により、$X$ の左双対が存在すれば、
それは一意な同型を除いて一意である。
逆に、関手 $Z \mapsto Z \otimes Y$ が
関手 $Z' \mapsto Z' \otimes X$ の右随伴であると仮定する。
すなわち、関手的な全単射
$$
\Mor(Z' \otimes X, Z) \longrightarrow \Mor(Z', Z \otimes Y)
$$
が $Z$ と $Z'$ の両方について与えられているとする。
随伴の単位は写像
$$
\eta_Z : Z \to Z \otimes X \otimes Y
$$
を生成し、これは $Z$ に関して関手的である。また随伴の余単位は写像
$$
\epsilon_{Z'} : Z' \otimes Y \otimes X \to Z'
$$
を生成し、これは $Z'$ に関して関手的である。
特に
$\eta = \eta_\mathbf{1} : \mathbf{1} \to X \otimes Y$ および
$\epsilon = \epsilon_\mathbf{1} : Y \otimes X \to \mathbf{1}$ を得る。
単位、余単位、随伴同型の関係についての練習として、次が成り立つことを示せる。
$$
(\epsilon \otimes \text{id}_Y) \circ \eta_Y = \text{id}_Y
\quad\text{and}\quad
\epsilon_X \circ (\eta \otimes \text{id}_X) = \text{id}_X
$$
しかし、一般には $\eta_Y = \text{id}_Y \otimes \eta$ であることも、
$\epsilon_X = \epsilon \otimes \text{id}_X$ であることも知らないので、これは
$(\epsilon \otimes \text{id}_Y) \circ (\text{id}_Y \otimes \eta) =
\text{id}_Y$
および
$(\text{id}_X \otimes \epsilon) \circ (\eta \otimes \text{id}_X) =
\text{id}_X$
を示すには不十分である。このためには、随伴同型が次の性質をもつことを知れば十分である。
任意の $W, Z, Z'$ に対して図式
$$
\xymatrix{
\Mor(Z' \otimes X, Z) \ar[r] \ar[d]_{\text{id}_W \otimes -} &
\Mor(Z', Z \otimes Y) \ar[d]^{\text{id}_W \otimes -} \\
\Mor(W \otimes Z' \otimes X, W \otimes Z) \ar[r] &
\Mor(W \otimes Z', W \otimes Z \otimes Y)
}
$$
が可換となることである。
これが成り立つとき、{\it 随伴が与えられたテンソル構造と両立する}という。
したがって、$Z \mapsto Z \otimes Y$ が
$Z' \mapsto Z' \otimes X$ の右随伴であり、与えられたテンソル構造と両立するという条件は、
左双対であるという性質の同値な定式化である。
\end{remark}
\begin{lemma}
\label{categories-lemma-tensor-dual}
$\mathcal{C}$ をモノイド圏とする。$i = 1, 2$ に対して $Y_i$ が
$i = 1, 2$ に対する $X_i$ の左双対ならば、
$Y_2 \otimes Y_1$ は $X_1 \otimes X_2$ の左双対である。
\end{lemma}

\begin{proof}
随伴の一意性と注意 \ref{categories-remark-left-dual-adjoint} から従う。
\end{proof}

\noindent
$(\mathcal{C}, \otimes)$ に対する{\it 可換性制約}とは、関手的同型
$$
\psi : X \otimes Y \longrightarrow Y \otimes X
$$
であって、合成
$$
X \otimes Y \xrightarrow{\psi} Y \otimes X \xrightarrow{\psi} X \otimes Y
$$
が恒等であるものをいう。$\psi$ が与えられた結合性制約 $\phi$ と
{\it 両立する}とは、すべての対象 $X, Y, Z$ に対して図式
\begin{equation}
\label{categories-equation-hexagon}
\vcenter{
\xymatrix{
X \otimes (Y \otimes Z) \ar[r]_\phi \ar[d]^\psi &
(X \otimes Y) \otimes Z \ar[r]_\psi &
Z \otimes (X \otimes Y) \ar[d]^\phi \\
X \otimes (Z \otimes Y) \ar[r]^\phi &
(X \otimes Z) \otimes Y \ar[r]^\psi &
(Z \otimes X) \otimes Y
}
}
\end{equation}
が可換となることである。

\begin{definition}
\label{categories-definition-symmetric-monoidal-category}
$\mathcal{C}$ が圏であり、
$\otimes : \mathcal{C} \otimes \mathcal{C} \to \mathcal{C}$ が関手で、
$\phi$ が結合性制約、$\psi$ が $\phi$ と両立する可換性制約であるような
四つ組 $(\mathcal{C}, \otimes, \phi, \psi)$ を、
単位が存在するとき{\it 対称モノイド圏}と呼ぶ。
\end{definition}

\noindent
念のため、$(\mathcal{C}, \otimes, \phi, \psi)$ が対称モノイド圏ならば、
$(\mathcal{C}, \otimes, \phi)$ はモノイド圏であり、上で述べた言葉と記法を
用いることができる。

\begin{lemma}
\label{categories-lemma-symmetric-monoidal-coherence}
対称モノイド圏
$\mathcal{C}, \otimes, \phi, \psi, \mathbf{1}, 1$ において、
\begin{enumerate}
\item 射 $1 \circ \psi, 1 :
\mathbf{1} \otimes \mathbf{1} \to \mathbf{1}$ は一致し、
\item 射 $\text{id}_X \otimes l_Y,
(l_X \otimes \text{id}) \circ (\psi \otimes \text{id}_Y):
X \otimes \mathbf{1} \otimes Y \to X \otimes Y$ は一致する。
\end{enumerate}
対称モノイド圏は \cite[定理5.1]{associativity} の仮定を満たす。
\end{lemma}

\begin{proof}
一意な同型 $a : \mathbf{1} \to \mathbf{1}$ を用いて
$\psi = a \otimes \text{id}_\mathbf{1}$ と書ける。
補題 \ref{categories-lemma-monoidal-unit-unique} により
$a \otimes \text{id}_\mathbf{1} = \text{id}_\mathbf{1} \otimes a$ である。
$\psi$ の関手性は、図式
$$
\xymatrix{
(\mathbf{1} \otimes \mathbf{1}) \otimes \mathbf{1}
\ar[r]_\psi
\ar[d]_{1 \otimes \text{id}_\mathbf{1}} &
\mathbf{1} \otimes (\mathbf{1} \otimes \mathbf{1})
\ar[d]^{\text{id}_\mathbf{1} \otimes 1} \\
\mathbf{1} \otimes \mathbf{1}  \ar[r]^\psi &
\mathbf{1} \otimes \mathbf{1}
}
$$
が可換であることをいう。したがって上の射は
$a \otimes \text{id}_\mathbf{1} \otimes \text{id}_\mathbf{1}$ に等しい。
よって $X = Y = Z = \mathbf{1}$ に対する
(\ref{categories-equation-hexagon}) は、
$a \otimes \text{id}_\mathbf{1} \otimes \text{id}_\mathbf{1}$ が
自分自身の平方に等しいことをいう。したがって
$a = \text{id}_\mathbf{1}$ である。これで (1) が示された。

\medskip\noindent
(2) は、モノイド圏において始域と終域を $X$ と同一視すれば、
$\psi : X \otimes \mathbf{1} \to \mathbf{1} \otimes X$ が恒等である
という主張である。これは $\mathbf{1}, X, \mathbf{1}$ に対する
(\ref{categories-equation-hexagon}) の可換性、すなわち
$$
\xymatrix{
\mathbf{1} \otimes (X \otimes \mathbf{1}) \ar[r]_\phi \ar[d]^\psi &
(\mathbf{1} \otimes X) \otimes \mathbf{1} \ar[r]_\psi &
\mathbf{1} \otimes (\mathbf{1} \otimes X) \ar[d]^\phi \\
\mathbf{1} \otimes (\mathbf{1} \otimes X) \ar[r]^\phi &
(\mathbf{1} \otimes \mathbf{1}) \otimes X \ar[r]^\psi &
(\mathbf{1} \otimes \mathbf{1}) \otimes X
}
$$
から従う。この図式の $\psi$ 以外の射はすべて恒等に等しいことが分かっている。

\medskip\noindent
(1) と (2) に加えて、図式
(\ref{categories-equation-pentagon}), (\ref{categories-equation-triangle}),
(\ref{categories-equation-hexagon}) の可換性、および補題
\ref{categories-lemma-monoidal-coherence} の結果から、補題の最後の主張が従う。
\end{proof}

\noindent
\cite[定理5.1]{associativity} では、上の補題のような六つ組
$(\mathcal{C}, \otimes, \phi, \psi, \mathbf{1}, 1)$ がコヒーレントであることが説明されている。
言い換えると、対称モノイド圏（ここでの意味）では、すべての $n \geq 1$ と
$\{1, \ldots, n\}$ の置換 $\sigma$ に対して関手的同型
$$
X_1 \otimes \ldots \otimes X_n
\to
X_{\sigma(1)} \otimes \ldots \otimes X_{\sigma(n)}
$$
があり、これらの同型から作られるすべての図式が可換で、モノイド圏の構造
（例えばスロットに $\mathbf{1}$ を挿入したときの同型）と両立する。

\begin{lemma}
\label{categories-lemma-dual-symmetric}
$(\mathcal{C}, \otimes, \phi, \psi)$ を対称モノイド圏とする。
$X$ を $\mathcal{C}$ の対象とし、$Y$、
$\eta : \mathbf{1} \to X \otimes Y$、および
$\epsilon : Y \otimes X \to \mathbf{1}$ が、定義
\ref{categories-definition-dual} における $X$ の左双対であるとする。
このとき
$\eta' = \psi \circ \eta : \mathbf{1} \to Y \otimes X$ および
$\epsilon' = \epsilon \circ \psi : X \otimes Y \to \mathbf{1}$ は、
$X$ を $Y$ の左双対にする。
\end{lemma}

\begin{proof}
省略。ヒント：定義による楽しい練習問題である。
\end{proof}

\begin{definition}
\label{categories-definition-functor-symmetric-monoidal-categories}
$\mathcal{C}$ と $\mathcal{C}'$ を対称モノイド圏とする。
{\it 対称モノイド圏の関手}
$F : \mathcal{C} \to \mathcal{C}'$ とは、指定された関手 $F$ と、
$$
F(X) \otimes F(Y) \to F(X \otimes Y)
$$
$X$ と $Y$ に関して関手的な同型の組であって、
$F$ がモノイド圏の関手であり、すべての対象 $X$ と $Y$ に対して図式
$$
\xymatrix{
F(X) \otimes F(Y) \ar[r] \ar[d] &
F(X \otimes Y) \ar[d] \\
F(Y) \otimes F(X) \ar[r] &
F(Y \otimes X)
}
$$
が可換となるものをいう。
\end{definition}

\begin{remark}
\label{categories-remark-internal-hom-monoidal}
モノイド圏 $\mathcal{C}$ を考える。$\mathcal{C}$ が
{\it 内部Hom} を持つとは、$\mathcal{C}$ の任意の対象の組 $X, Y$ に対して、
次を満たす $\mathcal{C}$ の対象 $hom(X, Y)$ が存在することをいう。
$$
\Mor(X, hom(Y, Z)) = \Mor(X \otimes Y, Z)
$$
これは $X, Y, Z$ に関して関手的である。ヨネダの補題により、二変数関手
$(X, Y) \mapsto hom(X, Y)$ は、それが存在するなら一意な同型を除いて
定まる。内部Homが与えられると、次の標準的な射が得られる。
\begin{enumerate}
\item $hom(X, Y) \otimes X \to Y$、
\item $hom(Y, Z) \otimes hom(X, Y) \to hom(X, Z)$、
\item $Z \otimes hom(X, Y) \to hom(X, Z \otimes Y)$、
\item $Y \to hom(X, Y \otimes X)$、および
\item $hom(Y, Z) \otimes X \to hom(hom(X, Y), Z)$（
$\mathcal{C}$ が対称モノイド圏の場合）。
\end{enumerate}
すなわち、(1) の射は $\text{id}_{hom(X, Y)}$ を
次の写像
$\Mor(hom(X, Y), hom(X, Y)) \to \Mor(hom(X, Y) \otimes X, Y)$
で写したものである。
(2) の射を $hom(X, Z)$ の定義する性質によって構成するには、次の写像を
構成すればよい。
$$
hom(Y, Z) \otimes hom(X, Y) \otimes X \longrightarrow Z
$$
このような写像は、(1) により
$hom(X, Y) \otimes X \to Y$ および
$hom(Y, Z) \otimes Y \to Z$ の写像が得られるので存在する。
(3) の射を $hom(X, Z \otimes Y)$ の定義する性質によって構成するには、
次の写像を構成すればよい。
$$
Z \otimes hom(X, Y) \otimes X \to Z \otimes Y
$$
この写像には $\text{id}_Z \otimes a$ を用いる。
ここで $a$ は (1) の写像である。(4) の射を構成するため、
(3) により既に
$Y \otimes hom(X, X) \to hom(X, Y \otimes X)$ の写像を得ている。
したがって $\mathbf{1} \to hom(X, X)$ の写像を構成すれば十分であり、
これには $\Mor(\mathbf{1}, hom(X, X))$ において、
$\Mor(\mathbf{1} \otimes X, X)$ の標準同型
$\mathbf{1} \otimes X \to X$ に対応する元を取ればよい。
最後に (5) を考える。$hom(hom(X, Y), Z)$ の普遍性により、
次の写像を構成すれば十分である。
$$
hom(Y, Z) \otimes X \otimes hom(X, Y) \longrightarrow Z
$$
これは可換性制約を用いて最後の二つのテンソル積を交換し、その後
$hom(X, Y) \otimes X \to Y$ および
$hom(Y, Z) \otimes Y \to Z$ の写像を用いて構成する。
\end{remark}
\section{点線矢印の圏}
\label{categories-section-dotted-arrows}

\noindent
$(2,1)$-圏における「点線矢印の圏」について論じる。
これらは代数的スタックに対するさまざまな持ち上げ判定を定式化するときに現れる。
例えば \emph{スタックの射} の Section
\ref{stacks-morphisms-section-valuative} および
\emph{スタックの射についてさらに} の Section
\ref{stacks-more-morphisms-section-formally-smooth} を参照されたい。

\begin{definition}
\label{categories-definition-dotted-arrows}
$\mathcal{C}$ を $(2,1)$-圏とする。次の $2$-可換な
充実図式を考える。
\begin{equation}
\label{categories-equation-dotted-arrows}
\vcenter{
\xymatrix{
S \ar[r]_-x \ar[d]_j & X \ar[d]^f \\
T \ar[r]^-y \ar@{..>}[ru] & Y
}
}
\end{equation}
$\mathcal{C}$ において、図式の $2$-可換性を示す
$2$-同型
$$
\gamma : y \circ j \rightarrow f \circ x
$$
を固定する。(\ref{categories-equation-dotted-arrows}) と $\gamma$ が与えられたとき、
\emph{点線矢印} とは、三つ組 $(a, \alpha, \beta)$ であって、射
$a \colon T \to X$ と $2$-同型
$\alpha : a \circ j \to x$, $\beta : y \to f \circ a$
からなるものとする。
$\gamma = (\text{id}_f \star \alpha) \circ (\beta \star \text{id}_j)$
を満たすもの、すなわち次の図式
$$
\xymatrix{
& f \circ a \circ j \ar[rd]^{\text{id}_f \star \alpha} \\
y \circ j \ar[ru]^{\beta \star \text{id}_j} \ar[rr]^\gamma & &
f \circ x
}
$$
が可換となるものをいう。{\it 点線矢印の射}
$(a, \alpha, \beta) \to (a', \alpha', \beta')$ とは、
$2$-射 $\theta : a \to a'$ であって、
$\alpha = \alpha' \circ (\theta \star \text{id}_j)$ および
$\beta' = (\text{id}_f \star \theta) \circ \beta$
を満たすものをいう。
\end{definition}

\noindent
定義 \ref{categories-definition-dotted-arrows} の状況には、対応する
\emph{点線矢印の圏} がある。この圏はグルーポイドである。
一般には $\gamma$ に依存しうる。次の二つの補題は、
点線矢印の圏が $f$ に関する基底変換および合成に対して
良い振る舞いをすることを述べる。

\begin{lemma}
\label{categories-lemma-cat-dotted-arrows-base-change}
$\mathcal{C}$ を $(2,1)$-圏とする。次の $2$-可換図式が与えられているとする。
$$
\xymatrix{
S \ar[r]_-{x'} \ar[d]_j &
X' \ar[d]^p \ar[r]_q &
X \ar[d]^f \\
T \ar[r]^-{y'} &
Y' \ar[r]^g &
Y
}
$$
$\mathcal{C}$ において、ここで右側の正方形は、$2$-同型
$\phi \colon g \circ p \to f \circ q$ に関して $2$-カルテシアンである。
$2$-射
$\gamma' : y' \circ j \to p \circ x'$ を選ぶ。
$x = q \circ x'$, $y = g \circ y'$ と置き、
$2$-同型
$\gamma : y \circ j \to f \circ x$ を
$\gamma = (\phi \star \text{id}_{x'}) \circ (\text{id}_g \star \gamma')$
で定める。このとき、左の正方形と $\gamma'$ に対する点線矢印の圏
$\mathcal{D}'$ は、外側の長方形と $\gamma$ に対する点線矢印の圏
$\mathcal{D}$ と同値である。
\end{lemma}

\begin{proof}
$\mathcal{D}' \to \mathcal{D}$ という関手で、対象上では
$(a, \alpha, \beta) \mapsto (q \circ a, \text{id}_q \star \alpha,
(\phi \star \text{id}_a) \circ (\text{id}_g \star \beta))$
とし、射上では $\theta \mapsto \text{id}_q \star \theta$ とするものがある。
この関手 $\mathcal{D}' \to \mathcal{D}$ が同値であることは、
Section \ref{categories-section-2-fibre-products} のように
$2$-ファイバー積の普遍性から形式的に従う。詳細は省略する。
\end{proof}

\begin{lemma}
\label{categories-lemma-cat-dotted-arrows-composition}
$\mathcal{C}$ を $(2,1)$-圏とする。次の充実な $2$-可換図式が与えられているとする。
$$
\xymatrix{
S \ar[r]_-x \ar[dd]_j & X \ar[d]^f \\
& Y \ar[d]^g \\
T \ar[r]^-z \ar@{..>}[ruu] & Z
}
$$
$\mathcal{C}$ において、$2$-同型
$\gamma \colon z \circ j \to g \circ f \circ x$ を選ぶ。
$\mathcal{D}$ を外側の長方形と $\gamma$ に対する点線矢印の圏とする。
$\mathcal{D}'$ を、次の充実な正方形と $\gamma$ に対する点線矢印の圏とする。
$$
\xymatrix{
S \ar[r]_-{f \circ x} \ar[d]_j & Y \ar[d]^g \\
T \ar[r]^-z \ar@{..>}[ru] & Z
}
$$
このとき $\mathcal{D}$ は、次の性質を持つ圏 $\mathcal{D}''$ と同値である。
すなわち、$\mathcal{D}'' \to \mathcal{D}'$ という関手が存在し、
$\mathcal{D}''$ を $\mathcal{D}'$ 上のグルーポイドにファイバー化された圏にし、
そのファイバー圏が、次の形のある充実な正方形と
$y \circ j \to f \circ x$ という $2$-同型の選択に対する点線矢印の圏と
同型になる。
$$
\xymatrix{
S \ar[r]_-x \ar[d]_j & X \ar[d]^f \\
T \ar[r]^-y \ar@{..>}[ru] & Y
}
$$
\end{lemma}

\begin{proof}
対象が五つ組
$(a,\alpha,\beta,b,\eta)$
であり、$(a,\alpha,\beta)$ が $\mathcal{D}$ の対象、
$b \colon T \rightarrow Y$ が $1$-射、
$\eta \colon b \rightarrow f \circ a$ が $2$-同型となるような圏
$\mathcal{D}''$ を構成する。
$\mathcal{D}''$ における
$(a,\alpha,\beta,b,\eta) \rightarrow (a',\alpha',\beta',b',\eta')$
という射は組 $(\theta_1,\theta_2)$ である。ここで
$\theta_1 \colon a \rightarrow a'$ は $\mathcal{D}$ における
$(a, \alpha, \beta) \rightarrow (a', \alpha', \beta')$ という射を定め、
$\theta_2 \colon b \rightarrow b'$ は $2$-同型であり、両立条件
$\eta' \circ \theta_2 = (\text{id}_f \star \theta_1) \circ \eta$
を満たす。

\medskip\noindent
$\mathcal{D}'' \rightarrow \mathcal{D}'$ という関手で、対象上では
$(a, \alpha, \beta, b, \eta) \mapsto
(b, (\text{id}_f \star \alpha) \circ (\eta \star \text{id}_j),
(\text{id}_g \star \eta^{-1}) \circ \beta)$
とし、射上では $(\theta_1,\theta_2) \mapsto \theta_2$ とするものがある。
このとき $\mathcal{D}'' \rightarrow \mathcal{D}'$ はグルーポイドにファイバー化されている。

\medskip\noindent
$(y, \delta, \epsilon)$ が $\mathcal{D}'$ の対象ならば、最後に表示した図式に対し、
$y \circ j \rightarrow f \circ x$ が $\delta$ で与えられる点線矢印の圏を
$\mathcal{D}_{y,\delta}$ と書く。
$\mathcal{D}_{y,\delta} \rightarrow \mathcal{D}''$ という関手で、対象上では
$(a, \alpha, \eta) \mapsto
(a, \alpha, (\text{id}_g \star \eta) \circ \epsilon, y, \eta)$
とし、射上では $\theta \mapsto (\theta, \text{id}_y)$ とするものがある。
これは $\mathcal{D}_{y,\delta}$ から
$\mathcal{D}'' \rightarrow \mathcal{D}'$ の
$(y,\delta,\epsilon)$ 上のファイバー圏への同型を与える。

\medskip\noindent
さらに $\mathcal{D} \rightarrow \mathcal{D}''$ という関手で、対象上では
$(a,\alpha,\beta) \mapsto (a,\alpha,\beta,f \circ a, \text{id}_{f \circ a})$
とし、射上では $\theta \mapsto (\theta, \text{id}_f \star \theta)$ とするものがある。
この関手は充満忠実かつ本質的全射であり、したがって同値である。詳細は省略する。
\end{proof}
